\chapter{Geometric Deep Learning}\label{ch:gdl}

In this chapter we explore \gls{GDL}, a paradigm that extends \gls{DL} beyond traditional Euclidean domains toward arbitrary geometric structures. While classical neural networks process data on Euclidean domains, such as regular grids (images) or linear sequences (text), the real world presents a wealth of complex geometric structures: molecules are three-dimensional graphs with rotational symmetries, social networks form irregular graphs with emergent topological properties, and physical surfaces constitute Riemannian manifolds with intrinsic curvature. \gls{GDL} provides the mathematical framework needed to build neural architectures that respect and exploit these natural geometries and are not restricted solely to representations of Euclidean domains.

Another of the fundamental contributions of \gls{GDL} is the demonstration that the seemingly disparate architectures of \gls{DL}, such as \glspl{CNN}, \glspl{Transformer}, or graph networks, are particular instances of a unifying principle: the equivariant processing of signals over geometric domains with specific symmetries. \glspl{CNN} emerge as the case of regular graphs with translational invariance, \glspl{Transformer} as complete graphs where attention implements weighted message passing, and \glspl{GNN} as the general case of arbitrary graphs with permutation invariance. This unifying perspective is grounded in deep mathematical tools from group theory, differential geometry, and harmonic analysis, providing not only theoretical foundations for why existing architectures work, but also design principles for new architectures with stability and generalization guarantees.

The impact of \gls{GDL} goes beyond theoretical foundations. By respecting the intrinsic geometric structure of data, architectures based on \gls{GDL} have achieved revolutionary advances in critical domains: AlphaFold solved the protein folding problem through SE(3)-equivariant networks, transforming structural biology; molecular graph networks have accelerated drug discovery by orders of magnitude; and geometric 3D vision models have enabled robust understanding of complex scenes. These successes demonstrate that geometry is not an optional technical detail, but the fundamental organizing principle for effective learning in complex domains. The theoretical foundation of this paradigm is established rigorously through the Equivariant Universal Approximation Theorem (\autoref{theorem:universal_approximation_equivariant}), developed in \autoref{sec:grupos_simetrias}.

\section{The Geometric Deep Learning Paradigm}\label{sec:paradigma_gdl}

\subsection{The Five Fundamental G's}\label{subsec:cinco_gs}

The \gls{GDL} framework is structured around five fundamental concepts, known as the ``5 G's'' \citep{bronstein2017geometricdeeplearning}, which provide a unified framework for understanding and designing neural architectures that exploit the properties of the geometries in which data reside:

\begin{enumerate}
    \item \textbf{Geometry}: The intrinsic structure of the data domain and its metric properties. This includes Riemannian manifolds (such as spheres for omnidirectional data), hyperbolic spaces (for hierarchical data), and product spaces (for multi-structured data). The geometry determines which operations are natural and efficient in each domain (see \autoref{sec:geometria_dominios}).

    \item \textbf{Groups}: The symmetries that preserve the structure of the problem through transformations. The most common groups include translations $\mathbb{R}^d$ (\glspl{CNN}), rotations $SO(n)$ (rotation-equivariant networks), permutations $S_n$ (graph networks), and Lorentz groups (relativistic networks). The correct identification of the symmetry group is crucial for architectural design (see \autoref{sec:grupos_simetrias}).

    \item \textbf{Graphs}: Discretizations of continuous geometric domains that capture local neighborhood relations. Graphs can be regular (uniform meshes), irregular (social networks), directed (causal networks), or dynamic (evolving systems). They provide the computational substrate for many \gls{GDL} architectures (see \autoref{sec:grafos_discretizacion}).

    \item \textbf{Geodesics}: Optimal paths that respect the intrinsic geometry of the space. On Riemannian manifolds, geodesics minimize distance; on graphs, they correspond to shortest paths; in information spaces, they define directions of maximum gradient. Geodesics determine how information propagates in geometric architectures (see \autoref{sec:geodesicas_propagacion}).

    \item \textbf{Gauges}: Invariances under gauge transformations that preserve the observable physics of the system. They include electromagnetic invariances, coordinate freedoms on manifolds, and orientation choices in undirected graphs. Gauges ensure that the learned representations are intrinsic to the geometric domain (see \autoref{sec:gauges_invarianzas}).
\end{enumerate}

\begin{figure}[htbp]
\centering
\includegraphics[width=0.7\textwidth]{figures/five_gs_framework.pdf}
\caption[The five G's of Geometric Deep Learning]{The five fundamental G's of Geometric Deep Learning. Each component represents an essential aspect of the unifying framework: \textbf{Geometry} (domains), \textbf{Groups} (symmetries), \textbf{Graphs} (discretization), \textbf{Geodesics} (distances), and \textbf{Gauges} (local coordinates).}
\label{fig:five_gs_framework}
\end{figure}

\autoref{fig:jerarquia_5gs} complements this view by showing the hierarchical structure of the five G's, from the most fundamental base (the geometry of the domain) to the most specific levels of abstraction.

\begin{figure}[htbp]
\centering
\begin{tikzpicture}[scale=0.75]
% Definir estilos
\tikzset{
    level/.style={rectangle, rounded corners=3pt, draw, minimum width=2.8cm, minimum height=0.7cm, align=center, font=\small\bfseries},
    arrow/.style={->, thick, >=stealth, gray!70},
    desc/.style={font=\scriptsize, text=black!60, align=left}
}

% Niveles de la pirámide (de abajo hacia arriba)
\node[level, fill=gdlgreen!15, minimum width=10cm] (geom) at (0,0) {Geometry};
\node[level, fill=gdlgreen!25, minimum width=8cm] (grupos) at (0,1.2) {Groups};
\node[level, fill=gdlgreen!35, minimum width=6cm] (grafos) at (0,2.4) {Graphs};
\node[level, fill=gdlgreen!45, minimum width=4cm] (geod) at (0,3.6) {Geodesics};
\node[level, fill=gdlgreen!55, minimum width=2.5cm] (gauges) at (0,4.8) {Gauges};

% Flechas de relación
\draw[arrow] (geom.north) ++(0,0.1) -- ++(0,0.3);
\draw[arrow] (grupos.north) ++(0,0.1) -- ++(0,0.3);
\draw[arrow] (grafos.north) ++(0,0.1) -- ++(0,0.3);
\draw[arrow] (geod.north) ++(0,0.1) -- ++(0,0.3);

% Descripciones a la derecha (posiciones relativas al borde de cada nivel)
\node[desc, anchor=west] at (geom.east) {\quad Domain $\Omega$ where data live};
\node[desc, anchor=west] at (grupos.east) {\quad Symmetries $G$ that preserve structure};
\node[desc, anchor=west] at (grafos.east) {\quad Discretization of the domain};
\node[desc, anchor=west] at (geod.east) {\quad Information propagation};
\node[desc, anchor=west] at (gauges.east) {\quad Local invariances};

% Símbolos a la izquierda (posiciones relativas al borde de cada nivel)
\node[font=\scriptsize, anchor=east] at (geom.west) {$\mathbb{R}^n$, $\mathcal{M}$, $\mathcal{G}$\quad};
\node[font=\scriptsize, anchor=east] at (grupos.west) {$\text{SO}(n)$, $S_n$, $E(n)$\quad};
\node[font=\scriptsize, anchor=east] at (grafos.west) {$G=(V,E)$\quad};
\node[font=\scriptsize, anchor=east] at (geod.west) {$d(x,y)$, paths\quad};
\node[font=\scriptsize, anchor=east] at (gauges.west) {frames\quad};

\end{tikzpicture}
\caption[Hierarchy of the five G's of GDL]{Hierarchical structure of the five G's of Geometric Deep Learning. From the base, \textbf{Geometry} defines the domain where data reside; \textbf{Groups} capture the symmetries of the problem; \textbf{Graphs} discretize the domain; \textbf{Geodesics} determine how information propagates; and \textbf{Gauges} handle local invariances in coordinate systems.}
\label{fig:jerarquia_5gs}
\end{figure}

From this unifying perspective, \gls{GDL} establishes that effective processing of signals over geometric domains requires respecting three fundamental principles:

\begin{enumerate}
    \item \textbf{Geometric locality}: Operations must be local with respect to the metric of the domain. Geometric architectures exploit local structure through neighborhoods defined by the underlying geometry, allowing information to propagate consistently with the intrinsic distances of the space (see \autoref{sec:grafos_discretizacion} for the formalization through graphs).

    \item \textbf{Equivariance}: Transformations must preserve the relevant symmetries of the problem. An equivariant function commutes with the actions of the symmetry group, ensuring that equivalent transformations of the input produce corresponding transformations in the output (see Section~\ref{subsec:equivarianza_invarianza} for the complete theory).

    \item \textbf{Compositionality}: Complex representations are built hierarchically through composition of simpler operations. The composition of equivariant layers preserves equivariance end to end (\autoref{theorem:equivariant_composition}), enabling deep architectures that maintain the desired geometric properties. This property is fundamental to the Equivariant Universal Approximation Theorem (\autoref{theorem:universal_approximation_equivariant}), which guarantees the expressive power of equivariant networks.
\end{enumerate}

\section{Geometry: The Domains of Data}\label{sec:geometria_dominios}

To design effective neural architectures in \gls{GDL}, we must understand the \textbf{intrinsic geometry} of the domains (\autoref{def:dominio_datos}) where data reside. Different types of data naturally inhabit different geometric spaces, and recognizing and respecting these geometries is decisive for the success of Deep Learning. The formal concepts of symmetry and equivariance are developed in \autoref{sec:grupos_simetrias}.

The choice of the geometric domain is neither arbitrary nor merely technical: it reflects the intrinsic structure of the data and determines which operations are natural and efficient. While the Deep Learning tradition has favored standard vector representations (vectors, matrices, tensors in $\mathbb{R}^n$) in a Euclidean space, much real-world data possesses intrinsic geometries that these representations do not adequately capture. Molecules naturally live in spaces with symmetries of the Euclidean group $E(3)$, hierarchies are better represented in hyperbolic spaces $\mathbb{H}^n$, and omnidirectional directions inhabit spheres $S^n$. Forcing this data into Euclidean spaces destroys valuable structural information and hinders effective learning.

\subsection{Geometry of Data}\label{subsec:geometria_datos}

In \autoref{ch:neural_networks} we established the fundamental conceptual framework of signals over structured domains (see \autoref{sec:datos_estructura}), where we introduced the distinction between \textit{domain} ($\Omega$), \textit{structure} ($\mathcal{S}$), and \textit{signal} ($f$). This framework allowed us to formally characterize different types of data through the triple $(\Omega, \mathcal{S}, f)$ or, equivalently, through the geometry-attributes pair $(\mathcal{D}, f)$ where $\mathcal{D} = (\Omega, \mathcal{S})$ represents the geometry of the domain and $f$ the observed attributes.

The concept of \textit{geometry of data} deepens this perspective, referring to the natural geometric structures that emerge from the intrinsic nature of the data, beyond their superficial numerical representations. \gls{GDL} proposes that these intrinsic geometries $\mathcal{D} = (\Omega, \mathcal{S})$ are not merely representation details, but fundamental properties that determine which neural architectures are appropriate and effective for each type of data.

\paragraph{From Euclidean representations to complex geometric structures}

Traditionally, machine learning has operated under the paradigm that all data can be represented as:

\begin{itemize}
    \item \textbf{Vectors}: $\mathbb{R}^n$ (sequences, embeddings, extracted features)
    \item \textbf{Matrices}: $\mathbb{R}^{n \times m}$ (grayscale images, tabular data tables)
    \item \textbf{Higher-order tensors}: $\mathbb{R}^{n_1 \times n_2 \times \cdots \times n_k}$ (RGB images as $\mathbb{R}^{H \times W \times 3}$, video as $\mathbb{R}^{T \times H \times W \times 3}$)
\end{itemize}

This simplifying perspective has been tremendously successful. However, this view is fundamentally limited. Let us revisit some of the examples presented in \autoref{ch:neural_networks} to examine the limitations of Euclidean representations from the perspective of \gls{GDL}:

\begin{itemize}
    \item \textbf{Molecules as 3D graphs} (cf.~\autoref{ex:molecula}): A molecule is not simply a set of atomic coordinates in $\mathbb{R}^{3n}$. As we saw earlier, its structure combines \textit{discrete connectivity} (graph of covalent chemical bonds) with \textit{continuous three-dimensional geometry} (spatial positions of atoms). More fundamentally, the physicochemical properties are invariant under the Euclidean group $E(3)$ of translations, rotations, and reflections. Vectorizing the coordinates arbitrarily breaks these natural symmetries and forces neural networks to relearn fundamental physical invariances.

    \item \textbf{Social networks} (cf.~\autoref{ex:red_social}): A social graph $G = (V, E)$ lacks a natural ordering among its nodes. Assigning indices $\{1, 2, \ldots, n\}$ to users and attempting to represent the graph as an adjacency matrix $A \in \mathbb{R}^{n \times n}$ introduces a completely artificial ordering. The relevant structure is purely topological: who is connected to whom, not in which ``position'' each user is.

    \item \textbf{Physical surfaces} (cf.~\autoref{ex:superficie_riemanniana}): A 3D surface (the cerebral cortex, the surface of a protein, a scanned object) is naturally a 2D Riemannian manifold embedded in $\mathbb{R}^3$. The geodesic distances over the surface and the intrinsic curvature are fundamental geometric properties independent of how we parameterize or triangulate the surface. A triangular mesh with $n$ vertices is not simply a tensor $\mathbb{R}^{n \times 3}$ of coordinates: it is a discretization of an underlying continuous manifold.
\end{itemize}

These examples illustrate a fundamental principle: data possesses intrinsic geometries that standard vector representations do not fully capture. \gls{GDL} proposes a paradigm shift: instead of forcing all data into $\mathbb{R}^n$, we must design architectures that operate directly on the natural geometries of the data.

\subsection{Traditional Euclidean Geometry}\label{subsec:geometria_euclidiana}

In order to understand the differences between geometries we must first understand traditional Euclidean geometry. The Euclidean space $\mathbb{R}^n$ constitutes the most familiar and widely used geometric domain in classical Deep Learning. Its ubiquity is not accidental: the Euclidean structure provides computational simplicity, well-developed analytical tools, and uniform representations that have proven effective in many applications.

\paragraph{Euclid's postulates}

Euclidean geometry is grounded in the five postulates established by Euclid in his \textit{Elements} \citep{euclid1956elements}. These axioms define the fundamental properties of flat space:

\begin{enumerate}
    \item \textbf{Postulate of the line}: Through two distinct points passes exactly one line.
    \item \textbf{Postulate of the segment}: Every segment can be extended indefinitely into a straight line.
    \item \textbf{Postulate of the circle}: Given a point and a distance, a circle can be drawn with center at that point and radius equal to that distance.
    \item \textbf{Postulate of right angles}: All right angles are equal to one another.
    \item \textbf{Postulate of parallels}: Through a point external to a line, exactly one line parallel to the first passes.
\end{enumerate}

The fifth postulate, known as the \textit{parallel postulate} or \textit{Euclid's postulate}, is historically the most controversial. For more than two millennia, mathematicians attempted to derive it from the other four, until in the 19th century Gauss, Bolyai, and Lobachevsky independently demonstrated that negating it produces consistent geometries: the \textit{non-Euclidean geometries} \citep{gray2007worlds}.

\paragraph{Modern mathematical formalization}

Intuitively, \textit{Euclidean geometry} is that which uses the Euclidean metric, or $\ell^2$ norm, to measure distances

\[
d(\mathbf{x}, \mathbf{y}) = \|\mathbf{x} - \mathbf{y}\|_2 = \sqrt{\sum_{i=1}^n (x_i - y_i)^2}
\]

This metric defines notions of distance, neighborhood, and continuity that are uniform and well behaved: it satisfies the properties we expect of everyday ``flat'' space: shortest paths are straight lines, triangles have angles that sum to $\pi$, and parallel lines remain equidistant. Moreover, the fundamental operations (inner products, projections, convolutions) have efficient closed forms and well-understood analytical properties.

In the language of modern differential geometry, the Euclidean space $\mathbb{R}^n$ is a Riemannian manifold (\autoref{def:metrica_riemanniana} and \autoref{subsec:variedades_riemannianas}) with the following characteristic properties (the referenced concepts are developed formally later in this chapter):

\begin{itemize}
    \item \textbf{Flat Riemannian metric}: The Riemannian metric $g$ is constant and equal to the identity ($g_{ij} = \delta_{ij}$), which induces the standard Euclidean distance (see \autoref{subsec:variedades_riemannianas}).
    \item \textbf{Zero curvature}: The sectional curvature (\autoref{def:curvatura_preliminar}) is $K = 0$ at every point and direction.
    \item \textbf{Straight geodesics}: The geodesics (\autoref{def:geodesica}) are straight lines, formalizing Euclid's postulates 1 and 2.
    \item \textbf{Symmetry group}: The group of isometries is the \textit{Euclidean group} $E(n) = \mathbb{R}^n \rtimes O(n)$ (\autoref{def:grupo}, \autoref{def:producto_semidirecto}), which includes rotations, reflections, and translations combined through the semidirect product.
\end{itemize}

The deep connection between Euclid's classical geometry and modern differential geometry is established through the following fundamental result (see \autoref{sec:gauss_bonnet} of Appendix~\ref{ch:teoremas_complementarios} for the generalized theorem and its proof):

\begin{theorem}[Gauss-Bonnet for geodesic triangles]\label{thm:gauss_bonnet_triangulo}
Let $\mathcal{M}$ be a Riemannian surface with Gaussian curvature $K$, and let $T \subset \mathcal{M}$ be a geodesic triangle with interior angles $\alpha$, $\beta$, $\gamma$. Then:
\begin{equation}\label{eq:gauss_bonnet_triangulo}
\alpha + \beta + \gamma = \pi + \int_T K \, dA
\end{equation}
where $dA$ denotes the area element induced by the Riemannian metric.
\end{theorem}

This theorem \citep{lee2018introduction} has a fundamental consequence: when $K = 0$ (zero curvature), we recover the classical Euclidean result that the angles of a triangle sum to exactly $\pi$. Moreover, the theorem demonstrates that Euclid's fifth postulate (parallels) is \textit{equivalent} to the condition of Gaussian curvature $K = 0$.

\paragraph{Limitations of Euclidean geometry}

However, Euclidean geometry is not universal. When data has intrinsically non-Euclidean structure, forcing it into standard vector/matrix representations destroys important structural information:

\begin{itemize}
    \item \textbf{Spherical images} ($360^\circ$ panoramas): They naturally live on the sphere $S^2$, but the equirectangular projection to $\mathbb{R}^{H \times W}$ introduces extreme distortions at the poles. The $\ell^2$ distance between pixels does not respect the geodesic distances over $S^2$. Rotations in $SO(3)$ become complex and nonlinear transformations in the matrix representation.

    \item \textbf{3D point clouds} (cf.~\autoref{ex:nube_puntos}): Sets $\{p_1, \ldots, p_n\} \subset \mathbb{R}^3$ without a canonical ordering. Forcibly vectorizing them into $\mathbb{R}^{3n}$ imposes an arbitrary ordering that is not intrinsic to the geometric object, violating the natural permutation invariance.

    \item \textbf{Molecular structures}: They have complex three-dimensional geometry with symmetries of the Euclidean group $E(3)$. Representing them as vectors in $\mathbb{R}^d$ breaks the fundamental physical invariances under rotations and translations.

    \item \textbf{Hierarchical data}: They naturally live in hyperbolic spaces $\mathbb{H}^n$ where distance captures hierarchy. Embeddings in Euclidean $\mathbb{R}^n$ require exponentially higher dimensionality to preserve distances \citep{nickel2017poincare}.
\end{itemize}

\subsection{Foundations of Manifolds}\label{subsec:fundamentos_variedades}

To formalize the idea that data resides in spaces with intrinsic geometry, we need the mathematical concept of \textit{manifold}. This subsection develops the necessary foundations, beginning with topological preliminaries and progressing toward differentiable manifolds.

\subsubsection{Topological preliminaries}\label{subsubsec:preliminares_topologicos}

Before introducing the concept of manifold, we need to establish the necessary topological foundations. For a complete exposition of these concepts, see~\citet{munkres2000topology}.

\begin{definition}[Topological space]\label{def:espacio_topologico}
A \textit{topological space} is a pair $(X, \tau)$ where $X$ is a set and $\tau$ is a collection of subsets of $X$ (called \textbf{open sets}) that satisfies:
\begin{enumerate}
	\item $\emptyset, X \in \tau$ (the empty set and the whole set are open)
	\item The arbitrary union of open sets is open: if $\{U_\alpha\}_{\alpha \in I} \subset \tau$, then $\bigcup_{\alpha \in I} U_\alpha \in \tau$
	\item The finite intersection of open sets is open: if $U_1, \ldots, U_n \in \tau$, then $\bigcap_{i=1}^n U_i \in \tau$
\end{enumerate}
The collection $\tau$ is called a \textbf{topology} on $X$.
\end{definition}

\begin{definition}[Hausdorff space]\label{def:espacio_hausdorff}
A topological space $(X, \tau)$ is a \textit{Hausdorff space} (or $T_2$) if for any two distinct points $x, y \in X$ with $x \neq y$, there exist disjoint open sets $U, V \in \tau$ such that $x \in U$, $y \in V$ and $U \cap V = \emptyset$.

Intuitively, this means that points can be ``separated'' by open sets.
\end{definition}

\begin{definition}[Countable base]\label{def:base_numerable}
A topological space $(X, \tau)$ has a \textit{countable base} (or satisfies the second axiom of countability) if there exists a countable collection $\mathcal{B} = \{B_n\}_{n \in \mathbb{N}}$ of open sets such that every open set of $\tau$ can be expressed as a union of elements of $\mathcal{B}$.

Formally, $\mathcal{B} \subset \tau$ is a base if for every $U \in \tau$ and every $x \in U$, there exists $B \in \mathcal{B}$ such that $x \in B \subseteq U$.
\end{definition}

\begin{definition}[Homeomorphism]\label{def:homeomorfismo}
Let $(X, \tau_X)$ and $(Y, \tau_Y)$ be topological spaces. A function $f: X \to Y$ is a \textit{homeomorphism} if:
\begin{enumerate}
	\item $f$ is bijective
	\item $f$ is continuous: for every $V \in \tau_Y$, we have $f^{-1}(V) \in \tau_X$
	\item $f^{-1}$ is continuous: for every $U \in \tau_X$, we have $(f^{-1})^{-1}(U) = f(U) \in \tau_Y$
\end{enumerate}
Two spaces are \textit{homeomorphic} if there exists a homeomorphism between them. Intuitively, they are ``topologically identical''.
\end{definition}

\subsubsection{Topological manifolds}\label{subsubsec:variedades_topologicas}

With these concepts established, we can now rigorously define topological manifolds. For a complete development of manifold theory, we recommend~\citet{lee2012introduction}.

\begin{definition}[Topological manifold]\label{def:variedad_topologica}
A \textit{topological manifold} of dimension $m$ is a topological space $\mathcal{M}$ that satisfies:
\begin{enumerate}
	\item $\mathcal{M}$ is a Hausdorff space
	\item $\mathcal{M}$ has a countable base
	\item $\mathcal{M}$ is locally Euclidean of dimension $m$: for every point $p \in \mathcal{M}$, there exists an open neighborhood $U \ni p$ and a homeomorphism $\phi: U \to V$ where $V$ is an open set of $\mathbb{R}^m$
\end{enumerate}
\end{definition}

Intuitively, a manifold is a space that locally ``resembles'' Euclidean space, although globally it may have a very different structure. We can illustrate it with the following examples:

\begin{itemize}
    \item \textbf{The Earth as a 2D manifold}: Although we know the Earth is approximately spherical, when we walk through a city, we locally perceive the ground as flat. This is precisely the concept of a manifold: a space that can be globally curved, but locally behaves like the familiar Euclidean space.
    \item \textbf{A crumpled sheet of paper}: If we crumple a sheet, it remains locally two-dimensional (we can place a finger at any point and move it in two independent directions), but globally it has a complex shape in three-dimensional space.
    \item \textbf{The configuration space of a robot}: A robotic arm with multiple joints has a configuration space that is not Euclidean, but locally we can parameterize nearby positions with standard coordinates.
\end{itemize}

This intuition is crucial because many high-dimensional datasets (images, text, signals) actually ``live'' on manifolds of much lower dimension, and exploiting this significantly improves learning.

\begin{definition}[Atlas and charts]\label{def:atlas_cartas}
An \textit{atlas} of dimension $m$ for a topological space $\mathcal{M}$ is a collection $\mathcal{A} = \{(U_\alpha, \phi_\alpha)\}_{\alpha \in I}$ where:
\begin{itemize}
	\item $\{U_\alpha\}_{\alpha \in I}$ is an open cover of $\mathcal{M}$
	\item Each $\phi_\alpha: U_\alpha \to V_\alpha \subset \mathbb{R}^m$ is a homeomorphism (called a \textit{chart})
\end{itemize}

The transition functions $\phi_\beta \circ \phi_\alpha^{-1}: \phi_\alpha(U_\alpha \cap U_\beta) \to \phi_\beta(U_\alpha \cap U_\beta)$ determine the structure of the manifold.
\end{definition}

\begin{figure}[htbp]
\centering
\includegraphics[width=0.75\textwidth]{figures/charts_atlas.pdf}
\caption[Atlas and charts of a manifold]{Atlas $\mathcal{A} = \{(U_\alpha, \phi_\alpha)\}_{\alpha \in I}$ of a manifold $\mathcal{M}$: each chart $(U_\alpha, \phi_\alpha)$ is a homeomorphism $\phi_\alpha: U_\alpha \to V_\alpha \subset \mathbb{R}^m$ that defines local coordinates on an open set of $\mathcal{M}$. The transition function $\phi_\beta \circ \phi_\alpha^{-1}: \phi_\alpha(U_\alpha \cap U_\beta) \to \phi_\beta(U_\alpha \cap U_\beta)$ connects the representations of overlapping charts.}
\label{fig:charts_atlas}
\end{figure}

\subsubsection{Differentiable manifolds and tangent space}\label{subsubsec:variedades_diferenciables}

\begin{definition}[Smooth function]\label{def:funcion_suave}
Let $U \subseteq \mathbb{R}^n$ be an open set. A function $f: U \rightarrow \mathbb{R}^m$ is \textit{smooth} (or of \textit{class $C^\infty$}) if it possesses continuous partial derivatives of all orders.
\end{definition}

More generally, a function between differentiable manifolds is smooth if its expression in local coordinates is smooth.

\begin{definition}[Differentiable manifold]\label{def:variedad_diferenciable}
A topological manifold $\mathcal{M}$ is \textit{differentiable} (or smooth) if it possesses an atlas $\mathcal{A} = \{(U_\alpha, \phi_\alpha)\}$ such that all its transition maps $\phi_\beta \circ \phi_\alpha^{-1}$ are of class $C^\infty$.
\end{definition}

The differentiable structure allows defining concepts from calculus (derivatives, integrals, differential equations) on the manifold. The most fundamental concept is that of the \textit{tangent space}, which captures the notion of ``infinitesimal directions'' at each point. We present two equivalent constructions that offer complementary perspectives.

\begin{definition}[Tangent space via curves]\label{def:espacio_tangente_curvas}
Let $\mathcal{M}$ be a differentiable manifold and $p \in \mathcal{M}$. Two curves $\gamma_1, \gamma_2: (-\epsilon, \epsilon) \to \mathcal{M}$ with $\gamma_1(0) = \gamma_2(0) = p$ are said to be \textit{tangent at $p$} if for some (and therefore every) chart $(U, \phi)$ with $p \in U$:
\[
\frac{d}{dt}\bigg|_{t=0} (\phi \circ \gamma_1)(t) = \frac{d}{dt}\bigg|_{t=0} (\phi \circ \gamma_2)(t)
\]
The \textit{tangent space} $T_p\mathcal{M}$ is the set of equivalence classes $[\gamma]$ of curves tangent at $p$.
\end{definition}

This definition captures the geometric intuition: a tangent vector represents a ``direction of motion'' from the point $p$, determined by the initial velocity of any curve passing through $p$ in that direction.

\begin{remark}[Vector space structure]\label{rem:estructura_vectorial_curvas}
The set $T_p^{\text{curvas}}\mathcal{M}$ admits a vector space structure induced by any chart $(U, \phi)$ with $p \in U$. The map $[\gamma] \mapsto (\phi \circ \gamma)'(0) \in \mathbb{R}^m$ is a bijection between $T_p^{\text{curvas}}\mathcal{M}$ and $\mathbb{R}^m$, and we \textit{transport} the vector structure of $\mathbb{R}^m$ to $T_p^{\text{curvas}}\mathcal{M}$ through this bijection. That is, we define:
\[
\alpha[\gamma_1] + [\gamma_2] \;=\; \bigl[\gamma\bigr], \quad \text{where $(\phi \circ \gamma)'(0) = \alpha(\phi \circ \gamma_1)'(0) + (\phi \circ \gamma_2)'(0)$}.
\]
This structure is \textit{independent of the chosen chart}: if $(V, \psi)$ is another chart with $p \in V$, the chain rule applied to the transition function $\psi \circ \phi^{-1}$ guarantees that the isomorphism induced by $\psi$ is compatible with the one induced by $\phi$.
\end{remark}

\begin{definition}[Tangent space via derivations]\label{def:espacio_tangente_derivaciones}
Let $\mathcal{M}$ be a differentiable manifold and $p \in \mathcal{M}$. A \textit{derivation at $p$} is a linear map $v: C^\infty(\mathcal{M}) \to \mathbb{R}$ that satisfies the Leibniz rule for all $f, g \in C^\infty(\mathcal{M})$:
\[
v(fg) = v(f)g(p) + f(p)v(g)
\]
The \textit{tangent space} $T_p\mathcal{M}$ is the vector space of all derivations at $p$.
\end{definition}

This algebraic definition captures the idea that a tangent vector is a ``directional derivative operator'': given a vector $v$ and a function $f$, $v(f)$ represents the rate of change of $f$ in the direction $v$.

\begin{lemma}[Decomposition of derivations in coordinates]\label{lem:descomposicion_derivaciones}
Let $\mathcal{M}$ be a differentiable manifold of dimension $m$, $p \in \mathcal{M}$ and $(U, \phi)$ a chart with $\phi(p) = 0$ and coordinate functions $x^1, \ldots, x^m$. Then:
\begin{enumerate}
	\item Derivations are local operators: if $f = g$ in a neighborhood of $p$, then $v(f) = v(g)$ for every derivation $v$ at $p$. This allows evaluating $v$ on functions defined only on $U$, extending them to $\mathcal{M}$ through bump functions, with a result independent of the extension.
	\item Every derivation $v$ at $p$ is expressed as:
	\[
	v = \sum_{j=1}^{m} v(x^j) \frac{\partial}{\partial x^j}\bigg|_p
	\]
	where $\frac{\partial}{\partial x^j}\big|_p(f) = \frac{\partial(f \circ \phi^{-1})}{\partial x^j}\big|_{\phi(p)}$. In particular, $\dim T_p^{\text{deriv}}\mathcal{M} = m$.
\end{enumerate}
\end{lemma}

\begin{proof}
\textbf{(1) Locality.} Let $\chi \in C^\infty(\mathcal{M})$ be a bump function with $\chi \equiv 1$ in a neighborhood $W \subset U$ of $p$ and $\operatorname{supp}(\chi) \subset U$. If $f = g$ in $U$, then $\chi(f - g) = 0$ on all of $\mathcal{M}$. By the Leibniz rule, $0 = v(\chi(f-g)) = v(\chi)(f-g)(p) + \chi(p) \, v(f-g) = v(f-g)$, since $\chi(p) = 1$ and $(f-g)(p) = 0$.

\textbf{(2) Decomposition.} First, every derivation annihilates constants: if $c \in \mathbb{R}$, then $v(c) = v(c \cdot 1) = c \, v(1)$ and $v(1) = v(1 \cdot 1) = v(1) \cdot 1 + 1 \cdot v(1) = 2v(1)$, whence $v(1) = 0$ and therefore $v(c) = 0$.

Let $f \in C^\infty(\mathcal{M})$. We write $\hat{f} = f \circ \phi^{-1}: \phi(U) \to \mathbb{R}$. By the fundamental theorem of calculus with integral remainder, for $\mathbf{x} = (x^1, \ldots, x^m)$ in a convex neighborhood of $0 \in \phi(U)$:
\[
\hat{f}(\mathbf{x}) = \hat{f}(0) + \sum_{j=1}^{m} x^j \, g_j(\mathbf{x}), \quad \text{where } g_j(\mathbf{x}) = \int_0^1 \frac{\partial \hat{f}}{\partial x^j}(t\mathbf{x}) \, dt
\]
The functions $g_j$ are smooth and $g_j(0) = \frac{\partial \hat{f}}{\partial x^j}(0)$. Composing with $\phi$ and using locality (item 1), we obtain on $U$:
\[
f = f(p) + \sum_{j=1}^{m} x^j \cdot (g_j \circ \phi)
\]
where we identify $x^j$ with the coordinate function $x^j = \pi^j \circ \phi$. Applying $v$ and using that $v$ annihilates constants, $x^j(p) = 0$ and the Leibniz rule:
\[
v(f) = \sum_{j=1}^{m} \bigl[v(x^j) \cdot g_j(\phi(p)) + x^j(p) \cdot v(g_j \circ \phi)\bigr] = \sum_{j=1}^{m} v(x^j) \frac{\partial(f \circ \phi^{-1})}{\partial x^j}\bigg|_{0}
\]
This shows that $v = \sum_j v(x^j) \frac{\partial}{\partial x^j}\big|_p$. Since $\frac{\partial}{\partial x^j}\big|_p(x^k) = \delta_j^k$, the derivations $\{\frac{\partial}{\partial x^j}\big|_p\}_{j=1}^m$ are linearly independent and generate $T_p^{\text{deriv}}\mathcal{M}$ \citep[Proposition 3.2]{lee2012introduction}.
\end{proof}

\begin{proposition}[Equivalence of the definitions of tangent space]\label{prop:equivalencia_espacio_tangente}
The two definitions of tangent space are canonically isomorphic (the isomorphism is independent of the choice of coordinates). The isomorphism $\Phi: T_p^{\text{curvas}}\mathcal{M} \to T_p^{\text{deriv}}\mathcal{M}$ is given by:
\[
\Phi([\gamma])(f) = \frac{d}{dt}\bigg|_{t=0} (f \circ \gamma)(t)
\]
for every function $f \in C^\infty(\mathcal{M})$.
\end{proposition}

\begin{proof}
We proceed in four steps: we verify that $\Phi$ is well defined, that its image consists of derivations, and that it is bijective.

\textbf{Step 1: $\Phi$ is well defined.} Let $\gamma_1, \gamma_2$ be two curves with $\gamma_1(0) = \gamma_2(0) = p$ and $\gamma_1 \sim \gamma_2$ (that is, they are tangent at $p$). We must show that $\Phi([\gamma_1]) = \Phi([\gamma_2])$. Let $(U, \phi)$ be a local chart at $p$ with $\phi = (x^1, \ldots, x^m)$. By the definition of tangency, $(\phi \circ \gamma_1)'(0) = (\phi \circ \gamma_2)'(0)$. For any $f \in C^\infty(\mathcal{M})$, we factor $f \circ \gamma_i = (f \circ \phi^{-1}) \circ (\phi \circ \gamma_i)$ and apply the chain rule:
\[
(f \circ \gamma_i)'(0) = \sum_{j=1}^{m} \frac{\partial (f \circ \phi^{-1})}{\partial x^j}\bigg|_{\phi(p)} \cdot (x^j \circ \gamma_i)'(0)
\]
Since $(x^j \circ \gamma_i)'(0)$ is the $j$-th component of $(\phi \circ \gamma_i)'(0)$ and $(\phi \circ \gamma_1)'(0) = (\phi \circ \gamma_2)'(0)$, we conclude that $(f \circ \gamma_1)'(0) = (f \circ \gamma_2)'(0)$.

\textbf{Step 2: $\Phi([\gamma])$ is a derivation.} Let $v_\gamma = \Phi([\gamma])$. We verify the two properties:
\begin{itemize}
    \item \textit{Linearity}: For $f, g \in C^\infty(\mathcal{M})$ and $\alpha, \beta \in \mathbb{R}$:
    \begin{align*}
    v_\gamma(\alpha f + \beta g) &= \frac{d}{dt}\bigg|_{t=0} (\alpha f + \beta g)(\gamma(t)) \\
    &= \alpha \frac{d}{dt}\bigg|_{t=0} f(\gamma(t)) + \beta \frac{d}{dt}\bigg|_{t=0} g(\gamma(t)) = \alpha v_\gamma(f) + \beta v_\gamma(g)
    \end{align*}
    \item \textit{Leibniz rule}: For $f, g \in C^\infty(\mathcal{M})$, we apply the product rule for real functions of a real variable $t \mapsto f(\gamma(t))$ and $t \mapsto g(\gamma(t))$:
    \begin{align*}
    v_\gamma(fg) &= \frac{d}{dt}\bigg|_{t=0} \bigl[f(\gamma(t)) \cdot g(\gamma(t))\bigr] \\
    &= \left(\frac{d}{dt}\bigg|_{t=0} f(\gamma(t))\right) g(\gamma(0)) + f(\gamma(0)) \left(\frac{d}{dt}\bigg|_{t=0} g(\gamma(t))\right) \\
    &= v_\gamma(f) \, g(p) + f(p) \, v_\gamma(g)
    \end{align*}
\end{itemize}

\textbf{Step 3: $\Phi$ is injective.} Suppose that $\Phi([\gamma_1]) = \Phi([\gamma_2])$, that is, that for every $f \in C^\infty(\mathcal{M})$ we have $(f \circ \gamma_1)'(0) = (f \circ \gamma_2)'(0)$. Let $(U, \phi)$ be a chart with $p \in U$ and coordinate functions $x^j = \pi^j \circ \phi$, where $\pi^j: \mathbb{R}^m \to \mathbb{R}$ is the $j$-th projection. The functions $x^j$ are defined on $U$, but by the locality of $\Phi([\gamma])$ (see \autoref{lem:descomposicion_derivaciones}), we can evaluate them. Taking $f = x^j$, we obtain $(x^j \circ \gamma_1)'(0) = (x^j \circ \gamma_2)'(0)$ for all $j = 1, \ldots, m$. This implies $(\phi \circ \gamma_1)'(0) = (\phi \circ \gamma_2)'(0)$, that is, $\gamma_1 \sim \gamma_2$.

\textbf{Step 4: $\Phi$ is surjective.} Let $v \in T_p^{\text{deriv}}\mathcal{M}$ be an arbitrary derivation. Let $(U, \phi)$ be a chart with $\phi(p) = 0$ and coordinate functions $x^1, \ldots, x^m$. We define the vector $\mathbf{a} = (a^1, \ldots, a^m) \in \mathbb{R}^m$ by $a^j = v(x^j)$ and consider the curve $\gamma(t) = \phi^{-1}(t\mathbf{a}) = \phi^{-1}(ta^1, \ldots, ta^m)$. Then $\gamma(0) = \phi^{-1}(0) = p$ and for any $f \in C^\infty(\mathcal{M})$, the chain rule gives:
\[
\Phi([\gamma])(f) = \frac{d}{dt}\bigg|_{t=0} (f \circ \phi^{-1})(t\mathbf{a}) = \sum_{j=1}^{m} a^j \frac{\partial (f \circ \phi^{-1})}{\partial x^j}\bigg|_{0} = \sum_{j=1}^{m} v(x^j) \frac{\partial}{\partial x^j}\bigg|_p (f)
\]
By \autoref{lem:descomposicion_derivaciones}, $v = \sum_{j=1}^{m} v(x^j) \frac{\partial}{\partial x^j}\big|_p$, whence $\Phi([\gamma]) = v$.
\end{proof}

\textbf{Intuition about the two constructions.} The two constructions offer complementary perspectives:
\begin{itemize}
	\item \textbf{Curves}: They capture the geometric idea of ``directions of motion''. They are intuitive for visualizing tangent vectors as velocities of particles moving over the manifold.
	\item \textbf{Derivations}: They capture the analytical idea of ``directions of change of functions''. They are convenient for computations and generalize naturally to more abstract spaces.
\end{itemize}
In practice, we use both as convenient: curves for geometric intuition, derivations for algebraic manipulations.

\begin{example}[Tangent space to the sphere]\label{ex:espacio_tangente_esfera}
Consider the sphere $S^2 = \{(x,y,z) \in \mathbb{R}^3 : x^2 + y^2 + z^2 = 1\}$. At the north pole $p = (0,0,1)$:
\begin{itemize}
	\item \textbf{Via curves}: A tangent vector is the initial velocity of any curve over $S^2$ passing through $p$. For example, $\gamma(t) = (\sin t, 0, \cos t)$ has $\gamma(0) = p$ and $\gamma'(0) = (1, 0, 0)$.
	\item \textbf{Via derivations}: In stereographic coordinates $(u,v)$ centered at $p$, the derivations $\partial/\partial u$ and $\partial/\partial v$ form a basis.
\end{itemize}
The tangent space $T_p S^2$ is the plane $\{(x,y,0) : x,y \in \mathbb{R}\}$, which is horizontal and tangent to the sphere at the north pole.
\end{example}

\begin{figure}[htbp]
\centering
\includegraphics[width=0.7\textwidth]{figures/tangent_space.pdf}
\caption[Tangent space to a manifold]{Tangent space $T_p\mathcal{M}$ at a point $p$ of a manifold. The tangent vectors represent initial velocities of curves passing through $p$, forming a vector space that linearly approximates the manifold at each point.}
\label{fig:tangent_space}
\end{figure}

\begin{definition}[Differential of a smooth map]\label{def:diferencial_aplicacion_suave}
Let $F: \mathcal{M} \to \mathcal{N}$ be a smooth map between differentiable manifolds. For each $p \in \mathcal{M}$, the \textit{differential} (or \textit{pushforward}) of $F$ at $p$ is the linear map $dF_p: T_p\mathcal{M} \to T_{F(p)}\mathcal{N}$ defined, in the formulation via derivations (\autoref{def:espacio_tangente_derivaciones}), by:
\[
dF_p(v)(g) = v(g \circ F) \quad \text{for all } g \in C^\infty(\mathcal{N}).
\]
If $(U, \phi)$ and $(V, \psi)$ are local charts at $p$ and $F(p)$ respectively, the expression of $dF_p$ in coordinates is the Jacobian matrix $\left(\frac{\partial (\psi \circ F \circ \phi^{-1})^i}{\partial x^j}\right)_{i,j}$.
\end{definition}

\begin{definition}[Smooth immersion]\label{def:inmersion_suave}
A smooth map $F: \mathcal{M} \to \mathcal{N}$ is an \textit{immersion} if its differential $dF_p: T_p\mathcal{M} \to T_{F(p)}\mathcal{N}$ is injective for every $p \in \mathcal{M}$. Equivalently, $\operatorname{rank}(dF_p) = \dim \mathcal{M}$ at each point.
\end{definition}

\begin{definition}[Smooth submersion]\label{def:submersion_suave}
A smooth map $F: \mathcal{M} \to \mathcal{N}$ is a \textit{submersion} if its differential $dF_p: T_p\mathcal{M} \to T_{F(p)}\mathcal{N}$ is surjective for every $p \in \mathcal{M}$. Equivalently, $\operatorname{rank}(dF_p) = \dim \mathcal{N}$ at each point, which requires $\dim \mathcal{M} \geq \dim \mathcal{N}$.
\end{definition}

\begin{remark}\label{rem:inmersion_submersion_dualidad}
Immersion and submersion are dual notions: the former requires $dF_p$ to be injective ($\operatorname{rank} = \dim \mathcal{M}$, needing $\dim \mathcal{M} \leq \dim \mathcal{N}$), while the latter requires surjectivity ($\operatorname{rank} = \dim \mathcal{N}$, needing $\dim \mathcal{M} \geq \dim \mathcal{N}$). The Regular Level Set Theorem (\autoref{thm:nivel_regular} of \autoref{ch:teoremas_complementarios}) guarantees that, if $F$ is a submersion, each fiber $F^{-1}(q)$ is a smooth submanifold of $\mathcal{M}$ of dimension $\dim \mathcal{M} - \dim \mathcal{N}$. This property of fibers as submanifolds will be fundamental in the theory of principal bundles (\autoref{subsec:fibrados_principales}), where the projection $\pi: P \to \mathcal{M}$ is a smooth submersion.
\end{remark}

\begin{definition}[Smooth embedding]\label{def:embebimiento_suave}
An \textit{immersion} $F: \mathcal{M} \to \mathcal{N}$ is an \textit{embedding} if in addition it is injective and is a homeomorphism onto its image $F(\mathcal{M})$ with the subspace topology.
\end{definition}

\begin{remark}\label{rem:compacta_inmersion_embebimiento}
If $\mathcal{M}$ is compact, every injective immersion $F: \mathcal{M} \to \mathcal{N}$ is automatically an embedding. Indeed, $F$ is a continuous, injective, and closed function (every continuous map from a compact space into a Hausdorff space is closed), whence $F$ is a homeomorphism onto its image.
\end{remark}

\subsubsection{Cotangent space and differential forms}\label{subsubsec:espacio_cotangente}

The tangent space $T_p\mathcal{M}$ captures the \textit{directions} at a point of the manifold; its dual, the \textit{cotangent space}, captures how \textit{functions change} along those directions. This step from vectors to covectors is essential for the differential geometry that underlies \gls{GDL}: the Riemannian metric (\autoref{subsec:variedades_riemannianas}), connections on bundles (\autoref{subsec:conexiones_fibrados}), and curvature (\autoref{def:curvatura_formal}) are expressed in terms of differential forms.

\begin{definition}[Cotangent space]\label{def:espacio_cotangente}
Let $\mathcal{M}$ be a differentiable manifold of dimension $m$ and $p \in \mathcal{M}$. The \textit{cotangent space} of $\mathcal{M}$ at $p$ is the algebraic dual of the tangent space:
\[
T_p^*\mathcal{M} = (T_p\mathcal{M})^* = \{\alpha: T_p\mathcal{M} \to \mathbb{R} \mid \alpha \text{ is linear}\}.
\]
The elements of $T_p^*\mathcal{M}$ are called \textit{covectors} (or \textit{pointwise 1-forms}). Since $\dim T_p\mathcal{M} = m$, we have $\dim T_p^*\mathcal{M} = m$.
\end{definition}

\begin{definition}[Differential of a function]\label{def:diferencial_funcion}
Let $f \in C^\infty(\mathcal{M})$ be a smooth function. The \textit{differential} of $f$ at $p$ is the covector $(df)_p \in T_p^*\mathcal{M}$ defined by:
\[
(df)_p(v) = v(f) \quad \text{for all } v \in T_p\mathcal{M},
\]
where $v(f)$ denotes the derivation $v$ applied to $f$ (\autoref{def:espacio_tangente_derivaciones}). In local coordinates $(x^1, \ldots, x^m)$, the differential is expressed as:
\[
df = \sum_{i=1}^m \frac{\partial f}{\partial x^i} \, dx^i.
\]
The covectors $\{dx^1, \ldots, dx^m\}$ form the \textit{dual basis} of $T_p^*\mathcal{M}$, determined by the duality relation with the basis $\{\partial/\partial x^j\}$ of $T_p\mathcal{M}$:
\[
dx^i\!\left(\frac{\partial}{\partial x^j}\right) = \delta^i_j.
\]
\end{definition}

\begin{remark}\label{rem:notacion_base_dual}
The notation $dx^i$ does not denote an ``infinitesimal'', but a well-defined linear functional: the covector that extracts the $i$-th component of a tangent vector with respect to the chart $(x^1, \ldots, x^m)$. This functional is precisely the differential of the coordinate function $x^i: \mathcal{M} \supset U \to \mathbb{R}$.
\end{remark}

\begin{definition}[Tensor product of covectors]\label{def:producto_tensorial_covectores}
Given $\alpha, \beta \in T_p^*\mathcal{M}$, their \textit{tensor product} is the bilinear form $\alpha \otimes \beta: T_p\mathcal{M} \times T_p\mathcal{M} \to \mathbb{R}$ defined by:
\[
(\alpha \otimes \beta)(v, w) = \alpha(v) \cdot \beta(w).
\]
The products $\{dx^i \otimes dx^j\}_{1 \leq i,j \leq m}$ form a basis of the space of $(0,2)$ tensors at $p$. In particular, the Riemannian metric (\autoref{def:metrica_riemanniana}) is expressed as $g_p = \sum_{i,j} g_{ij}(p) \, dx^i \otimes dx^j$.
\end{definition}

\begin{definition}[Differential 1-form]\label{def:1_forma_diferencial}
A \textit{differential 1-form} (or \textit{covector field}) on $\mathcal{M}$ is a smooth assignment $\omega: p \mapsto \omega_p \in T_p^*\mathcal{M}$. In local coordinates it is written $\omega = \sum_{i=1}^m \omega_i(x) \, dx^i$ with $\omega_i \in C^\infty(U)$. The space of all differential 1-forms is denoted $\Omega^1(\mathcal{M})$. Connections on principal bundles (\autoref{def:conexion_fibrado}) are 1-forms with values in a Lie algebra.
\end{definition}

\begin{definition}[Exterior product and $k$-forms]\label{def:producto_exterior_k_formas}
A \textit{differential $k$-form} on $\mathcal{M}$ is a smooth assignment that associates to each $p \in \mathcal{M}$ an antisymmetric tensor of type $(0, k)$. The space of $k$-forms is denoted $\Omega^k(\mathcal{M})$, with $\Omega^0(\mathcal{M}) = C^\infty(\mathcal{M})$ and $\Omega^k(\mathcal{M}) = 0$ for $k > m = \dim \mathcal{M}$; the 1-forms of \autoref{def:1_forma_diferencial} correspond to $k = 1$.

The \textit{exterior product} (or \textit{wedge}) is the operation
\[
\wedge: \Omega^k(\mathcal{M}) \times \Omega^l(\mathcal{M}) \longrightarrow \Omega^{k+l}(\mathcal{M})
\]
defined through the \textit{alternation operator} $\operatorname{Alt}$, which to a tensor $T$ of type $(0,n)$ assigns its totally antisymmetric part:
\[
\operatorname{Alt}(T)(v_1, \ldots, v_n) = \frac{1}{n!} \sum_{\sigma \in S_n} \operatorname{sgn}(\sigma)\, T(v_{\sigma(1)}, \ldots, v_{\sigma(n)}),
\]
where $S_n$ is the permutation group and $\operatorname{sgn}(\sigma) \in \{+1, -1\}$ the sign of~$\sigma$. The exterior product is defined as the normalized antisymmetrization of the tensor product (\autoref{def:producto_tensorial_covectores}):
\begin{equation}\label{eq:wedge_alt}
\alpha \wedge \beta = \frac{(k+l)!}{k!\,l!}\, \operatorname{Alt}(\alpha \otimes \beta),
\end{equation}
where $(\alpha \otimes \beta)(v_1, \ldots, v_{k+l}) = \alpha(v_1, \ldots, v_k)\,\beta(v_{k+1}, \ldots, v_{k+l})$. For 1-forms $\alpha, \beta \in \Omega^1(\mathcal{M})$, equation~\eqref{eq:wedge_alt} reduces to:
\[
(\alpha \wedge \beta)(v, w) = \alpha(v)\beta(w) - \alpha(w)\beta(v),
\]
which is equivalent to $\alpha \wedge \beta = \alpha \otimes \beta - \beta \otimes \alpha$. In particular, on the dual basis:
\[
dx^i \wedge dx^j = -dx^j \wedge dx^i, \qquad dx^i \wedge dx^i = 0.
\]
In general, the exterior product satisfies \textit{graded anticommutativity}: for $\alpha \in \Omega^k(\mathcal{M})$ and $\beta \in \Omega^l(\mathcal{M})$,
\[
\alpha \wedge \beta = (-1)^{kl}\, \beta \wedge \alpha.
\]
In local coordinates, every $k$-form is written:
\[
\omega = \sum_{i_1 < \cdots < i_k} \omega_{i_1 \cdots i_k}(x) \, dx^{i_1} \wedge \cdots \wedge dx^{i_k}.
\]
The volume element (\autoref{def:elemento_volumen}) is an $m$-form.
\end{definition}

\begin{proposition}[Algebraic properties of the exterior product]\label{prop:propiedades_producto_exterior}
The exterior product $\wedge$ satisfies the following properties for $\alpha \in \Omega^k(\mathcal{M})$, $\beta \in \Omega^l(\mathcal{M})$, $\gamma \in \Omega^r(\mathcal{M})$ and $f \in C^\infty(\mathcal{M})$:
\begin{enumerate}
    \item \textbf{Bilinearity}: $(\alpha + \alpha') \wedge \beta = \alpha \wedge \beta + \alpha' \wedge \beta$ and $\alpha \wedge (\beta + \beta') = \alpha \wedge \beta + \alpha \wedge \beta'$.
    \item \textbf{Associativity}: $(\alpha \wedge \beta) \wedge \gamma = \alpha \wedge (\beta \wedge \gamma)$.
    \item \textbf{Compatibility with scalars}: $(f\alpha) \wedge \beta = f(\alpha \wedge \beta) = \alpha \wedge (f\beta)$.
    \item \textbf{Graded anticommutativity}: $\alpha \wedge \beta = (-1)^{kl}\, \beta \wedge \alpha$.
\end{enumerate}
\end{proposition}

\begin{proof}
All the properties derive from definition~\eqref{eq:wedge_alt}: $\alpha \wedge \beta = \frac{(k+l)!}{k!\,l!}\, \operatorname{Alt}(\alpha \otimes \beta)$.

\textbf{(1) Bilinearity.} The tensor product $\otimes$ is bilinear and $\operatorname{Alt}$ is linear (being a linear combination of evaluations with fixed coefficients $\operatorname{sgn}(\sigma)/n!$). The composition~\eqref{eq:wedge_alt} preserves bilinearity.

\textbf{(3) Compatibility with scalars.} For $f \in C^\infty(\mathcal{M})$, we have $(f\alpha) \otimes \beta = f(\alpha \otimes \beta)$ pointwise, since $f(p)$ is a scalar that does not depend on the tangent vectors permuted by $\operatorname{Alt}$. By the linearity of $\operatorname{Alt}$:
\[
(f\alpha) \wedge \beta = \tfrac{(k+l)!}{k!\,l!}\,\operatorname{Alt}((f\alpha) \otimes \beta) = \tfrac{(k+l)!}{k!\,l!}\,f\cdot\operatorname{Alt}(\alpha \otimes \beta) = f(\alpha \wedge \beta).
\]
Analogously, $\alpha \wedge (f\beta) = f(\alpha \wedge \beta)$.

\textbf{(4) Graded anticommutativity.} Let $\tau \in S_{k+l}$ be the permutation $(1, \ldots, k+l) \mapsto (k{+}1, \ldots, k{+}l, 1, \ldots, k)$ that swaps the first $k$ indices with the last $l$. This permutation decomposes into $kl$ transpositions (each of the first $k$ elements crosses the last $l$), so $\operatorname{sgn}(\tau) = (-1)^{kl}$. The relation between the tensor products is:
\[
(\beta \otimes \alpha)(v_1, \ldots, v_{k+l}) = \beta(v_1, \ldots, v_l)\,\alpha(v_{l+1}, \ldots, v_{l+k}) = (\alpha \otimes \beta)(v_{\tau(1)}, \ldots, v_{\tau(k+l)}).
\]
Substituting into $\operatorname{Alt}$ and making the change $\rho = \tau \circ \sigma$ (which ranges over $S_{k+l}$ as $\sigma$ varies):
\begin{align*}
\operatorname{Alt}(\beta \otimes \alpha)(v_1, \ldots, v_{k+l})
  &= \frac{1}{(k{+}l)!}\sum_{\sigma \in S_{k+l}} \operatorname{sgn}(\sigma)\,(\alpha \otimes \beta)(v_{\tau\sigma(1)}, \ldots, v_{\tau\sigma(k+l)}) \\
  &= \frac{(-1)^{kl}}{(k{+}l)!}\sum_{\rho \in S_{k+l}} \operatorname{sgn}(\rho)\,(\alpha \otimes \beta)(v_{\rho(1)}, \ldots, v_{\rho(k+l)})
  = (-1)^{kl}\,\operatorname{Alt}(\alpha \otimes \beta),
\end{align*}
where we have used $\operatorname{sgn}(\tau\sigma) = (-1)^{kl}\operatorname{sgn}(\sigma)$. Multiplying by $\frac{(k+l)!}{k!\,l!}$ we obtain $\beta \wedge \alpha = (-1)^{kl}\,\alpha \wedge \beta$.

\textbf{(2) Associativity.} Let $n = k+l+r$. By~\eqref{eq:wedge_alt}:
\[
(\alpha \wedge \beta) \wedge \gamma = \frac{n!}{(k{+}l)!\,r!}\,\operatorname{Alt}_n\!\big((\alpha \wedge \beta) \otimes \gamma\big) = \frac{n!}{k!\,l!\,r!}\,\operatorname{Alt}_n\!\big(\operatorname{Alt}_{k+l}(\alpha \otimes \beta) \otimes \gamma\big).
\]
We identify $S_{k+l}$ with the subgroup of $S_n$ that fixes the last $r$ indices. For $\tau$ in this subgroup, $\tau$ acts only on the first $k+l$ arguments of $\alpha \otimes \beta$ and leaves $\gamma$ invariant. Evaluating explicitly:
\begin{align*}
&\operatorname{Alt}_n\!\big(\operatorname{Alt}_{k+l}(\alpha \otimes \beta) \otimes \gamma\big)(v_1, \ldots, v_n) \\
  &\quad= \frac{1}{n!\,(k{+}l)!}\sum_{\sigma \in S_n}\sum_{\tau \in S_{k+l}}\operatorname{sgn}(\sigma)\,\operatorname{sgn}(\tau)\,(\alpha \otimes \beta \otimes \gamma)(v_{\sigma(\tau(1))}, \ldots, v_{\sigma(\tau(n))}).
\end{align*}
Substituting $\rho = \sigma \circ \tau$: for each $\rho \in S_n$, the equation $\sigma \circ \tau = \rho$ with $\tau \in S_{k+l}$ admits exactly $(k{+}l)!$ solutions (since $\sigma = \rho \circ \tau^{-1}$ is determined). Since $\operatorname{sgn}(\sigma)\operatorname{sgn}(\tau) = \operatorname{sgn}(\rho)$:
\[
= \frac{(k{+}l)!}{n!\,(k{+}l)!}\sum_{\rho \in S_n}\operatorname{sgn}(\rho)\,(\alpha \otimes \beta \otimes \gamma)(v_{\rho(1)}, \ldots, v_{\rho(n)}) = \operatorname{Alt}_n(\alpha \otimes \beta \otimes \gamma).
\]
Therefore, $(\alpha \wedge \beta) \wedge \gamma = \frac{n!}{k!\,l!\,r!}\,\operatorname{Alt}_n(\alpha \otimes \beta \otimes \gamma)$. The computation for $\alpha \wedge (\beta \wedge \gamma)$ is identical, using $S_{l+r} \hookrightarrow S_n$ in the last positions, and produces the same expression.
\end{proof}

\begin{definition}[Exterior derivative]\label{def:derivada_exterior}
The \textit{exterior derivative} is the unique linear operator $d: \Omega^k(\mathcal{M}) \to \Omega^{k+1}(\mathcal{M})$ that satisfies:
\begin{enumerate}
    \item \textbf{Extension of the differential}: for $f \in \Omega^0(\mathcal{M}) = C^\infty(\mathcal{M})$, it coincides with the differential of \autoref{def:diferencial_funcion}.
    \item \textbf{Graded Leibniz rule}: $d(\alpha \wedge \beta) = d\alpha \wedge \beta + (-1)^k \alpha \wedge d\beta$ for $\alpha \in \Omega^k$.
    \item \textbf{Nilpotency}: $d \circ d = 0$.
\end{enumerate}
Nilpotency implies that $\operatorname{Im}(d: \Omega^{k-1} \to \Omega^k) \subseteq \ker(d: \Omega^k \to \Omega^{k+1})$, giving rise to the de Rham complex $0 \to \Omega^0 \xrightarrow{d} \Omega^1 \xrightarrow{d} \cdots \xrightarrow{d} \Omega^m \to 0$. In the curvature equation (\autoref{def:curvatura_formal}), the exterior derivative $d\omega$ and the Bianchi identity $d\Omega_\omega = [\omega \wedge \Omega_\omega]$ are direct applications of this operator.
\end{definition}

\begin{example}[Differential forms on $S^2$]\label{ex:formas_esfera}
On the sphere $S^2$ with spherical coordinates $(\theta, \phi)$:
\begin{itemize}
    \item The function $f = \cos\theta$ is a 0-form ($f \in \Omega^0(S^2)$).
    \item Its differential $df = -\sin\theta \, d\theta$ is a 1-form.
    \item The area element $\sin\theta \, d\theta \wedge d\phi$ is the unique 2-form (up to smooth scalar multiples) on $S^2$, since $\dim S^2 = 2$.
\end{itemize}
The area 2-form appears in spherical convolution, where signals over $S^2$ are integrated with respect to this volume form~\citep{cohen2018sphericalcnns}.
\end{example}

\begin{remark}[Relevance for GDL]\label{rem:formas_gdl}
The constructions of this section reappear in the following chapters under four concrete forms:
\begin{enumerate}
    \item The \textbf{Riemannian metric} is a smooth field of tensor products of covectors (\autoref{def:producto_tensorial_covectores}).
    \item The \textbf{connections} on principal bundles are 1-forms with values in the Lie algebra of the structure group (\autoref{def:1_forma_diferencial}).
    \item The \textbf{volume form} $\sqrt{\det(g_{ij})} \, dx^1 \wedge \cdots \wedge dx^m$ is an $m$-form that allows integrating signals over manifolds (\autoref{def:producto_exterior_k_formas}).
    \item The \textbf{de Rham complex} $(\Omega^\bullet, d)$ is the continuous version of simplicial chains (\autoref{subsec:tdl_grafos}), connecting differential geometry with algebraic topology.
\end{enumerate}
\end{remark}

\subsection{The Manifold Hypothesis}\label{subsec:hipotesis_variedad}

A fundamental perspective that connects Euclidean representations with richer geometries is the \textbf{manifold hypothesis}. This hypothesis, central to modern machine learning, postulates that although data is represented as high-dimensional vectors in $\mathbb{R}^d$, it actually concentrates near a manifold of much lower intrinsic dimension.

The foundations of topological and differentiable manifolds have been developed in the previous subsection (see Definitions~\ref{def:variedad_topologica} and~\ref{def:variedad_diferenciable}). In this subsection we focus on the implications of the manifold hypothesis for \gls{GDL} and present fundamental theoretical results on the structure and complexity of manifolds in high-dimensional spaces.

\paragraph{Preliminary Riemannian notions}

The quantitative results of this section and the classification of geometries in \autoref{subsec:espacios_no_euclidianos} employ notions from Riemannian geometry that we will formalize completely in \autoref{subsec:variedades_riemannianas}. We present here the necessary intuitive definitions.

\begin{definition}[Inner product]\label{def:producto_interno}
Let $V$ be a real vector space. An \textit{inner product} on $V$ is a bilinear form $\langle \cdot, \cdot \rangle: V \times V \to \mathbb{R}$ that satisfies:
\begin{enumerate}
	\item \textbf{Symmetry}: $\langle u, v \rangle = \langle v, u \rangle$ for all $u, v \in V$
	\item \textbf{Linearity}: $\langle \alpha u + \beta v, w \rangle = \alpha \langle u, w \rangle + \beta \langle v, w \rangle$ for all $u, v, w \in V$ and $\alpha, \beta \in \mathbb{R}$
	\item \textbf{Positive definiteness}: $\langle v, v \rangle > 0$ for all $v \neq 0$
\end{enumerate}
A vector space equipped with an inner product is called an \textit{inner product space}. If it is also complete with respect to the induced norm $\|v\| = \sqrt{\langle v, v \rangle}$, it is called a \textit{Hilbert space}.
\end{definition}

\begin{definition}[Proper subspace]\label{def:subespacio_propio}
Let $V$ be a vector space. A subspace $W \subseteq V$ is a \textit{proper subspace} if $W \neq \{0\}$ and $W \neq V$, that is, $0 < \dim(W) < \dim(V)$.
\end{definition}

\begin{definition}[Riemannian manifold (preliminary)]\label{def:variedad_riemanniana_preliminar}
A \textit{Riemannian manifold} is a pair $(\mathcal{M}, g)$ where $\mathcal{M}$ is a differentiable manifold and $g$ assigns to each point $p \in \mathcal{M}$ an inner product $g_p: T_p\mathcal{M} \times T_p\mathcal{M} \to \mathbb{R}$ on the tangent space, which varies smoothly with $p$. The tensor $g$ is called the \textit{Riemannian metric} and allows measuring lengths, angles, and volumes over the manifold. The complete formalization is developed in \autoref{def:variedad_riemanniana}.
\end{definition}

\begin{definition}[Geodesic (preliminary)]\label{def:geodesica_preliminar}
Let $(\mathcal{M}, g)$ be a Riemannian manifold. A \textit{geodesic} is a curve $\gamma: I \to \mathcal{M}$ that locally minimizes the length between nearby points. That is, for each $t_0 \in I$ there exists $\varepsilon > 0$ such that $\gamma|_{[t_0-\varepsilon, t_0+\varepsilon]}$ is the shortest path between $\gamma(t_0 - \varepsilon)$ and $\gamma(t_0 + \varepsilon)$.
\end{definition}

Geodesics generalize straight lines to curved spaces: in $\mathbb{R}^n$ they are the straight lines, and on the sphere $S^2$ they are the great circles. The differential characterization through the equation $\nabla_{\dot{\gamma}} \dot{\gamma} = 0$ is developed in \autoref{def:geodesica}.

\begin{definition}[Sectional curvature (preliminary)]\label{def:curvatura_preliminar}
Let $(\mathcal{M}, g)$ be a Riemannian manifold. The \textit{sectional curvature} $K$ at a point $p \in \mathcal{M}$ and a 2-plane $\sigma \subset T_p\mathcal{M}$ measures the rate at which initially parallel geodesics in directions within $\sigma$ separate or approach:
\begin{itemize}
    \item $K = 0$: The geodesics remain equidistant (Euclidean geometry).
    \item $K > 0$: The geodesics converge (spherical/elliptic geometry).
    \item $K < 0$: The geodesics diverge (hyperbolic geometry).
\end{itemize}
The formalization through the Riemann curvature tensor is developed in \autoref{def:tensor_riemann}.
\end{definition}

\begin{figure}[htbp]
\centering
\includegraphics[width=0.9\textwidth]{figures/manifold_hypothesis.png}
\caption[Manifold hypothesis]{Illustration of the manifold hypothesis. Left: the torus $\mathbb{T}^2$ as a two-dimensional manifold embedded in $\mathbb{R}^3$. Right: hyperbolic paraboloid (saddle) with data distributed over the surface, illustrating how high-dimensional data concentrates on low-dimensional manifolds.}
\label{fig:manifold_hypothesis}
\end{figure}

\begin{hypothesis}[Manifold hypothesis]\label{hyp:hipotesis_variedad_gdl}
High-dimensional data in real applications typically lies near a manifold $\mathcal{M}$ of intrinsic dimension $m \ll d$ immersed in $\mathbb{R}^d$. Formally, there exists a differentiable manifold $\mathcal{M}$ such that the probability density of the data $p(\mathbf{x})$ concentrates in a neighborhood of $\mathcal{M}$:
\[
p(\mathbf{x}) \approx 0 \quad \text{for } \mathbf{x} \notin \mathcal{U}_\epsilon(\mathcal{M})
\]
where $\mathcal{U}_\epsilon(\mathcal{M})$ is an $\epsilon$-tubular neighborhood of $\mathcal{M}$.
\end{hypothesis}

This hypothesis is not merely speculative: it has a solid theoretical foundation and abundant empirical evidence. The following fundamental results characterize the structure of manifolds in high-dimensional spaces.

\begin{theorem}[Whitney Embedding Theorem]\label{thm:whitney_embedding}
Let $\mathcal{M}$ be a compact differentiable manifold of dimension $m$. Then $\mathcal{M}$ can be smoothly embedded in $\mathbb{R}^{2m+1}$ and, if $m > 1$, immersed in $\mathbb{R}^{2m-1}$.
\end{theorem}

\begin{proof}[Sketch]
The proof proceeds in two phases: first an embedding into a Euclidean space of arbitrarily large dimension is constructed, and then the dimension is iteratively reduced until reaching $\mathbb{R}^{2m+1}$.

\textbf{Step 1: Immersion in high dimension.}
By compactness, $\mathcal{M}$ admits a finite atlas $\{(U_i, \phi_i)\}_{i=1}^N$. Let $\{\rho_i\}_{i=1}^N$ be a partition of unity subordinate to this cover, that is, smooth functions $\rho_i: \mathcal{M} \to [0,1]$ with $\operatorname{supp}(\rho_i) \subset U_i$ and $\sum_i \rho_i \equiv 1$. We define $F: \mathcal{M} \to \mathbb{R}^{N(m+1)}$ by:
\[
F(p) = \bigl(\rho_1(p)\phi_1(p), \ldots, \rho_N(p)\phi_N(p), \rho_1(p), \ldots, \rho_N(p)\bigr)
\]
where $\rho_i(p)\phi_i(p) = 0$ when $p \notin U_i$.

\textbf{Step 2: $F$ is an embedding.}
The map $F$ is smooth, being a composition of smooth functions. It is injective: if $F(p) = F(q)$, the last $N$ components give $\rho_i(p) = \rho_i(q)$ for all $i$; choosing $i$ with $\rho_i(p) > 0$, the first components give $\rho_i(p)\phi_i(p) = \rho_i(q)\phi_i(q)$, whence $\phi_i(p) = \phi_i(q)$ and therefore $p = q$. To see that $dF_p$ (\autoref{def:diferencial_aplicacion_suave}) is injective, we fix $i$ with $\rho_i(p) > 0$: the corresponding block of the differential contains $\rho_i(p) \cdot d(\phi_i)_p$, which has rank $m$ since $\phi_i$ is a local diffeomorphism. Since $\mathcal{M}$ is compact, by \autoref{rem:compacta_inmersion_embebimiento} this injective immersion is an embedding.

\textbf{Step 3: Dimension reduction.}
One iteratively projects from $\mathbb{R}^k$ to $\mathbb{R}^{k-1}$ by choosing a direction $\mathbf{v} \in S^{k-1}$ and projecting orthogonally onto $\mathbf{v}^\perp \cong \mathbb{R}^{k-1}$. The projection $\pi_\mathbf{v}$ restricted to $F(\mathcal{M})$ fails to be an embedding only if $\mathbf{v}$ is tangent to $F(\mathcal{M})$ (loss of injectivity of the differential) or if $\mathbf{v}$ is a secant direction, that is, $\mathbf{v} \propto F(p) - F(q)$ for some pair $p \neq q$ (loss of global injectivity). The set of tangent directions has dimension $\leq 2m - 1$ in $S^{k-1}$ (image of the unit tangent bundle, of dimension $2m-1$) and that of secant directions has dimension $\leq 2m$ (image of $\mathcal{M} \times \mathcal{M} \setminus \Delta$, of dimension $2m$). By Sard's Theorem (\autoref{thm:sard}, see \autoref{sec:sard} of Appendix~\ref{ch:teoremas_complementarios}), the set of critical values of the maps that parameterize these directions has measure zero in $S^{k-1}$. Since $\dim S^{k-1} = k - 1 > 2m$ when $k > 2m + 1$, there exist ``good'' directions and the projection preserves the embedding. Iterating, one obtains an embedding in $\mathbb{R}^{2m+1}$.

\textbf{Strong version.}
The improvement from $\mathbb{R}^{2m+1}$ to $\mathbb{R}^{2m}$ requires additional techniques (the ``Whitney trick'') to eliminate residual self-intersections; see \citep[Chapter 6]{lee2012introduction}.
\end{proof}

Whitney's theorem guarantees that any compact $m$-dimensional manifold can be faithfully represented in $\mathbb{R}^{2m+1}$. This theoretically justifies why high-dimensional representations work: if the data really resides on a manifold of dimension $m$, then embeddings in $\mathbb{R}^{2m+1}$ or higher are sufficient to capture all of its geometric structure. However, the theorem does not provide an efficient constructive algorithm, which motivates the use of dimensionality reduction techniques such as variational autoencoders and \gls{GDL}.

\begin{definition}[Reach of a manifold]\label{def:reach}
Let $\mathcal{M} \subset \mathbb{R}^d$ be an embedded closed manifold. The \textit{reach} of $\mathcal{M}$, denoted $\tau(\mathcal{M})$, is the maximum value of $r > 0$ such that any point $\mathbf{x} \in \mathbb{R}^d$ with $d(\mathbf{x}, \mathcal{M}) < r$ has a unique projection onto $\mathcal{M}$.
\end{definition}

The reach characterizes how ``regular'' the geometry of the manifold is: large reach indicates smooth curvature and absence of self-proximity, while small reach indicates complex geometry with regions of high curvature.

\begin{proposition}[Sample complexity on manifolds]\label{prop:complejidad_muestral_variedades}
Let $\mathcal{M}$ be a compact Riemannian manifold (\autoref{def:espacio_compacto}) of dimension $m$, with reach $\tau > 0$ and sectional curvature bounded in absolute value by $\kappa$. Then, to uniformly approximate an $L$-Lipschitz function $f: \mathcal{M} \to \mathbb{R}$ with error $\epsilon$ from samples, the following number of samples is required:
\[
N = \Omega\left(\left(\frac{1}{\epsilon}\right)^m \cdot \operatorname{poly}(\tau^{-1}, \kappa)\right)
\]
\citep{niyogi2008finding}.
\end{proposition}

For a sketch of the proof based on covering numbers and volume comparison, see \autoref{sec:complejidad_muestral} of Appendix~\ref{ch:teoremas_complementarios}.

This proposition reveals that the sample complexity depends on the \textbf{intrinsic dimension} $m$ of the manifold, not on the ambient dimension $d$. This is the ``blessing of dimensionality'': if the high-dimensional data really resides on a low-dimensional manifold, then the number of samples needed for learning is exponentially smaller than what the ambient dimension would suggest. This explains why neural networks work effectively in high-dimensional spaces such as images ($d \sim 10^6$) or text sequences ($d \sim 10^4$): the effective intrinsic dimension $m$ is much smaller.

\paragraph{Examples of manifolds in real data}

The manifold hypothesis is not abstract: it manifests concretely in real application domains:

\begin{itemize}
    \item \textbf{Natural images}: Although a $256 \times 256$ RGB image is formally a vector in $\mathbb{R}^{196608}$, natural images do not uniformly fill this space. They concentrate on a manifold of estimated dimension between 10 and 100, determined by semantic structures (objects, textures, illumination) and natural statistics.

    \item \textbf{Configurations of articulated robots}: A robotic arm with $n$ joints has configurations that reside in the product of circles $(S^1)^n$, a compact manifold of dimension $n$. Although we can represent each configuration as a vector of angles in $\mathbb{R}^n$, the intrinsic topology is non-Euclidean: the angle $2\pi$ is identical to the angle $0$.

    \item \textbf{Brain diffusion MRI data}: The neural fibers in diffusion magnetic resonance imaging live in the tangent bundle of the brain manifold. At each voxel, the diffusion direction is a vector tangent to the cortical surface, not an arbitrary vector in $\mathbb{R}^3$.
\end{itemize}

\subsection{Non-Euclidean Spaces}\label{subsec:espacios_no_euclidianos}

The manifold hypothesis establishes that data resides on low-dimensional manifolds, but these manifolds are, in general, non-Euclidean spaces whose intrinsic geometry differs from the Euclidean one. Non-Euclidean spaces therefore provide the concrete geometric framework for modeling data manifolds: each type of curvature and topology gives rise to distinct geometric operations that, as we will see, determine the appropriate neural architecture. Notably, all the geometries we will study in this section (spheres, projective spaces, Grassmann manifolds, hyperbolic spaces) share a common algebraic structure that allows building equivariant architectures in a unified way; this structure is formalized in \autoref{subsec:espacios_homogeneos}.

\paragraph{Origin of non-Euclidean geometries}

Non-Euclidean geometries arise from modifying Euclid's fifth postulate. In the language of modern differential geometry (\autoref{subsec:variedades_riemannianas}), this modification corresponds to considering Riemannian manifolds with nonzero curvature:

\begin{itemize}
    \item \textbf{Hyperbolic geometry} \citep{gray2007worlds}: Through a point external to a line pass \textit{infinitely many} parallel lines. Developed independently by Lobachevsky and Bolyai around 1830, it corresponds to manifolds with constant sectional curvature $K < 0$. The canonical model is the hyperbolic space $\mathbb{H}^n$ with $K = -1$.
    \item \textbf{Elliptic/spherical geometry} \citep{riemann1854hypothesen}: Through a point external to a line passes \textit{no} parallel line (all geodesics intersect). Formalized by Riemann in 1854, it corresponds to manifolds with constant sectional curvature $K > 0$. The canonical model is the sphere $S^n$ with $K = +1$.
\end{itemize}

In the context of \gls{GDL}, these geometric properties (see \autoref{rem:significado_curvatura}) have direct implications for the design of neural architectures: the curvature of the domain determines how information propagates between neighboring points, since in spaces of positive curvature information converges (facilitating global aggregation), while in negative curvature the representational capacity grows exponentially with distance (facilitating hierarchical representations). Choosing the appropriate geometric space amounts, in this sense, to incorporating a geometric inductive bias into the architecture. We next examine each concrete space, presenting its formal definition, its geometric properties, and its characteristic applications.


\begin{figure}[htbp]
\centering
\includegraphics[width=0.85\textwidth]{figures/curvature_comparison.pdf}
\caption[Comparison of geometries according to curvature]{Comparison of the three classical geometries according to the sign of the sectional curvature $K$. In positive curvature ($K > 0$), geodesics converge, facilitating the global aggregation of information. In zero curvature ($K = 0$), geodesics maintain constant distance (Euclidean geometry). In negative curvature ($K < 0$), geodesics diverge exponentially, providing greater representational capacity for hierarchical data.}
\label{fig:curvature_comparison}
\end{figure}

\paragraph{Spheres and projective spaces: directional data}

\begin{definition}[$n$-dimensional sphere]\label{def:esfera_n}
The \textit{$n$-dimensional sphere} is the manifold:
\[
S^n = \{x \in \mathbb{R}^{n+1} : \|x\| = 1\}
\]
The Riemannian metric induced by the ambient Euclidean metric has constant sectional curvature $+1$ (see \autoref{thm:curvatura_esfera} of Appendix~\ref{ch:teoremas_complementarios}).
\end{definition}

The spheres $S^{n-1} \subset \mathbb{R}^n$ are fundamental for directional data where only the orientation matters, not the magnitude, with intrinsic metric given by the angular distance:
\[
d_{S^{n-1}}(\mathbf{u}, \mathbf{v}) = \arccos(\mathbf{u}^T \mathbf{v})
\]

\textbf{Key applications}:
\begin{itemize}
    \item \textbf{Omnidirectional vision}: $360^\circ$ panoramic images naturally live on $S^2$. \glspl{SphericalCNN} \citep{cohen2018sphericalcnns} operate directly on the sphere using convolutions equivariant under $SO(3)$, avoiding the distortions of planar projections (see \autoref{subsec:spherical_cnns_so3}).

    \item \textbf{Normalized embeddings}: In metric learning, normalizing embeddings to the unit sphere $S^{d-1}$ improves generalization capacity by focusing on directionality rather than magnitude \citep{wang2017normface} (see \autoref{subsec:atencion_esfera}).

    \item \textbf{Molecular orientations}: Directions of chemical bonds and dipole moments are naturally elements of $S^2$, with rotational symmetries under $SO(3)$.

    \item \textbf{Architectural implication}: The constant positive curvature of $S^n$ implies that convolutions must replace Euclidean translations with rotations of the group $SO(n+1)$. \glspl{SphericalCNN} exploit this by defining convolution as cross-correlation over $SO(3)$, guaranteeing exact rather than approximate rotational equivariance (see \autoref{subsec:spherical_cnns_so3}).
\end{itemize}

In many applications, the data represents \textit{directions} but does not distinguish between opposite senses: a line in an image, an axis of rotation, or a principal direction in \gls{PCA}. This data does not live on the sphere $S^{n-1}$ (where $\mathbf{v}$ and $-\mathbf{v}$ are distinct points), but in the \textbf{real projective space}.

\begin{definition}[Real Projective Space]\label{def:espacio_proyectivo}
The \textit{real projective space} $\mathbb{RP}^n$ is the space of lines passing through the origin in $\mathbb{R}^{n+1}$, that is, the quotient of $S^n$ by the antipodal identification $x \sim -x$:
\[
\mathbb{RP}^n = S^n / \{\pm 1\}
\]
As a manifold, $\dim \mathbb{RP}^n = n$, since the antipodal projection $\pi: S^n \to \mathbb{RP}^n$ is a double cover (local diffeomorphism). The space $\mathbb{RP}^n$ is compact and connected (continuous image of $S^n$), but not simply connected: its fundamental group is $\pi_1(\mathbb{RP}^n) \cong \mathbb{Z}/2\mathbb{Z}$, with $S^n$ ($n \geq 2$) being its universal cover.

Concrete cases: $\mathbb{RP}^1 \cong S^1$ (the circle, since the identification of diametrically opposite points on $S^1$ again produces a circle); $\mathbb{RP}^2$ is the first non-orientable closed manifold (the Möbius band is an open subset of $\mathbb{RP}^2$).
\end{definition}

The canonical metric of $\mathbb{RP}^n$ is inherited from the round metric of $S^n$ through the antipodal projection, which is a local isometry (Riemannian cover):
\[
d_{\mathbb{RP}^n}([x], [y]) = \arccos(|x^T y|)
\]
The absolute value in the argument of the arccosine distinguishes this metric from the spherical one $d_{S^n}(x,y) = \arccos(x^T y)$, reflecting that $[x] = [-x]$ in $\mathbb{RP}^n$. Since $\pi: S^n \to \mathbb{RP}^n$ is a Riemannian cover (local isometry), the sectional curvature of $\mathbb{RP}^n$ is constant $+1$, equal to that of $S^n$; the geometric differences between the two spaces are purely global (orientability, fundamental group).

\textbf{Key applications} of projective spaces:
\begin{itemize}
    \item \textbf{Computer vision}: Lines and edges in images are elements of $\mathbb{RP}^2$, since a line has no ``preferred direction''. Networks that process this data must be invariant under the antipodal identification $\mathbf{v} \sim -\mathbf{v}$; projective geometric algebra provides a natural framework for encoding this symmetry \citep{dehaan2024euclideanprojectiveconformalchoosing}.

    \item \textbf{Principal component analysis}: The principal axes of \gls{PCA} are unsigned directions, elements of $\mathbb{RP}^{n-1}$. Optimization methods over $\mathbb{RP}^{n-1}$ require Riemannian rather than Euclidean gradients, since the antipodal identification turns the space into a nontrivial manifold.

    \item \textbf{Architectural implication}: The quotient structure $\mathbb{RP}^{n-1} = S^{n-1}/\mathbb{Z}_2$ implies that the neural layers must be equivariant under the antipodal action. A projective layer can be constructed by taking a spherical layer and symmetrizing: $f([x]) = g(x) + g(-x)$. As a homogeneous space, $\mathbb{RP}^{n-1} \cong O(n)/(O(n-1) \times \mathbb{Z}_2)$ (\autoref{prop:proyectivo_cociente}), and this representation guides the construction of equivariant networks over the projective space.
\end{itemize}

\begin{definition}[Grassmann manifold]\label{def:variedad_grassmann}
The \textit{Grassmann manifold} $Gr(k,n)$ is the space of all $k$-dimensional subspaces of $\mathbb{R}^n$. It is a compact manifold of dimension $\dim Gr(k,n) = k(n-k)$. Each point of $Gr(k,n)$ can be uniquely represented by the orthogonal projection matrix $P \in \mathbb{R}^{n \times n}$ onto the corresponding subspace, characterized by:
\[
P^2 = P, \qquad P^T = P, \qquad \operatorname{tr}(P) = k
\]
The real projective space is the special case $\mathbb{RP}^{n-1} = Gr(1,n)$ (lines through the origin). As concrete cases: $Gr(1,3) = \mathbb{RP}^2$ parameterizes the lines in $\mathbb{R}^3$, and $Gr(2,4)$ (planes in $\mathbb{R}^4$) has dimension $2 \cdot 2 = 4$.
\end{definition}

The natural metric of $Gr(k,n)$ is based on the \textit{principal angles} between subspaces. Given $U, V \in Gr(k,n)$, the principal angles $\theta_1, \ldots, \theta_k \in [0, \pi/2]$ are the angles between the most aligned pairs of vectors of $U$ and $V$ (defined recursively as $\cos \theta_i = \max_{u \in U, v \in V} u^T v$ subject to orthonormality with the previous pairs). The geodesic distance is:
\[
d_{Gr}(U, V) = \sqrt{\sum_{i=1}^{k} \theta_i^2}
\]
When $k = 1$, this metric reduces to the Fubini-Study metric over $\mathbb{RP}^{n-1}$. The Grassmann manifold also possesses a natural symmetry $Gr(k,n) \cong Gr(n-k,n)$ (duality by orthogonal complement: to each subspace $V$ corresponds $V^\perp$), which in dimension $k(n-k) = (n-k)k$ is reflected in the equality of dimensions.

Grassmann manifolds are relevant for data that represents linear subspaces rather than points. Since subspaces do not admit standard addition or scaling, the neural network operations must respect the curved geometry of $Gr(k,n)$.

\textbf{Key applications}:
\begin{itemize}
    \item \textbf{Action recognition in video}: Sequences of frames are modeled as subspaces of $\mathbb{R}^n$, and classification operates over $Gr(k,n)$ with geodesic metrics between subspaces. The GrassmannNet of \citet{huang2017grassmann} define layers that operate through projections and re-orthogonalizations that preserve the subspace structure.

    \item \textbf{Local principal component analysis}: In many vision tasks, the local descriptors of an image generate subspaces whose comparison requires the Grassmann metric, not the Euclidean one. The geodesic distance based on principal angles provides a similarity measure intrinsic to the space of subspaces.

    \item \textbf{Architectural implication}: The standard operations (addition, scaling) do not preserve the subspace structure in $Gr(k,n)$, so neural networks must redefine their basic operations: propagation uses exponential and logarithmic maps between the tangent space and the manifold, and linear layers are replaced by transformations that preserve membership in $Gr(k,n)$ \citep{huang2017grassmann}. As a homogeneous space, $Gr(k,n) \cong O(n)/(O(k) \times O(n-k))$ (\autoref{prop:grassmann_stiefel_cociente}), which allows deriving the network operations from the representation theory of the orthogonal group.
\end{itemize}


\paragraph{Hyperbolic spaces: hierarchical data}

\begin{definition}[Hyperbolic Space]\label{def:espacio_hiperbolico}
The \textit{hyperbolic space} $\mathbb{H}^n$ is the unique simply connected space of constant sectional curvature $-1$ (the proof that $K = -1$ is presented in \autoref{thm:curvatura_hiperbolico} of Appendix~\ref{ch:teoremas_complementarios}). Geometrically, it is the $n$-dimensional analogue of the Lobachevsky-Bolyai hyperbolic plane.
\end{definition}

\textbf{Models of hyperbolic space.}\label{rem:modelos_hiperbolico} There exist several isometric models of $\mathbb{H}^n$:
\begin{itemize}
    \item \textbf{Poincaré disk model}: $\mathbb{D}^n = \{x \in \mathbb{R}^n : \|x\| < 1\}$
          with metric $ds^2 = \frac{4}{(1 - \|\mathbf{x}\|^2)^2} \sum_{i=1}^n (dx^i)^2$
    \item \textbf{Upper half-plane model}: $\mathbb{H}^2 = \{z \in \mathbb{C} : \Im(z) > 0\}$
          with metric $ds^2 = \frac{dx^2 + dy^2}{y^2}$
    \item \textbf{Hyperboloid model}: Surface $x_0^2 - x_1^2 - \cdots - x_n^2 = 1$ in $\mathbb{R}^{n+1}$
\end{itemize}

These three models are isometric as Riemannian manifolds; the explicit proof of the isometries between them (through stereographic projection and the Cayley transform) is presented in \autoref{sec:isometrias_hiperbolico} of Appendix~\ref{ch:teoremas_complementarios}.

The hyperbolic spaces $\mathbb{H}^n$ possess constant negative curvature, which confers on them a distinctive geometric property: the volume of geodesic balls grows exponentially with the radius. This characteristic makes them ideal for representing hierarchical and tree-structured data.

\textbf{Key applications}:
\begin{itemize}
    \item \textbf{Knowledge graphs and lexical semantics}: Taxonomic hierarchies (WordNet, biological ontologies) and hyperonymy (is-a) relations in natural language are embedded with less distortion in $\mathbb{H}^n$ than in $\mathbb{R}^n$ \citep{nickel2017poincare}. The required dimensionality is logarithmic rather than exponential (see Section~\ref{subsubsec:embeddings_hiperbolicos_grafos} for the mathematical justification and its connection with hyperbolic attention mechanisms).

    \item \textbf{Social networks}: Hierarchical communities in social networks exhibit a tree structure that is naturally captured in hyperbolic geometry \citep{krioukov2010hyperbolic}.
\end{itemize}

\textbf{Architectural implication}: The negative curvature of $\mathbb{H}^n$ requires redefining the basic operations of neural networks. Vector addition is replaced by \textit{Möbius addition} in the Poincaré disk model, and matrix multiplication is adapted through \textit{exponential and logarithmic maps} that transport vectors between the (Euclidean) tangent space and the (hyperbolic) manifold. The Poincaré disk model $\mathbb{D}^n$ is the most used computationally for its differentiability and compatibility with gradient-based optimization \citep{nickel2017poincare}.

\begin{figure}[htbp]
\centering
\includegraphics[width=\textwidth]{figures/hyperbolic_embedding}
\caption[Hierarchical embedding in the Poincaré disk]{Embedding of a hierarchical structure in the Poincaré disk $\mathbb{D}^2$. The exponential growth of volume in hyperbolic geometry allows representing trees and hierarchies with minimal distortion, in contrast to Euclidean space where the required dimensionality grows exponentially.}
\label{fig:hyperbolic_embedding}
\end{figure}

\paragraph{Product and composite spaces: multi-structured data}

Much real-world data combines multiple types of geometry. \textbf{Product spaces} allow modeling this structural heterogeneity.

\begin{definition}[Product Manifold]\label{def:variedad_producto}
Given two Riemannian manifolds $(\mathcal{M}, g_\mathcal{M})$ and $(\mathcal{N}, g_\mathcal{N})$, the \textit{product manifold} $\mathcal{M} \times \mathcal{N}$ is the manifold whose underlying set is the Cartesian product, equipped with the \textit{product metric}:
\[
g_{(p,q)}((v_1, w_1), (v_2, w_2)) = g_{\mathcal{M},p}(v_1, v_2) + g_{\mathcal{N},q}(w_1, w_2)
\]
for $p \in \mathcal{M}$, $q \in \mathcal{N}$, $v_i \in T_p\mathcal{M}$, $w_i \in T_q\mathcal{N}$.
\end{definition}

\begin{proposition}[Properties of product manifolds]\label{prop:propiedades_producto}
The product manifold $\mathcal{M} \times \mathcal{N}$ satisfies:
\begin{enumerate}
    \item $\dim(\mathcal{M} \times \mathcal{N}) = \dim \mathcal{M} + \dim \mathcal{N}$
    \item The geodesics are products of geodesics: $\gamma(t) = (\gamma_\mathcal{M}(t), \gamma_\mathcal{N}(t))$
    \item The sectional curvatures of planes tangent to each factor are preserved
\end{enumerate}
\end{proposition}

\begin{proof}
\leavevmode
\begin{enumerate}
    \item \textbf{Dimension.} The tangent spaces satisfy $T_{(p,q)}(\mathcal{M} \times \mathcal{N}) \cong T_p\mathcal{M} \oplus T_q\mathcal{N}$, hence $\dim(\mathcal{M} \times \mathcal{N}) = \dim T_p\mathcal{M} + \dim T_q\mathcal{N} = \dim \mathcal{M} + \dim \mathcal{N}$.

    \item \textbf{Geodesics.} The product metric $g = g_\mathcal{M} + g_\mathcal{N}$ has block-diagonal structure on the tangent spaces. The geodesic equation $\nabla_{\dot\gamma}\dot\gamma = 0$ decouples into the components of each factor, since the mixed Christoffel symbols vanish. Hence $\gamma(t) = (\gamma_\mathcal{M}(t), \gamma_\mathcal{N}(t))$ is a geodesic if and only if each component is a geodesic in its factor.

    \item \textbf{Sectional curvatures.} Let $\sigma$ be a plane tangent to $\mathcal{M}$ at $(p,q)$, that is, $\sigma \subset T_p\mathcal{M} \oplus \{0\}$. The Riemann tensor of the product metric is block-diagonal (the mixed terms vanish), hence $K_{\mathcal{M}\times\mathcal{N}}(\sigma) = K_\mathcal{M}(\sigma)$. \qedhere
\end{enumerate}
\end{proof}

\begin{definition}[$n$-dimensional torus]\label{def:toro_n}
The \textit{$n$-dimensional torus} is the product of $n$ circles:
\[
T^n = \underbrace{S^1 \times S^1 \times \cdots \times S^1}_{n \text{ times}} = (\mathbb{R}/\mathbb{Z})^n
\]
The torus has zero curvature (it is flat) and admits periodic coordinates $(\theta_1, \ldots, \theta_n) \in [0, 2\pi)^n$.
\end{definition}

\textbf{Applications of product spaces:}

\begin{itemize}
    \item \textbf{Torus} $T^n = (S^1)^n$: For periodic data such as cyclic time series (seasonal climate patterns, circadian rhythms) and Fourier analysis.

    \item \textbf{Products of manifolds}: Graphs with geometric node attributes can be modeled as products $G \times \mathcal{M}$, where $G$ is the discrete structure of the graph and $\mathcal{M}$ is the space of continuous attributes.

    \item \textbf{Mixed-curvature embeddings}: Data with heterogeneous structure (partially hierarchical, partially cyclic) is represented in products $\mathbb{H}^{d_1} \times S^{d_2} \times \mathbb{R}^{d_3}$, where each component captures a different geometric aspect of the data. The dimensionality of each factor is learned automatically \citep{bronstein2017geometricdeeplearning}.
\end{itemize}

\paragraph{Graphs as discrete geometry}

Graphs $G = (V, E)$ represent a special case of \textit{discrete} non-Euclidean geometry. As developed extensively in \autoref{sec:grafos_discretizacion}, graphs provide a more flexible representation than Euclidean spaces by capturing arbitrary connectivity without imposing rigid metrics.

\paragraph{Synthesis: geometry as inductive bias}

The choice of the geometric space is not a mere representation detail: it constitutes a fundamental \textit{inductive bias} that determines the appropriate neural architecture. Each geometry imposes specific restrictions on the admissible operations:

\begin{table}[htbp]
\centering
\small
\caption{Correspondence between geometries, inductive biases, and operations in GDL.}
\label{tab:geometria_sesgo_operacion}
\begin{tabular}{@{}lll@{}}
    \toprule
    \textbf{Geometry} & \textbf{Inductive bias} & \textbf{GDL operation} \\
    \midrule
    $\mathbb{R}^n$ (Euclidean) & Translational symmetry & Standard convolution \\
    $S^n$ (spherical, $K>0$) & Rotational symmetry & Convolution over $SO(n{+}1)$ \\
    $\mathbb{H}^n$ (hyperbolic, $K<0$) & Exponential hierarchy & Möbius addition, exp/log maps \\
    $Gr(k,n)$ (Grassmann) & Data as subspaces & Geodesic metrics in $Gr$ \\
    $\mathcal{M} \times \mathcal{N}$ (product) & Heterogeneous structure & Component-wise operations \\
    Graphs & Discrete connectivity & Message passing \\
    \bottomrule
\end{tabular}
\end{table}

This correspondence between geometry and architecture is the central organizing principle of \gls{GDL}: the symmetries of the geometric domain determine the equivariance constraints that the neural network must satisfy. Notably, all the spaces presented in this section admit a unified description as homogeneous spaces $G/H$, as formalized and demonstrated in \autoref{subsec:espacios_homogeneos}; the geometric interpretation of each decomposition (how $G$ and $H$ determine curvature, compactness, and volume growth) is detailed alongside each proposition in \autoref{subsec:espacios_homogeneos}. The algebraic foundations for formalizing these symmetries are developed in \autoref{sec:grupos_simetrias}, and the complete table with specific architectures and references is presented in \autoref{rem:espacios_homogeneos_arquitecturas}.
\section{Groups: The Symmetries of Data}\label{sec:grupos_simetrias}

The previous section identified the geometric domains where data reside and pointed out that each domain possesses natural symmetries: translations in $\mathbb{R}^n$, rotations in $S^n$, isometries in $\mathbb{H}^n$. But what does it mean exactly for a neural architecture to ``respect'' these symmetries? The answer requires a precise algebraic language. \textbf{Group theory} provides this language: a group captures the collection of all transformations that preserve the structure of a space, and \textbf{representation theory} translates these abstract symmetries into concrete constraints on the functions that the network can learn. From these foundations, we formalize the concepts of \textbf{equivariance} and \textbf{invariance} (\autoref{subsec:equivarianza_invarianza}) and show that the geometric domains studied earlier admit a unified description as \textbf{homogeneous spaces} (\autoref{subsec:espacios_homogeneos}).

\subsection{Group Theory}\label{subsec:teoria_grupos_representaciones}

We begin with the algebraic foundations: the definition of group, subgroups, group actions, and Lie groups. These structures allow us to model both discrete symmetries (permutations $S_n$) and continuous ones (rotations $SO(n)$, translations $\mathbb{R}^d$). The natural progression leads from groups to their linear representations, which are the mechanism by which symmetries act on the feature spaces of a neural network.

\begin{definition}[Internal Operation]\label{def:operacion_interna}
An \textit{internal operation} (or binary operation) on a set $G$ is a function:
\begin{equation}
\cdot: G \times G \rightarrow G
\end{equation}
that assigns to each ordered pair $(a,b) \in G \times G$ a unique element $a \cdot b \in G$.
\end{definition}

\begin{definition}[Group]\label{def:grupo}
A \textit{group} is a set $G$ equipped with an internal operation $\cdot: G \times G \rightarrow G$ that satisfies:
\begin{enumerate}
    \item \textbf{Associativity}: $(a \cdot b) \cdot c = a \cdot (b \cdot c)$ for all $a, b, c \in G$
    \item \textbf{Identity element}: There exists $e \in G$ such that $e \cdot a = a \cdot e = a$ for all $a \in G$
    \item \textbf{Inverse element}: For each $a \in G$, there exists $a^{-1} \in G$ such that $a \cdot a^{-1} = a^{-1} \cdot a = e$
\end{enumerate}
\end{definition}

\textbf{Abelian Groups.}\label{rem:grupos_abelianos} A group $(G, \cdot)$ is \textbf{abelian} (or commutative) if $a \cdot b = b \cdot a$ for all $a, b \in G$. Translations in \glspl{CNN} form abelian groups, where the order in which the translations are applied does not affect the final result: $(x,y) + (a,b) + (c,d) = (x,y) + (c,d) + (a,b)$.

\begin{definition}[Subgroup]\label{def:subgrupo}
Let $(G, \cdot)$ be a group. A subset $H \subseteq G$ is a \textit{subgroup} if $(H, \cdot)$ is a group under the inherited operation. Notation: $H \leq G$.
\end{definition}

\begin{proposition}[Subgroup Criterion]\label{prop:criterio_subgrupo}
$H \leq G$ if and only if:
\begin{enumerate}
    \item $e \in H$ (contains the identity element)
    \item $h_1, h_2 \in H \Rightarrow h_1 h_2 \in H$ (closed under the operation)
    \item $h \in H \Rightarrow h^{-1} \in H$ (contains inverses)
\end{enumerate}
\end{proposition}

\begin{proof}
$(\Rightarrow)$ If $H \leq G$, the three conditions are immediate from the group axioms.

$(\Leftarrow)$ Suppose $H$ satisfies (1)--(3). The operation on $H$ is the one inherited from $G$, which is associative since it is associative in $G$. Condition (1) provides the identity, (2) guarantees closure, and (3) the existence of inverses. Hence $(H, \cdot)$ is a group.
\end{proof}

\begin{definition}[Normal Subgroup]\label{def:subgrupo_normal}
A subgroup $H \leq G$ is \textit{normal} (notation: $H \triangleleft G$) if:
\begin{equation}
gHg^{-1} = H \quad \text{for all } g \in G
\end{equation}
Equivalently, $ghg^{-1} \in H$ for all $g \in G$ and $h \in H$.
\end{definition}

\begin{definition}[Order of a Group]\label{def:orden_grupo}
The \textit{order} of a group $G$, denoted $|G|$, is its cardinality. The \textbf{index} of a subgroup $H \leq G$, denoted $[G:H]$, is the number of cosets of $H$ in $G$.
\end{definition}

For finite groups, the order of the group and that of any subgroup are related precisely by Lagrange's Theorem (\autoref{thm:lagrange}, see \autoref{sec:lagrange} of Appendix~\ref{ch:teoremas_complementarios}).

The concepts of normal subgroup and group order lead us naturally toward a fundamental algebraic construction: the quotient group. When a subgroup $H$ is normal in $G$, the cosets $gH$ not only partition the group, but also inherit a well-defined group structure. This construction is essential in \gls{GDL} because it allows us to model systematic reductions of symmetries in hierarchical architectures.

\paragraph{Quotient Group}

\begin{definition}[Quotient Group]\label{def:grupo_cociente}
Let $G$ be a group and $H \triangleleft G$ a normal subgroup. The \textit{quotient group} (or factor group) $G/H$ is the set of left cosets:
\begin{equation}
G/H = \{gH : g \in G\}
\end{equation}
equipped with the operation:
\begin{equation}
(g_1H) \cdot (g_2H) = (g_1g_2)H
\end{equation}
This operation is well defined, as the following proposition shows.
\end{definition}

\begin{proposition}[Well-definedness of the quotient operation]\label{prop:cociente_bien_definido}
Let $G$ be a group and $H \leq G$ a subgroup. The operation $(g_1H)(g_2H) = (g_1g_2)H$ on the set of cosets $G/H$ is well defined if and only if $H \triangleleft G$.
\end{proposition}

\begin{proof}
$(\Leftarrow)$ Suppose $H \triangleleft G$. Let $g_1H = g_1'H$ and $g_2H = g_2'H$, that is, there exist $h_1, h_2 \in H$ such that $g_1' = g_1h_1$ and $g_2' = g_2h_2$. Then:
\[
g_1'g_2' = g_1h_1g_2h_2 = g_1g_2\underbrace{(g_2^{-1}h_1g_2)}_{\in H}h_2
\]
Since $H \triangleleft G$, we have $g_2^{-1}h_1g_2 \in H$, hence $(g_2^{-1}h_1g_2)h_2 \in H$ and therefore $(g_1'g_2')H = (g_1g_2)H$.

$(\Rightarrow)$ Suppose the operation is well defined. Let $g \in G$ and $h \in H$ be arbitrary. Since $h \in H$, we have $ghH = gH$. Then:
\[
(ghg^{-1})H = (ghH)(g^{-1}H) = (gH)(g^{-1}H) = (gg^{-1})H = eH = H
\]
where the second equality uses well-definedness with the representatives $gh$ and $g$ of the same coset $gH$. Thus, $ghg^{-1} \in H$ for all $g \in G$ and $h \in H$, which means $H \triangleleft G$.
\end{proof}

Quotient groups are fundamental in \gls{GDL} because they model symmetry reductions in hierarchical architectures. For example, in networks that process 3D data, we can start with full $SE(3)$ symmetries and gradually reduce to $SO(3)$ or $SO(2)$ in higher layers via natural morphisms $SE(3) \rightarrow SO(3)$ whose kernel (translations) forms a normal subgroup.

The construction of quotient groups depends crucially on the existence of functions that preserve the group structure. This concept leads us directly to homomorphisms, functions that respect the algebraic operations and constitute the bridge between different groups. In \gls{GDL}, homomorphisms translate abstract symmetries into computable transformations, providing the formal mechanism to implement equivariance in neural architectures.

\paragraph{Homomorphisms}

\begin{definition}[Group Homomorphism]\label{def:homomorfismo_grupos}
Let $(G, \cdot_G)$ and $(H, \cdot_H)$ be groups. A function $\phi: G \rightarrow H$ is a \textit{homomorphism} if:
\begin{equation}
\phi(g_1 \cdot_G g_2) = \phi(g_1) \cdot_H \phi(g_2) \quad \forall g_1, g_2 \in G
\end{equation}
\end{definition}

\begin{definition}[Group Isomorphism]\label{def:isomorfismo_grupos}
Let $(G, \cdot_G)$ and $(H, \cdot_H)$ be groups. A homomorphism $\phi: G \rightarrow H$ is an \textit{isomorphism} if it is bijective. In this case, we say that $G$ and $H$ are \textit{isomorphic} and write $G \cong H$.
\end{definition}

Two isomorphic groups have the same algebraic structure and differ only in the symbols used to denote the elements and the operation. Isomorphisms allow us to transfer properties between structurally identical groups, facilitating the analysis of complex symmetries through simpler or more familiar groups.

\begin{definition}[Kernel and Image]\label{def:nucleo_imagen}
Let $\phi: G \rightarrow H$ be a group homomorphism. The \textit{kernel} of $\phi$ is:
\begin{equation}
\ker(\phi) = \{g \in G : \phi(g) = e_H\}
\end{equation}
where $e_H$ is the identity of $H$. The \textit{image} of $\phi$ is:
\begin{equation}
\mathrm{Im}(\phi) = \{\phi(g) : g \in G\} \subseteq H
\end{equation}
The kernel is always a normal subgroup of $G$, that is, $\ker(\phi) \triangleleft G$.
\end{definition}

\begin{theorem}[First Isomorphism Theorem]\label{thm:primer_isomorfismo}
Let $\phi: G \rightarrow H$ be a group homomorphism. Then:
\begin{equation}
G/\ker(\phi) \cong \mathrm{Im}(\phi)
\end{equation}
The isomorphism is given by $\bar{\phi}: G/\ker(\phi) \rightarrow \mathrm{Im}(\phi)$, $\bar{\phi}(g\ker(\phi)) = \phi(g)$.
\end{theorem}

\begin{proof}
\textbf{Step 1: $\bar{\phi}$ is well defined.} If $g_1\ker(\phi) = g_2\ker(\phi)$, then $g_1^{-1}g_2 \in \ker(\phi)$, so $\phi(g_1^{-1}g_2) = e_H$. This implies $\phi(g_1)^{-1}\phi(g_2) = e_H$, that is, $\phi(g_1) = \phi(g_2)$.

\textbf{Step 2: $\bar{\phi}$ is a homomorphism.}
\[
\bar{\phi}(g_1\ker(\phi) \cdot g_2\ker(\phi)) = \bar{\phi}((g_1g_2)\ker(\phi)) = \phi(g_1g_2) = \phi(g_1)\phi(g_2) = \bar{\phi}(g_1\ker(\phi)) \cdot \bar{\phi}(g_2\ker(\phi))
\]

\textbf{Step 3: $\bar{\phi}$ is injective.} If $\bar{\phi}(g\ker(\phi)) = e_H$, then $\phi(g) = e_H$, so $g \in \ker(\phi)$, and $g\ker(\phi) = \ker(\phi)$ is the identity in $G/\ker(\phi)$.

\textbf{Step 4: $\bar{\phi}$ is surjective onto $\mathrm{Im}(\phi)$ by definition.}
\end{proof}

This theorem is fundamental because it allows us to identify quotient groups with images of homomorphisms, simplifying the analysis of complex algebraic structures.

\begin{definition}[Group Representation]\label{def:representacion_grupo}
Let $G$ be a group and $V$ a vector space over a field $\mathbb{K}$. A \textit{representation} of $G$ on $V$ is a group homomorphism:
\begin{equation}
\rho: G \rightarrow GL(V)
\end{equation}
where $GL(V)$ denotes the group of invertible linear transformations of $V$. That is, to each element $g \in G$ an invertible matrix $\rho(g)$ is assigned so that composition in the group translates into matrix product: $\rho(g_1 g_2) = \rho(g_1)\rho(g_2)$.
\end{definition}

Representations are fundamental in \gls{GDL} because they translate abstract symmetries into computable operations on neural feature spaces: each symmetry of the domain becomes a concrete linear transformation that can be implemented and optimized.

Homomorphisms establish relations between abstract groups, but to apply groups to concrete problems we need them to act on specific data spaces. This need leads us to the concept of group action, which formalizes how a group transforms elements of a set. Group actions translate abstract symmetries into explicit transformations on concrete geometric domains such as images, graphs, or manifolds.

\paragraph{Group Action}

\begin{definition}[Group Action]\label{def:accion_grupo}
Let $G$ be a group and $X$ a set. A \textit{group action} of $G$ on $X$ is a function:
\begin{equation}
\rho: G \times X \rightarrow X, \quad (g, x) \mapsto g \cdot x
\end{equation}
that satisfies:
\begin{enumerate}
    \item \textbf{Compatibility}: $(g_1 g_2) \cdot x = g_1 \cdot (g_2 \cdot x)$ for all $g_1, g_2 \in G$ and $x \in X$
    \item \textbf{Identity}: $e \cdot x = x$ for all $x \in X$, where $e$ is the identity element of $G$
\end{enumerate}
\end{definition}

\begin{proposition}[Actions as Homomorphisms]\label{prop:acciones_homomorfismos}
Every group action $\rho: G \times X \rightarrow X$ induces a group homomorphism:
\begin{equation}
\tilde{\rho}: G \rightarrow \text{Sym}(X), \quad \tilde{\rho}(g)(x) = g \cdot x
\end{equation}
where $\text{Sym}(X)$ denotes the group of bijections of $X$ with itself.
\end{proposition}

\begin{proof}
For each $g \in G$, we define $\tilde{\rho}(g): X \rightarrow X$ as the map $x \mapsto g \cdot x$. We must verify:
\begin{enumerate}
    \item \textbf{$\tilde{\rho}(g)$ is bijective}: Its inverse is $\tilde{\rho}(g^{-1})$, since $(g^{-1}) \cdot (g \cdot x) = (g^{-1}g) \cdot x = e \cdot x = x$.
    \item \textbf{$\tilde{\rho}$ preserves the operation}: For $g_1, g_2 \in G$,
    \begin{equation}
    \tilde{\rho}(g_1 g_2)(x) = (g_1 g_2) \cdot x = g_1 \cdot (g_2 \cdot x) = \tilde{\rho}(g_1)(\tilde{\rho}(g_2)(x)) = (\tilde{\rho}(g_1) \circ \tilde{\rho}(g_2))(x)
    \end{equation}
\end{enumerate}
Therefore, $\tilde{\rho}$ is a homomorphism from $G$ into $\text{Sym}(X)$.
\end{proof}

This correspondence is fundamental: studying actions of $G$ on $X$ is equivalent to studying homomorphisms from $G$ into $\text{Sym}(X)$. In the context of \gls{GDL}, this means that the symmetries of the data translate directly into concrete transformations on the input space.

\begin{example}[Group Actions in Deep Learning]\label{ex:acciones_grupo_dl}
\leavevmode
\begin{itemize}
    \item \textbf{Translations in images}: The group $\mathbb{Z}^2$ acts on discrete images by shifting pixels
    \item \textbf{Rotations in 3D data}: $SO(3)$ acts on point clouds via rotation matrices
    \item \textbf{Permutations in graphs}: $S_n$ acts on graphs with $n$ nodes by reordering the node indices
\end{itemize}
\end{example}

Before moving on to Lie groups, we introduce a fundamental algebraic construction that we will use to describe symmetry groups such as the Euclidean group.

\begin{definition}[Semidirect Product]\label{def:producto_semidirecto}
Let $N$ and $H$ be groups, and let $\phi: H \to \mathrm{Aut}(N)$ be a homomorphism (where $\mathrm{Aut}(N)$ is the group of automorphisms of $N$). The \textit{semidirect product} of $N$ by $H$ with respect to $\phi$, denoted $N \rtimes_\phi H$ (or simply $N \rtimes H$), is the set $N \times H$ equipped with the operation:
\begin{equation}\label{eq:producto_semidirecto}
(n_1, h_1) \cdot (n_2, h_2) = (n_1 \cdot \phi(h_1)(n_2), h_1 \cdot h_2)
\end{equation}
where $\phi(h_1)(n_2)$ denotes the action of the automorphism $\phi(h_1)$ on $n_2$.
\end{definition}

Intuitively, the semidirect product generalizes the direct product $N \times H$ by allowing $H$ to ``act'' on $N$ in a nontrivial way. In the direct product, the elements of $N$ and $H$ commute; in the semidirect product, $H$ ``twists'' $N$ via automorphisms.

\begin{example}[Euclidean Group as a Semidirect Product]\label{ex:grupo_euclidiano_semidirecto}
The \textit{Euclidean group} $E(n) = \mathbb{R}^n \rtimes O(n)$ is the semidirect product of the translations $\mathbb{R}^n$ by the rotations and reflections $O(n)$. The action $\phi: O(n) \to \mathrm{Aut}(\mathbb{R}^n)$ is matrix multiplication: $\phi(R)(t) = Rt$. Thus, the group operation is:
\begin{equation}
(t_1, R_1) \cdot (t_2, R_2) = (t_1 + R_1 t_2, R_1 R_2)
\end{equation}
This structure reflects how rotations ``affect'' translations: first one rotates, then one translates.
\end{example}

So far we have developed group theory from a purely algebraic perspective, treating groups as abstract structures defined by binary operations. However, the applications of \gls{GDL} frequently require continuous symmetries: smooth rotations in $SO(3)$, translations in $\mathbb{R}^d$, or Euclidean transformations in $SE(3)$. This transition from the discrete to the continuous leads us naturally toward Lie groups, where the algebraic structure is combined with differential geometry.

\subsubsection{Lie Groups}

Differentiable manifolds (see \autoref{def:variedad_diferenciable}) constitute the natural geometry on which many \gls{GDL} architectures operate, especially when dealing with continuous data such as 3D point clouds or surfaces. For this reason, it is essential to consider Lie groups instead of abstract groups, since continuous symmetries require differentiable structure.

\begin{definition}[Lie Group]\label{def:grupo_lie}
A \textit{Lie group} is a group $G$ that is also a differentiable manifold, where the group operations:
\begin{itemize}
    \item \textbf{Multiplication}: $\mu: G \times G \rightarrow G, \quad (g,h) \mapsto gh$
    \item \textbf{Inversion}: $\iota: G \rightarrow G, \quad g \mapsto g^{-1}$
\end{itemize}
are smooth functions (of class $C^\infty$).
\end{definition}

Lie groups constitute the most important class of groups in \gls{GDL}, since most continuous symmetries in real applications (rotations, translations, Euclidean transformations) possess a natural differentiable structure.

\begin{definition}[Smooth Action]\label{def:accion_suave}
Let $G$ be a Lie group and $M$ a differentiable manifold. A \textit{smooth action} of $G$ on $M$ is a group action $\rho: G \times M \rightarrow M$ that is moreover a differentiable map (of class $C^\infty$).
\end{definition}

The smoothness condition guarantees that the group transformations preserve the differentiable structure of the manifold, allowing the use of tools from differential calculus such as derivatives and flows. In \gls{GDL}, this is essential for defining gradients and performing optimization over spaces of transformations.

\begin{definition}[Lie Bracket]\label{def:corchete_lie}
Let $\mathfrak{g}$ be a vector space. A \textit{Lie bracket} is a binary operation $[\cdot, \cdot]: \mathfrak{g} \times \mathfrak{g} \rightarrow \mathfrak{g}$ that satisfies:
\begin{enumerate}
    \item \textbf{Bilinearity}: $[aX + bY, Z] = a[X,Z] + b[Y,Z]$ for all $a, b \in \mathbb{R}$
    \item \textbf{Antisymmetry}: $[X,Y] = -[Y,X]$
    \item \textbf{Jacobi identity}: $[X,[Y,Z]] + [Y,[Z,X]] + [Z,[X,Y]] = 0$
\end{enumerate}
\end{definition}

\begin{definition}[Lie Algebra]\label{def:algebra_lie}
A \textit{Lie algebra} is a vector space $\mathfrak{g}$ equipped with a Lie bracket $[\cdot, \cdot]: \mathfrak{g} \times \mathfrak{g} \to \mathfrak{g}$.
\end{definition}

Every Lie group $G$ carries an associated canonical Lie algebra: the tangent space at the identity $T_e G$, with the bracket induced by the left-invariant vector fields. The advantage of this construction is that it allows us to work in a linear vector space instead of a curved manifold, which is fundamental for optimization and differentiation in neural architectures. The bridge between the (nonlinear) group and its (linear) algebra is the \textit{exponential map}.

\begin{definition}[Exponential Map]\label{def:mapa_exponencial}
Let $G$ be a Lie group with Lie algebra $\mathfrak{g}$. The \textit{exponential map} is the smooth map:
\begin{equation}
\exp: \mathfrak{g} \rightarrow G
\end{equation}
that sends each $X \in \mathfrak{g}$ to the point $\gamma(1)$, where $\gamma: \mathbb{R} \rightarrow G$ is the unique one-parameter subgroup with $\gamma'(0) = X$.
\end{definition}

For matrix Lie groups, the exponential map coincides with the matrix exponential:
\begin{equation}
\exp(A) = \sum_{k=0}^{\infty} \frac{A^k}{k!} = I + A + \frac{A^2}{2!} + \frac{A^3}{3!} + \cdots
\end{equation}

\begin{example}[Exponential Map in $SO(2)$]\label{ex:mapa_exp_so2}
For $SO(2)$, the Lie algebra $\mathfrak{so}(2)$ consists of matrices of the form $\theta J$ where $J = \begin{psmallmatrix} 0 & -1 \\ 1 & 0 \end{psmallmatrix}$. The exponential map gives:
\begin{equation}
\exp(\theta J) = \begin{pmatrix} \cos\theta & -\sin\theta \\ \sin\theta & \cos\theta \end{pmatrix} = R_\theta
\end{equation}
Thus, the exponential converts an angle (element of the algebra) into a rotation matrix (element of the group).
\end{example}

\begin{figure}[htbp]
\centering
\includegraphics[width=0.7\textwidth]{figures/exponential_map.pdf}
\caption[Exponential map of Lie groups]{Exponential map $\exp: \mathfrak{g} \to G$. It transforms vectors of the Lie algebra (tangent space at the identity $e$) into elements of the Lie group, connecting the linear structure of the algebra with the nonlinear structure of the group.}
\label{fig:exponential_map}
\end{figure}

\paragraph{Representations of Lie Groups}

Representation theory extends the concept of group representation (\autoref{def:representacion_grupo}) to the differentiable setting.

\begin{definition}[Lie Group Representation]\label{def:representacion_grupo_lie}
Let $G$ be a Lie group and $V$ a finite-dimensional vector space. A \textit{representation of $G$} is a Lie group homomorphism:
\begin{equation}
\rho: G \rightarrow GL(V)
\end{equation}
that is, $\rho$ is a group homomorphism that is moreover a smooth map.
\end{definition}

\begin{proposition}[Induced Lie Algebra Representation]\label{prop:representacion_algebra_lie}
Every representation $\rho: G \rightarrow GL(V)$ of a Lie group induces a \textit{Lie algebra representation} $d\rho: \mathfrak{g} \rightarrow \mathfrak{gl}(V)$ defined by:
\begin{equation}
d\rho(X) = \frac{d}{dt}\Big|_{t=0} \rho(\exp(tX))
\end{equation}
This representation satisfies $d\rho([X,Y]) = [d\rho(X), d\rho(Y)]$ (it preserves the Lie bracket).
\end{proposition}

\begin{proof}
\textbf{Well-definedness.} Since $\rho$ is smooth and $\exp: \mathfrak{g} \to G$ is smooth, the composition $t \mapsto \rho(\exp(tX))$ is a smooth curve in $GL(V)$. Its derivative at $t=0$ exists and belongs to $\mathfrak{gl}(V) = T_{\text{Id}} GL(V)$.

\textbf{Linearity.} For $X, Y \in \mathfrak{g}$ and $\alpha \in \mathbb{R}$:
\[
d\rho(\alpha X) = \frac{d}{dt}\Big|_{t=0} \rho(\exp(t\alpha X)) = \alpha \frac{d}{ds}\Big|_{s=0} \rho(\exp(sX)) = \alpha\, d\rho(X)
\]
For additivity, since $\exp(t(X+Y)) = \exp(tX)\exp(tY)\exp(-\tfrac{t^2}{2}[X,Y] + O(t^3))$ by the Baker-Campbell-Hausdorff formula, differentiating at $t=0$ the terms of order $\geq 2$ vanish, and using that $\rho$ is a homomorphism:
\[
d\rho(X+Y) = \frac{d}{dt}\Big|_{t=0} \rho(\exp(t(X+Y))) = \frac{d}{dt}\Big|_{t=0} \rho(\exp(tX))\rho(\exp(tY)) = d\rho(X) + d\rho(Y)
\]
where the last equality is obtained by the product rule for derivatives.

\textbf{Preservation of the bracket.} The Baker-Campbell-Hausdorff formula gives:
\[
\exp(tX)\exp(tY)\exp(-tX)\exp(-tY) = \exp(t^2[X,Y] + O(t^3))
\]
Applying $\rho$ (which is a homomorphism) and differentiating twice at $t=0$:
\[
\frac{d^2}{dt^2}\Big|_{t=0} \rho(\exp(t^2[X,Y] + O(t^3))) = d\rho([X,Y])
\]
On the other hand, expanding the left-hand side with the product rule yields $d\rho(X)d\rho(Y) - d\rho(Y)d\rho(X) = [d\rho(X), d\rho(Y)]$. See \citep[Theorem~3.28]{Kirillov2008} for the complete details.
\end{proof}

This correspondence is fundamental in \gls{GDL}: it allows us to optimize in the linear space $\mathfrak{g}$ (using standard gradients) and then apply transformations in the group $G$ via the exponential map.

\begin{definition}[Irreducible Representation]\label{def:representacion_irreducible}
A representation $\rho: G \rightarrow GL(V)$ is \textit{irreducible} if $V \neq \{0\}$ and the only subspaces of $V$ invariant under the action of $G$ are $\{0\}$ and $V$ itself. That is, there is no proper subspace $W \subset V$ such that $\rho(g)(W) \subseteq W$ for all $g \in G$.
\end{definition}

Irreducible representations are the ``fundamental building blocks'' from which all other representations are constructed.

\begin{example}[Irreducible Representations of $SO(2)$]\label{ex:irreps_so2}
The irreducible representations of $SO(2)$ over $\mathbb{C}$ are one-dimensional, indexed by $n \in \mathbb{Z}$:
\begin{equation}
\rho_n: SO(2) \rightarrow \mathbb{C}^*, \quad \rho_n(R_\theta) = e^{in\theta}
\end{equation}
Each $\rho_n$ corresponds to a ``frequency'' in circular Fourier analysis (see \autoref{sec:analisis_armonico} of \autoref{ch:cnn}). In G-CNNs over rotated images, the feature channels are organized according to these representations.
\end{example}

The following fundamental results of representation theory require the group to be \textit{compact}:

\begin{definition}[Compact Group]\label{def:grupo_compacto}
A Lie group $G$ is \textit{compact} if it is compact as a topological space. Equivalently, $G$ is compact if and only if it is closed and bounded in any embedding into $\mathbb{R}^n$.
\end{definition}

The most important compact groups in \gls{GDL} include $SO(n)$, $O(n)$, $SU(n)$, $U(n)$, and all finite groups.

The results we develop next constitute the harmonic analysis of compact groups, generalizing classical Fourier analysis (where the underlying group is $SO(2) \cong S^1$) to arbitrary symmetry groups. The natural sequence of construction is: invariant integration over the group (Haar measure), decomposition of representations into irreducible blocks, uniqueness of these blocks (Schur's lemma), coordinate functions and their orthogonality relations, and finally the completeness of the system (Peter-Weyl theorem). In \gls{GDL}, this machinery provides the theoretical basis for implementing equivariant convolutions in the spectral domain, transforming the problem of designing equivariant neural networks into a well-structured algebraic problem.

\begin{definition}[Haar Measure]\label{def:medida_haar_gdl}
Let $G$ be a compact Lie group. The \textit{Haar measure} is the unique measure $\mu$ on $G$ that satisfies:
\begin{enumerate}
    \item \textbf{Invariance}: $\mu(gA) = \mu(Ag) = \mu(A)$ for all $g \in G$ and measurable set $A$
    \item \textbf{Normalization}: $\mu(G) = 1$
\end{enumerate}
This measure allows us to define integrals over the group: $\int_G f(g) \, d\mu(g)$.
\end{definition}

The Haar measure provides a canonical way to ``average over all the symmetries'' of the group. In \gls{GDL}, this averaging is the fundamental mechanism underlying group convolutions: the convolution integrates a filter against the signal using precisely the Haar measure, and the invariance of this measure under group translations is what guarantees the equivariance of the resulting operator \citep{cohen2016groupequivariantconvolutional, kondor2018generalizationconvolutionalneural}.

\begin{theorem}[Decomposition of Representations]\label{thm:descomposicion_representaciones}
Every representation of a compact group $G$ on a finite-dimensional Hilbert space can be decomposed as a direct sum of irreducible representations:
\begin{equation}
V = \bigoplus_{i} n_i V_i
\end{equation}
where each $V_i$ is an irreducible representation and $n_i \in \mathbb{N}$ is its multiplicity.
\end{theorem}

\begin{proof}
Let $\rho: G \rightarrow GL(V)$ be a representation of a compact group $G$ on a finite-dimensional vector space $V$. The proof proceeds in three steps \citep{Fulton1991}.

\textbf{Step 1: $G$-invariant inner product.}
Let $\langle \cdot, \cdot \rangle_0$ be an arbitrary inner product on $V$. We define:
\begin{equation}\label{eq:producto_invariante}
\langle u, v \rangle = \int_G \langle \rho(g)u, \rho(g)v \rangle_0 \, d\mu(g)
\end{equation}
where $\mu$ is the normalized Haar measure on $G$ (\autoref{def:medida_haar_gdl}). Let us verify that $\langle \cdot, \cdot \rangle$ is a $G$-invariant inner product:
\begin{itemize}
    \item \textit{Linearity and symmetry}: inherited from $\langle \cdot, \cdot \rangle_0$ and the linearity of the integral.
    \item \textit{Positivity}: if $u \neq 0$, then $\rho(g)u \neq 0$ for all $g$ (since $\rho(g)$ is invertible), hence $\langle \rho(g)u, \rho(g)u \rangle_0 > 0$ for all $g$, and the integral is strictly positive.
    \item \textit{$G$-invariance}: for any $h \in G$,
    \[
    \langle \rho(h)u, \rho(h)v \rangle = \int_G \langle \rho(g)\rho(h)u, \rho(g)\rho(h)v \rangle_0 \, d\mu(g) = \int_G \langle \rho(gh)u, \rho(gh)v \rangle_0 \, d\mu(g).
    \]
    The change of variable $g' = gh$ and the invariance of the Haar measure give $d\mu(g) = d\mu(g')$, hence this integral equals $\langle u, v \rangle$.
\end{itemize}

\textbf{Step 2: Invariant orthogonal complement.}
Let $W \subseteq V$ be a $G$-invariant subspace. Consider the orthogonal complement $W^\perp$ with respect to the $G$-invariant inner product constructed in Step~1. For any $v \in W^\perp$, $g \in G$ and $w \in W$:
\[
\langle \rho(g)v, w \rangle = \langle v, \rho(g^{-1})w \rangle = 0
\]
where the first equality uses the $G$-invariance of the inner product (with $h = g^{-1}$) and the second follows from $\rho(g^{-1})w \in W$ (since $W$ is $G$-invariant) and $v \in W^\perp$. Therefore $\rho(g)v \in W^\perp$, which proves that $W^\perp$ is $G$-invariant.

\textbf{Step 3: Induction on $\dim(V)$.}
We proceed by strong induction on $n = \dim(V)$.

\textit{Base case}: if $n = 1$, every representation is irreducible (there are no nontrivial proper subspaces) and $V$ is its own decomposition.

\textit{Inductive step}: suppose the result is proved for every space of dimension less than $n$. If $\rho$ is irreducible, the decomposition is $V = V$ and there is nothing to prove. If $\rho$ is not irreducible, there exists a proper $G$-invariant subspace $W \subset V$ with $0 < \dim(W) < n$. By Step~2, $W^\perp$ is also $G$-invariant, and $V = W \oplus W^\perp$ with $\dim(W) < n$ and $\dim(W^\perp) < n$. By the inductive hypothesis, both $W$ and $W^\perp$ decompose as direct sums of irreducible representations, and their union provides the desired decomposition of $V$.
\end{proof}

This theorem guarantees that any representation, however complex, can be analyzed in terms of its irreducible components. In \gls{GDL}, the feature maps of an equivariant network are representations of the symmetry group, and the decomposition theorem guarantees that they can always be organized as direct sums of irreducible ``frequency channels''. Steerable networks (see \autoref{subsubsec:steerable_cnns_so2}) exploit this principle explicitly, parameterizing the features as collections of irreducible representations of different types \citep{weiler2018learningsteerable, cohen2016steerablecnns}.

The following lemma, due to Issai Schur, is fundamental in representation theory:

\begin{lemma}[Schur]\label{lem:schur}
Let $(\rho, V)$ and $(\sigma, W)$ be irreducible representations of a group $G$, and let $T: V \to W$ be a morphism of representations (that is, $T \circ \rho(g) = \sigma(g) \circ T$ for all $g \in G$). Then:
\begin{enumerate}
    \item $T$ is an isomorphism or $T = 0$.
    \item If $V = W$ and the base field is $\mathbb{C}$, then $T = \lambda \cdot \mathrm{Id}$ for some $\lambda \in \mathbb{C}$.
\end{enumerate}
\end{lemma}

\begin{proof}
\textbf{Part 1.} Consider $\ker(T) \subseteq V$. For any $v \in \ker(T)$ and $g \in G$:
\[
T(\rho(g)v) = \sigma(g)T(v) = \sigma(g) \cdot 0 = 0
\]
Therefore $\rho(g)v \in \ker(T)$, which shows that $\ker(T)$ is a $G$-invariant subspace of $V$. Since $\rho$ is irreducible, $\ker(T) = \{0\}$ or $\ker(T) = V$.

Analogously, $\mathrm{Im}(T) \subseteq W$ is $G$-invariant: for $w = T(v) \in \mathrm{Im}(T)$,
\[
\sigma(g)w = \sigma(g)T(v) = T(\rho(g)v) \in \mathrm{Im}(T)
\]
Since $\sigma$ is irreducible, $\mathrm{Im}(T) = \{0\}$ or $\mathrm{Im}(T) = W$.

If $\ker(T) = V$, then $T = 0$. If $\ker(T) = \{0\}$ and $\mathrm{Im}(T) = W$, then $T$ is an isomorphism.

\textbf{Part 2.} Let $\lambda$ be an eigenvalue of $T$ (it exists because we work over $\mathbb{C}$). Then $T - \lambda \cdot \mathrm{Id}$ is a morphism of representations with nontrivial kernel. By Part 1, $T - \lambda \cdot \mathrm{Id} = 0$.
\end{proof}

Schur's lemma severely restricts the form of equivariant linear maps between irreducible representations: between distinct types, every equivariant map is necessarily zero; between copies of the same type, it must be proportional to the identity. For \gls{GDL}, this implies that the learnable parameters in equivariant layers are strongly restricted by the group symmetry, drastically reducing the parameter space compared to unconstrained architectures \citep{cohen2016steerablecnns}. This algebraic restriction is precisely the reason for the parameter efficiency exhibited by equivariant networks.

\begin{definition}[Matrix Coefficients]\label{def:coeficientes_matriciales}
Let $(\rho, V)$ be a finite-dimensional unitary representation of a compact group $G$, with orthonormal basis $\{e_1, \ldots, e_d\}$ where $d = \dim(V)$. The \textit{matrix coefficients} of $\rho$ are the functions $\rho_{ij}: G \to \mathbb{C}$ defined by:
\begin{equation}
\rho_{ij}(g) = \langle e_i, \rho(g) e_j \rangle
\end{equation}
Equivalently, $\rho_{ij}(g)$ is the $(i,j)$ entry of the matrix that represents $\rho(g)$ in the basis $\{e_i\}$.
\end{definition}

The matrix coefficients are the ``coordinate functions'' of the representations. For $SO(2)$, the coefficients of $\rho_n(R_\theta) = e^{in\theta}$ are the functions $e^{in\theta}$, which form the Fourier basis. In general, the matrix coefficients play the role of ``basis functions'' for the generalized Fourier transform over the group, analogous to the complex exponentials $e^{in\theta}$ in classical Fourier analysis. In spectral methods of \gls{GDL}, equivariant convolutions over the group become pointwise multiplications in this spectral basis \citep{kondor2018generalizationconvolutionalneural}, enabling computationally efficient implementations.

\begin{proposition}[Schur Orthogonality Relations]\label{prop:ortogonalidad_schur}
Let $(\rho, V)$ and $(\sigma, W)$ be irreducible unitary representations of a compact group $G$ with normalized Haar measure $\mu$. Then:
\begin{equation}\label{eq:ortogonalidad_schur}
\int_G \overline{\rho_{ij}(g)} \sigma_{kl}(g) \, d\mu(g) = \frac{\delta_{\rho\sigma} \delta_{ik} \delta_{jl}}{d_\rho}
\end{equation}
where $d_\rho = \dim(V)$ and $\delta_{\rho\sigma} = 1$ if $\rho \cong \sigma$, and $0$ otherwise.
\end{proposition}

\begin{proof}
Define the operator $T: W \to V$ by:
\[
T = \int_G \rho(g) \circ P \circ \sigma(g^{-1}) \, d\mu(g)
\]
where $P: W \to V$ is an arbitrary linear operator. By the invariance of the Haar measure, $T$ is a morphism of representations: $\rho(h) T = T \sigma(h)$ for all $h \in G$.

By Schur's Lemma~\ref{lem:schur}: if $\rho \not\cong \sigma$, then $T = 0$; if $\rho = \sigma$, then $T = \lambda(P) \cdot \mathrm{Id}$.

Choosing $P$ as the operator $|e_l\rangle\langle e_j|$ and computing the trace, we obtain $\lambda = \delta_{jl}/d_\rho$, which gives the result.
\end{proof}

The orthogonality relations guarantee that the spectral components associated with distinct irreducible representations are perfectly separated, with no cross-interference. In equivariant architectures, this independence allows different ``frequency types'' to be processed in parallel, analogously to band filtering in classical signal processing.

\begin{corollary}\label{cor:base_ortonormal_coeficientes}
The matrix coefficients $\{\sqrt{d_\lambda} \rho^\lambda_{ij}\}$, where $\lambda$ ranges over $\widehat{G}$ and $1 \leq i,j \leq d_\lambda$, form an orthonormal system in $L^2(G)$.
\end{corollary}

\begin{proof}
A direct consequence of the Schur orthogonality relations. By the previous proposition:
\[
\langle \rho^\lambda_{ij}, \rho^\mu_{kl} \rangle_{L^2(G)} = \frac{1}{d_\lambda} \delta_{\lambda\mu} \delta_{ik} \delta_{jl}
\]
Normalizing by $\sqrt{d_\lambda}$, we obtain:
\[
\langle \sqrt{d_\lambda}\, \rho^\lambda_{ij},\; \sqrt{d_\mu}\, \rho^\mu_{kl} \rangle = \sqrt{d_\lambda d_\mu} \cdot \frac{1}{d_\lambda} \delta_{\lambda\mu} \delta_{ik} \delta_{jl} = \delta_{\lambda\mu} \delta_{ik} \delta_{jl}
\]
which proves the orthonormality of the system.
\end{proof}

The previous results establish that the normalized matrix coefficients form an orthonormal system in $L^2(G)$. The following theorem completes the picture by showing that this system is moreover \textit{complete}: every square-integrable function over the group can be recovered from its spectral components, ensuring that the generalized Fourier decomposition (which extends the classical transform, \autoref{def:transformada_fourier_continua}) loses no information.

\begin{theorem}[Peter-Weyl]\label{thm:peter_weyl}
Let $G$ be a compact Lie group. The space $L^2(G)$ of square-integrable functions over $G$ decomposes as:
\begin{equation}
L^2(G) = \bigoplus_{\lambda \in \widehat{G}} d_\lambda \cdot V_\lambda
\end{equation}
where $\widehat{G}$ is the set of equivalence classes of irreducible representations, $V_\lambda$ is the space of the representation $\lambda$, and $d_\lambda = \dim(V_\lambda)$ is its dimension.
\end{theorem}

\begin{proof}
The proof proceeds in two steps.

\textbf{Step 1: Orthogonality.} By \autoref{cor:base_ortonormal_coeficientes}, the normalized matrix coefficients $\{\sqrt{d_\lambda} \rho^\lambda_{ij}\}$ form an orthonormal system in $L^2(G)$.

\textbf{Step 2: Completeness.} Let $f \in L^2(G)$ be orthogonal to all matrix coefficients. Define the operator $T_f: V_\lambda \to V_\lambda$ by:
\[
T_f = \int_G f(g) \rho^\lambda(g) \, d\mu(g)
\]
For each $v, w \in V_\lambda$:
\[
\langle w, T_f v \rangle = \int_G f(g) \langle w, \rho^\lambda(g) v \rangle \, d\mu(g) = \int_G f(g) \rho^\lambda_{wv}(g) \, d\mu(g) = 0
\]
where the last equality uses that $f$ is orthogonal to all matrix coefficients. Since this holds for all $v, w$, we have $T_f = 0$ for every representation $\lambda$.

The key step is to show that $T_f = 0$ for every $\lambda$ implies $f = 0$. This follows from the fact that the representations of $G$ separate points: if $g \neq h$, there exists $\lambda$ such that $\rho^\lambda(g) \neq \rho^\lambda(h)$. Combining this with the Stone-Weierstrass theorem for compact groups, the matrix coefficients are dense in $C(G)$, and therefore in $L^2(G)$.

Thus, the orthonormal system of matrix coefficients is complete, which establishes the decomposition.
\end{proof}

The Peter-Weyl theorem is the generalization of the Fourier series to non-abelian groups. For $G = SO(2)$, we recover the classical Fourier decomposition: the irreducible representations are $\rho_n(R_\theta) = e^{in\theta}$ for $n \in \mathbb{Z}$, and the decomposition becomes the Fourier series $f(\theta) = \sum_{n \in \mathbb{Z}} c_n e^{in\theta}$ (for a detailed development of Fourier series and the Fourier transform, see Section~\ref{subsec:fundamentos_fourier} of \autoref{ch:cnn}).

\textbf{Implications for GDL.} The Peter-Weyl theorem has direct consequences for the design of equivariant neural architectures:
\begin{itemize}
    \item \textbf{Channel structure}: Each irreducible representation $V_\lambda$ corresponds to a ``type'' of channel in G-CNNs.
    \item \textbf{Determined dimensions}: The dimensions $d_\lambda$ of the layers are not arbitrary, but are fixed by the structure of the group.
    \item \textbf{Convolutions in frequency}: Equivariant convolutions are implemented efficiently in the generalized Fourier domain, multiplying coefficients in each $V_\lambda$.
\end{itemize}

\begin{corollary}[Channel structure in G-CNNs]\label{cor:peter_weyl_gcnn}
	Let $G$ be a compact Lie group and $\Phi: L^2(G) \to L^2(G)$ a linear $G$-equivariant operator. By the Peter-Weyl theorem, $\Phi$ preserves the decomposition into irreducibles:
	\[
		\Phi = \bigoplus_{\lambda \in \widehat{G}} \Phi_\lambda, \quad \text{where } \Phi_\lambda: d_\lambda V_\lambda \to d_\lambda V_\lambda
	\]
	Each $\Phi_\lambda$ acts as $\Phi_\lambda = I_{d_\lambda} \otimes W_\lambda$ for some matrix $W_\lambda \in \mathbb{R}^{d_\lambda \times d_\lambda}$.

	Consequences for $G$-equivariant architectures (developed in \autoref{ch:cnn}):
	\begin{enumerate}
		\item The ``channels'' of an equivariant layer are indexed by irreducible representations $\lambda \in \widehat{G}$
		\item The number of learnable parameters per channel is $d_\lambda^2$, independent of the input resolution
		\item Equivariance is guaranteed automatically by working in the spectral domain
	\end{enumerate}
\end{corollary}

\begin{proof}
By \autoref{thm:peter_weyl}, $L^2(G) = \bigoplus_{\lambda \in \widehat{G}} d_\lambda V_\lambda$. Let $\Phi: L^2(G) \to L^2(G)$ be a linear $G$-equivariant operator, that is, $\Phi \circ L_g = L_g \circ \Phi$ for all $g \in G$, where $L_g$ is the left regular representation. Since $\Phi$ commutes with the action of $G$, it must preserve each isotypic summand $d_\lambda V_\lambda$, since these are the irreducible $G$-invariant subspaces of $L^2(G)$. Thus, $\Phi = \bigoplus_\lambda \Phi_\lambda$ with $\Phi_\lambda: d_\lambda V_\lambda \to d_\lambda V_\lambda$.

Within the summand $d_\lambda V_\lambda \cong V_\lambda \otimes \mathbb{R}^{d_\lambda}$, the action of $G$ is $\rho^\lambda \otimes I_{d_\lambda}$. By Schur's Lemma~\ref{lem:schur}, every $G$-equivariant endomorphism of $V_\lambda \otimes \mathbb{R}^{d_\lambda}$ has the form $I_{d_\lambda} \otimes W_\lambda$ for some matrix $W_\lambda \in \mathbb{R}^{d_\lambda \times d_\lambda}$, which establishes the channel structure described.
\end{proof}

\begin{example}[Fourier Decomposition in $SO(2)$: Connection with classical CNNs]\label{ex:peter_weyl_so2_detallado}
	For $G = SO(2) \cong S^1$, the abelian group of rotations in the plane, the Peter-Weyl decomposition reduces to the classical Fourier series.

	\textbf{Irreducible representations:} All irreducibles of $SO(2)$ are one-dimensional ($d_n = 1$ for all $n$), given by:
	\[
		\rho_n: SO(2) \to \mathbb{C}^*, \quad \rho_n(R_\theta) = e^{in\theta}, \quad n \in \mathbb{Z}
	\]

	\textbf{Decomposition:} Every function $f \in L^2(SO(2)) \cong L^2([0, 2\pi])$ decomposes as:
	\[
		f(\theta) = \sum_{n \in \mathbb{Z}} \hat{f}_n e^{in\theta}, \quad \hat{f}_n = \frac{1}{2\pi} \int_0^{2\pi} f(\theta) e^{-in\theta} d\theta
	\]

	\textbf{Connection with cyclic CNNs:} For discrete signals over $\mathbb{Z}_N$ (the discrete version of $SO(2)$), the circular convolution is computed in the Fourier domain via the DFT:
	\[
		(f \star g)[k] = \text{IDFT}(\text{DFT}(f) \odot \text{DFT}(g))
	\]
	where $\odot$ is the elementwise product. This is exactly the efficient implementation of cyclic CNNs using the FFT, a special case of the Convolution Theorem in the Fourier Domain (\autoref{thm:convolucion_fourier} of \autoref{ch:cnn}).
\end{example}

\subsubsection{Fundamental Lie Groups}\label{subsubsec:grupos_lie_fundamentales}

We now develop the Lie groups most relevant to \gls{GDL}, presenting their formal definitions, explicit parameterizations, and associated Lie algebras. For each group, we derive its dimension as an application of a general tool that we present first.

\begin{remark}[Dimension of a Lie group]\label{rem:dimension_grupo_lie}
The \textbf{dimension} of a Lie group $G$ is its dimension as a differentiable manifold (\autoref{def:variedad_diferenciable}), which coincides with the dimension of its Lie algebra: $\dim G = \dim \mathfrak{g} = \dim T_eG$.

This dimension has a direct interpretation in the context of \gls{GDL}, in connection with the manifold hypothesis (\autoref{subsec:hipotesis_variedad}): the dimension of the symmetry group quantifies the degrees of freedom of the transformations that preserve the structure of the data. If a signal in $\mathbb{R}^d$ is invariant under the action of a group $G$ of dimension $k$, the signal effectively resides on a manifold of dimension $\leq d - k$, since $k$ degrees of freedom are fixed by the symmetry. The more symmetries the data respect, the greater the dimensional reduction.
\end{remark}

\begin{proposition}[Dimension of matrix subgroups]\label{prop:dimension_subgrupos_matriciales}
Let $F: \mathbb{R}^{n \times n} \to \mathbb{R}^k$ be a smooth map and $c \in \mathbb{R}^k$ a regular value of $F$ (see \autoref{def:valor_regular} of \autoref{ch:teoremas_complementarios}). Then $F^{-1}(c)$ is a differentiable submanifold of $\mathbb{R}^{n \times n}$ of dimension $n^2 - k$.
\end{proposition}

\begin{proof}
This is a direct application of the Regular Level Set Theorem (\autoref{thm:nivel_regular} of \autoref{ch:teoremas_complementarios}; see also \citep[Chapter~5]{lee2012introduction}). Since $c$ is a regular value, the differential $dF_A$ has rank $k$ at every point $A \in F^{-1}(c)$, so $F^{-1}(c)$ is a smooth submanifold of dimension $n^2 - k$. The same argument applies over $\mathbb{C}^{n \times n} \cong \mathbb{R}^{2n^2}$ for unitary groups.
\end{proof}

\paragraph{Special Orthogonal Group $SO(n)$}

\begin{definition}[Special Orthogonal Group]\label{def:grupo_so_n}
The \textit{special orthogonal group} $SO(n)$ is the group of orthogonal matrices with determinant $1$:
\begin{equation}
SO(n) = \{R \in \mathbb{R}^{n \times n} : R^T R = I_n, \det(R) = 1\}
\end{equation}
with the operation of matrix multiplication. It is a compact Lie group of dimension $\frac{n(n-1)}{2}$.
\end{definition}

We derive this dimension as an application of \autoref{prop:dimension_subgrupos_matriciales}. Consider the map $F: \mathbb{R}^{n \times n} \to \operatorname{Sym}(n)$ defined by $F(R) = R^T R - I_n$, where $\operatorname{Sym}(n)$ denotes the space of symmetric $n \times n$ matrices, of dimension $\frac{n(n+1)}{2}$ ($n$ diagonal entries and $\binom{n}{2}$ entries above the diagonal). Then $O(n) = F^{-1}(0)$.

To verify that $0$ is a regular value, we compute the differential: $dF_R(H) = H^T R + R^T H$. Given any $S \in \operatorname{Sym}(n)$, the choice $H = \frac{1}{2}RS$ satisfies
\[
	dF_R(H) = \tfrac{1}{2}S^T R^T R + \tfrac{1}{2}R^T R S = \tfrac{1}{2}S + \tfrac{1}{2}S = S
\]
using $R^T R = I_n$ and $S^T = S$, which proves the surjectivity of $dF_R$. Therefore:
\[
	\dim O(n) = n^2 - \frac{n(n+1)}{2} = \frac{n(n-1)}{2}
\]
The additional condition $\det(R) = 1$ selects the connected component of the identity within $O(n)$: since $\det: O(n) \to \{-1, +1\}$ is a discrete function, it imposes no differential restriction, and $\dim SO(n) = \dim O(n) = \frac{n(n-1)}{2}$. As a check, the Lie algebra $\mathfrak{so}(n) = \{A \in \mathbb{R}^{n \times n} : A^T + A = 0\}$ of antisymmetric matrices has exactly $\binom{n}{2} = \frac{n(n-1)}{2}$ free entries (those located strictly above the diagonal).

\begin{definition}[Orthogonal Group]\label{def:grupo_o_n}
The \textit{orthogonal group} $O(n)$ is the group of all orthogonal matrices:
\begin{equation}
O(n) = \{R \in \mathbb{R}^{n \times n} : R^T R = I_n\}
\end{equation}
It includes both rotations ($\det R = +1$) and reflections ($\det R = -1$), so that $SO(n) \subset O(n)$ is the index-2 subgroup formed by the orientation-preserving transformations. It is a compact Lie group of dimension $\frac{n(n-1)}{2}$.
\end{definition}

\begin{example}[$SO(2)$: Rotations in the Plane]\label{ex:so2_rotaciones}
The group $SO(2)$ parameterizes the rotations in $\mathbb{R}^2$. Each element is represented by an angle $\theta \in [0, 2\pi)$:
\begin{equation}
R_\theta = \begin{pmatrix}
\cos\theta & -\sin\theta \\
\sin\theta & \cos\theta
\end{pmatrix}
\end{equation}
For example, a rotation of $90^\circ$ corresponds to:
\begin{equation}
R_{\pi/2} = \begin{pmatrix}
0 & -1 \\
1 & 0
\end{pmatrix}
\end{equation}
Topologically, $SO(2) \cong S^1$ (the circle). Its Lie algebra $\mathfrak{so}(2)$ is one-dimensional, generated by:
\begin{equation}
J = \begin{pmatrix}
0 & -1 \\
1 & 0
\end{pmatrix}, \quad \text{with } e^{\theta J} = R_\theta
\end{equation}
\end{example}

\begin{figure}[htbp]
\centering
\includegraphics[width=0.6\textwidth]{figures/so2_rotations.pdf}
\caption[Rotations in the plane: $SO(2)$]{The group $SO(2)$ of rotations in the plane. Each element $R_\theta$ rotates vectors by an angle $\theta$. Topologically, $SO(2)$ is isomorphic to the circle $S^1$.}
\label{fig:so2_rotations}
\end{figure}

\begin{example}[$SO(3)$: Rotations in Space]\label{ex:so3_rotaciones}
The group $SO(3)$ describes the rotations in $\mathbb{R}^3$, fundamental for processing 3D data. The elementary rotations about the coordinate axes are:
\begin{align}
R_x(\theta) &= \begin{pmatrix}
1 & 0 & 0 \\
0 & \cos\theta & -\sin\theta \\
0 & \sin\theta & \cos\theta
\end{pmatrix}, \quad
R_y(\theta) = \begin{pmatrix}
\cos\theta & 0 & \sin\theta \\
0 & 1 & 0 \\
-\sin\theta & 0 & \cos\theta
\end{pmatrix}, \\
R_z(\theta) &= \begin{pmatrix}
\cos\theta & -\sin\theta & 0 \\
\sin\theta & \cos\theta & 0 \\
0 & 0 & 1
\end{pmatrix}
\end{align}
The Lie algebra $\mathfrak{so}(3)$ consists of antisymmetric $3 \times 3$ matrices, with basis:
\begin{equation}
L_x = \begin{pmatrix}
0 & 0 & 0 \\
0 & 0 & -1 \\
0 & 1 & 0
\end{pmatrix}, \quad
L_y = \begin{pmatrix}
0 & 0 & 1 \\
0 & 0 & 0 \\
-1 & 0 & 0
\end{pmatrix}, \quad
L_z = \begin{pmatrix}
0 & -1 & 0 \\
1 & 0 & 0 \\
0 & 0 & 0
\end{pmatrix}
\end{equation}
The Lie bracket satisfies $[L_i, L_j] = \epsilon_{ijk} L_k$, where $\epsilon_{ijk}$ is the Levi-Civita symbol.
\end{example}

\begin{example}[Exponential map in $SO(3)$: Rodrigues' Formula]\label{ex:mapa_exp_so3}
	The exponential map $\exp: \mathfrak{so}(3) \to SO(3)$ has a closed form known as \textbf{Rodrigues' formula}. For a vector $\boldsymbol{\omega} = (\omega_1, \omega_2, \omega_3)^T \in \mathbb{R}^3$ with $\theta = \|\boldsymbol{\omega}\|$, we define the antisymmetric matrix:
	\[
		[\boldsymbol{\omega}]_\times = \begin{pmatrix}
			0 & -\omega_3 & \omega_2 \\
			\omega_3 & 0 & -\omega_1 \\
			-\omega_2 & \omega_1 & 0
		\end{pmatrix} \in \mathfrak{so}(3)
	\]
	Then the exponential map is:
	\begin{equation}\label{eq:rodrigues}
		\exp([\boldsymbol{\omega}]_\times) = I + \frac{\sin\theta}{\theta} [\boldsymbol{\omega}]_\times + \frac{1 - \cos\theta}{\theta^2} [\boldsymbol{\omega}]_\times^2
	\end{equation}

	\textbf{Numerical example:} Consider a rotation of $90^\circ$ about the $z$ axis, that is, $\boldsymbol{\omega} = (0, 0, \pi/2)^T$. Then $\theta = \pi/2$ and:
	\[
		[\boldsymbol{\omega}]_\times = \frac{\pi}{2} \begin{pmatrix}
			0 & -1 & 0 \\
			1 & 0 & 0 \\
			0 & 0 & 0
		\end{pmatrix}
	\]
	Applying Rodrigues' formula with $\sin(\pi/2) = 1$ and $\cos(\pi/2) = 0$:
	\begin{align*}
		\exp([\boldsymbol{\omega}]_\times) &= I + \frac{2}{\pi} \cdot \frac{\pi}{2} \begin{pmatrix}
			0 & -1 & 0 \\
			1 & 0 & 0 \\
			0 & 0 & 0
		\end{pmatrix} + \frac{4}{\pi^2} \cdot \frac{\pi^2}{4} \begin{pmatrix}
			-1 & 0 & 0 \\
			0 & -1 & 0 \\
			0 & 0 & 0
		\end{pmatrix} \\
		&= \begin{pmatrix}
			0 & -1 & 0 \\
			1 & 0 & 0 \\
			0 & 0 & 1
		\end{pmatrix} = R_z(\pi/2)
	\end{align*}

	\textbf{Relevance for GDL:} Rodrigues' formula is computationally efficient and avoids the singularity of Euler angles. It is fundamental for:
	\begin{itemize}
		\item $SE(3)$-equivariant networks that process 3D point clouds
		\item Smooth interpolation of rotations in animation and robotics
		\item Parameterization of the space of rotations in optimization
	\end{itemize}
\end{example}

\begin{figure}[htbp]
\centering
\includegraphics[width=0.75\textwidth]{figures/so3_rotations.pdf}
\caption[Rotations in space: $SO(3)$]{The group $SO(3)$ of rotations in three-dimensional space. The elementary rotation matrices $R_x$, $R_y$ and $R_z$ rotate about the coordinate axes. Topologically, $SO(3) \cong \mathbb{RP}^3$.}
\label{fig:so3_rotations}
\end{figure}

\paragraph{Special Euclidean Group $SE(n)$}

\begin{definition}[Special Euclidean Group]\label{def:grupo_se_n}
The \textit{special Euclidean group} $SE(n)$ is the group of rigid transformations in $\mathbb{R}^n$ (rotation + translation). It is defined as the semidirect product (\autoref{def:producto_semidirecto}):
\begin{equation}
SE(n) = \mathbb{R}^n \rtimes SO(n)
\end{equation}
where the group operation is $(t_1, R_1) \cdot (t_2, R_2) = (t_1 + R_1 t_2, R_1 R_2)$.
It is a Lie group of dimension $\frac{n(n-1)}{2} + n = \frac{n(n+1)}{2}$.
\end{definition}

\begin{definition}[Euclidean Group]\label{def:grupo_e_n}
The \textit{Euclidean group} $E(n)$ is the group of all isometries of $\mathbb{R}^n$ (rotations, reflections, and translations). It is defined as the semidirect product:
\begin{equation}
E(n) = \mathbb{R}^n \rtimes O(n)
\end{equation}
where $O(n)$ is the orthogonal group (\autoref{def:grupo_o_n}). Unlike $SE(n)$, which only includes rotations and translations, $E(n)$ also includes reflections. It is a Lie group of dimension $\frac{n(n+1)}{2}$.
\end{definition}

\begin{example}[$SE(3)$ in Homogeneous Coordinates]\label{ex:se3_homogeneas}
The transformations of $SE(3)$ are conveniently represented as $4 \times 4$ matrices in homogeneous coordinates:
\begin{equation}
T = \begin{pmatrix}
R & t \\
0 & 1
\end{pmatrix} \in \mathbb{R}^{4 \times 4}, \quad R \in SO(3), \, t \in \mathbb{R}^3
\end{equation}
The action on a point $\mathbf{x} \in \mathbb{R}^3$ is:
\begin{equation}
\begin{pmatrix}
R & t \\
0 & 1
\end{pmatrix}
\begin{pmatrix}
\mathbf{x} \\
1
\end{pmatrix}
= \begin{pmatrix}
R\mathbf{x} + t \\
1
\end{pmatrix}
\end{equation}
The Lie algebra $\mathfrak{se}(3)$ has dimension 6, with generators for rotations $(L_x, L_y, L_z)$ and translations $(P_x, P_y, P_z)$:
\begin{equation}
P_x = \begin{pmatrix}
0 & 0 & 0 & 1 \\
0 & 0 & 0 & 0 \\
0 & 0 & 0 & 0 \\
0 & 0 & 0 & 0
\end{pmatrix}, \quad
P_y = \begin{pmatrix}
0 & 0 & 0 & 0 \\
0 & 0 & 0 & 1 \\
0 & 0 & 0 & 0 \\
0 & 0 & 0 & 0
\end{pmatrix}, \quad
P_z = \begin{pmatrix}
0 & 0 & 0 & 0 \\
0 & 0 & 0 & 0 \\
0 & 0 & 0 & 1 \\
0 & 0 & 0 & 0
\end{pmatrix}
\end{equation}
In \gls{GDL}, $SE(3)$ is essential for 3D point clouds and molecular data, where the predictions must be invariant to rigid transformations.
\end{example}

\paragraph{Linear Groups $GL(n)$ and $SL(n)$}

\begin{definition}[General Linear Group]\label{def:grupo_gl_n}
The \textit{general linear group} $GL(n, \mathbb{R})$ is the group of invertible $n \times n$ matrices:
\begin{equation}
GL(n, \mathbb{R}) = \{A \in \mathbb{R}^{n \times n} : \det(A) \neq 0\}
\end{equation}
It is a non-compact Lie group of dimension $n^2$. Its Lie algebra $\mathfrak{gl}(n, \mathbb{R})$ is the space of all $n \times n$ matrices, with the bracket $[A, B] = AB - BA$.
\end{definition}

\begin{definition}[Special Linear Group]\label{def:grupo_sl_n}
The \textit{special linear group} $SL(n, \mathbb{R})$ is the subgroup of matrices with determinant $1$:
\begin{equation}
SL(n, \mathbb{R}) = \{A \in GL(n, \mathbb{R}) : \det(A) = 1\}
\end{equation}
It is a Lie group of dimension $n^2 - 1$. Its Lie algebra $\mathfrak{sl}(n, \mathbb{R})$ consists of trace-zero matrices.
\end{definition}

The dimensions of these groups are obtained by different methods. For $GL(n, \mathbb{R})$: since $\det: \mathbb{R}^{n \times n} \to \mathbb{R}$ is a polynomial in the entries (and therefore continuous), $GL(n, \mathbb{R}) = \det^{-1}(\mathbb{R} \setminus \{0\})$ is the preimage of an open set, hence it is an open subset of $\mathbb{R}^{n \times n}$. An open subset of $\mathbb{R}^{n^2}$ inherits the manifold structure with the same dimension $n^2$.

For $SL(n, \mathbb{R})$, we apply \autoref{prop:dimension_subgrupos_matriciales} with $F = \det$ and $c = 1$. The differential of the determinant is $d(\det)_A(H) = \det(A) \cdot \operatorname{tr}(A^{-1}H)$. For $A \in SL(n)$, $\det(A) = 1$, and the map $H \mapsto \operatorname{tr}(A^{-1}H)$ is surjective onto $\mathbb{R}$: it suffices to take $H = tA$ to obtain $\operatorname{tr}(I_n) \cdot t = nt$. Thus $1$ is a regular value and $\dim SL(n) = n^2 - 1$. As a check, $\mathfrak{sl}(n, \mathbb{R}) = \{A : \operatorname{tr}(A) = 0\}$ imposes a linear restriction on $\mathfrak{gl}(n, \mathbb{R}) \cong \mathbb{R}^{n^2}$, leaving dimension $n^2 - 1$.

\paragraph{Unitary Groups $U(n)$ and $SU(n)$}

\begin{definition}[Unitary Group]\label{def:grupo_u_n}
The \textit{unitary group} $U(n)$ is the group of unitary $n \times n$ matrices:
\begin{equation}
U(n) = \{U \in \mathbb{C}^{n \times n} : U^\dagger U = I_n\}
\end{equation}
where $U^\dagger = \overline{U}^T$ denotes the conjugate transpose. It is a compact Lie group of dimension $n^2$.
\end{definition}

\begin{definition}[Special Unitary Group]\label{def:grupo_su_n}
The \textit{special unitary group} $SU(n)$ is the subgroup of unitary matrices with determinant $1$:
\begin{equation}
SU(n) = \{U \in U(n) : \det(U) = 1\}
\end{equation}
It is a compact Lie group of dimension $n^2 - 1$.
\end{definition}

We derive the dimensions of $U(n)$ and $SU(n)$. For $U(n)$, we define $F: \mathbb{C}^{n \times n} \to \operatorname{Herm}(n)$ by $F(U) = U^\dagger U - I_n$, where $\operatorname{Herm}(n)$ denotes the space of Hermitian $n \times n$ matrices. This space has real dimension $n^2$: the $n$ diagonal entries are real, and the $\binom{n}{2}$ entries above the diagonal are complex ($n(n-1)$ real parameters), totaling $n + n(n-1) = n^2$. The ambient space $\mathbb{C}^{n \times n}$ has real dimension $2n^2$. The differential is $dF_U(H) = H^\dagger U + U^\dagger H$; for any $S \in \operatorname{Herm}(n)$, the choice $H = \frac{1}{2}US$ satisfies $dF_U(H) = S$ (analogously to the orthogonal case, using $U^\dagger U = I_n$ and $S^\dagger = S$). By \autoref{prop:dimension_subgrupos_matriciales}:
\[
	\dim U(n) = 2n^2 - n^2 = n^2
\]

For $SU(n)$, the additional restriction $\det(U) = 1$ imposes a smooth condition on $U(n)$. For unitary $U$, $|\det(U)| = 1$, so $\det: U(n) \to S^1$. The differential $d(\det)_U$ restricted to $T_U U(n)$ is surjective onto $T_{\det(U)}S^1 \cong i\mathbb{R}$: for $X \in \mathfrak{u}(n)$, $d(\det)_I(X) = \operatorname{tr}(X)$, which is surjective (it suffices to take $X = itI_n$). Since $S^1$ has real dimension $1$, the condition $\det = 1$ defines a submanifold of codimension $1$, hence $\dim SU(n) = n^2 - 1$.

\begin{example}[$SU(2)$ and its Relation with $SO(3)$]\label{ex:su2_so3}
The group $SU(2)$ consists of $2 \times 2$ matrices of the form:
\begin{equation}
U = \begin{pmatrix}
\alpha & -\overline{\beta} \\
\beta & \overline{\alpha}
\end{pmatrix}, \quad |\alpha|^2 + |\beta|^2 = 1
\end{equation}
Topologically, $SU(2) \cong S^3$ (the 3-sphere). We will explicitly construct a surjective homomorphism $\pi: SU(2) \to SO(3)$ and prove that $SO(3) \cong SU(2)/\mathbb{Z}_2$.

\textbf{The Lie algebra $\mathfrak{su}(2)$ and its inner product.}
The Lie algebra $\mathfrak{su}(2)$ consists of the trace-zero anti-Hermitian matrices:
\[
	\mathfrak{su}(2) = \{X \in M_2(\mathbb{C}) : X^\dagger = -X, \; \operatorname{tr}(X) = 0\}
\]
The \textit{Pauli matrices} provide a natural basis:
\begin{equation}
\sigma_x = \begin{pmatrix}
0 & 1 \\
1 & 0
\end{pmatrix}, \quad
\sigma_y = \begin{pmatrix}
0 & -i \\
i & 0
\end{pmatrix}, \quad
\sigma_z = \begin{pmatrix}
1 & 0 \\
0 & -1
\end{pmatrix}
\end{equation}
The generators $\tau_k = \frac{i}{2}\sigma_k$ for $k \in \{x, y, z\}$ form a basis of $\mathfrak{su}(2)$ and satisfy $[\tau_i, \tau_j] = \epsilon_{ijk}\tau_k$. We endow $\mathfrak{su}(2)$ with the inner product:
\begin{equation}\label{eq:inner_product_su2}
	\langle X, Y \rangle = -2\operatorname{tr}(XY)
\end{equation}
To verify that $\{\tau_1, \tau_2, \tau_3\}$ is an orthonormal basis, we compute $\langle \tau_j, \tau_k \rangle = -2\operatorname{tr}(\tau_j \tau_k) = \frac{1}{2}\operatorname{tr}(\sigma_j \sigma_k) = \delta_{jk}$, where the last equality uses the Pauli identity $\sigma_j \sigma_k = \delta_{jk} I + i\epsilon_{jkl}\sigma_l$ and $\operatorname{tr}(\sigma_l) = 0$.

\textbf{Construction of $\pi$ via the adjoint action.}
For each $U \in SU(2)$, we define the adjoint representation $\operatorname{Ad}_U: \mathfrak{su}(2) \to \mathfrak{su}(2)$ by:
\begin{equation}\label{eq:adjoint_su2}
	\operatorname{Ad}_U(X) = UXU^\dagger
\end{equation}
Let us see that $\operatorname{Ad}_U$ preserves $\mathfrak{su}(2)$: if $X^\dagger = -X$ and $\operatorname{tr}(X) = 0$, then $(UXU^\dagger)^\dagger = UX^\dagger U^\dagger = -UXU^\dagger$ and $\operatorname{tr}(UXU^\dagger) = \operatorname{tr}(X) = 0$. Moreover, $\operatorname{Ad}_U$ preserves the inner product:
\[
	\langle UXU^\dagger, UYU^\dagger \rangle = -2\operatorname{tr}(UXU^\dagger \cdot UYU^\dagger) = -2\operatorname{tr}(UXYU^\dagger) = -2\operatorname{tr}(XY) = \langle X, Y \rangle
\]
We define $\pi(U)$ as the matrix of $\operatorname{Ad}_U$ in the orthonormal basis $\{\tau_1, \tau_2, \tau_3\}$, that is, $[\pi(U)]_{jk} = \langle \tau_j, \operatorname{Ad}_U(\tau_k) \rangle$. Since $\operatorname{Ad}_U$ preserves the inner product, $\pi(U) \in O(3)$. The group $SU(2)$ is connected (it is homeomorphic to $S^3$), $\pi$ is continuous and $\pi(I) = I_3 \in SO(3)$. Since $\det \circ\, \pi: SU(2) \to \{-1, +1\}$ is continuous on a connected space, it must be constant, hence $\det(\pi(U)) = 1$ for all $U$. We conclude that $\pi: SU(2) \to SO(3)$ is a group homomorphism, since $\operatorname{Ad}_{UV} = \operatorname{Ad}_U \circ \operatorname{Ad}_V$.

\textbf{Computation of the kernel.}
We determine $\ker(\pi) = \{U \in SU(2) : UXU^\dagger = X \text{ for all } X \in \mathfrak{su}(2)\}$. If $U \in \ker(\pi)$, then $U$ commutes with $\tau_1, \tau_2, \tau_3$, and therefore with $\sigma_x, \sigma_y, \sigma_z$. Since $\{I, \sigma_x, \sigma_y, \sigma_z\}$ is a basis of $M_2(\mathbb{C})$, the matrix $U$ commutes with every $2 \times 2$ matrix, hence $U = \lambda I$ for some $\lambda \in \mathbb{C}$. The unitarity ($|\lambda|^2 = 1$) and determinant ($\lambda^2 = 1$) conditions imply $\lambda = \pm 1$. Therefore:
\begin{equation}\label{eq:kernel_pi}
	\ker(\pi) = \{I, -I\} \cong \mathbb{Z}_2
\end{equation}

\textbf{Surjectivity of $\pi$.}
The kernel $\{I, -I\}$ is discrete, so that $\operatorname{Lie}(\ker \pi) = \{0\}$. Since $\ker(d\pi_I) = \operatorname{Lie}(\ker \pi)$ for Lie group homomorphisms, the differential $d\pi_I: \mathfrak{su}(2) \to \mathfrak{so}(3)$ is injective. Since $\dim \mathfrak{su}(2) = \dim \mathfrak{so}(3) = 3$, the differential $d\pi_I$ is a linear isomorphism. By the inverse function theorem (\autoref{thm:funcion_inversa}), $\pi$ is a local diffeomorphism, so $\operatorname{Im}(\pi)$ is open in $SO(3)$. Since $SU(2)$ is compact, $\operatorname{Im}(\pi)$ is compact and therefore closed in $SO(3)$ (which is Hausdorff). Since $SO(3)$ is connected and $\operatorname{Im}(\pi)$ is a nonempty open and closed subset, we conclude that $\operatorname{Im}(\pi) = SO(3)$.

\textbf{Quotient isomorphism.}
By the First Isomorphism Theorem (\autoref{thm:primer_isomorfismo}), since $\pi: SU(2) \to SO(3)$ is a surjective homomorphism:
\begin{equation}\label{eq:su2_z2_so3}
	SO(3) \cong SU(2)/\ker(\pi) = SU(2)/\{I, -I\} = SU(2)/\mathbb{Z}_2
\end{equation}
Thus, $SU(2)$ is the \textbf{double cover} of $SO(3)$.

\textbf{Isomorphism of Lie algebras.}
Since $d\pi_I: \mathfrak{su}(2) \to \mathfrak{so}(3)$ is a linear isomorphism and the differential of a Lie group homomorphism preserves the bracket ($d\pi_I([X,Y]) = [d\pi_I(X), d\pi_I(Y)]$), we obtain $\mathfrak{su}(2) \cong \mathfrak{so}(3)$ as Lie algebras. This explains why both admit bases with the same commutation relation $[e_i, e_j] = \epsilon_{ijk}e_k$ (compare with \autoref{ex:so3_rotaciones}).

\textbf{Relevance.}
This relation is fundamental in quantum physics and in neural networks that use quaternions, since the unit quaternions are isomorphic to $SU(2)$ (see Section~\ref{sec:cuaterniones} of Appendix~\ref{ch:teoremas_complementarios}) and provide a singularity-free parameterization of 3D rotations.
\end{example}

\subsection{Equivariance and Invariance}\label{subsec:equivarianza_invarianza}

The concepts of equivariance and invariance constitute the theoretical core of \gls{GDL}. Having developed the theory of Lie groups and representations, we are now in a position to formalize these fundamental concepts, which capture how functions can respect the symmetries of a problem.

\begin{definition}[Equivariant Function]\label{def:equivariance_gdl}
Let $\phi: X \rightarrow Y$ be a function, $G$ a group with actions $\rho_X: G \times X \rightarrow X$ and $\rho_Y: G \times Y \rightarrow Y$ on the spaces $X$ and $Y$ respectively. The function $\phi$ is \textit{$G$-equivariant} if:
\begin{equation}
\phi(g \cdot x) = g \cdot \phi(x)
\label{eq:equivariance_gdl_formal}
\end{equation}
for all $g \in G$ and $x \in X$, where $g \cdot x$ denotes $\rho_X(g, x)$ and $g \cdot \phi(x)$ denotes $\rho_Y(g, \phi(x))$.
\end{definition}

\begin{definition}[Invariant Function]\label{def:invariance_gdl}
A function $\phi: X \rightarrow Y$ is \textit{$G$-invariant} if:
\begin{equation}
\phi(g \cdot x) = \phi(x)
\label{eq:invariance_gdl_formal}
\end{equation}
for all $g \in G$ and $x \in X$. This is equivalent to $G$-equivariance when the action of $G$ on $Y$ is trivial (the identity).
\end{definition}

Invariance is a special case of equivariance in which the group acts trivially on the output space. In many \gls{GDL} applications, we seek equivariant intermediate layers that preserve structure, followed by a final invariant layer.

\begin{figure}[htbp]
\centering
\includegraphics[width=0.65\textwidth]{figures/equivariance_diagram.pdf}
\caption[Commutative diagram of equivariance]{Commutative diagram of equivariance. A function $\phi: X \to Y$ is $G$-equivariant if the diagram commutes: applying first the transformation $\rho_X(g)$ and then $\phi$ yields the same result as applying first $\phi$ and then $\rho_Y(g)$.}
\label{fig:equivariance_diagram}
\end{figure}

Traditional \glspl{ANN} lack equivariance by design. To approximate invariance, they require \textit{data augmentation} (\autoref{def:data_augmentation}) with all the transformations of the group $G$, increasing the size of the training set by a factor of $|G|$ with no guarantee of generalization. \gls{GDL} solves this problem by incorporating symmetries directly into the architecture, propagating equivariance through all layers.

\begin{theorem}[Composition of Equivariant Layers]\label{theorem:equivariant_composition}
If $f: X \rightarrow Y$ and $g: Y \rightarrow Z$ are equivariant functions with respect to the transformations $T_X$, $T_Y$ and $T_Z$ on the spaces $X$, $Y$ and $Z$ respectively, then their composition $g \circ f: X \rightarrow Z$ is also equivariant.
\end{theorem}

\begin{proof}
Let $h \in G$ be arbitrary and $x \in X$. By the $G$-equivariance of $f$:
\[
(g \circ f)(T_X(h) \cdot x) = g(f(T_X(h) \cdot x)) = g(T_Y(h) \cdot f(x))
\]
By the $G$-equivariance of $g$:
\[
= T_Z(h) \cdot g(f(x)) = T_Z(h) \cdot (g \circ f)(x)
\]
Therefore, $g \circ f$ is $G$-equivariant.
\end{proof}

This theorem guarantees that deep networks built by composing equivariant layers maintain equivariance end to end, a fundamental property for the design of geometric architectures.

\subsubsection{Equivariant Universal Approximation Theorem}\label{sec:teorema_equivariante}

The following result, proved by \citet{yarotsky2022} and \citet{kondor2018generalizationconvolutionalneural}, extends the classical universal approximation theorem to the equivariant case:

\begin{theorem}[Equivariant Universal Approximation]\label{theorem:universal_approximation_equivariant}
Let $G$ be a compact group (\autoref{def:grupo_compacto}) acting continuously on compact metric spaces $X$ and $Y$, and let $\mathcal{F}^G(X,Y)$ be the space of all continuous $G$-equivariant functions $f: X \rightarrow Y$. Then, for any function $f \in \mathcal{F}^G(X,Y)$ and any $\epsilon > 0$, there exists a $G$-equivariant neural network $\mathcal{N}$ such that:
\begin{equation}
\sup_{x \in X} d_Y(f(x), \mathcal{N}(x)) < \epsilon
\label{eq:universal_approximation_equivariant}
\end{equation}
where $d_Y$ is the metric on $Y$.
\end{theorem}

\begin{proof}
The proof proceeds in four steps, following the ideas of \citet{yarotsky2022}.

\textbf{Step 1: Non-equivariant approximation.}
By the classical universal approximation theorem (\autoref{thm:aproximacion_universal}), for any continuous function $f: X \to Y$ and $\epsilon > 0$, there exists a neural network $\tilde{f}: X \to Y$ such that $\sup_{x \in X} d_Y(f(x), \tilde{f}(x)) < \epsilon$.

\textbf{Step 2: Construction of the equivariant averaging.}
We define the \textit{Reynolds operator} by averaging over the group:
\[
\overline{f}(x) = \int_G \rho_Y(g^{-1}) \tilde{f}(g \cdot x) \, d\mu(g)
\]
where $\rho_Y: G \to GL(Y)$ is the representation of $G$ on $Y$ (\autoref{def:equivariance_gdl}), and $\mu$ is the normalized Haar measure on $G$ (\autoref{def:medida_haar_gdl}).

\textbf{Step 3: Verification of equivariance.}
We verify that $\overline{f}$ is $G$-equivariant. For any $h \in G$:
\begin{align*}
\overline{f}(h \cdot x) &= \int_G \rho_Y(g^{-1}) \tilde{f}(g \cdot (h \cdot x)) \, d\mu(g) \\
&= \int_G \rho_Y(g^{-1}) \tilde{f}((gh) \cdot x) \, d\mu(g)
\end{align*}
Performing the change of variables $g' = gh$, and using the invariance of the Haar measure ($d\mu(g) = d\mu(g')$):
\begin{align*}
&= \int_G \rho_Y((g'h^{-1})^{-1}) \tilde{f}(g' \cdot x) \, d\mu(g') \\
&= \int_G \rho_Y(h)\rho_Y(g'^{-1}) \tilde{f}(g' \cdot x) \, d\mu(g') \\
&= \rho_Y(h) \int_G \rho_Y(g'^{-1}) \tilde{f}(g' \cdot x) \, d\mu(g') = \rho_Y(h) \overline{f}(x)
\end{align*}
Therefore, $\overline{f}$ is $G$-equivariant according to \autoref{def:equivariance_gdl}.

\textbf{Step 4: Control of the approximation error.}
Let $f$ be the target $G$-equivariant function. By the linearity of the integral and the equivariance of $f$:
\begin{align*}
d_Y(\overline{\tilde{f}}(x), f(x)) &\leq \int_G d_Y(\rho_Y(g^{-1}) \tilde{f}(g \cdot x), \rho_Y(g^{-1}) f(g \cdot x)) \, d\mu(g)
\end{align*}
For compact groups, the representations can be chosen unitary, so that $\rho_Y(g^{-1})$ preserves distances:
\[
\leq \int_G d_Y(\tilde{f}(g \cdot x), f(g \cdot x)) \, d\mu(g) \leq \sup_{x' \in X} d_Y(\tilde{f}(x'), f(x')) < \epsilon
\]
The last inequality uses the fact that the action of $G$ on the compact set $X$ preserves compactness.
\end{proof}

\textbf{Implications for architecture design.} This theorem is crucial because it shows that equivariance constraints do not limit the expressive capacity of neural networks. Specifically:

\begin{itemize}
    \item $G$-equivariant networks can approximate any continuous $G$-equivariant function
    \item It is not necessary to sacrifice expressive power for equivariance
    \item Equivariance acts as an inductive bias that reduces the function space without losing approximation capacity in the relevant subspace
\end{itemize}

\begin{example}[CNNs as a special case]
For the group of translations $G = \mathbb{Z}^2$ acting on discrete images, this theorem guarantees that \glspl{CNN} (which are translation-equivariant, see \autoref{ch:cnn}) can approximate any continuous function that is translation-invariant or translation-equivariant. This provides a theoretical justification for the success of \glspl{CNN} in computer vision tasks.
\end{example}

\subsection{Homogeneous Spaces and Geometric Domains}\label{subsec:espacios_homogeneos}

So far we have studied equivariance as an algebraic property of functions between signal spaces. However, many important data domains in \gls{GDL} (the sphere $S^2$, projective space, Grassmann manifolds) share a common geometric structure: they can be expressed as quotients $G/H$ of a Lie group by a closed subgroup. This observation makes it possible to unify the design of equivariant architectures for seemingly distinct domains under a single theoretical framework.

\begin{definition}[Homogeneous Space]\label{def:espacio_homogeneo}
Let $G$ be a Lie group and $H$ a closed subgroup. The \textit{homogeneous space} $G/H$ is the set of left cosets:
\begin{equation}
G/H = \{gH : g \in G\}
\end{equation}
The group $G$ acts transitively on $G/H$ via $g' \cdot (gH) = (g'g)H$.
\end{definition}

\subsubsection{Theoretical foundations}

To formalize the quotient structure $G/H$, we need two concepts from the theory of group actions: the \textit{orbit} of an element (all the points reachable by the action) and its \textit{stabilizer} (the subgroup that leaves it fixed). The orbit-stabilizer theorem will connect both notions and justify why every homogeneous space admits a representation as a quotient.

\begin{definition}[Orbit]\label{def:orbita}
Let $G$ be a group acting on a set $X$. For any element $x \in X$, the \textit{orbit} of $x$ under the action of $G$ is:
\begin{equation}
\text{Orb}(x) = \{g \cdot x : g \in G\}
\end{equation}
That is, the set of all elements of $X$ that can be reached by applying elements of $G$ to $x$.
\end{definition}

\begin{definition}[Stabilizer Subgroup]\label{def:estabilizador}
Let $G$ be a group acting on a set $X$. For any element $x \in X$, the \textit{stabilizer subgroup} (or isotropy subgroup) of $x$ is:
\begin{equation}
\text{Stab}(x) = \{g \in G : g \cdot x = x\}
\end{equation}
That is, the subgroup of all elements of $G$ that leave $x$ invariant.
\end{definition}

The fundamental relation between these concepts is captured by the orbit-stabilizer theorem, which is central to understanding the structure of homogeneous spaces:

\begin{theorem}[Orbit-Stabilizer Theorem]\label{thm:orbita_estabilizador}
Let $G$ be a group acting on a set $X$, and let $x \in X$. Then there exists a natural isomorphism:
\begin{equation}
G/\text{Stab}(x) \cong \text{Orb}(x)
\end{equation}
\end{theorem}

\begin{proof}
We construct a direct bijection $\psi: G/\mathrm{Stab}(x) \to \mathrm{Orb}(x)$ defined by $\psi(g\,\mathrm{Stab}(x)) = g \cdot x$.

\textbf{Well-definedness.} If $g\,\mathrm{Stab}(x) = h\,\mathrm{Stab}(x)$, then $h^{-1}g \in \mathrm{Stab}(x)$, that is, $(h^{-1}g) \cdot x = x$. Hence $g \cdot x = h \cdot ((h^{-1}g) \cdot x) = h \cdot x$, and $\psi$ is well defined.

\textbf{Injectivity.} If $\psi(g\,\mathrm{Stab}(x)) = \psi(h\,\mathrm{Stab}(x))$, then $g \cdot x = h \cdot x$, from which $(h^{-1}g) \cdot x = x$ and $h^{-1}g \in \mathrm{Stab}(x)$. Therefore $g\,\mathrm{Stab}(x) = h\,\mathrm{Stab}(x)$.

\textbf{Surjectivity.} Every $y \in \mathrm{Orb}(x)$ is, by definition, of the form $y = g \cdot x$ for some $g \in G$, hence $y = \psi(g\,\mathrm{Stab}(x))$.
\end{proof}

This theorem is fundamental for \gls{GDL} because it establishes that any geometric domain admitting transitive symmetries can be represented as a homogeneous space $G/H$, where $G$ is the symmetry group and $H$ is the subgroup that fixes a reference point.

\begin{figure}[htbp]
\centering
\includegraphics[width=0.7\textwidth]{figures/orbit_stabilizer.pdf}
\caption[Orbit-Stabilizer Theorem]{Visualization of the Orbit-Stabilizer Theorem. The orbit $\mathrm{Orb}(x)$ under the action of $G$ is isomorphic to the quotient space $G/\mathrm{Stab}(x)$: each coset $g\,\mathrm{Stab}(x)$ corresponds to exactly one point of the orbit.}
\label{fig:orbit_stabilizer}
\end{figure}

\begin{example}[Application of the Theorem in \gls{GDL}]\label{ex:aplicacion_teorema_gdl}
Consider the sphere $S^2$ as a domain for directional data, with $G = SO(3)$ acting by rotations.

\textbf{Identification of the quotient.}
\begin{itemize}
    \item \emph{Transitivity}: $SO(3)$ acts transitively on $S^2$. Indeed, given $\mathbf{x}, \mathbf{y} \in S^2$, each can be completed to a positively oriented orthonormal basis of $\mathbb{R}^3$, and the rotation sending one basis to the other satisfies $R\mathbf{x} = \mathbf{y}$. Therefore $\mathrm{Orb}(\mathbf{p}) = S^2$ for any reference point $\mathbf{p}$.

    \item \emph{Stabilizer}: Take the north pole $\mathbf{p} = (0,0,1)^{\!\top}$. A rotation $R \in SO(3)$ fixes $\mathbf{p}$ if and only if its axis of rotation is the $z$-axis, that is, $R = R_z(\theta)$ for some $\theta \in [0, 2\pi)$ (see \autoref{ex:so3_rotaciones}). These rotations form a subgroup isomorphic to $SO(2)$, hence $\mathrm{Stab}(\mathbf{p}) \cong SO(2)$.

    \item \emph{Interpretation of the cosets}: Two rotations $g_1, g_2 \in SO(3)$ produce the same point on $S^2$ (that is, $g_1 \cdot \mathbf{p} = g_2 \cdot \mathbf{p}$) if and only if $g_1^{-1}g_2 \in SO(2)$, which means that they differ by a rotation about the axis passing through that point. Each coset $gSO(2)$ corresponds to exactly one point of the sphere.
\end{itemize}
Applying Theorem~\ref{thm:orbita_estabilizador} yields
\begin{equation}\label{eq:esfera_cociente_ejemplo}
    S^2 \cong SO(3)/SO(2).
\end{equation}
The rigorous proof for $S^{n-1}$ in arbitrary dimension is presented in Proposition~\ref{prop:esfera_cociente}.

\textbf{Interpretation for signal processing.}
A signal on the sphere $f\colon S^2 \to \mathbb{R}^C$ can be ``lifted'' to a function on $SO(3)$ that is right-invariant under $SO(2)$: it suffices to define $\tilde{f}(g) = f(g \cdot \mathbf{p})$, so that $\tilde{f}(gh) = \tilde{f}(g)$ for all $h \in SO(2)$. As a consequence, the equivariant convolution on $S^2$ must be defined as a correlation over the full group $SO(3)$, rather than as a convolution on a flat projection (which would introduce distortions). In this way, the first convolutional layer transforms signals on $S^2$ into functions on $SO(3)$, and the subsequent layers operate directly on $SO(3)$. The technical details of this construction are developed in \autoref{subsec:spherical_cnns_so3}.

\textbf{Applications.}
\begin{enumerate}
    \item \emph{Climate and geophysical data}: signals on the Earth's surface (temperature, pressure, wind) are functions on $S^2$. Equivariance under $SO(3)$ guarantees that predictions do not depend on the choice of geographic coordinates \citep{cohen2018sphericalcnns}.

    \item \emph{Omnidirectional vision}: $360^\circ$ images are naturally represented as signals on $S^2$. Processing these signals directly on the sphere avoids the distortions inherent in the equirectangular or cubic projections used by planar \glspl{CNN}.

    \item \emph{Molecular orientations}: molecular surfaces lack a preferred orientation. Using $SO(3)$-equivariant architectures guarantees physically meaningful representations, independent of the arbitrary orientation in which the molecule is presented.
\end{enumerate}

\textbf{General methodology.}
This example illustrates a systematic pattern for designing \gls{GDL} architectures: (1)~identify the domain $\Omega$ of the data; (2)~determine the relevant symmetry group $G$; (3)~compute the stabilizer $H$ of a reference point; (4)~recognize $\Omega \cong G/H$ via the Orbit-Stabilizer Theorem; and (5)~build the equivariant architecture using the harmonic analysis of $G$ (Peter-Weyl Theorem~\ref{thm:peter_weyl}). Propositions~\ref{prop:esfera_cociente}--\ref{prop:grassmann_stiefel_cociente} apply this methodology systematically to the fundamental domains of \gls{GDL}.
\end{example}

\paragraph{Fundamental geometric properties}

Homogeneous spaces possess distinctive geometric properties that make them especially well suited for Deep Learning:

\begin{proposition}[Transitivity and Connectedness]\label{prop:transitividad_espacios_homogeneos}
Let $G/H$ be a homogeneous space. Then:
\begin{enumerate}
    \item $G$ acts transitively on $G/H$: given $g_1H, g_2H \in G/H$, there exists $g \in G$ such that $g \cdot (g_1H) = g_2H$
    \item If $G$ is connected, then $G/H$ is connected
    \item The action of $G$ preserves any $G$-invariant metric on $G/H$
\end{enumerate}
\end{proposition}

\begin{proof}
\leavevmode
\begin{enumerate}
    \item \textbf{Transitivity.} Given $g_1H, g_2H \in G/H$, let $g = g_2 g_1^{-1} \in G$. Then:
    \[
    g \cdot (g_1H) = (g_2 g_1^{-1}) g_1 H = g_2 H
    \]

    \item \textbf{Connectedness.} The canonical projection $\pi: G \to G/H$, $\pi(g) = gH$, is continuous by definition of the quotient topology and surjective by definition. Since $G$ is connected, $G/H = \pi(G)$ is connected (the continuous image of a connected space is connected).

    \item \textbf{Preservation of the metric.} Let $g$ be a $G$-invariant metric on $G/H$ (\autoref{def:metrica_invariante}). For all $\gamma \in G$, the map $L_\gamma: p \mapsto \gamma \cdot p$ satisfies, by definition of $G$-invariance:
    \[
    g_p(u, v) = g_{\gamma \cdot p}(d(L_\gamma)_p\, u,\; d(L_\gamma)_p\, v)
    \]
    for all $p \in G/H$ and $u, v \in T_p(G/H)$. This means that $L_\gamma^* g = g$, that is, $L_\gamma$ is an isometry. \qedhere
\end{enumerate}
\end{proof}

\begin{definition}[Invariant Metric]\label{def:metrica_invariante}
A Riemannian metric $g$ on a homogeneous space $G/H$ is \textit{$G$-invariant} if for all $\gamma \in G$ and tangent vectors $u, v \in T_p(G/H)$:
\begin{equation}
g_p(u, v) = g_{\gamma \cdot p}(d\gamma_p u, d\gamma_p v)
\end{equation}
where $d\gamma_p$ denotes the differential of the action of $\gamma$ at the point $p$.
\end{definition}

\paragraph{Correspondence between geometry and equivariance}

The theory of homogeneous spaces establishes a fundamental connection between the geometric structure of data domains and the equivariance properties of the functions defined on them.

\begin{theorem}[Determination of equivariant operators by the base point]\label{thm:correspondencia_geometria_equivarianza}
Let $\mathcal{M} = G/H$ be a homogeneous space and let $\rho_1, \rho_2$ be representations of $G$ on vector spaces $V_1, V_2$. Let $\mathcal{F}(\mathcal{M}, V_i) = \{f: \mathcal{M} \to V_i\}$ be the space of functions on $\mathcal{M}$ with values in $V_i$, with the left regular action $(L_g f)(x) = \rho_i(g)\, f(g^{-1} \cdot x)$. Every $G$-equivariant operator $\Phi: \mathcal{F}(\mathcal{M}, V_1) \to \mathcal{F}(\mathcal{M}, V_2)$ (that is, $\Phi \circ L_g = L_g \circ \Phi$ for all $g \in G$) is completely determined by its values at the base point $o = eH$: if $\Phi_1$ and $\Phi_2$ are $G$-equivariant and $(\Phi_1 f)(o) = (\Phi_2 f)(o)$ for all $f$, then $\Phi_1 = \Phi_2$.
\end{theorem}

\begin{proof}
Let $x \in \mathcal{M}$ be an arbitrary point. Since $G$ acts transitively on $G/H$, there exists $g \in G$ with $g \cdot o = x$. For any $f \in \mathcal{F}(\mathcal{M}, V_1)$, $G$-equivariance gives:
\[
(\Phi f)(x) = (\Phi f)(g \cdot o) = \rho_2(g)\,(\Phi f)(g^{-1} \cdot (g \cdot o)) \quad \text{(evaluating } L_{g^{-1}}(\Phi f) \text{ at } o\text{)}
\]
More precisely, by equivariance $L_g(\Phi f) = \Phi(L_g f)$, from which:
\[
(\Phi f)(x) = (L_{g^{-1}}^{-1}(\Phi f))(o) = \rho_2(g)\,(\Phi(L_{g^{-1}} f))(o)
\]
since $(L_{g^{-1}} h)(o) = \rho_2(g^{-1})\,h(g \cdot o) = \rho_2(g^{-1})\,h(x)$, which implies $h(x) = \rho_2(g)\,(L_{g^{-1}} h)(o)$. Thus, $(\Phi f)(x) = \rho_2(g)\,(\Phi(L_{g^{-1}} f))(o)$, and the value of $\Phi f$ at any point $x$ is determined by the evaluation of $\Phi$ at the base point $o$.
\end{proof}

\begin{corollary}[Equivariant Linear Operators]\label{cor:operadores_equivariantes}
Let $G$ be a compact Lie group acting transitively on itself by left translation, and let $\Phi: L^2(G) \to L^2(G)$ be a continuous $G$-equivariant linear operator (that is, $\Phi \circ L_g = L_g \circ \Phi$ for all $g \in G$). Then $\Phi$ is a generalized convolution operator on $G$.
\end{corollary}

\begin{proof}
By \autoref{thm:correspondencia_geometria_equivarianza}, $\Phi$ is determined by its values at the identity $e \in G$. Equivariance and linearity imply that $\Phi$ admits a representation as an integral operator with kernel $k$:
\[
(\Phi f)(g) = \int_{G} k(g, g') f(g')\, d\mu(g')
\]
where $\mu$ is the normalized Haar measure. The equivariance condition $\Phi(L_h f) = L_h(\Phi f)$ implies, by a direct computation with a change of variables, that $k(hg, hg') = k(g, g')$ for all $h \in G$. Taking $h = g^{-1}$, we obtain $k(g, g') = k(e, g^{-1}g') = \kappa(g^{-1}g')$ for the function $\kappa = k(e, \cdot): G \to \mathbb{R}$. The operator $\Phi$ then takes the convolution form:
\[
(\Phi f)(g) = \int_{G} \kappa(g^{-1}g') f(g')\, d\mu(g') = (\kappa * f)(g) \qedhere
\]
\end{proof}

This result mathematically justifies why convolutions are the natural operation in equivariant architectures: they are precisely the linear operators that respect the symmetry structure of the domain.

\textbf{Consequences for \gls{GDL}.}
\begin{enumerate}
    \item \textbf{Principle of necessary invariance}: Only the genuine symmetries of the problem should be incorporated into the architecture. Artificial symmetries can degrade performance.

    \item \textbf{Hierarchy of equivariances}: In homogeneous spaces with multiple levels of symmetry, architectures can implement approximate equivariance in early layers and exact invariance in final layers.

    \item \textbf{Geometric transferability}: Models trained on a homogeneous space $G/H$ can be transferred to other homogeneous spaces $G/H'$ if there exists a natural homomorphism between the stabilizers $H$ and $H'$.
\end{enumerate}

\paragraph{Fundamental representations as quotients}

The Orbit-Stabilizer Theorem~\ref{thm:orbita_estabilizador} and \autoref{ex:aplicacion_teorema_gdl} showed how $S^2 \cong SO(3)/SO(2)$. We now apply the same technique systematically to the non-Euclidean spaces introduced in \autoref{subsec:espacios_no_euclidianos}, rigorously justifying each representation $G/H$.

\begin{proposition}[Sphere as a Homogeneous Space]\label{prop:esfera_cociente}
The sphere $S^{n-1}$ (\autoref{def:esfera_n}) is a homogeneous space:
\[
S^{n-1} \cong SO(n)/SO(n-1).
\]
\end{proposition}

\begin{proof}
Let $\mathbf{e}_n = (0,\ldots,0,1) \in S^{n-1}$ be the north pole. The group $SO(n)$ (\autoref{def:grupo_so_n}) acts on $S^{n-1}$ by matrix multiplication: if $R \in SO(n)$ and $\mathbf{x} \in S^{n-1}$, then $\|R\mathbf{x}\| = \|\mathbf{x}\| = 1$, so $R\mathbf{x} \in S^{n-1}$.

\textit{Transitivity.} Given $\mathbf{x} \in S^{n-1}$, we complete $\mathbf{x}$ to a positively oriented orthonormal basis $\{\mathbf{v}_1, \ldots, \mathbf{v}_{n-1}, \mathbf{x}\}$ of $\mathbb{R}^n$ by means of the Gram-Schmidt process. The matrix $R = [\mathbf{v}_1 \mid \cdots \mid \mathbf{v}_{n-1} \mid \mathbf{x}]$ satisfies $R \in SO(n)$ and $R\mathbf{e}_n = \mathbf{x}$.

\textit{Stabilizer.} If $R\mathbf{e}_n = \mathbf{e}_n$, the last column of $R$ is $\mathbf{e}_n$ and orthogonality forces $R$ to have the form:
\[
R = \begin{pmatrix} A & 0 \\ 0 & 1 \end{pmatrix}, \quad A \in SO(n-1).
\]
Hence $\mathrm{Stab}(\mathbf{e}_n) \cong SO(n-1)$.

By the Orbit-Stabilizer Theorem~\ref{thm:orbita_estabilizador}, $S^{n-1} = \mathrm{Orb}(\mathbf{e}_n) \cong SO(n)/SO(n-1)$. This result generalizes \autoref{ex:aplicacion_teorema_gdl} of $S^2 \cong SO(3)/SO(2)$ to arbitrary dimension.
\end{proof}

The compactness of $SO(n)$ algebraically reflects the constant positive curvature $K=+1$ of the sphere; this homogeneous space is fundamental for directional data, omnidirectional vision and molecular orientations.

\begin{proposition}[Euclidean Space as a Homogeneous Space]\label{prop:euclidiano_cociente}
Euclidean space is a homogeneous space:
\[
\mathbb{R}^n \cong E(n)/O(n) \cong SE(n)/SO(n).
\]
\end{proposition}

\begin{proof}
The Euclidean group $E(n) = \mathbb{R}^n \rtimes O(n)$ (\autoref{def:grupo_e_n}) acts on $\mathbb{R}^n$ by $(t, R) \cdot x = Rx + t$.

\textit{Transitivity.} For all $x \in \mathbb{R}^n$, $(x, I_n) \cdot \mathbf{0} = x$, so the orbit of the origin is all of $\mathbb{R}^n$.

\textit{Stabilizer.} If $(t, R) \cdot \mathbf{0} = t = \mathbf{0}$, then $t = \mathbf{0}$ and $R$ is arbitrary in $O(n)$. Hence $\mathrm{Stab}(\mathbf{0}) \cong O(n)$.

The Orbit-Stabilizer Theorem~\ref{thm:orbita_estabilizador} gives $\mathbb{R}^n \cong E(n)/O(n)$. Restricting to the special Euclidean group $SE(n) = \mathbb{R}^n \rtimes SO(n)$ (\autoref{def:grupo_se_n}), which also acts transitively, one obtains analogously $\mathbb{R}^n \cong SE(n)/SO(n)$.
\end{proof}

The curvature $K=0$ is reflected in the fact that the translational component of $E(n)$ is abelian; this is the base case of classical \glspl{CNN}, where translational equivariance is implemented through standard convolutions.

\begin{proposition}[Projective Space as a Homogeneous Space]\label{prop:proyectivo_cociente}
Real projective space (\autoref{def:espacio_proyectivo}) is a homogeneous space:
\[
\mathbb{RP}^{n-1} \cong O(n)/(O(n-1) \times \mathbb{Z}_2).
\]
\end{proposition}

\begin{proof}
The orthogonal group $O(n)$ (\autoref{def:grupo_o_n}) acts on the set of lines through the origin in $\mathbb{R}^n$ via $R \cdot [\mathbf{x}] = [R\mathbf{x}]$. The action is well defined: if $[\mathbf{x}] = [\mathbf{y}]$, then $\mathbf{y} = \lambda \mathbf{x}$ with $\lambda \neq 0$, and $[R\mathbf{y}] = [\lambda R\mathbf{x}] = [R\mathbf{x}]$.

\textit{Transitivity.} Every line through the origin contains a unit vector. Since $O(n)$ acts transitively on $S^{n-1}$ (because both rotations and reflections preserve the norm), it acts transitively on the lines.

\textit{Stabilizer.} The line $[\mathbf{e}_n]$ is fixed by $R$ if and only if $R\mathbf{e}_n = \pm\mathbf{e}_n$. In both cases, orthogonality forces:
\[
R = \begin{pmatrix} A & 0 \\ 0 & \pm 1 \end{pmatrix}, \quad A \in O(n-1).
\]
Hence $\mathrm{Stab}([\mathbf{e}_n]) \cong O(n-1) \times \mathbb{Z}_2$, and the result follows from the Orbit-Stabilizer Theorem~\ref{thm:orbita_estabilizador}.
\end{proof}

The factor $\mathbb{Z}_2$ in the stabilizer reflects the antipodal identification that defines projective space; the curvature is inherited from $S^{n-1}$ via the double cover $S^{n-1} \to \mathbb{RP}^{n-1}$.

\begin{proposition}[Hyperbolic Space as a Homogeneous Space]\label{prop:hiperbolico_cociente}
Hyperbolic space (\autoref{def:espacio_hiperbolico}) is a homogeneous space:
\[
\mathbb{H}^n \cong SO^+(n,1)/SO(n).
\]
\end{proposition}

\begin{proof}
We consider $\mathbb{H}^n$ in the hyperboloid model: $\mathbb{H}^n = \{x \in \mathbb{R}^{n+1} : x_0^2 - x_1^2 - \cdots - x_n^2 = 1, \; x_0 > 0\}$. The group $SO^+(n,1)$, the identity-connected component of the Lorentz group, preserves the Lorentz quadratic form $\langle x, y \rangle_{n,1} = x_0 y_0 - \sum_{i=1}^n x_i y_i$ and the condition $x_0 > 0$.

\textit{Transitivity.} Given two points $p, q \in \mathbb{H}^n$, the composition of a (hyperbolic) Lorentz boost taking $p$ to the base point $\mathbf{e}_0 = (1,0,\ldots,0)$, followed by another boost taking $\mathbf{e}_0$ to $q$, provides an element of $SO^+(n,1)$ sending $p$ to $q$.

\textit{Stabilizer.} If $g \in SO^+(n,1)$ fixes $\mathbf{e}_0 = (1,0,\ldots,0)$, then $g$ preserves the orthogonal complement of $\mathbf{e}_0$ with respect to the Lorentz form, which is the spatial subspace $\{0\} \times \mathbb{R}^n$, where the Lorentz form restricts to $-\sum_{i=1}^n x_i^2$. Therefore, $g$ acts as a rotation on the spatial coordinates, hence $\mathrm{Stab}(\mathbf{e}_0) \cong SO(n)$.

By the Orbit-Stabilizer Theorem~\ref{thm:orbita_estabilizador}, $\mathbb{H}^n \cong SO^+(n,1)/SO(n)$.
\end{proof}

Unlike the spherical case, $SO^+(n,1)$ is non-compact, which algebraically reflects the negative curvature $K=-1$ and the exponential growth of the volume of balls in $\mathbb{H}^n$.

\begin{proposition}[Grassmann and Stiefel Manifolds as Homogeneous Spaces]\label{prop:grassmann_stiefel_cociente}
The Grassmann and Stiefel manifolds (\autoref{def:variedad_grassmann}) are homogeneous spaces:
\[
Gr(k,n) \cong O(n)/(O(k) \times O(n-k)), \qquad St(k,n) \cong O(n)/O(n-k).
\]
\end{proposition}

\begin{proof}
\textit{Grassmann.} The group $O(n)$ acts on the set of $k$-dimensional subspaces of $\mathbb{R}^n$ by $R \cdot V = \{Rv : v \in V\}$.

The action is transitive: given a subspace $V$ of dimension $k$, we choose an orthonormal basis $\{v_1,\ldots,v_k\}$ of $V$ and complete it to an orthonormal basis $\{v_1,\ldots,v_n\}$ of $\mathbb{R}^n$. The matrix $R = [v_1 \mid \cdots \mid v_n] \in O(n)$ satisfies $R \cdot \mathrm{span}(\mathbf{e}_1,\ldots,\mathbf{e}_k) = V$.

The stabilizer of $V_0 = \mathrm{span}(\mathbf{e}_1,\ldots,\mathbf{e}_k)$ consists of the matrices $R \in O(n)$ that preserve $V_0$ and its orthogonal complement $V_0^\perp$. These have the block-diagonal form:
\[
R = \begin{pmatrix} A & 0 \\ 0 & B \end{pmatrix}, \quad A \in O(k), \; B \in O(n-k).
\]
Hence $\mathrm{Stab}(V_0) \cong O(k) \times O(n-k)$ and $Gr(k,n) \cong O(n)/(O(k) \times O(n-k))$.

\textit{Stiefel.} The Stiefel manifold $St(k,n)$ is the space of orthonormal $k$-frames in $\mathbb{R}^n$. Taking the canonical frame $(\mathbf{e}_1,\ldots,\mathbf{e}_k)$ as the base point, the stabilizer consists of the matrices in $O(n)$ that fix the first $k$ canonical vectors, that is, matrices of the form $\mathrm{diag}(I_k, B)$ with $B \in O(n-k)$. Hence $St(k,n) \cong O(n)/O(n-k)$.
\end{proof}

The stabilizer $O(k) \times O(n{-}k)$ separates the rotations within the subspace from those of the orthogonal complement, reflecting the metric of principal angles. The natural projection $\pi\colon St(k,n) \to Gr(k,n)$ forgets the specific basis and retains only the spanned subspace.

In the preceding propositions, the pair $(G, H)$ encodes all the geometric information: $G$ captures the global symmetries and $H$ the local symmetry at a point, determining curvature, compactness and volume growth. We now complete the picture with two families of homogeneous spaces whose quotients involve unitary and general linear groups.

\paragraph{Complex Spaces}

\begin{definition}[Complex Projective Space]\label{def:cp_homogeneo}
The \textit{complex projective space} $\mathbb{CP}^{n-1}$ is the space of complex lines in $\mathbb{C}^n$:
\begin{equation}
\mathbb{CP}^{n-1} = SU(n)/S(U(n-1) \times U(1)) = U(n)/(U(n-1) \times U(1))
\end{equation}
It is a compact Kähler manifold, fundamental in quantum mechanics (the space of pure states) and complex signal processing.
\end{definition}

\paragraph{Matrix Spaces}

\begin{definition}[Manifold of Symmetric Positive Definite Matrices]\label{def:spd_homogeneo}
The space $\mathrm{Sym}^+(n)$ of $n \times n$ symmetric positive definite matrices is:
\begin{equation}
\mathrm{Sym}^+(n) = GL(n)/O(n)
\end{equation}
Equipped with the affine-invariant metric, it is a symmetric space of nonpositive curvature. It is fundamental in geometric statistics for covariance matrices and diffusion tensors.
\end{definition}

\begin{remark}[Homogeneous spaces and GDL architectures]\label{rem:espacios_homogeneos_arquitecturas}
	The classification of homogeneous spaces has direct implications for the design of neural architectures in \gls{GDL}. \autoref{tab:espacios_homogeneos_arquitecturas} summarizes the main connections:

	\begin{table}[htbp]
	\centering
	\caption{Correspondence between homogeneous spaces, GDL architectures and applications.}
	\label{tab:espacios_homogeneos_arquitecturas}
	\footnotesize
	\begin{tabular}{@{}llll@{}}
		\toprule
		\textbf{Space} & \textbf{Quotient $G/H$} & \textbf{Architecture} & \textbf{Application} \\
		\midrule
		$S^{n-1}$ & $SO(n)/SO(n{-}1)$ & \makecell[l]{Spherical CNN \\  \citealp{cohen2018sphericalcnns}} & Climate, planetary data \\
		$\mathbb{R}^n$ & $E(n)/O(n)$ & \makecell[l]{Classical CNN \\ \citealp{cohen2016groupequivariantconvolutional}} & Images, signals \\
		$\mathbb{RP}^{n-1}$ & $O(n)/(O(n{-}1) \times \mathbb{Z}_2)$ & \makecell[l]{Projective CNN} & Unoriented axes \\
		$\mathbb{H}^n$ & $SO^+(n,1)/SO(n)$ & \makecell[l]{Hyperbolic NN \\  \citealp{nickel2017poincare}} & Knowledge graphs \\
		$SE(3)/SO(3)$ & $\cong \mathbb{R}^3$ & \makecell[l]{SE(3)-equivariant \\  \citealp{fuchs2020se3transformers}} & Molecules, proteins \\
		$Gr(k,n)$ & $O(n)/(O(k) \times O(n{-}k))$ & \makecell[l]{Grassmann NN \\ \scriptsize \citealp{huang2017grassmann}} & Subspace clustering \\
		$\mathbb{CP}^{n-1}$ & $U(n)/(U(n{-}1) \times U(1))$ & \makecell[l]{Complex NN} & Quantum mechanics \\
		$\mathrm{Sym}^+(n)$ & $GL(n)/O(n)$ & \makecell[l]{SPD-Net \\  \citealp{huang2017spd}} & Covariances, EEG \\
		\bottomrule
	\end{tabular}
	\end{table}

	Beyond the classification, the algebraic nature of $G$ determines the computational strategy. When $G$ is compact (such as $SO(n)$ for the sphere), the Peter-Weyl Theorem~\ref{thm:peter_weyl} provides a complete basis of irreducible representations that makes it possible to build equivariant filters in the spectral domain. When $G$ is non-compact (such as $SO^+(n,1)$ for hyperbolic space), this harmonic decomposition is not directly available, and one resorts to operations in the tangent space by means of the exponential and logarithmic maps.

	This correspondence illustrates a fundamental principle of \gls{GDL}: the choice of homogeneous space determines the symmetry group $G$ and the stabilizer $H$, which in turn dictates the structure of the available equivariant representations.
\end{remark}


\section{Graphs: The Discretization of Geometries}\label{sec:grafos_discretizacion}

Neural networks operate on finite, discrete structures, but real-world data often resides on continuous manifolds: curved surfaces, three-dimensional Euclidean space, the terrestrial sphere. Graphs constitute the natural bridge between both worlds. Through the construction of \textbf{proximity graphs}, a continuous manifold is approximated by a finite set of points connected by edges that encode neighborhood relations, preserving essential topological and spectral properties.

This approximation has profound implications for \gls{GDL}. On the one hand, it theoretically justifies the use of \glspl{GNN} for data that naturally reside on manifolds (molecules, 3D point clouds, climate data). On the other hand, it establishes fundamental limits: any geometric information of the continuous structure that is not captured by the discretization will be inaccessible to graph-based architectures.

\subsection{Graphs as the Discretization of Manifolds}\label{subsec:grafos_discretizacion_variedades}

Formally, the discrete approximation of a manifold $\mathcal{M}$ is built from a finite sampling of points and a connectivity rule based on distances. Recall that the fundamental definitions of graphs (graph, directed and undirected graph, paths, and cycles) were introduced in \autoref{sec:redes_grafos} of \autoref{ch:neural_networks}. Here we extend these concepts to the geometric context, endowing graphs with weights that encode metric information of the underlying manifold.

\begin{definition}[Weighted graph]\label{def:grafo_ponderado}
	A \textit{weighted graph} is a graph $G = (V, E)$ equipped with a weight function $w: E \to \mathbb{R}$ that assigns a numerical value to each edge. When the graph is undirected, it is required that $w(u,v) = w(v,u)$.
\end{definition}

\begin{definition}[Neighborhood]\label{def:vecindad_grafo}
	For a vertex $v \in V$, the \textit{neighborhood} of $v$ is defined as:
	\[\mathcal{N}(v) = \{u \in V : (v,u) \in E\}\]
	that is, the set of all vertices that are directly connected to $v$ by an edge.
\end{definition}

\begin{definition}[Neighborhood graph]\label{def:grafo_vecindad_epsilon}
	Given a set of points $\{x_i\}_{i=1}^n$ on a manifold $\mathcal{M}$, we can construct a \textit{neighborhood graph} $G = (V,E)$ where:
	\begin{itemize}
		\item $V = \{x_1, \ldots, x_n\}$ (the vertices are the sampled points)
		\item $(x_i, x_j) \in E$ if $d_\mathcal{M}(x_i, x_j) < \epsilon$ for some threshold $\epsilon > 0$
	\end{itemize}
	where $d_\mathcal{M}$ denotes the geodesic distance on $\mathcal{M}$ (\autoref{def:variedad_riemanniana}).
\end{definition}

This simple construction has profound consequences: if the points $\{x_i\}$ are sufficiently dense in $\mathcal{M}$ and the radius $\epsilon$ is chosen appropriately, the resulting graph $G$ captures fundamental topological and geometric properties of the original manifold.

\begin{figure}[htbp]
\centering
\includegraphics[width=0.7\textwidth]{figures/epsilon_neighborhood_graph.pdf}
\caption[$\epsilon$-neighborhood graph]{Construction of an $\epsilon$-neighborhood graph. Points sampled from a manifold (left) are connected when their distance is less than $\epsilon$, forming a graph that approximates the local structure of the original manifold (right).}
\label{fig:epsilon_neighborhood_graph}
\end{figure}

\paragraph{Regular Grids as a Special Case}

The regular grids on which traditional \glspl{CNN} operate are particular cases of this general construction:

\begin{example}[2D image as a regular graph]\label{ex:imagen_grafo_grilla}
	A digital image of dimension $H \times W$ can be interpreted as a graph $G = (V,E)$ where:
	\begin{itemize}
		\item $V = \{(i,j) : 1 \leq i \leq H, 1 \leq j \leq W\}$ (pixels)
		\item $(i,j) \sim (i',j')$ if $|i-i'| + |j-j'| = 1$ (4-neighborhood connectivity)
	\end{itemize}

	This graph is a \textbf{uniform discretization} of the continuous domain $[0,1] \times [0,1]$ with the Euclidean metric. The regularity of the grid translates into special properties:
	\begin{enumerate}
		\item All interior nodes have exactly degree 4
		\item The graph possesses discrete translational symmetry
		\item The geodesic distance in the graph approximates the $L^1$ distance in the continuous domain
	\end{enumerate}
\end{example}

\begin{example}[Triangulated 3D surface as an irregular graph]\label{ex:superficie_triangulada}
	Consider a three-dimensional surface $\mathcal{S} \subset \mathbb{R}^3$ (for example, a 3D model of an object). A \textit{triangulation} of $\mathcal{S}$ produces a graph $G = (V,E)$ where:
	\begin{itemize}
		\item $V$ are the vertices of the triangulation
		\item $E$ are the edges of the triangles
	\end{itemize}

	Unlike regular grids:
	\begin{enumerate}
		\item The node degrees are \textbf{heterogeneous} (typically between 4 and 8 for well-formed meshes)
		\item There is no global translational symmetry
		\item The local curvature of $\mathcal{S}$ is reflected in the density of vertices and the irregularity of the graph
	\end{enumerate}

	This construction is fundamental in computer graphics, physical simulations, and 3D shape analysis. Since $\mathcal{S}$ is a surface (compact 2-manifold), the existence of such a triangulation is guaranteed by \autoref{thm:rado}. \glspl{GNN} on these triangulations must respect \textbf{gauge equivariance} (see \autoref{sec:gauges_invarianzas}): the predictions cannot depend on the specific orientation of the mesh or on the chosen parametrization.
\end{example}

\paragraph{Simplices, Triangulations, and Simplicial Complexes}

To formalize triangulations and their generalizations, we need to define precisely the fundamental building blocks of combinatorial topology.

\begin{definition}[$k$-simplex]\label{def:simplicio_gdl}
	A \textit{$k$-simplex} $\sigma$ is the convex hull (the set of all convex combinations) of $k+1$ affinely independent points $v_0, v_1, \ldots, v_k$ in $\mathbb{R}^n$:
	\[
		\sigma = [v_0, v_1, \ldots, v_k] = \left\{ \sum_{i=0}^k t_i v_i : \sum_{i=0}^k t_i = 1, \, t_i \geq 0 \right\}
	\]
	The points $v_i$ are called the \textit{vertices} of the simplex, and $k$ is its \textit{dimension}. The fundamental cases are: a 0-simplex is a point, a 1-simplex an edge, a 2-simplex a triangle, a 3-simplex a tetrahedron, and so on for higher dimensions.
\end{definition}

\begin{definition}[Face of a simplex]\label{def:cara_simplicio_gdl}
	A \textit{face} of dimension $j < k$ of the $k$-simplex $\sigma = [v_0, v_1, \ldots, v_k]$ is any simplex determined by a subset of $j+1$ vertices:
	\[
		\tau = [v_{i_0}, v_{i_1}, \ldots, v_{i_j}] \quad \text{with } \{i_0, i_1, \ldots, i_j\} \subseteq \{0, 1, \ldots, k\}
	\]
	We write $\tau \leq \sigma$ when $\tau$ is a face of $\sigma$.
\end{definition}

\begin{definition}[Simplicial complex]\label{def:complejo_simplicial}
	A \textit{simplicial complex} $\mathcal{K}$ is a finite collection of simplices such that:
	\begin{enumerate}
		\item If $\sigma \in \mathcal{K}$ and $\tau \leq \sigma$, then $\tau \in \mathcal{K}$ (closed under faces)
		\item If $\sigma, \tau \in \mathcal{K}$, then $\sigma \cap \tau$ is a face of both (well-defined intersection)
	\end{enumerate}

	The \textit{1-skeleton} of $\mathcal{K}$ is the graph $G = (V,E)$ formed by the 0-simplices (vertices) and 1-simplices (edges).
\end{definition}

\begin{figure}[htbp]
\centering
\includegraphics[width=0.85\textwidth]{figures/simplicial_complex.pdf}
\caption[Simplices and simplicial complex]{Top: simplices of increasing dimension: a 0-simplex (point), 1-simplex (edge), 2-simplex (filled triangle), and 3-simplex (tetrahedron). Bottom: a simplicial complex $\mathcal{K}$ assembled from simplices, whose 1-skeleton forms the underlying graph $G = (V,E)$.}
\label{fig:simplicial_complex}
\end{figure}

\begin{definition}[Triangulation]\label{def:triangulacion}
	A \textit{triangulation} of a compact topological manifold $M$ is a pair $(\mathcal{K}, h)$ where $\mathcal{K}$ is a simplicial complex and $h: |\mathcal{K}| \to M$ is a homeomorphism from the \textit{geometric realization} $|\mathcal{K}|$ (the union of the simplices of $\mathcal{K}$ with the induced topology) to $M$.
\end{definition}

\begin{theorem}[Radó's theorem]\label{thm:rado}
	Every compact topological 2-manifold admits a triangulation.
\end{theorem}

For a proof sketch, see Section~\ref{sec:rado} of Appendix~\ref{ch:teoremas_complementarios}. This classical result, proved by \citet{rado1925aufgabe}, guarantees that every compact surface can be decomposed into triangles in a manner compatible with its topology, which is fundamental both for the Gauss-Bonnet theorem and for the applications of \gls{GDL} on triangulated meshes. Note that triangulability depends crucially on dimension: in dimension 3 it was extended by \citet{moise1977geometric}, but in dimension $\geq 4$ there exist non-triangulable manifolds (see Remark~\ref{rem:triangulabilidad_dimensiones} of Appendix~\ref{ch:teoremas_complementarios}).

Simplicial complexes allow capturing higher-order topological information (not just pairwise connectivity, but also cycles, holes, higher-dimensional components), which motivates \textbf{\gls{TDL}}. The algebraic tools to quantify this topology and their applications to neural networks are developed in \autoref{subsec:tdl_grafos}.

\subsection{Topological Geometry of Graphs}\label{subsec:geometria_topologica_grafos}

Graphs possess a rich topological and spectral structure that allows defining discrete analogues of continuous geometric concepts. This section develops the \textbf{calculus on graphs}: discrete differential operators, harmonic analysis, and spectral properties.

\paragraph{Basic Topological Properties}

\begin{definition}[Path in a graph]\label{def:camino_grafo}
	A \textit{path} from $u$ to $v$ in a graph $G = (V,E)$ is a sequence of vertices $u = v_0, v_1, \ldots, v_k = v$ where $(v_{i-1}, v_i) \in E$ for all $i = 1, \ldots, k$. The integer $k$ is called the \textit{length} of the path.
\end{definition}

\begin{definition}[Connected graph]\label{def:grafo_conexo}
	A graph $G = (V,E)$ is \textit{connected} if for every pair of vertices $u, v \in V$ there exists a path from $u$ to $v$.
\end{definition}

\begin{definition}[Connected component]\label{def:componente_conexa}
	A \textit{connected component} of a graph $G = (V,E)$ is a maximal connected subgraph of $G$.
\end{definition}

\begin{definition}[Cycle]\label{def:ciclo_grafo}
	A \textit{cycle} in a graph $G = (V,E)$ is a closed path $v_0, v_1, \ldots, v_k = v_0$ with $k \geq 3$ and no repetition of intermediate vertices.
\end{definition}

\begin{definition}[Tree]\label{def:arbol_grafo}
	A \textit{tree} is a connected graph without cycles.
\end{definition}

These combinatorial properties encode fundamental aspects of the topology of the discretized space. For example, the number of connected components corresponds to the \textbf{Betti number} $\beta_0$ of the underlying manifold, while the fundamental cycles relate to $\beta_1$.

\paragraph{The Graph Laplacian}

The central operator of graph analysis is the \textbf{Laplacian}, which discretizes the Laplace-Beltrami operator on manifolds. To construct it, we first need the adjacency and degree matrices.

\begin{definition}[Adjacency matrix]\label{def:matriz_adyacencia}
	For an undirected weighted graph $G = (V,E,W)$ with $n = |V|$ vertices, the \textit{adjacency matrix} is $A \in \mathbb{R}^{n \times n}$ where $A_{ij} = W_{ij}$ if $(i,j) \in E$ and $A_{ij} = 0$ otherwise. For unweighted graphs, $A_{ij} \in \{0, 1\}$.
\end{definition}

\begin{definition}[Degree matrix]\label{def:matriz_grados}
	For a graph with adjacency matrix $A$, the \textit{degree matrix} is the diagonal matrix $D \in \mathbb{R}^{n \times n}$ where $D_{ii} = \deg(i) = \sum_{j=1}^{n} A_{ij}$.
\end{definition}

\begin{definition}[Graph Laplacian]\label{def:laplaciano_grafo}
	The \textit{Laplacian} of a graph with adjacency matrix $A$ and degree matrix $D$ is
	\[
		L_G = D - A.
	\]
\end{definition}

\begin{definition}[Normalized Laplacian]\label{def:laplaciano_normalizado}
	The \textit{normalized Laplacian} of a graph with Laplacian $L_G$ and degree matrix $D$ (with $D_{ii} > 0$ for all $i$) is
	\[
		\mathcal{L} = D^{-1/2} L_G \, D^{-1/2} = I - D^{-1/2} A \, D^{-1/2}.
	\]
\end{definition}

\begin{figure}[htbp]
\centering
\includegraphics[width=0.8\textwidth]{figures/graph_laplacian.pdf}
\caption[Graph Laplacian matrices]{Concrete example of the matrices associated with a graph: adjacency matrix $A$, degree matrix $D$, and Laplacian $L_G = D - A$. The entries of the Laplacian encode the local connectivity: $L_{ii} = \deg(i)$ and $L_{ij} = -W_{ij}$ for adjacent vertices.}
\label{fig:graph_laplacian}
\end{figure}

\begin{definition}[Positive semidefinite matrix]\label{def:semidefinida_positiva}
	A symmetric matrix $A \in \mathbb{R}^{n \times n}$ is \textit{positive semidefinite} (written $A \succeq 0$) if
	\[
		\mathbf{x}^T A \mathbf{x} \geq 0 \quad \text{for all } \mathbf{x} \in \mathbb{R}^n.
	\]
	If the inequality is strict for all $\mathbf{x} \neq \mathbf{0}$, then $A$ is said to be \textit{positive definite} ($A \succ 0$). By \autoref{thm:espectral_matrices_simetricas} (Appendix~\ref{ch:teoremas_complementarios}), $A \succeq 0$ if and only if all its eigenvalues are nonnegative.
\end{definition}

\begin{proposition}[Algebraic properties of the Laplacian]\label{prop:propiedades_algebraicas_laplaciano}
	The Laplacian $L_G = D - A$ of an undirected weighted graph satisfies:
	\begin{enumerate}
		\item \textbf{Symmetry}: $L_G = L_G^T$
		\item \textbf{Rows sum to zero}: $L_G \mathbf{1} = \mathbf{0}$
		\item \textbf{Positive semidefiniteness}: for all $\mathbf{x} \in \mathbb{R}^n$,
		\[
			\mathbf{x}^T L_G \mathbf{x} = \frac{1}{2}\sum_{(i,j)\in E} W_{ij}(x_i - x_j)^2 \geq 0
		\]
		\item \textbf{Discrete operator}: for functions $f: V \to \mathbb{R}$,
		\[
			(L_G f)(i) = \sum_{j\sim i} W_{ij}(f(i)-f(j))
		\]
	\end{enumerate}
\end{proposition}

\begin{proof}
	\textbf{(1) Symmetry}: By definition, $D$ is diagonal (hence symmetric) and $A$ is symmetric since the graph is undirected, so $L_G = D - A$ is symmetric.

	\textbf{(2) Rows sum to zero}: The $i$-th row of $L_G$ has diagonal entry $D_{ii} = \deg(i) = \sum_j A_{ij}$ and off-diagonal entries $-A_{ij}$. Thus, $\sum_j (L_G)_{ij} = \deg(i) - \sum_j A_{ij} = 0$.

	\textbf{(3) Positive semidefiniteness}: For any $\mathbf{x} \in \mathbb{R}^n$:
	\begin{align}
		\mathbf{x}^T L_G \mathbf{x} &= \mathbf{x}^T (D - A) \mathbf{x} = \mathbf{x}^T D \mathbf{x} - \mathbf{x}^T A \mathbf{x}\\
		&= \sum_{i=1}^n D_{ii} x_i^2 - \sum_{i,j=1}^n A_{ij} x_i x_j\\
		&= \sum_{i=1}^n x_i^2 \sum_{j} A_{ij} - \sum_{i,j=1}^n A_{ij} x_i x_j\\
		&= \sum_{i,j=1}^n A_{ij} (x_i^2 - x_i x_j)\\
		&= \frac{1}{2} \sum_{i,j=1}^n A_{ij} (x_i^2 - 2x_i x_j + x_j^2)\\
		&= \frac{1}{2} \sum_{(i,j) \in E} W_{ij} (x_i - x_j)^2 \geq 0
	\end{align}

	\textbf{(4) Discrete operator}: By definition:
	\[
		(L_G f)(i) = (Df)(i) - (Af)(i) = \sum_{j} A_{ij} f(i) - \sum_{j} A_{ij} f(j) = \sum_{j \sim i} W_{ij}(f(i) - f(j)) \qedhere
	\]
\end{proof}

\begin{proposition}[Spectral properties of the Laplacian]\label{prop:propiedades_espectrales_laplaciano_grafos}
	Let $L_G$ be the Laplacian of an undirected graph with $n$ vertices. Then:
	\begin{enumerate}
		\item The eigenvalues satisfy $0 = \lambda_1 \leq \lambda_2 \leq \cdots \leq \lambda_n$
		\item The multiplicity of $\lambda_1 = 0$ equals the number of connected components of $G$
	\end{enumerate}
\end{proposition}

\begin{proof}
	\textbf{(1)}: By \autoref{prop:propiedades_algebraicas_laplaciano}, $L_G$ is symmetric and positive semidefinite (\autoref{def:semidefinida_positiva}). \autoref{thm:espectral_matrices_simetricas} (Appendix~\ref{ch:teoremas_complementarios}) guarantees that all eigenvalues are real, and positive semidefiniteness implies $\lambda_i \geq 0$ by \autoref{cor:caracterizacion_semidefinida}. Moreover, $L_G \mathbf{1} = \mathbf{0}$ (item 2 of \autoref{prop:propiedades_algebraicas_laplaciano}), so $\lambda_1 = 0$ with eigenvector $\mathbf{1}$.

	\textbf{(2)}: $L_G \mathbf{x} = \mathbf{0}$ if and only if $\mathbf{x}^T L_G \mathbf{x} = 0$, which by the quadratic form of \autoref{prop:propiedades_algebraicas_laplaciano} is equivalent to $(x_i - x_j)^2 = 0$ for every edge $(i,j) \in E$, that is, $\mathbf{x}$ is constant on each connected component. The space of such vectors has dimension equal to the number of connected components.
\end{proof}

The quadratic form $f^T L_G f$ measures the \textbf{smoothness} of the function $f$ over the graph: it is small when $f$ varies little between adjacent vertices, large when there are abrupt changes. This interpretation connects with the \textbf{Dirichlet energy} in functional analysis \citep[Ch.~1]{Chung1997}.

\paragraph{Connection with Differential Geometry}

The graph Laplacian $L_G$ is the discrete analogue of the Laplace-Beltrami operator $\Delta_\mathcal{M}$ on Riemannian manifolds (\autoref{def:laplace_beltrami}), which acts on smooth functions $f: \mathcal{M} \to \mathbb{R}$ as $\Delta_\mathcal{M} f = \operatorname{div}(\operatorname{grad} f)$.

The analogy is profound:
\begin{itemize}
	\item \textbf{On continuous manifolds}: $\Delta_\mathcal{M} f = \text{div}(\text{grad } f)$ measures the divergence of the gradient
	\item \textbf{On discrete graphs}: $(L_G f)(i) = \sum_{j \sim i} (f(i) - f(j))$ measures the difference between $f(i)$ and the average of its neighbors
\end{itemize}

Both operators generalize the Euclidean Laplacian $\nabla^2 f$ (\autoref{def:laplaciano_euclidiano}) to non-Euclidean domains.

\subsection{Topological Deep Learning}\label{subsec:tdl_grafos}

The Laplacian $L_G = D - A$ from \autoref{subsec:geometria_topologica_grafos} operates exclusively on node signals (0-chains). However, the simplicial complexes introduced in \autoref{def:complejo_simplicial} possess a much richer algebraic structure: signals defined on edges, triangles, and higher-order simplices. To quantify this structure, we introduce the tools of algebraic topology.

\begin{definition}[Chains]\label{def:cadenas_gdl}
	Let $\mathcal{K}$ be a simplicial complex. A \textit{$k$-chain} is a formal linear combination of $k$-simplices with coefficients in a field (typically $\mathbb{R}$ or $\mathbb{Z}_2$):
	\[
		C_k(\mathcal{K}) = \left\{ \sum_{i} \alpha_i \sigma_i^{(k)} : \alpha_i \in \mathbb{R}, \, \sigma_i^{(k)} \in \mathcal{K} \text{ is a } k\text{-simplex} \right\}
	\]
\end{definition}

\begin{definition}[Boundary operator]\label{def:operador_frontera_gdl}
	Let $\mathcal{K}$ be a simplicial complex. The \textit{boundary operator} $\partial_k: C_k(\mathcal{K}) \to C_{k-1}(\mathcal{K})$ is defined on an oriented $k$-simplex $[v_0, v_1, \ldots, v_k]$ as:
	\[
		\partial_k [v_0, v_1, \ldots, v_k] = \sum_{i=0}^{k} (-1)^i [v_0, \ldots, \widehat{v_i}, \ldots, v_k]
	\]
	where $\widehat{v_i}$ indicates the omission of vertex $v_i$.
\end{definition}

\begin{proposition}[Fundamental property of the boundary operator]\label{prop:frontera_cuadrado_cero}
	The boundary operator satisfies $\partial_{k-1} \circ \partial_k = 0$ for all $k$. That is, the boundary of a boundary is empty.
\end{proposition}

\begin{proof}
	It suffices to verify the identity on an oriented $k$-simplex $[v_0, \ldots, v_k]$:
	\begin{align*}
		\partial_{k-1}(\partial_k [v_0, \ldots, v_k]) &= \partial_{k-1}\!\left(\sum_{i=0}^{k} (-1)^i [v_0, \ldots, \widehat{v_i}, \ldots, v_k]\right) \\
		&= \sum_{i=0}^{k} (-1)^i \sum_{\substack{j=0 \\ j \neq i}}^{k} (-1)^{j'} [v_0, \ldots, \widehat{v_i}, \ldots, \widehat{v_j}, \ldots, v_k]
	\end{align*}
	where $j' = j$ if $j < i$ and $j' = j-1$ if $j > i$. Each codimension-2 face, obtained by omitting two vertices $v_i$ and $v_j$ with $i < j$, appears exactly twice: once with sign $(-1)^{i+j}$ (omitting first $v_i$ and then $v_j$) and another with sign $(-1)^{i+j-1}$ (omitting first $v_j$ and then $v_i$). Since $(-1)^{i+j} + (-1)^{i+j-1} = 0$, all terms cancel.
\end{proof}

\begin{definition}[Homology groups]\label{def:homologia_betti}
	The \textit{homology groups} of a simplicial complex $\mathcal{K}$ are defined as:
	\[
		H_k(\mathcal{K}) = \ker(\partial_k) / \mathrm{im}(\partial_{k+1}) = Z_k(\mathcal{K}) / B_k(\mathcal{K})
	\]
	where $Z_k = \ker(\partial_k)$ are the \textit{$k$-cycles} (chains without boundary) and $B_k = \mathrm{im}(\partial_{k+1})$ are the \textit{$k$-boundaries} (chains that are the boundary of something).
\end{definition}

\begin{definition}[Betti numbers]\label{def:numeros_betti_gdl}
	The \textit{$k$-th Betti number} of a simplicial complex $\mathcal{K}$ is $\beta_k = \dim H_k(\mathcal{K})$, which counts:
	\begin{itemize}
		\item $\beta_0$: number of connected components
		\item $\beta_1$: number of ``holes'' or independent cycles
		\item $\beta_2$: number of three-dimensional cavities
	\end{itemize}
\end{definition}

\begin{example}[Betti numbers of surfaces]\label{ex:betti_superficies}
\leavevmode
	\begin{itemize}
		\item \textbf{Sphere $S^2$}: $\beta_0 = 1$ (connected), $\beta_1 = 0$ (no holes), $\beta_2 = 1$ (one cavity)
		\item \textbf{Torus $T^2$}: $\beta_0 = 1$, $\beta_1 = 2$ (two independent cycles), $\beta_2 = 1$
		\item \textbf{Klein bottle}: $\beta_0 = 1$, $\beta_1 = 1$ (over $\mathbb{Z}_2$)
	\end{itemize}
\end{example}

\textbf{\gls{TDL}} extends message passing (\autoref{def:mpnn_unificador}) and graph analysis to these higher-order structures via the \textit{Hodge Laplacian}, which generalizes $L_G$ to arbitrary $k$-chains.

\subsubsection{The Hodge Laplacian}\label{subsubsec:hodge_laplacian_gdl}

The boundary operator $\partial_k$ (\autoref{def:operador_frontera_gdl}) and the property $\partial_{k-1} \circ \partial_k = 0$ (\autoref{prop:frontera_cuadrado_cero}) provide the sequence of chain spaces:
\[
	\cdots \xrightarrow{\partial_{k+2}} C_{k+1}(\mathcal{K}) \xrightarrow{\partial_{k+1}} C_k(\mathcal{K}) \xrightarrow{\partial_k} C_{k-1}(\mathcal{K}) \xrightarrow{\partial_{k-1}} \cdots
\]

\begin{definition}[Coboundary operator]\label{def:operador_cofrontera_gdl}
	We equip each chain space $C_k(\mathcal{K})$ with the \textit{standard inner product} in which the oriented $k$-simplices form an orthonormal basis. The \textit{coboundary operator} $\partial_k^T : C_{k-1}(\mathcal{K}) \to C_k(\mathcal{K})$ is the formal adjoint of the boundary operator (\autoref{def:operador_frontera_gdl}) with respect to this inner product, that is, the unique linear operator satisfying:
	\[
		\langle \partial_k^T c, \, \sigma \rangle = \langle c, \, \partial_k \sigma \rangle \quad \text{for all } c \in C_{k-1}(\mathcal{K}),\; \sigma \in C_k(\mathcal{K}).
	\]

	\textbf{Combinatorial interpretation.} For a $(k{-}1)$-simplex $\tau$, the coboundary operator aggregates all $k$-simplices of $\mathcal{K}$ that contain $\tau$ as a face, with the inherited incidence signs:
	\[
		\partial_k^T(\tau) = \sum_{\substack{\sigma \in \mathcal{K} \\ \dim \sigma = k}} [\partial_k]_{\tau\sigma} \, \sigma,
	\]
	where $[\partial_k]_{\tau\sigma}$ denotes the coefficient of $\tau$ in $\partial_k(\sigma)$, which is $(-1)^i$ if $\tau$ is obtained from $\sigma$ by omitting its $i$-th vertex, and $0$ if $\tau$ is not a face of $\sigma$.

	\textbf{Cochain complex.} The coboundary operators reverse the direction of the chain complex, forming the \textit{cochain complex}:
	\[
		\cdots \xleftarrow{\partial_{k+2}^T} C_{k+1}(\mathcal{K}) \xleftarrow{\partial_{k+1}^T} C_k(\mathcal{K}) \xleftarrow{\partial_k^T} C_{k-1}(\mathcal{K}) \xleftarrow{\partial_{k-1}^T} \cdots
	\]

	\textbf{Nilpotency.} Transposing $\partial_{k} \circ \partial_{k+1} = 0$ (\autoref{prop:frontera_cuadrado_cero}) yields $\partial_{k+1}^T \circ \partial_k^T = 0$, which guarantees that the previous sequence is indeed a cochain complex.
\end{definition}

\begin{definition}[Hodge Laplacian]\label{def:hodge_laplacian_gdl}
	Let $\mathcal{K}$ be a finite simplicial complex with boundary operators $\partial_k$. The \textit{Hodge Laplacian of order $k$} is the operator $L_k: C_k(\mathcal{K}) \to C_k(\mathcal{K})$ defined by:
	\[
		L_k = \underbrace{\partial_k^T \partial_k}_{L_k^{\mathrm{down}}} + \underbrace{\partial_{k+1} \partial_{k+1}^T}_{L_k^{\mathrm{up}}}
	\]
	where:
	\begin{itemize}
		\item $L_k^{\mathrm{down}} = \partial_k^T \partial_k$ captures \textit{lower} adjacency: two $k$-simplices are adjacent if they share a $(k{-}1)$-face.
		\item $L_k^{\mathrm{up}} = \partial_{k+1} \partial_{k+1}^T$ captures \textit{upper} adjacency: two $k$-simplices are adjacent if they belong to the same $(k{+}1)$-simplex.
	\end{itemize}
\end{definition}

\begin{proposition}[Properties of the Hodge Laplacian]\label{prop:propiedades_hodge_gdl}
	The Hodge Laplacian $L_k$ satisfies:
	\begin{enumerate}
		\item \textbf{Symmetry and positive semidefiniteness}: $L_k$ is symmetric and $\langle x, L_k x \rangle \geq 0$ for every $k$-chain $x$, since $L_k$ is a sum of Gram matrices: $\langle x, \partial_k^T \partial_k x \rangle = \|\partial_k x\|^2 \geq 0$ and $\langle x, \partial_{k+1} \partial_{k+1}^T x \rangle = \|\partial_{k+1}^T x\|^2 \geq 0$.
		\item \textbf{Recovery of the graph Laplacian}: For $k = 0$, since $\partial_0 = 0$ by convention, we have $L_0 = \partial_1 \partial_1^T$, which coincides with the Laplacian $L_G = D - A$ of \autoref{def:laplaciano_grafo}.
		\item \textbf{Connection with homology}: $\ker(L_k) \cong H_k(\mathcal{K})$, that is, the harmonic $k$-chains correspond to the $k$-th homology group. This result is an immediate corollary of \autoref{thm:hodge_decomposition_gdl}.
	\end{enumerate}
\end{proposition}

\begin{proof}
	\textbf{(1)}: $L_k$ is symmetric since it is a sum of Gram matrices: $(\partial_k^T \partial_k)^T = \partial_k^T \partial_k$ and $(\partial_{k+1} \partial_{k+1}^T)^T = \partial_{k+1} \partial_{k+1}^T$. For positive semidefiniteness, let $x \in C_k(\mathcal{K})$:
	\[
		\langle x, L_k x \rangle = \langle x, \partial_k^T \partial_k x \rangle + \langle x, \partial_{k+1} \partial_{k+1}^T x \rangle = \|\partial_k x\|^2 + \|\partial_{k+1}^T x\|^2 \geq 0.
	\]

	\textbf{(2)}: For $k = 0$, the space $C_{-1}(\mathcal{K})$ is trivial, so $\partial_0 = 0$ and $L_0 = \partial_1 \partial_1^T$. If $\partial_1$ is the oriented incidence matrix of the graph, with entries $(\partial_1)_{e,v} = +1$ if $v$ is the origin of $e$ and $-1$ if it is the destination, then the $(i,j)$ entry of $\partial_1 \partial_1^T$ is: $d_i$ if $i = j$ (vertex degree), $-1$ if $(i,j) \in E$, and $0$ otherwise. That is, $\partial_1 \partial_1^T = D - A = L_G$.

	\textbf{(3)}: By (1), $L_k x = 0$ if and only if $\langle x, L_k x \rangle = \|\partial_k x\|^2 + \|\partial_{k+1}^T x\|^2 = 0$, which is equivalent to $\partial_k x = 0$ and $\partial_{k+1}^T x = 0$ simultaneously. The first condition says $x \in Z_k = \ker(\partial_k)$; the second says $x \in \mathrm{im}(\partial_{k+1})^\perp$, since $\langle \partial_{k+1} a, x \rangle = \langle a, \partial_{k+1}^T x \rangle = 0$ for all $a$. Thus, $x$ is a cycle orthogonal to all boundaries, and the class $[x] \in Z_k / B_k = H_k(\mathcal{K})$ establishes the isomorphism $\ker(L_k) \cong H_k(\mathcal{K})$.
\end{proof}

The central result of discrete Hodge theory provides an orthogonal decomposition of the chain space that generalizes Fourier analysis on graphs (see \autoref{subsec:analisis_espectral_grafos_gdl}). For an exhaustive treatment of the Hodge Laplacian on graphs, see \citet{lim2020hodge}.

\begin{theorem}[Hodge decomposition]\label{thm:hodge_decomposition_gdl}
	For a finite simplicial complex $\mathcal{K}$, the space of $k$-chains admits the orthogonal decomposition:
	\[
		C_k(\mathcal{K}) = \underbrace{\mathrm{im}(\partial_{k+1})}_{\text{boundaries (``gradients'')}} \oplus \underbrace{\ker(L_k)}_{\text{harmonics}} \oplus \underbrace{\mathrm{im}(\partial_k^T)}_{\text{coboundaries (``curls'')}}
	\]
	The three components capture complementary aspects of a simplicial signal: the boundaries represent irrotational flow, the harmonics encode the global topology ($H_k$), and the coboundaries represent pure circulation.
\end{theorem}

\begin{proof}
	We must show that the three subspaces $\mathrm{im}(\partial_{k+1})$, $\ker(L_k)$, and $\mathrm{im}(\partial_k^T)$ are mutually orthogonal and that their direct sum is $C_k(\mathcal{K})$.

	\textbf{Step 1: Mutual orthogonality.} Let $x \in \mathrm{im}(\partial_{k+1})$ and $y \in \mathrm{im}(\partial_k^T)$. Then $x = \partial_{k+1} a$ for some $a \in C_{k+1}$ and $y = \partial_k^T b$ for some $b \in C_{k-1}$. Using the property $\partial_k \circ \partial_{k+1} = 0$ (\autoref{prop:frontera_cuadrado_cero}):
	\[
		\langle x, y \rangle = \langle \partial_{k+1} a, \partial_k^T b \rangle = \langle \partial_k \partial_{k+1} a, b \rangle = \langle 0, b \rangle = 0.
	\]
	Now, let $z \in \ker(L_k)$. From the proof of \autoref{prop:propiedades_hodge_gdl}(3), $L_k z = 0$ implies $\partial_k z = 0$ and $\partial_{k+1}^T z = 0$. Then:
	\[
		\langle z, \partial_{k+1} a \rangle = \langle \partial_{k+1}^T z, a \rangle = 0, \qquad
		\langle z, \partial_k^T b \rangle = \langle \partial_k z, b \rangle = 0.
	\]
	The three subspaces are therefore mutually orthogonal.

	\textbf{Step 2: Completeness.} For finite simplicial complexes, all spaces are finite-dimensional, so it suffices to verify that $C_k = V_1 \oplus V_2 \oplus V_3$ where $V_1 = \mathrm{im}(\partial_{k+1})$, $V_2 = \mathrm{im}(\partial_k^T)$, and $V_3 = \ker(L_k)$. Let $w \in (V_1 \oplus V_2)^\perp$. Then $w \perp \mathrm{im}(\partial_{k+1})$ and $w \perp \mathrm{im}(\partial_k^T)$, which implies:
	\[
		\partial_{k+1}^T w = 0 \quad \text{and} \quad \partial_k w = 0.
	\]
	Therefore $L_k w = \partial_k^T \partial_k w + \partial_{k+1} \partial_{k+1}^T w = 0$, that is, $w \in \ker(L_k)$. This proves $(V_1 \oplus V_2)^\perp = V_3$, and thus $C_k = V_1 \oplus V_2 \oplus V_3$.
\end{proof}

The Hodge decomposition is of central importance for \gls{GDL} for the following reasons:
\begin{itemize}
	\item \textbf{Generalization of Fourier analysis}: while the diagonalization of $L_0$ provides the Fourier transform on graphs (see \autoref{subsec:analisis_espectral_grafos_gdl}), the Hodge decomposition extends this analysis to signals on simplices of any order $k$.
	\item \textbf{Physical interpretation of the components}: every simplicial signal decomposes into a gradient component $\mathrm{im}(\partial_{k+1})$ (irrotational flow), a harmonic component $\ker(L_k)$ (global topology), and a curl component $\mathrm{im}(\partial_k^T)$ (pure circulation), analogously to the Helmholtz decomposition of vector fields.
	\item \textbf{Detection of topological features}: the harmonic component satisfies $\ker(L_k) \cong H_k(\mathcal{K})$, which allows simplicial neural networks to directly access the homological information of the complex.
\end{itemize}
The extensions of this decomposition to simplicial attention mechanisms are developed in \autoref{ch:attention}.

\begin{example}[Hodge Laplacian $L_1$ of a triangular complex]\label{ex:hodge_laplacian_triangular}
	Consider a simplicial complex $\mathcal{K}$ with vertices $\{a, b, c, d\}$, edges $e_1 = [a,b]$, $e_2 = [b,c]$, $e_3 = [a,c]$, $e_4 = [c,d]$, and a single 2-simplex $\sigma = [a,b,c]$. The matrix of the boundary operator $\partial_2: C_2 \to C_1$ is obtained from $\partial_2[a,b,c] = [b,c] - [a,c] + [a,b] = e_2 - e_3 + e_1$:
	\[
		\partial_2 = \begin{pmatrix} 1 \\ 1 \\ -1 \\ 0 \end{pmatrix} \quad \text{(column corresponding to $\sigma$)}
	\]
	and the matrix $\partial_1: C_1 \to C_0$ is:
	\[
		\partial_1 = \begin{pmatrix} -1 & 0 & -1 & 0 \\ 1 & -1 & 0 & 0 \\ 0 & 1 & 1 & -1 \\ 0 & 0 & 0 & 1 \end{pmatrix}
	\]
	It can be verified that $\partial_1 \partial_2 = 0$, as required by \autoref{prop:frontera_cuadrado_cero}. Then:
	\[
		L_1^{\mathrm{down}} = \partial_1^T \partial_1 = \begin{pmatrix} 2 & -1 & 1 & 0 \\ -1 & 2 & 1 & -1 \\ 1 & 1 & 2 & -1 \\ 0 & -1 & -1 & 2 \end{pmatrix}, \quad
		L_1^{\mathrm{up}} = \partial_2 \partial_2^T = \begin{pmatrix} 1 & 1 & -1 & 0 \\ 1 & 1 & -1 & 0 \\ -1 & -1 & 1 & 0 \\ 0 & 0 & 0 & 0 \end{pmatrix}
	\]
	Note how $L_1^{\mathrm{down}}$ connects edges that share a vertex (lower adjacency), while $L_1^{\mathrm{up}}$ connects edges that coexist in the triangle $\sigma$ (upper adjacency). The edge $e_4$, which does not belong to any triangle, has $L_1^{\mathrm{up}}$ null in its row/column.
\end{example}

\subsubsection{Message Passing on Simplicial Complexes}\label{subsubsec:message_passing_simplicial_gdl}

The decomposition $L_k = L_k^{\mathrm{down}} + L_k^{\mathrm{up}}$ reveals that each $k$-simplex receives information from two geometrically distinct sources: its $(k{-}1)$-faces (lower adjacency) and the $(k{+}1)$-simplices that contain it (upper adjacency). This dual structure motivates a natural generalization of message passing on graphs to simplicial complexes, whose formalization as an instance of the \gls{MPNN} framework is developed in \autoref{subsubsec:tdl_mpnn_general}.

\subsubsection{Architectural Variants and Extensions}\label{subsubsec:variantes_tdl}

The \gls{TDL} framework presented in this section admits various specializations and extensions, treated in detail in the following chapters:

\begin{itemize}
	\item \textbf{Topological convolutions}: Spectral filters on the Hodge Laplacian, parametrized via polynomials of $L_k$. See \autoref{def:hodge_laplacians_cnn} of \autoref{ch:cnn}.
	\item \textbf{Simplicial attention}: Attention mechanisms that weight simplicial messages according to orientation and topological relevance. See \autoref{def:mensaje_simplicial} of \autoref{ch:attention} and \citep{giusti2022simplicial}.
	\item \textbf{Cell complexes}: Generalization of simplicial complexes where the cells of dimension $k$ can have arbitrary shape (not only simplices). See \autoref{def:complejo_celular} of \autoref{ch:cnn} and \citep{bodnar2021cellular}.
	\item \textbf{Simplicial diffusion}: Extension of the diffusion equation to simplicial complexes via $L_k$. See \autoref{def:ecuacion_difusion_simplicial} of \autoref{ch:attention}.
	\item \textbf{Higher-order attention networks}: Architectures that combine multi-head attention with simplicial topology \citep{hajij2022higherorder, hajij2023topological}.
\end{itemize}

\subsection{Convergence of Graphs to Manifolds}\label{subsec:convergencia_grafos_variedades}

The fundamental question that connects graph theory with differential geometry is: \textbf{under what conditions does a discrete graph faithfully approximate a continuous manifold?} The following theorem, one of the central results in manifold learning theory, connects concepts from manifold analysis, random graph theory, and asymptotic analysis to provide a rigorous answer.

\begin{theorem}[Convergence of graphs to manifolds]\label{thm:convergencia_grafos_variedades_detallado}
	Let $(\mathcal{M}, g)$ be a compact Riemannian manifold of dimension $d$ immersed in $\mathbb{R}^N$. Let $\{x_1, \ldots, x_n\}$ be a sample of points on $\mathcal{M}$ distributed according to a density $p(x)$ bounded below: $p(x) \geq p_{\min} > 0$. We construct a graph $G_n$ with weight matrix $W_{ij} = k_\epsilon(\|x_i - x_j\|)$ where $k_\epsilon(t) = \frac{1}{\epsilon^d} k\left(\frac{t}{\epsilon}\right)$ and $k: \mathbb{R}_+ \to \mathbb{R}_+$ is a smooth radial kernel with compact support.

	Then, under the scaling conditions:
	\begin{itemize}
		\item $\epsilon = \epsilon_n \to 0$ as $n \to \infty$
		\item $n\epsilon^{d+2} \to \infty$ (enough neighbors per node)
	\end{itemize}

	the normalized graph Laplacian converges to the weighted Laplace-Beltrami operator:
	\[
		\lim_{n \to \infty} \frac{1}{p(x_i)\,\epsilon^2} L_n f(x_i) = \Delta_\mathcal{M} f(x_i) + \langle \nabla \log p(x_i), \nabla f(x_i) \rangle
	\]
	where $L_n$ is the graph Laplacian operator defined in Step~1 of the proof and the convergence is in probability uniformly for $f \in C^2(\mathcal{M})$.
\end{theorem}

\begin{proof}
	The proof proceeds in seven steps, combining techniques from differential geometry, functional analysis, and probability theory.

	\textbf{Step 1: Normalization of the graph Laplacian.}

	We define the density-normalized Laplacian as:
	\[
		L_n f(x_i) = \frac{1}{n \epsilon^d} \sum_{j=1}^n k_\epsilon(\|x_i - x_j\|) [f(x_j) - f(x_i)]
	\]
	where $k_\epsilon(t) = \frac{1}{\epsilon^d} k\left(\frac{t}{\epsilon}\right)$ normalizes the kernel appropriately. The factor $n\epsilon^d$ represents the expected local volume of the neighborhood.

	\textbf{Step 2: Local approximation via normal coordinates.}

	For $x_i \in \mathcal{M}$, we work in geodesic normal coordinates centered at $x_i$. In these coordinates, the Riemannian metric satisfies:
	\[
		g_{jk}(y) = \delta_{jk} + O(\|y\|^2)
	\]
	that is, the manifold is locally approximated by its tangent space $T_{x_i}\mathcal{M} \cong \mathbb{R}^d$.

	For points $x_j$ close to $x_i$ with $\|x_i - x_j\| < \epsilon$, we can write:
	\[
		x_j = \exp_{x_i}(\epsilon v_j)
	\]
	where $v_j \in T_{x_i}\mathcal{M}$ with $\|v_j\| \leq C$ for some constant $C$, and $\exp$ denotes the exponential map.

	\textbf{Step 3: Taylor expansion of $f$.}

	Using the Taylor expansion of $f$ around $x_i$ on the manifold:
	\begin{align}
		f(x_j) - f(x_i) &= f(\exp_{x_i}(\epsilon v_j)) - f(x_i)\\
		&= \epsilon \langle \nabla f(x_i), v_j \rangle + \frac{\epsilon^2}{2} \text{Hess}_f(x_i)(v_j, v_j) + O(\epsilon^3)
	\end{align}
	where:
	\begin{itemize}
		\item $\nabla f(x_i) \in T_{x_i}\mathcal{M}$ is the Riemannian gradient of $f$
		\item $\text{Hess}_f(x_i): T_{x_i}\mathcal{M} \times T_{x_i}\mathcal{M} \to \mathbb{R}$ is the intrinsic Hessian of $f$
	\end{itemize}

	\textbf{Step 4: Computation of the expected limit.}

	Substituting the expansion into $L_n f(x_i)$:
	\begin{align}
		\frac{1}{\epsilon^2} L_n f(x_i) &= \frac{1}{n \epsilon^{d+2}} \sum_{j=1}^n k_\epsilon(\|x_i - x_j\|) [f(x_j) - f(x_i)]\\
		&= \frac{1}{n \epsilon^{d+2}} \sum_{j=1}^n k_\epsilon(\|x_i - x_j\|) \left[\epsilon \langle \nabla f(x_i), v_j \rangle + \frac{\epsilon^2}{2} \text{Hess}_f(x_i)(v_j,v_j)\right]\\
		&\quad + O(\epsilon)
	\end{align}

	\textbf{Step 5: Asymptotic averaging and the law of large numbers.}

	As $n \to \infty$, by the strong law of large numbers (the points $\{x_j\}$ are i.i.d. with density $p$):
	\[
		\frac{1}{n} \sum_{j=1}^n F(x_j) \xrightarrow{\text{a.s.}} \mathbb{E}[F] = \int_{\mathcal{M}} F(y) \, p(y) \, dV_g(y)
	\]

	Applying this to the first-order term:
	\begin{align}
		&\frac{1}{n\epsilon^{d+1}} \sum_{j=1}^n k_\epsilon(\|x_i - x_j\|) \langle \nabla f(x_i), v_j \rangle\\
		&\approx \frac{1}{\epsilon^{d+1}} \int_{\mathcal{M}} k_\epsilon(\|x_i - y\|) \langle \nabla f(x_i), v(y) \rangle \, p(y) \, dV_g(y)
	\end{align}

	Changing variables to $y = \exp_{x_i}(\epsilon v)$ and expanding the density $p(\exp_{x_i}(\epsilon v)) = p(x_i) + \epsilon \langle \nabla p(x_i), v \rangle + O(\epsilon^2)$:
	\[
		\approx \int_{T_{x_i}\mathcal{M}} k(\|v\|) \langle \nabla f(x_i), v \rangle \, \bigl[p(x_i) + \epsilon \langle \nabla p(x_i), v \rangle\bigr] \, dv
	\]

	\textbf{Step 6: Analysis of the first-order term.}

	The term $p(x_i) \int k(\|v\|) \langle \nabla f, v \rangle \, dv = 0$ vanishes by the radial symmetry of $k$. However, the cross term of order $\epsilon$:
	\[
		\epsilon \int_{T_{x_i}\mathcal{M}} k(\|v\|) \langle \nabla f(x_i), v \rangle \langle \nabla p(x_i), v \rangle \, dv
	\]
	does not vanish in general, since the integrand contains products $v_\alpha v_\beta$ that survive the integration when $\alpha = \beta$. By radial symmetry, $\int k(\|v\|) v_\alpha v_\beta \, dv = \frac{1}{d}\left(\int k(\|v\|) \|v\|^2 \, dv\right) \delta_{\alpha\beta}$, so this term contributes proportionally to $\langle \nabla f(x_i), \nabla p(x_i) \rangle$.

	\textbf{Step 7: Leading second-order term and final result.}

	The second-order term gives:
	\begin{align}
		\frac{p(x_i)}{2} \int_{T_{x_i}\mathcal{M}} k(\|v\|) \text{Hess}_f(x_i)(v,v) \, dv
	\end{align}

	Using the relation $\text{tr}(\text{Hess}_f) = \Delta_\mathcal{M} f$ (the trace of the Hessian is the Laplacian) and radial symmetry:
	\begin{align}
		&= \frac{p(x_i)}{2d} \left(\int_{\mathbb{R}^d} k(\|v\|) \|v\|^2 \, dv\right) \Delta_\mathcal{M} f(x_i)
	\end{align}

	Choosing the kernel $k$ such that $\int_{\mathbb{R}^d} k(\|v\|) \|v\|^2 \, dv = 2d$, and combining the terms from Steps~6 and~7:
	\[
		\lim_{n \to \infty} \frac{1}{\epsilon^2} L_n f(x_i) = p(x_i) \Delta_\mathcal{M} f(x_i) + \langle \nabla p(x_i), \nabla f(x_i) \rangle
	\]

	Dividing by $p(x_i)$ and using that $\nabla p / p = \nabla \log p$:
	\[
		\lim_{n \to \infty} \frac{1}{p(x_i)\,\epsilon^2} L_n f(x_i) = \Delta_\mathcal{M} f(x_i) + \langle \nabla \log p(x_i), \nabla f(x_i) \rangle
	\]

	The term $\langle \nabla \log p, \nabla f \rangle$ reflects the influence of the non-uniform sampling density and can be eliminated with more sophisticated constructions of the Laplacian \citep{coifman2006diffusion, hein2007graph}.

	\textbf{Technical conditions:} The convergence requires that $\epsilon \to 0$ (local radius tends to zero) but $n\epsilon^{d+2} \to \infty$ (enough points in each neighborhood). The optimal rate is:
	\[
		\epsilon \sim \left(\frac{\log n}{n}\right)^{1/(d+2+\alpha)}
	\]
	for some small $\alpha > 0$, balancing the bias of the approximation with the variance of the sampling.
\end{proof}

\paragraph{Implications for Geometric Deep Learning}

This fundamental theoretical result has profound consequences for \gls{GDL}:

\begin{enumerate}
	\item \textbf{Theoretical justification of GNNs}: Graphs can faithfully capture the geometry of manifolds when we have enough data. This justifies the use of \glspl{GNN} for data that naturally reside on manifolds.

	\item \textbf{Guidance for graph construction}: The conditions of the theorem ($\epsilon \to 0$, $n\epsilon^{d+2} \to \infty$) provide principles for choosing the neighborhood radius $\epsilon$ as a function of the number of samples $n$ and the intrinsic dimension $d$.

	\item \textbf{Fundamental limitations}: The data are finite, so we only have approximations. The graph construction (choice of $\epsilon$, type of kernel) affects the quality of the approximation. Information about the smooth geometry between sampled points is lost.

	\item \textbf{Computational trade-off}: Working directly with continuous manifolds is mathematically elegant but computationally intractable. Graphs offer a practical compromise between geometric expressiveness and computational efficiency.

	\item \textbf{Architecture design}: \glspl{GNN} that approximate spectral filters of the Laplacian inherit geometric properties of the continuous domain, motivating architectures such as ChebNet and GCN.
\end{enumerate}

The unifying perspective is clear: \textbf{graphs are an intermediate step between Euclidean representations (which ignore geometry) and continuous manifolds (which are mathematically elegant but computationally costly)}. Both graphs and manifolds preserve the intrinsic structure of the data better than traditional vector representations, but with different trade-offs between geometric precision and computational feasibility.

\begin{remark}[Convergence rates and scaling regimes]\label{rem:tasas_convergencia_laplaciano}
	The convergence of the graph Laplacian to the Laplace-Beltrami operator depends critically on the relationship between $n$ (number of points) and $\epsilon$ (connectivity radius). Three regimes are identified:

	\begin{center}
	\small
	\begin{tabular}{@{}llp{5cm}@{}}
		\toprule
		\textbf{Regime} & \textbf{Condition} & \textbf{Behavior} \\
		\midrule
		Underconnected & $n\epsilon^d \to 0$ & Disconnected graph; no convergence \\
		\addlinespace
		Optimal & $n\epsilon^d \sim \log n$ & Convergence with rate $O(\epsilon^2 + 1/\sqrt{n\epsilon^d})$ \\
		\addlinespace
		Overconnected & $n\epsilon^d \to \infty$ very fast & Convergence, but loss of locality \\
		\bottomrule
	\end{tabular}
	\end{center}

	\textbf{Numerical example (sphere $S^2$, $d=2$):} For $n = 10000$ points uniformly distributed on $S^2$:
	\begin{itemize}
		\item $\epsilon = 0.05$: Error in the first nontrivial eigenvalue $\approx 15\%$ (underconnected)
		\item $\epsilon = 0.15$: Error $\approx 2\%$ (optimal regime)
		\item $\epsilon = 0.40$: Error $\approx 5\%$ (overconnected, loses locality)
	\end{itemize}
	The optimal choice $\epsilon^* \approx (\log n / n)^{1/4} \approx 0.14$ coincides with the theoretical predictions.
\end{remark}

\subsection{Spectral Analysis on Graphs}\label{subsec:analisis_espectral_grafos_gdl}

\textbf{Graph spectral analysis} extends the ideas of classical Fourier analysis to the discrete domain (the foundations of harmonic analysis and the Fourier transform are developed in \autoref{sec:analisis_armonico} of \autoref{ch:cnn}). The eigenvectors of the Laplacian form an orthogonal basis analogous to the sinusoidal waves in Euclidean Fourier analysis, allowing functions to be decomposed into ``frequencies'' over the graph.

\paragraph{Graph Fourier Transform}

\begin{definition}[Graph Fourier transform]\label{def:transformada_fourier_grafos}
	Let $G = (V,E)$ be a graph with normalized Laplacian $\mathcal{L}$ and spectral decomposition $\mathcal{L} = U \Lambda U^T$ where:
	\begin{itemize}
		\item $\Lambda = \text{diag}(\lambda_1, \ldots, \lambda_n)$ with $0 = \lambda_1 \leq \lambda_2 \leq \cdots \leq \lambda_n \leq 2$
		\item $U = [u_1, \ldots, u_n]$ with $u_i \in \mathbb{R}^n$ orthonormal eigenvectors
	\end{itemize}

	The \textit{graph Fourier transform} of a function $f: V \to \mathbb{R}$ (viewed as $f \in \mathbb{R}^n$) is:
	\[
		\hat{f} = U^T f \in \mathbb{R}^n
	\]

	The \textit{inverse transform} is:
	\[
		f = U \hat{f} = \sum_{i=1}^n \hat{f}(i) \, u_i
	\]
\end{definition}

This transformation decomposes $f$ into the ``harmonic bases'' $\{u_i\}$ of the graph, where $\hat{f}(i)$ represents the ``frequency coefficient'' associated with the eigenvalue $\lambda_i$. Note the direct analogy with \autoref{def:transformada_fourier_continua}: the eigenvectors $u_i$ of the Laplacian play the role of the complex exponentials $e^{2\pi i \omega \cdot x}$, and the eigenvalues $\lambda_i$ that of the frequencies $\omega$.

\begin{figure}[htbp]
\centering
\includegraphics[width=0.95\textwidth]{figures/laplacian_eigenvectors}
\caption[Graph Laplacian eigenvectors]{Eigenvectors of the Laplacian on a grid graph. Small eigenvalues correspond to smooth functions (low frequency), while large eigenvalues correspond to oscillatory functions (high frequency).}
\label{fig:laplacian_eigenvectors}
\end{figure}

\begin{theorem}[Frequency interpretation of eigenvectors]\label{thm:interpretacion_frecuencial_autovectores}
	The eigenvectors of the Laplacian $\{u_i\}$ satisfy:
	\begin{enumerate}
		\item \textbf{Small eigenvalues $\Leftrightarrow$ smooth functions}: If $\mathcal{L} u_i = \lambda_i u_i$ with $\lambda_i$ small, then:
		\[
			u_i^T \mathcal{L} u_i = \lambda_i \|u_i\|^2 = \lambda_i
		\]
		is small, which by \autoref{prop:propiedades_algebraicas_laplaciano} implies that $u_i$ varies smoothly over the graph.

		\item \textbf{Large eigenvalues $\Leftrightarrow$ oscillatory functions}: If $\lambda_i \approx 2$, then $u_i$ oscillates rapidly between positive and negative values at adjacent nodes.

		\item \textbf{Analogy with frequencies}: $\lambda_i$ plays the role of ``frequency squared'': $\lambda_i \sim \omega_i^2$ in classical Fourier analysis (\autoref{sec:analisis_armonico}).
	\end{enumerate}
\end{theorem}

\begin{proof}
	\textbf{(1)}: Analogously to the quadratic form of the unnormalized Laplacian (\autoref{prop:propiedades_algebraicas_laplaciano}), for the normalized Laplacian $\mathcal{L} = D^{-1/2} L_G D^{-1/2}$ we have:
	\[
		u_i^T \mathcal{L} u_i = \frac{1}{2} \sum_{(j,k) \in E} W_{jk} \left(\frac{u_i(j)}{\sqrt{d_j}} - \frac{u_i(k)}{\sqrt{d_k}}\right)^2
	\]

	If $\lambda_i$ is small, this sum must be small, which requires $|u_i(j) - u_i(k)|$ to be small for edges $(j,k)$, that is, $u_i$ is smooth.

	\textbf{(2)}: For the normalized Laplacian, $\lambda_{\max} \leq 2$. Values close to 2 are attained when $u_i$ changes sign frequently between neighbors, maximizing $(u_i(j) - u_i(k))^2$.

	\textbf{(3)}: In the Euclidean Laplacian $\nabla^2$, the eigenfunctions are plane waves $e^{i\omega \cdot x}$ with eigenvalues $-\omega^2$. Changing sign (convention), $\lambda \sim \omega^2$.
\end{proof}

\begin{figure}[htbp]
\centering
\includegraphics[width=0.95\textwidth]{figures/graph_fourier_transform}
\caption[Graph Fourier transform]{Graph Fourier transform. Top: a signal on a cyclic graph, its Fourier coefficients, and the frequency spectrum. Bottom: partial reconstruction using different numbers of modes, showing how the low-frequency components capture smooth variations.}
\label{fig:graph_fourier_transform}
\end{figure}

\paragraph{Spectral Filters}

The graph Fourier transform allows defining \textbf{spectral filters}, fundamental for spectral \glspl{GNN}:

\begin{definition}[Spectral filter on graphs]\label{def:filtro_espectral_grafo}
	A \textit{spectral filter} is a linear operator $g_\theta: \mathbb{R}^n \to \mathbb{R}^n$ defined via a transfer function $g_\theta: \mathbb{R} \to \mathbb{R}$ applied to the eigenvalues of the Laplacian:
	\[
		g_\theta(\mathcal{L}) = U \, g_\theta(\Lambda) \, U^T = U \, \text{diag}(g_\theta(\lambda_1), \ldots, g_\theta(\lambda_n)) \, U^T
	\]

	The \textit{spectral convolution} of a signal $f$ with the filter $g_\theta$ is:
	\[
		f \star_G g_\theta = g_\theta(\mathcal{L}) f = U \, g_\theta(\Lambda) \, U^T f
	\]
\end{definition}

This definition generalizes the Euclidean convolution, where the convolution theorem establishes that convolution in the spatial domain is equivalent to multiplication in the Fourier domain (see \autoref{subsec:convoluciones_espectrales_cnns} of \autoref{ch:cnn} for the detailed treatment in the Euclidean case). The \textbf{Spectral Graph Convolutional Networks} \citep{bruna2014spectralnetworkslocallyconnected} use this idea, parametrizing $g_\theta(\lambda)$ as a learnable function of the frequencies; for the complete mathematical development of spectral convolutions on graphs, including polynomial filters and the Chebyshev approximation (ChebNet), see \autoref{subsec:analisis_espectral_grafos} of \autoref{ch:cnn}.

\paragraph{Connection with Continuous Harmonic Analysis}

The analogy between spectral analysis on graphs and harmonic analysis on manifolds is precise in the limit:

\begin{table}[htbp]
\centering
\begin{tabular}{|l|l|}
\hline
\textbf{Continuous Manifold} & \textbf{Discrete Graph} \\
\hline
Laplace-Beltrami $\Delta_\mathcal{M}$ & Graph Laplacian $L_G$ \\
Eigenfunctions $\Delta_\mathcal{M} \phi_i = \lambda_i \phi_i$ & Eigenvectors $L_G u_i = \lambda_i u_i$ \\
Fourier transform $\hat{f}(\lambda) = \langle f, \phi_\lambda \rangle$ & $\hat{f} = U^T f$ \\
Smooth functions (low frequency) & Eigenvectors with $\lambda_i$ small \\
Oscillatory functions (high frequency) & Eigenvectors with $\lambda_i$ large \\
\hline
\end{tabular}
\caption[Analogy between spectral analysis on manifolds and graphs]{Analogy between spectral analysis on Riemannian manifolds and discrete graphs.}
\label{tab:analogia_espectral_variedades_grafos}
\end{table}

This correspondence justifies the use of spectral techniques in \glspl{GNN} and connects graph learning with the theory of signals on manifolds.
\section{Geodesics: Optimal Paths and Information Propagation}\label{sec:geodesicas_propagacion}

The preceding sections have established two pillars of \gls{GDL}: the \textbf{geometries} where data reside and the \textbf{graphs} that discretize them. In particular, \autoref{thm:convergencia_grafos_variedades_detallado} proved that the Laplacian of a proximity graph converges to the Laplace-Beltrami operator $\Delta_\mathcal{M}$ of the underlying manifold, and \autoref{subsec:analisis_espectral_grafos_gdl} exploited this connection to define Fourier transforms and spectral filters on graphs. However, the operator $\Delta_\mathcal{M}$ presupposes that the manifold is endowed with a \textit{Riemannian metric}, a concept we have not yet formalized: the differentiable manifolds introduced in \autoref{subsec:fundamentos_variedades} possess topological and differential structure, but lack an intrinsic notion of distance.

\textit{Geodesics} fill this gap. In a Euclidean space, the ``shortest path'' between two points is trivially a straight line; on a curved manifold, geodesics generalize this notion as curves that locally minimize length. They are the piece that connects the three preceding G's: they provide the intrinsic distance that justifies the weights of proximity graphs, define the local coordinate systems needed for geometric convolutions, and determine how to propagate information in a manner compatible with the curvature of the domain. To formalize these ideas, we first endow differentiable manifolds with metric structure (\autoref{subsec:variedades_riemannianas}), introduce the connections and curvature that measure how the geometry deviates from the flat case (\autoref{subsec:conexiones_curvatura}), define geodesics as solutions of a differential equation (\autoref{subsec:geodesicas_variedades}), and study their discretization on graphs (\autoref{subsec:geodesicas_grafos}).

\subsection{Riemannian Manifolds}\label{subsec:variedades_riemannianas}

To define geodesics rigorously, we need to endow differentiable manifolds (introduced in \autoref{subsec:fundamentos_variedades}) with metric structure. A \textit{Riemannian metric} provides the machinery needed to measure lengths, angles, and volumes intrinsically.

\begin{definition}[Riemannian metric]\label{def:metrica_riemanniana}
A \textit{Riemannian metric} on a differentiable manifold $\mathcal{M}$ is a smooth assignment that associates to each point $p \in \mathcal{M}$ an inner product $g_p: T_p\mathcal{M} \times T_p\mathcal{M} \to \mathbb{R}$ on the tangent space. In local coordinates $(x^1, \ldots, x^m)$, it is expressed via the dual basis $\{dx^i\}$ (\autoref{def:diferencial_funcion}) and the tensor product of covectors (\autoref{def:producto_tensorial_covectores}) as:
\[
g_p = \sum_{i,j=1}^m g_{ij}(p) \, dx^i \otimes dx^j
\]
where $g_{ij}(p)$ are smooth functions that form a symmetric positive-definite matrix at each point.
\end{definition}

\begin{definition}[Riemannian manifold]\label{def:variedad_riemanniana}
A \textit{Riemannian manifold} is the pair $(\mathcal{M}, g)$ where $\mathcal{M}$ is a differentiable manifold and $g$ is a Riemannian metric on $\mathcal{M}$.
\end{definition}

\begin{example}[Metrics on common manifolds]\label{ex:metricas_variedades}
\leavevmode
\begin{itemize}
	\item \textbf{Euclidean space} $\mathbb{R}^n$: The standard metric is $g_{ij} = \delta_{ij}$ (identity), giving $ds^2 = dx_1^2 + \cdots + dx_n^2$.

	\item \textbf{Sphere} $S^2$: In spherical coordinates $(\theta, \phi)$, the metric induced by the immersion in $\mathbb{R}^3$ is $ds^2 = d\theta^2 + \sin^2\theta \, d\phi^2$.

	\item \textbf{Hyperbolic space} $\mathbb{H}^n$ (Poincaré model): On the disk $\mathbb{D}^n = \{x \in \mathbb{R}^n : \|x\| < 1\}$, the metric is $g_{ij} = \frac{4\delta_{ij}}{(1-\|x\|^2)^2}$, which ``stretches'' distances near the boundary.
\end{itemize}
\end{example}

The Riemannian metric provides all the geometric structure needed to measure lengths and distances on a manifold, through a three-step conceptual chain. First, since $g_p$ is an inner product on each tangent space $T_p\mathcal{M}$, it induces a pointwise norm $\|v\|_g = \sqrt{g_p(v,v)}$ for all $v \in T_p\mathcal{M}$. For a curve $\gamma(t)$ on the manifold, this norm measures the ``instantaneous velocity'' $\|\dot\gamma(t)\|_g$ at each point of the curve. Second, integrating this instantaneous velocity along the entire curve yields the total length $L(\gamma)$. Third, taking the infimum of the lengths over all smooth curves joining two points $p$ and $q$ yields the geodesic distance $d_g(p,q)$. We now formalize these concepts:

\begin{definition}[Length of curves]\label{def:longitud_curvas}
For a curve $\gamma: [a,b] \to \mathcal{M}$, its \textit{length} is defined as:
\[
L(\gamma) = \int_a^b \sqrt{g_{\gamma(t)}(\dot{\gamma}(t), \dot{\gamma}(t))} \, dt
\]
\end{definition}

\begin{definition}[Geodesic distance]\label{def:distancia_geodesica}
The \textit{geodesic distance} between two points $p, q \in \mathcal{M}$ is:
\[
d_g(p,q) = \inf\{L(\gamma) : \gamma \text{ is a smooth curve from } p \text{ to } q\}
\]
This function $d_g: \mathcal{M} \times \mathcal{M} \to \mathbb{R}_{\geq 0}$ is a metric in the topological sense (see \autoref{prop:dg_metrica}).
\end{definition}

\begin{example}[Geodesic distance on the sphere $S^2$]\label{ex:longitud_esfera}
Consider the unit sphere $S^2 \subset \mathbb{R}^3$ with the metric $ds^2 = d\theta^2 + \sin^2\theta\,d\phi^2$ from \autoref{ex:metricas_variedades}. Let $p = (0,0,1)$ (north pole) and $q = (\sin\alpha, 0, \cos\alpha)$ for some $\alpha \in (0, \pi)$. The great-circle arc joining them is parametrized as
\[
\gamma(t) = (\sin t, 0, \cos t), \quad t \in [0, \alpha],
\]
which in spherical coordinates corresponds to $\theta(t) = t$, $\phi(t) = 0$. The tangent vector is $\dot\theta(t) = 1$, $\dot\phi(t) = 0$, so that
\[
g_{\gamma(t)}(\dot\gamma, \dot\gamma) = \dot\theta^2 + \sin^2\theta\,\dot\phi^2 = 1.
\]
Integrating, we obtain $L(\gamma) = \int_0^\alpha \sqrt{1}\,dt = \alpha$. Since the great-circle arcs are geodesics of $S^2$ (they minimize length between their endpoints), we conclude that
\[
d_g(p, q) = \alpha = \arccos\langle p, q \rangle,
\]
where $\langle \cdot, \cdot \rangle$ denotes the Euclidean inner product of $\mathbb{R}^3$.
\end{example}

\begin{proposition}[The geodesic distance is a metric]\label{prop:dg_metrica}
Let $(\mathcal{M}, g)$ be a connected Riemannian manifold. The function $d_g: \mathcal{M} \times \mathcal{M} \to \mathbb{R}_{\geq 0}$ defined by the infimum of curve lengths is a metric in the topological sense: for all $p, q, r \in \mathcal{M}$,
\begin{enumerate}
	\item $d_g(p,q) \geq 0$, with $d_g(p,q) = 0$ if and only if $p = q$;
	\item $d_g(p,q) = d_g(q,p)$ (symmetry);
	\item $d_g(p,r) \leq d_g(p,q) + d_g(q,r)$ (triangle inequality).
\end{enumerate}
\end{proposition}

\begin{proof}
\textbf{Non-degeneracy.} If $p = q$, the constant curve $\gamma(t) = p$ has zero length, so $d_g(p,p) = 0$. Conversely, suppose $p \neq q$ and let us show that $d_g(p,q) > 0$. Let $(U, \varphi)$ be a chart with $\varphi(p) = 0 \in \mathbb{R}^m$ and choose $r > 0$ such that $\overline{B_r(0)} \subset \varphi(U)$ and $q \notin \varphi^{-1}(\overline{B_r(0)})$. On the compact set $K = \varphi^{-1}(\overline{B_r(0)})$, the Riemannian metric $g$ is expressed in coordinates as $g_x(v,v) = \sum_{i,j} g_{ij}(\varphi(x))\, v^i v^j$. The function $(x,v) \mapsto g_x(v,v)$ is continuous and strictly positive for $v \neq 0$, so restricted to the compact set $K \times S^{m-1}_{\text{euc}}$ it attains a minimum $c > 0$. By homogeneity,
\[
g_x(v,v) \geq c\,\|v\|_{\text{euc}}^2 \quad \text{for all } x \in K,\; v \in T_x\mathcal{M}.
\]
Every smooth curve $\gamma: [a,b] \to \mathcal{M}$ from $p$ to $q$ must leave $K$, that is, there exists $t^* \in [a,b]$ such that $\gamma([a,t^*]) \subset K$ and $\varphi(\gamma(t^*)) \in \partial B_r(0)$. Its length satisfies
\[
L(\gamma) \geq \int_a^{t^*} \|\dot\gamma(t)\|_g\, dt \geq \int_a^{t^*} \sqrt{c}\,\|\dot\gamma(t)\|_{\text{euc}}\, dt \geq \sqrt{c}\,\|\varphi(\gamma(t^*)) - \varphi(\gamma(a))\|_{\text{euc}} = \sqrt{c}\, r,
\]
where the last inequality uses that the Euclidean length of a curve is an upper bound for the distance between its endpoints. Since this bound $\sqrt{c}\, r > 0$ is independent of $\gamma$, we conclude $d_g(p,q) \geq \sqrt{c}\, r > 0$.

\textbf{Symmetry.} Given a curve $\gamma: [a,b] \to \mathcal{M}$ from $p$ to $q$, we define its reverse reparametrization $\bar\gamma(t) = \gamma(a+b-t)$, which goes from $q$ to $p$. Then $\dot{\bar\gamma}(t) = -\dot\gamma(a+b-t)$, so $\|\dot{\bar\gamma}(t)\|_g = \|\dot\gamma(a+b-t)\|_g$. Via the change of variable $s = a+b-t$,
\[
L(\bar\gamma) = \int_a^b \|\dot{\bar\gamma}(t)\|_g\, dt = \int_a^b \|\dot\gamma(s)\|_g\, ds = L(\gamma).
\]
The map $\gamma \mapsto \bar\gamma$ establishes a bijection between the curves from $p$ to $q$ and those from $q$ to $p$ that preserves lengths, hence both infima coincide.

\textbf{Triangle inequality.} For any $\varepsilon > 0$, there exist curves $\gamma_1$ from $p$ to $q$ and $\gamma_2$ from $q$ to $r$ with $L(\gamma_1) < d_g(p,q) + \varepsilon$ and $L(\gamma_2) < d_g(q,r) + \varepsilon$. The concatenation $\gamma_1 \cdot \gamma_2$ is a curve from $p$ to $r$ with $L(\gamma_1 \cdot \gamma_2) = L(\gamma_1) + L(\gamma_2)$, hence $d_g(p,r) \leq L(\gamma_1) + L(\gamma_2) < d_g(p,q) + d_g(q,r) + 2\varepsilon$. Since $\varepsilon > 0$ is arbitrary, the result follows.
\end{proof}

\begin{definition}[Angles on manifolds]\label{def:angulos_variedades}
The \textit{angle} between two vectors $v, w \in T_p\mathcal{M}$ is determined by:
\[
\cos\theta = \frac{g_p(v,w)}{\sqrt{g_p(v,v)}\sqrt{g_p(w,w)}}
\]
\end{definition}

\begin{definition}[Volume element]\label{def:elemento_volumen}
The \textit{volume element} induced by the metric $g$ is expressed in local coordinates via the exterior product (\autoref{def:producto_exterior_k_formas}) as:
\[
dV_g = \sqrt{\det(g_{ij})} \, dx^1 \wedge \cdots \wedge dx^m
\]
This element is a differential $m$-form that allows integrating functions over the manifold intrinsically.
\end{definition}

\begin{definition}[Euclidean Laplacian]\label{def:laplaciano_euclidiano}
Let $f: \mathbb{R}^m \to \mathbb{R}$ be a function of class $C^2$. The \textit{Euclidean Laplacian} (or Laplace operator) is the second-order differential operator
\[
\nabla^2 f = \sum_{i=1}^{m} \frac{\partial^2 f}{\partial (x^i)^2},
\]
equivalently $\nabla^2 f = \operatorname{div}(\operatorname{grad} f)$. It measures the deviation of the value of $f$ at a point from its local average: $\nabla^2 f(p) > 0$ indicates that $f(p)$ is below the average in a neighborhood of $p$, and $\nabla^2 f(p) < 0$ indicates that it is above.
\end{definition}

\begin{definition}[Laplace-Beltrami operator]\label{def:laplace_beltrami}
For a Riemannian manifold $(\mathcal{M}, g)$, the \textit{Laplace-Beltrami operator} $\Delta_\mathcal{M}$ acts on smooth functions $f: \mathcal{M} \to \mathbb{R}$ as:
\[
\Delta_\mathcal{M} f = \text{div}(\text{grad } f) = \frac{1}{\sqrt{\det(g)}} \partial_i\left(\sqrt{\det(g)} \, g^{ij} \partial_j f\right)
\]
where $g_{ij}$ are the components of the Riemannian metric and $g^{ij}$ are the components of the inverse metric. This operator generalizes the Euclidean Laplacian (\autoref{def:laplaciano_euclidiano}) to arbitrary manifolds: when $(\mathcal{M}, g) = (\mathbb{R}^m, \delta_{ij})$ we have $\det(g) = 1$ and $g^{ij} = \delta^{ij}$, so that the previous expression reduces to $\Delta_{\mathbb{R}^m} f = \sum_{i=1}^{m} \frac{\partial^2 f}{\partial (x^i)^2} = \nabla^2 f$.
\end{definition}

\begin{remark}[Role of the Laplace-Beltrami operator in GDL]\label{rem:laplace_beltrami_gdl}
The operator $\Delta_\mathcal{M}$ is the central object that connects differential geometry with geometric deep learning, for the following reasons:

\begin{enumerate}
    \item \textbf{Diffusion on manifolds.} $\Delta_\mathcal{M}$ governs the heat equation
    \[
    \frac{\partial u}{\partial t} = \Delta_\mathcal{M} u,
    \]
    which models the propagation of information over the manifold. In GDL, the diffusion of features in a neural network over $\mathcal{M}$ is analogous to this process: each layer ``diffuses'' the signal toward geodesic neighborhoods, in a manner controlled by the local geometry.

    \item \textbf{Intrinsic spectral analysis.} The eigenfunctions $\{\varphi_k\}_{k \geq 0}$ of $\Delta_\mathcal{M}$ (with $\Delta_\mathcal{M} \varphi_k = \lambda_k \varphi_k$) form an orthonormal basis of $L^2(\mathcal{M})$, generalizing the classical Fourier basis. This allows defining an intrinsic Fourier transform $\hat{f}(\lambda_k) = \langle f, \varphi_k \rangle$ and, from it, spectral convolutions over the manifold, in perfect analogy with spectral filters on graphs (\autoref{def:filtro_espectral_grafo}).

    \item \textbf{Invariance under isometries.} For every isometry $\phi: \mathcal{M} \to \mathcal{M}$, the Laplace-Beltrami operator satisfies
    \[
    \Delta_\mathcal{M}(f \circ \phi) = (\Delta_\mathcal{M} f) \circ \phi,
    \]
    that is, $\Delta_\mathcal{M}$ commutes with isometries. This makes it a canonical operator compatible with the principle of geometric equivariance: any architecture built from $\Delta_\mathcal{M}$ automatically inherits invariance under the symmetries of the manifold.

    \item \textbf{Discretization and connection with graphs.} The graph Laplacian $L_G$ (\autoref{def:laplaciano_grafo}) is precisely the discrete analogue of $\Delta_\mathcal{M}$, and under appropriate sampling conditions it converges to the continuous operator (\autoref{thm:convergencia_grafos_variedades_detallado}). \autoref{tab:analogia_espectral_variedades_grafos} summarizes this correspondence in detail.
\end{enumerate}
\end{remark}

\subsection{Connections and Curvature}\label{subsec:conexiones_curvatura}

To define geodesics and quantify the deviation of a manifold from Euclidean space, we need two fundamental tools. First, a \textit{connection} provides a mechanism for differentiating vector fields and transporting vectors along curves: in $\mathbb{R}^n$, we can compare vectors at distinct points directly because all tangent spaces are canonically identified with $\mathbb{R}^n$; on a general manifold, $T_p\mathcal{M}$ and $T_q\mathcal{M}$ are distinct vector spaces and there is no canonical way to ``transport'' a vector from one to the other. Second, the \textit{curvature tensor} quantifies the non-commutativity of the resulting covariant derivatives, measuring the intrinsic deviation from flat geometry.

\subsubsection{Affine connection and Levi-Civita connection}\label{subsec:conexion_levi_civita}

\begin{definition}[Affine connection]\label{def:conexion_afin}
Let $\mathcal{M}$ be a differentiable manifold. An \textit{affine connection} (or covariant connection) on $\mathcal{M}$ is a map\footnote{The notation $\Gamma(T\mathcal{M})$ denotes the space of \textit{smooth vector fields} on $\mathcal{M}$, that is, smooth assignments $X\colon p \mapsto X(p) \in T_p\mathcal{M}$ that associate a tangent vector to each point. This notation will be formalized in \autoref{subsec:fibrados_principales} as the space of sections of the tangent bundle.}
\[
\nabla: \Gamma(TM) \times \Gamma(TM) \to \Gamma(TM), \quad (X, Y) \mapsto \nabla_X Y
\]
satisfying the following properties for any vector fields $X, Y, Z \in \Gamma(TM)$ and functions $f, g \in C^\infty(\mathcal{M})$:
\begin{enumerate}
	\item \textbf{$C^\infty(\mathcal{M})$-linearity in $X$}: $\nabla_{fX + gY} Z = f\nabla_X Z + g\nabla_Y Z$
	\item \textbf{$\mathbb{R}$-linearity in $Y$}: $\nabla_X(Y + Z) = \nabla_X Y + \nabla_X Z$
	\item \textbf{Leibniz rule in $Y$}: $\nabla_X(fY) = (Xf)Y + f\nabla_X Y$
\end{enumerate}
\end{definition}

\begin{definition}[Torsion tensor]\label{def:tensor_torsion}
Let $\nabla$ be an affine connection on $\mathcal{M}$. The \textit{torsion tensor} of $\nabla$ is the map $T\colon \Gamma(TM) \times \Gamma(TM) \to \Gamma(TM)$ defined by:
\[
T(X,Y) = \nabla_X Y - \nabla_Y X - [X,Y]
\]
Unlike the covariant derivative $\nabla_X Y$ (which is $\mathbb{R}$-linear, but not $C^\infty(\mathcal{M})$-linear, in $Y$), the torsion tensor $T$ is $C^\infty(\mathcal{M})$-bilinear: for every function $f \in C^\infty(\mathcal{M})$,
\[
T(fX, Y) = fT(X,Y) = T(X, fY).
\]
This guarantees that $T$ is a genuine tensor: its value $T(X,Y)|_p$ depends only on $X(p)$ and $Y(p)$, not on the values of $X$ and $Y$ at neighboring points.

A connection is said to be \textit{torsion-free} (or symmetric) if $T \equiv 0$, that is, if $\nabla_X Y - \nabla_Y X = [X,Y]$ for any fields $X, Y$. Geometrically, $T$ measures the closure defect of infinitesimal parallelograms: when transporting two vectors $X, Y$ one along the other, the difference between the endpoints is proportional to $T(X,Y)$. For coordinate fields $\partial_i, \partial_j$ (where $[\partial_i, \partial_j] = 0$), the condition $T = 0$ is equivalent to the symmetry of the Christoffel symbols: $\Gamma^k_{ij} = \Gamma^k_{ji}$.
\end{definition}

A manifold admits infinitely many affine connections. However, when the manifold is endowed with a Riemannian metric $g$, there exists a unique connection that is simultaneously compatible with $g$ and torsion-free (\autoref{def:tensor_torsion}). This is the Levi-Civita connection, whose existence and uniqueness constitute the fundamental theorem of Riemannian geometry.

\begin{definition}[Levi-Civita connection]\label{def:conexion_levi_civita_gdl}
Let $(\mathcal{M}, g)$ be a Riemannian manifold. The \textit{Levi-Civita connection} is the unique connection $\nabla$ on $T\mathcal{M}$ that satisfies:
\begin{enumerate}
	\item \textbf{Compatibility with the metric}: For any vector fields $X, Y, Z$:
	\[
	X(g(Y,Z)) = g(\nabla_X Y, Z) + g(Y, \nabla_X Z)
	\]
	\item \textbf{Symmetry (zero torsion)}: $\nabla_X Y - \nabla_Y X = [X, Y]$
\end{enumerate}
In local coordinates, it is determined by the \textit{Christoffel symbols}:
\[
\Gamma^k_{ij} = \frac{1}{2} g^{kl}\left(\frac{\partial g_{il}}{\partial x^j} + \frac{\partial g_{jl}}{\partial x^i} - \frac{\partial g_{ij}}{\partial x^l}\right)
\]
where $g^{kl}$ denotes the inverse matrix of $g_{ij}$.
\end{definition}

\begin{theorem}[Existence and uniqueness of the Levi-Civita connection]\label{thm:existencia_levi_civita}
Let $(\mathcal{M}, g)$ be a Riemannian manifold. There exists a unique affine connection $\nabla$ on $\mathcal{M}$ that satisfies:
\begin{enumerate}
	\item \textbf{Compatibility with the metric}: $Xg(Y,Z) = g(\nabla_X Y, Z) + g(Y, \nabla_X Z)$ for all $X, Y, Z \in \Gamma(TM)$.
	\item \textbf{Zero torsion}: $\nabla_X Y - \nabla_Y X = [X,Y]$ for all $X, Y \in \Gamma(TM)$.
\end{enumerate}
\end{theorem}

\begin{proof}
The proof proceeds in two steps: uniqueness (which fixes the formula) and existence (which verifies the axioms).

\textbf{Uniqueness.} Suppose $\nabla$ satisfies both properties. We cyclically expand the metric compatibility:
\begin{align*}
Xg(Y,Z) &= g(\nabla_X Y, Z) + g(Y, \nabla_X Z), \\
Yg(Z,X) &= g(\nabla_Y Z, X) + g(Z, \nabla_Y X), \\
Zg(X,Y) &= g(\nabla_Z X, Y) + g(X, \nabla_Z Y).
\end{align*}
Adding the first two and subtracting the third, and using the zero-torsion condition $\nabla_X Y - \nabla_Y X = [X,Y]$ to eliminate the cross terms, we obtain the \textit{Koszul formula}:
\begin{equation}\label{eq:koszul}
2g(\nabla_X Y, Z) = Xg(Y,Z) + Yg(Z,X) - Zg(X,Y) + g([X,Y],Z) - g([Y,Z],X) + g([Z,X],Y).
\end{equation}
Since $g$ is non-degenerate, this formula determines $\nabla_X Y$ uniquely.

\textbf{Existence.} Let $F(X,Y,Z)$ denote the right-hand side of~\eqref{eq:koszul}. We define $\nabla_X Y$ as the unique vector field such that $2g(\nabla_X Y, Z) = F(X,Y,Z)$ for all $Z \in \Gamma(TM)$. Well-definedness follows from the fact that $F(X,Y,\cdot)$ is $C^\infty(\mathcal{M})$-linear in $Z$ (a direct verification using $[X, fZ] = f[X,Z] + (Xf)Z$, $[fZ, X] = f[Z,X] - (Xf)Z$ and the symmetry of $g$ to cancel the terms in $Xf$ and $Yf$), so that it defines a $1$-form; by the non-degeneracy of $g$, there exists a unique vector field $\nabla_X Y$ representing it. We verify the five required properties.

\textit{(1) $C^\infty(\mathcal{M})$-linearity in $X$.} Substituting $X$ by $fX$ in $F$, we expand each term using $(fX)g(Y,Z) = f\,Xg(Y,Z)$, the Leibniz rules $Yg(Z,fX) = (Yf)g(Z,X) + f\,Yg(Z,X)$ and $Zg(fX,Y) = (Zf)g(X,Y) + f\,Zg(X,Y)$, and the identities $[fX,Y] = f[X,Y] - (Yf)X$, $[Z,fX] = f[Z,X] + (Zf)X$. The terms with factor $f$ reproduce $f \cdot F(X,Y,Z)$, while the remaining ones vanish by symmetry of $g$:
\[
(Yf)\bigl[g(Z,X) - g(X,Z)\bigr] + (Zf)\bigl[g(X,Y) - g(X,Y)\bigr] = 0.
\]
The $\mathbb{R}$-additivity $F(X_1+X_2,Y,Z) = F(X_1,Y,Z) + F(X_2,Y,Z)$ is immediate, whence $\nabla_{fX+gY}Z = f\nabla_X Z + g\nabla_Y Z$.

\textit{(2) $\mathbb{R}$-linearity in $Y$.} The additivity $F(X, Y_1+Y_2, Z) = F(X,Y_1,Z) + F(X,Y_2,Z)$ follows from the bilinearity of $g$ and of the Lie bracket in each term of~\eqref{eq:koszul}.

\textit{(3) Leibniz rule in $Y$.} Substituting $Y$ by $fY$ and using $[X,fY] = f[X,Y] + (Xf)Y$ and $[fY,Z] = f[Y,Z] - (Zf)Y$, the terms with factor $f$ reproduce $f \cdot F(X,Y,Z)$ and the additional terms are
\[
(Xf)g(Y,Z) + (Xf)g(Y,Z) + (Zf)\bigl[g(Y,X) - g(X,Y)\bigr] = 2(Xf)\,g(Y,Z),
\]
where the first $(Xf)g(Y,Z)$ comes from $Xg(fY,Z)$ and the second from $g([X,fY],Z)$, while the terms in $Zf$ cancel by symmetry of $g$. Thus, $2g(\nabla_X(fY),Z) = 2f\,g(\nabla_X Y, Z) + 2(Xf)\,g(Y,Z) = 2g(f\nabla_X Y + (Xf)Y, Z)$ for all $Z$, and the non-degeneracy of $g$ gives $\nabla_X(fY) = (Xf)Y + f\nabla_X Y$.

\textit{(4) Zero torsion.} We compute $F(X,Y,Z) - F(Y,X,Z)$. By the symmetry of $g$, the six terms in Lie derivatives cancel in pairs (for example, $Xg(Y,Z) - Xg(Z,Y) = 0$). For the bracket terms, the antisymmetry $[Y,X] = -[X,Y]$ gives $g([X,Y],Z) - g([Y,X],Z) = 2g([X,Y],Z)$, and the four remaining ones vanish: $-g([Y,Z],X) + g([X,Z],Y) + g([Z,X],Y) - g([Z,Y],X) = 0$, where we have used $[X,Z] = -[Z,X]$ and $[Z,Y] = -[Y,Z]$. Therefore,
\[
2g(\nabla_X Y - \nabla_Y X, Z) = F(X,Y,Z) - F(Y,X,Z) = 2g([X,Y], Z)
\]
for all $Z$, and the non-degeneracy of $g$ implies $\nabla_X Y - \nabla_Y X = [X,Y]$.

\textit{(5) Compatibility with the metric.} We add $F(X,Y,Z) + F(X,Z,Y)$. The terms in Lie derivatives give $Xg(Y,Z) + Xg(Z,Y) = 2Xg(Y,Z)$ by symmetry of $g$, while $Yg(Z,X) - Yg(X,Z) = 0$ and $Zg(Y,X) - Zg(X,Y) = 0$ vanish. The three pairs of bracket terms also cancel by antisymmetry: $g([X,Y],Z) + g([Y,X],Z) = 0$, and analogously for the others. In conclusion,
\[
2g(\nabla_X Y, Z) + 2g(Y, \nabla_X Z) = F(X,Y,Z) + F(X,Z,Y) = 2Xg(Y,Z),
\]
that is, $Xg(Y,Z) = g(\nabla_X Y, Z) + g(Y, \nabla_X Z)$.
\end{proof}

\begin{remark}[Geometric meaning]\label{rem:interpretacion_levi_civita}
The two properties that characterize the Levi-Civita connection have clear geometric interpretations:
\begin{itemize}
	\item \textbf{Compatibility with the metric}: parallel transport along any curve preserves lengths and angles between vectors. In particular, if $V$ and $W$ are parallel fields along a curve $\gamma$, then $g(V,W)$ is constant.
	\item \textbf{Zero torsion}: ``infinitesimal parallelograms'' close, that is, parallel transport does not introduce spurious rotations. Formally, the coordinate fields $\partial_i, \partial_j$ satisfy $\nabla_{\partial_i}\partial_j = \nabla_{\partial_j}\partial_i$.
\end{itemize}
\end{remark}

\subsubsection{Curvature}\label{subsec:curvatura_riemanniana}

In Euclidean space $\mathbb{R}^n$, partial derivatives commute: $\partial_i \partial_j f = \partial_j \partial_i f$ for every smooth function $f$. Analogously, the covariant derivatives of a constant vector field satisfy $\nabla_{\partial_i} \nabla_{\partial_j} Z = \nabla_{\partial_j} \nabla_{\partial_i} Z$. On a general Riemannian manifold, however, this commutativity no longer holds: the result of covariantly differentiating a vector field depends on the order in which the derivations are performed. The \textit{Riemann curvature tensor} quantifies exactly this non-commutativity, and constitutes the central geometric invariant that distinguishes curved spaces from flat ones: a Riemannian manifold is locally isometric to $\mathbb{R}^n$ if and only if its Riemann tensor vanishes identically.

\begin{definition}[Riemann curvature tensor]\label{def:tensor_riemann}
Let $(\mathcal{M}, g)$ be a Riemannian manifold with Levi-Civita connection $\nabla$. The \textit{Riemann curvature tensor} is the map $R\colon \Gamma(TM)^3 \to \Gamma(TM)$ defined by:
\[
R(X,Y)Z = \nabla_X \nabla_Y Z - \nabla_Y \nabla_X Z - \nabla_{[X,Y]} Z
\]
for any vector fields $X, Y, Z \in \Gamma(TM)$. The tensor $R$ is $C^\infty(\mathcal{M})$-multilinear in its three arguments, which guarantees that its value at a point depends only on the values of $X$, $Y$ and $Z$ at that point.
\end{definition}

\begin{definition}[Sectional curvature]\label{def:curvatura_seccional}
Let $(\mathcal{M}, g)$ be a Riemannian manifold and $\sigma \subset T_p\mathcal{M}$ a 2-plane spanned by linearly independent vectors $u, v \in T_p\mathcal{M}$. The \textit{sectional curvature}\footnote{Sectional curvature generalizes the Gaussian curvature of surfaces to arbitrary dimensions. For a surface ($\dim \mathcal{M} = 2$), there is a unique 2-plane at each tangent point, and $K$ coincides with the classical Gaussian curvature. The terminology ``sectional'' comes from the fact that $K(\sigma)$ is the Gaussian curvature of the surface generated by considering the ``section'' of $\mathcal{M}$ by geodesics in the directions of $\sigma$.} of $\sigma$ is defined as:
\begin{equation}\label{eq:curvatura_seccional}
K(\sigma) = K(u,v) = \frac{\langle R(u,v)v, u \rangle_g}{\|u\|^2\|v\|^2 - \langle u,v \rangle_g^2}
\end{equation}
where the denominator is the square of the area of the parallelogram spanned by $u$ and $v$. Geometrically, $K(\sigma)$ measures the rate at which initially parallel geodesics in directions within $\sigma$ converge ($K > 0$) or diverge ($K < 0$).
\end{definition}

\begin{definition}[Ricci tensor and scalar curvature]\label{def:ricci_escalar}
The \textit{Ricci tensor} $\mathrm{Ric}$ is the contraction of the Riemann tensor:
\[
R_{ij} = R^k_{\ ikj} = \sum_{k} \langle R(\partial_i, \partial_k)\partial_j, \partial_k \rangle_g
\]
Intuitively, $R_{ij}$ sums the sectional curvatures of all 2-planes containing the direction $\partial_i$.

The \textit{scalar curvature} $\mathcal{R} = g^{ij}R_{ij}$ provides an average measure of curvature at each point, obtained as the trace of the Ricci tensor (equivalently, as the sum of sectional curvatures over an orthonormal basis).
\end{definition}

\begin{remark}[Geometric meaning of sectional curvature]\label{rem:significado_curvatura}
The sectional curvature $K$ determines fundamental properties that generalize Euclid's postulates:
\begin{itemize}
    \item \textbf{Sum of angles (local Gauss-Bonnet theorem)}: In a geodesic triangle
          of area $A$: $\alpha + \beta + \gamma = \pi + K \cdot A$. This generalizes the
          Euclidean theorem that the sum of the angles of a triangle is $\pi$.
    \item \textbf{Area growth}: The area of a geodesic disk of radius $r$ on a surface of constant curvature $K$ satisfies:
          \[
          A(r) = 2\pi \left( \frac{1 - \cos(\sqrt{K} \cdot r)}{K} \right) \xrightarrow{K \to 0} \pi r^2
          \]
          For $K > 0$ it grows more slowly than $\pi r^2$; for $K < 0$ it grows exponentially (see \autoref{sec:area_disco_geodesico} of Appendix~\ref{ch:teoremas_complementarios} for the derivation).
    \item \textbf{Parallel geodesics} (see \autoref{subsec:geodesicas_variedades}):
          The deviation $\xi(t)$ between initially parallel geodesics satisfies the
          Jacobi equation: $\ddot{\xi} + K\xi = 0$, whose solutions are constant ($K=0$),
          oscillatory ($K>0$), or exponential ($K<0$).
\end{itemize}
\end{remark}

\subsection{Geodesics on Riemannian Manifolds}\label{subsec:geodesicas_variedades}

\textit{Geodesics} are the curves that (locally) minimize length, generalizing straight lines to curved spaces. The Levi-Civita connection introduced in \autoref{subsec:conexion_levi_civita} provides the covariant derivative needed to formalize them: a geodesic is a curve whose tangent vector is parallel along itself.

\subsubsection{Definition and properties of geodesics}\label{subsec:geodesicas_definicion}

\begin{definition}[Geodesic]\label{def:geodesica}
A curve $\gamma: I \to \mathcal{M}$ is a \textit{geodesic} if its tangent vector is parallel along itself:
\[
\nabla_{\dot{\gamma}} \dot{\gamma} = 0
\]
In local coordinates, this is equivalent to the differential equation:
\[
\ddot{\gamma}^k + \Gamma^k_{ij} \dot{\gamma}^i \dot{\gamma}^j = 0
\]
where we use Einstein's summation convention.
\end{definition}

\textbf{Physical and geometric interpretation of geodesics.}\label{rem:interpretacion_geodesicas} Geodesics generalize the concept of ``straight line'' to curved spaces:
\begin{itemize}
	\item \textbf{Variational principle}: They are extremals of the length functional $L(\gamma)$.
	\item \textbf{Principle of inertia}: A free particle (with no external forces) follows geodesics.
	\item \textbf{Local minimization}: Between sufficiently close points, geodesics are the shortest paths.
\end{itemize}

\textbf{Concrete examples}:
\begin{itemize}
	\item In $\mathbb{R}^n$ with the Euclidean metric: straight lines.
	\item In $S^2$: great circles (meridians, equator).
	\item In $\mathbb{H}^n$ (Poincaré model): arcs of circles perpendicular to the boundary, or diameters.
\end{itemize}

\subsubsection{Exponential map and normal coordinates}\label{subsec:mapa_exponencial}

The geodesics defined in the previous section provide the ``straight paths'' of a Riemannian manifold. The exponential map turns this notion into an analytical tool: it assigns to each tangent vector $v \in T_p\mathcal{M}$ the point reached by ``walking'' the geodesic in direction $v$ for unit time, thus establishing a local diffeomorphism $T_p\mathcal{M} \to \mathcal{M}$. Its inverse, the logarithmic map, ``projects'' nearby points to the tangent space. The resulting normal coordinates make the metric Euclidean at $p$ to first order, with errors governed by the curvature (\autoref{rem:coordenadas_normales_linealizacion}). These three tools are the operational basis of the geodesic convolutions and gauge-equivariant architectures that we will see in \autoref{subsec:geodesicas_relevancia_gdl}.

\begin{definition}[Riemannian exponential map]\label{def:mapa_exponencial_riemanniano}
Let $(\mathcal{M}, g)$ be a Riemannian manifold and $p \in \mathcal{M}$. The \textit{exponential map} $\exp_p: T_p\mathcal{M} \to \mathcal{M}$ assigns to each vector $v \in T_p\mathcal{M}$ the point:
\[
\exp_p(v) = \gamma_v(1)
\]
where $\gamma_v$ is the unique geodesic with $\gamma_v(0) = p$ and $\dot{\gamma}_v(0) = v$.
\end{definition}

\begin{proposition}[Properties of the exponential map]\label{prop:propiedades_exp}
Let $(\mathcal{M}, g)$ be a Riemannian manifold and $p \in \mathcal{M}$. The exponential map satisfies:
\begin{enumerate}
	\item $\exp_p(0) = p$.
	\item $\exp_p(tv) = \gamma_v(t)$ for all $t \in [0,1]$ and $v \in T_p\mathcal{M}$ such that $v$ belongs to the domain of $\exp_p$.
	\item $d(\exp_p)_0 = \operatorname{id}_{T_p\mathcal{M}}$.
	\item $\exp_p$ is a diffeomorphism from a neighborhood of $0 \in T_p\mathcal{M}$ to a neighborhood of $p$ in $\mathcal{M}$.
	\item $d_g(p, \exp_p(v)) = \|v\|_g$ for $v$ sufficiently small.
\end{enumerate}
\end{proposition}

\begin{proof}
\leavevmode
\begin{enumerate}
	\item The geodesic with initial velocity $v = 0$ is the constant curve $\gamma_0(t) = p$ for all $t$, since it trivially satisfies $\nabla_{\dot{\gamma}} \dot{\gamma} = 0$. Evaluating at $t = 1$: $\exp_p(0) = \gamma_0(1) = p$.

	\item We must show that $\gamma_{tv}(1) = \gamma_v(t)$. We define $\sigma(s) = \gamma_v(ts)$, which is a reparametrization of $\gamma_v$. Then $\dot{\sigma}(s) = t \dot{\gamma}_v(ts)$ and:
	\[
	\nabla_{\dot{\sigma}} \dot{\sigma}(s) = t^2 \nabla_{\dot{\gamma}_v} \dot{\gamma}_v(ts) = 0,
	\]
	hence $\sigma$ is a geodesic. Moreover $\sigma(0) = \gamma_v(0) = p$ and $\dot{\sigma}(0) = t \dot{\gamma}_v(0) = tv$. By uniqueness of geodesics (a consequence of the existence and uniqueness theorem for ODEs, see \citep[Theorem~4.10]{lee2018introduction}), $\sigma = \gamma_{tv}$, whence $\exp_p(tv) = \gamma_{tv}(1) = \sigma(1) = \gamma_v(t)$.

	\item Let $v \in T_p\mathcal{M}$ and consider the curve $c(t) = tv$ in $T_p\mathcal{M}$. Then:
	\[
	d(\exp_p)_0(v) = \frac{d}{dt}\bigg|_{t=0} \exp_p(tv) = \frac{d}{dt}\bigg|_{t=0} \gamma_v(t) = \dot{\gamma}_v(0) = v,
	\]
	where the second equality uses property (2). Since this holds for all $v$, we have $d(\exp_p)_0 = \operatorname{id}_{T_p\mathcal{M}}$.

	\item Property (3) shows that $d(\exp_p)_0$ is invertible (in fact, it is the identity). By the inverse function theorem (see \autoref{sec:funcion_inversa} of Appendix~\ref{ch:teoremas_complementarios}), $\exp_p$ is a diffeomorphism of class $C^\infty$ on a neighborhood of $0 \in T_p\mathcal{M}$.

	\item The geodesic $t \mapsto \exp_p(tv) = \gamma_v(t)$ has constant speed $\|\dot{\gamma}_v(t)\|_g$ (since the Levi-Civita connection preserves the norm of the tangent vector along a geodesic: $\frac{d}{dt}\|\dot{\gamma}_v\|_g^2 = 2g(\nabla_{\dot{\gamma}_v}\dot{\gamma}_v, \dot{\gamma}_v) = 0$). Thus $\|\dot{\gamma}_v(t)\|_g = \|\dot{\gamma}_v(0)\|_g = \|v\|_g$ for all $t$. The length of $\gamma_v$ on $[0,1]$ is:
	\[
	L(\gamma_v|_{[0,1]}) = \int_0^1 \|\dot{\gamma}_v(t)\|_g \, dt = \|v\|_g.
	\]
	For $v$ sufficiently small, $\gamma_v|_{[0,1]}$ is minimizing (see \citep[Theorem~6.15]{lee2018introduction}), hence $d_g(p, \exp_p(v)) = L(\gamma_v|_{[0,1]}) = \|v\|_g$. \qedhere
\end{enumerate}
\end{proof}

\begin{definition}[Logarithmic map]\label{def:mapa_logaritmico}
The \textit{logarithmic map} $\log_p: U \to T_p\mathcal{M}$, defined on a neighborhood $U$ of $p$ where $\exp_p$ is a diffeomorphism, is the local inverse of $\exp_p$. For $q \in U$:
\[
\log_p(q) = v \quad \text{where} \quad \exp_p(v) = q
\]
Intuitively, $\log_p(q)$ is the ``vector pointing from $p$ to $q$'' in the tangent space.
\end{definition}

\begin{example}[Exponential map on the sphere $S^2$]\label{ex:mapa_exp_esfera}
Consider the unit sphere $S^2 \subset \mathbb{R}^3$ with the north pole $p = (0, 0, 1)$. The tangent space $T_p S^2$ is identified with the $xy$ plane, that is, $T_p S^2 = \{(v_1, v_2, 0) : v_1, v_2 \in \mathbb{R}\}$. For $v \in T_p S^2$ with $r = \|v\| > 0$ and unit vector $\hat{v} = v/r$, the exponential map is given by:
\[
\exp_p(v) = \cos(r) \cdot p + \sin(r) \cdot \hat{v} = (\sin(r) \hat{v}_1, \, \sin(r) \hat{v}_2, \, \cos(r)).
\]
Geometrically, this formula describes a point at angular distance $r$ from the north pole in the direction $\hat{v}$, which corresponds to traversing the great circle (a geodesic in $S^2$, see \autoref{ex:longitud_esfera}) determined by $v$.

The inverse logarithmic map, for $q \neq -p$, satisfies $\log_p(q) = r \hat{v}$ where $r = \arccos(p \cdot q)$ and $\hat{v}$ is the tangent direction from $p$ to $q$. The \textit{injectivity radius} (the maximum $r$ for which $\exp_p$ restricted to $B(0, r) \subset T_p S^2$ is a diffeomorphism) is $\operatorname{inj}(p) = \pi$: the exponential map ceases to be injective exactly at the antipodal point $-p$, where all geodesics from $p$ meet again.
\end{example}

\begin{definition}[Geodesic normal coordinates]\label{def:coordenadas_normales}
Let $p \in \mathcal{M}$ and $\{e_1, \ldots, e_m\}$ an orthonormal basis of $T_p\mathcal{M}$. The \textit{normal coordinates} centered at $p$ are:
\[
\phi: U \subset \mathcal{M} \to \mathbb{R}^m, \quad \phi(\exp_p(v^i e_i)) = (v^1, \ldots, v^m)
\]
\end{definition}

\begin{proposition}[Properties of normal coordinates]\label{prop:coordenadas_normales}
In normal coordinates centered at $p$:
\begin{enumerate}
	\item Geodesics from $p$ are straight lines through the origin
	\item $g_{ij}(p) = \delta_{ij}$ (Euclidean metric at the origin)
	\item $\Gamma^k_{ij}(p) = 0$ (Christoffel symbols vanish at $p$)
	\item $\frac{\partial g_{ij}}{\partial x^k}(p) = 0$ (first derivatives of the metric vanish)
\end{enumerate}
These properties make normal coordinates ideal for local computations \citep[cf.][Proposition~5.24]{lee2018introduction}.
\end{proposition}

\begin{proof}
\leavevmode
\begin{enumerate}
	\item Let $v = v^i e_i \in T_p\mathcal{M}$ and $\gamma_v(t) = \exp_p(tv)$ the geodesic with initial velocity $v$. By \autoref{def:coordenadas_normales}, the normal coordinates of $\gamma_v(t)$ are:
	\[
	\phi(\gamma_v(t)) = \phi(\exp_p(t v^i e_i)) = (tv^1, \ldots, tv^m),
	\]
	which describes a straight line through the origin in $\mathbb{R}^m$ with direction $(v^1, \ldots, v^m)$.

	\item The coordinate vectors at $p$ are $\frac{\partial}{\partial x^i}\big|_p = d(\exp_p)_0(e_i)$. By \autoref{prop:propiedades_exp}(3), $d(\exp_p)_0 = \operatorname{id}_{T_p\mathcal{M}}$, hence $\frac{\partial}{\partial x^i}\big|_p = e_i$. Since $\{e_1, \ldots, e_m\}$ is an orthonormal basis:
	\[
	g_{ij}(p) = g\!\left(\frac{\partial}{\partial x^i}\bigg|_p, \frac{\partial}{\partial x^j}\bigg|_p\right) = g(e_i, e_j) = \delta_{ij}.
	\]

	\item By property (1), in normal coordinates every geodesic from $p$ is written $\gamma^k(t) = tv^k$, hence $\ddot{\gamma}^k(t) = 0$ and $\dot{\gamma}^i(0) = v^i$. Evaluating the geodesic equation (\autoref{def:geodesica}) at $t = 0$:
	\[
	\Gamma^k_{ij}(p)\, v^i v^j = 0 \quad \text{for all } v \in T_p\mathcal{M}.
	\]
	Substituting $v \to v + w$ and expanding:
	\[
	\Gamma^k_{ij}(p)(v^i + w^i)(v^j + w^j) = \Gamma^k_{ij}(p)\,v^iv^j + \Gamma^k_{ij}(p)\,v^iw^j + \Gamma^k_{ij}(p)\,w^iv^j + \Gamma^k_{ij}(p)\,w^iw^j = 0.
	\]
	The first and fourth terms vanish (by the identity applied to $v$ and $w$ respectively), so that:
	\[
	\Gamma^k_{ij}(p)\,v^iw^j + \Gamma^k_{ij}(p)\,w^iv^j = \bigl(\Gamma^k_{ij}(p) + \Gamma^k_{ji}(p)\bigr)v^i w^j = 0 \quad \text{for all } v, w.
	\]
	This implies $\Gamma^k_{ij}(p) + \Gamma^k_{ji}(p) = 0$. But the Levi-Civita connection has zero torsion (\autoref{def:conexion_levi_civita_gdl}), that is, $\Gamma^k_{ij} = \Gamma^k_{ji}$, hence $2\Gamma^k_{ij}(p) = 0$ and therefore $\Gamma^k_{ij}(p) = 0$.

	\item From the formula for the Christoffel symbols (\autoref{def:conexion_levi_civita_gdl}) and properties (2) and (3):
	\[
	0 = \Gamma^k_{ij}(p) = \frac{1}{2}g^{kl}(p)\left(\frac{\partial g_{il}}{\partial x^j}(p) + \frac{\partial g_{jl}}{\partial x^i}(p) - \frac{\partial g_{ij}}{\partial x^l}(p)\right) = \frac{1}{2}\delta^{kl}\left(\partial_j g_{il} + \partial_i g_{jl} - \partial_l g_{ij}\right)\!(p),
	\]
	where we use $g^{kl}(p) = \delta^{kl}$ (inverse of $\delta_{ij}$). Taking $l = k$:
	\begin{equation}\tag{A}
	\partial_j g_{ik}(p) + \partial_i g_{jk}(p) - \partial_k g_{ij}(p) = 0.
	\end{equation}
	Cyclically permuting $(i,j,k) \to (j,k,i)$ and using the symmetry $g_{ab} = g_{ba}$:
	\begin{equation}\tag{B}
	\partial_k g_{ji}(p) + \partial_j g_{ki}(p) - \partial_i g_{jk}(p) = 0.
	\end{equation}
	Adding (A) and (B), the terms $\partial_i g_{jk}$ and $-\partial_i g_{jk}$ cancel, and by the symmetry $g_{ik} = g_{ki}$ and $g_{ij} = g_{ji}$:
	\[
	2\,\partial_j g_{ik}(p) = 0,
	\]
	whence $\frac{\partial g_{ij}}{\partial x^k}(p) = 0$ for all $i, j, k$. \qedhere
\end{enumerate}
\end{proof}

\begin{remark}[Normal coordinates as a local linearization]\label{rem:coordenadas_normales_linealizacion}
\autoref{prop:coordenadas_normales} states that the metric is Euclidean at $p$ and its first derivatives vanish. The next term in the Taylor expansion of the metric in normal coordinates is:
\begin{equation}\label{eq:expansion_metrica_normales}
g_{ij}(x) = \delta_{ij} - \frac{1}{3} R_{ikjl}(p) \, x^k x^l + O(\|x\|^3),
\end{equation}
where $R_{ikjl}(p)$ are the components of the Riemann curvature tensor (\autoref{def:tensor_riemann}) evaluated at $p$ \citep[Lemma~5.10]{lee2018introduction}. This expansion has two fundamental consequences:
\begin{itemize}
	\item \textbf{Linearization}: in a geodesic neighborhood of radius $r$, the manifold behaves like $\mathbb{R}^m$ with errors of order $O(r^2)$. The ``rate'' at which the geometry deviates from the Euclidean case is governed by the curvature.
	\item \textbf{Bridge to GDL}: geodesic convolutions (\autoref{subsec:geodesicas_relevancia_gdl}) operate on normal-coordinate patches, where this expansion guarantees that the convolutional kernel ``sees'' an approximately flat geometry, with curvature corrections captured by the Jacobian factor.
\end{itemize}
\end{remark}

\subsubsection{Relevance for Geometric Deep Learning}\label{subsec:geodesicas_relevancia_gdl}

The exponential map, its logarithmic inverse, and normal coordinates constitute the operational tools that allow transferring deep learning techniques from Euclidean space to Riemannian manifolds. In this section we formalize the four main applications.

\begin{remark}[Geodesic tools in GDL]\label{rem:geodesicas_gdl}
\leavevmode
\begin{enumerate}
	\item \textbf{Intrinsic distances.} \autoref{prop:propiedades_exp}(5) provides the identity $d_g(p,q) = \|\log_p(q)\|_g$ for $q$ within the injectivity radius of $p$. Unlike the Euclidean distance in the ambient space $\mathbb{R}^N$, the geodesic distance respects the geometry of the domain: on a protein surface, for example, two residues may be 5~\AA{} apart in Euclidean distance but 30~\AA{} in geodesic distance if they are separated by a fold.

	\item \textbf{Geodesic convolutions.} Architectures such as \gls{GeoCNN}~\citep{masci2015geodesic} define convolutional kernels in normal coordinates. Let $\{e_1, \ldots, e_m\}$ be an orthonormal basis of $T_p\mathcal{M}$ and $\{(r_i, \theta_j)\}$ a polar discretization of the tangent space. The geodesic convolution of a signal $f: \mathcal{M} \to \mathbb{R}$ with a kernel $k$ is defined as:
	\begin{equation}\label{eq:convolucion_geodesica}
	(f *_g k)(p) = \sum_{i,j} k(r_i, \theta_j) \cdot f\!\left(\exp_p(r_i \, e_{\theta_j})\right) \cdot J_p(r_i, \theta_j),
	\end{equation}
	where $e_{\theta_j} = \cos\theta_j \, e_1 + \sin\theta_j \, e_2$ is the angular direction and $J_p(r, \theta)$ is a Jacobian factor that corrects the distortion between the tangent patch and the manifold. By the expansion~\eqref{eq:expansion_metrica_normales} of \autoref{rem:coordenadas_normales_linealizacion}, this factor satisfies:
	\[
	J_p(r, \theta) = 1 - \frac{1}{6} K(p) \, r^2 + O(r^3),
	\]
	where $K(p)$ is the Gaussian curvature, so that the correction is negligible for small patches and grows with the curvature.

	\item \textbf{Parallel transport and comparison of features.} To compare feature vectors $f(p) \in T_p\mathcal{M}$ and $f(q) \in T_q\mathcal{M}$ at distinct points, it is necessary to transport one to the tangent space of the other. The Levi-Civita connection (\autoref{def:conexion_levi_civita_gdl}) provides the parallel transport $\Pi_{p \to q}^{\gamma}: T_p\mathcal{M} \to T_q\mathcal{M}$ along the geodesic $\gamma$ joining $p$ and $q$, preserving norms and angles. This operation is fundamental for the gauge-equivariant networks of \autoref{sec:gauges_invarianzas}.

	\item \textbf{Geodesic interpolation and residual connections.} The exponential map allows interpolating between points of the manifold:
	\[
	\gamma(t) = \exp_p(t \cdot \log_p(q)), \quad t \in [0,1],
	\]
	which is the geodesic from $p$ to $q$. This idea generalizes to Riemannian residual connections: in a deep network whose states $x_\ell$ live in $\mathcal{M}$, the residual update
	\begin{equation}\label{eq:residual_riemanniana}
	x_{\ell+1} = \exp_{x_\ell}(F_\ell(x_\ell)),
	\end{equation}
	where $F_\ell(x_\ell) \in T_{x_\ell}\mathcal{M}$ is the output of layer $\ell$, generalizes the Euclidean skip connection $x_{\ell+1} = x_\ell + F_\ell(x_\ell)$ to the case where the state space is not linear.
\end{enumerate}
\end{remark}

\begin{remark}[Gauge ambiguity in normal coordinates]\label{rem:ambiguedad_gauge_normales}
The definition of normal coordinates (\autoref{def:coordenadas_normales}) requires choosing an orthonormal basis $\{e_1, \ldots, e_m\}$ of $T_p\mathcal{M}$. Any other orthonormal basis $\{e_1', \ldots, e_m'\}$ is related to the first by a rotation $O \in O(m)$: $e_i' = O_{ji} e_j$. This non-canonical choice introduces a \textit{gauge ambiguity}: the kernel of the geodesic convolution~\eqref{eq:convolucion_geodesica} transforms as $k'(r, \theta) = k(r, \theta - \alpha)$ under a rotation of the basis by angle $\alpha$. For the network to be invariant under this choice, the kernel must satisfy equivariance constraints that we will study in \autoref{sec:gauges_invarianzas}.
\end{remark}

\subsection{Geodesics on Graphs}\label{subsec:geodesicas_grafos}

Shortest paths on graphs represent the natural discretization of continuous geodesics, establishing a deep connection between differential geometry and graph theory.

\begin{definition}[Geodesic distance on graphs]\label{def:distancia_geodesica_grafo}
For a graph $G = (V, E)$, the \textit{discrete geodesic distance} between vertices $u, v \in V$ is:
\[
d_G(u,v) = \min\{k : \exists \text{ path } u = v_0, v_1, \ldots, v_k = v \text{ with } (v_i, v_{i+1}) \in E\}
\]
That is, the minimum number of edges in a path from $u$ to $v$.
\end{definition}

\begin{definition}[Geodesic neighborhood graph]\label{def:grafo_vecindad_geodesico}
Given a graph $G = (V, E)$ and a radius $r \in \mathbb{N}$, the \textit{geodesic $r$-neighborhood} of a vertex $v$ is:
\[
B_r(v) = \{u \in V : d_G(v, u) \leq r\}
\]
This neighborhood is the discrete analogue of a geodesic ball on a manifold.
\end{definition}

\textbf{Convergence of graphs to manifolds.}\label{rem:convergencia_grafos_variedades} When a graph $G_n$ is constructed by sampling $n$ points from a manifold $\mathcal{M}$ and connecting nearby points, under appropriate conditions:
\begin{enumerate}
	\item The discrete geodesic distances $d_{G_n}$ converge to the continuous geodesic distances $d_g$ as $n \to \infty$ and the connectivity parameter is scaled appropriately (\autoref{thm:convergencia_distancias_geodesicas}).

	\item The \textit{normalized graph Laplacian} $\mathcal{L}_{G_n}$ converges to the \textit{Laplace-Beltrami operator} $\Delta_g$ of the manifold, which is fundamental for defining diffusion and intrinsic spectral analysis (\autoref{thm:convergencia_grafos_variedades_detallado}).
\end{enumerate}
This convergence theoretically justifies the use of graphs as computationally tractable approximations of continuous manifolds in \gls{GDL} applications.

\subsubsection{Convergence Theorem for Geodesic Distances}

\begin{theorem}[Convergence of shortest-path distances to geodesic distances]\label{thm:convergencia_distancias_geodesicas}
	Let $(\mathcal{M}, g)$ be a compact, connected Riemannian manifold of dimension $d$, with bounded sectional curvature. Let $X_n = \{x_1, \ldots, x_n\}$ be a set of points sampled i.i.d. according to a distribution with density $p$ such that $p(x) \geq p_{\min} > 0$ for all $x \in \mathcal{M}$.

	We construct the graph $G_n = (X_n, E_n)$ where $(x_i, x_j) \in E_n$ if and only if $\|x_i - x_j\| \leq \epsilon_n$, with weights $w_{ij} = \|x_i - x_j\|$ (Euclidean distance in the ambient space).

	Let $d_{G_n}(x_i, x_j)$ be the shortest-path distance in the graph (sum of weights along the shortest path) and $d_g(x_i, x_j)$ the geodesic distance in $\mathcal{M}$.

	If $\epsilon_n \to 0$ and $n \epsilon_n^d / \log n \to \infty$ as $n \to \infty$, then with high probability:
	\[
		\sup_{x_i, x_j \in X_n} \left| d_{G_n}(x_i, x_j) - d_g(x_i, x_j) \right| \to 0 \quad \text{as } n \to \infty
	\]
\end{theorem}

\begin{proof}
	The proof, following the argument of \citet{bernstein2000graph}, proceeds in five steps.

	\textbf{Step 1: Local metric comparison.}
	Let $|K| \leq \kappa$ be the sectional curvature bound of $\mathcal{M}$. For two points $x, y \in \mathcal{M}$ with $d_g(x, y) \leq \epsilon$, the relationship between the Euclidean distance in the ambient space $\mathbb{R}^N$ and the geodesic distance satisfies:
	\begin{equation}\label{eq:comparacion_metrica_local}
	d_g(x,y)\bigl(1 - C_\kappa \epsilon^2\bigr) \leq \|x - y\| \leq d_g(x,y),
	\end{equation}
	where $C_\kappa > 0$ is a constant depending only on $\kappa$ and the dimension $d$. The \emph{upper bound} is immediate: the straight segment in $\mathbb{R}^N$ is shorter than any curve in $\mathcal{M}$, in particular the geodesic. For the \emph{lower bound}, we consider normal coordinates centered at $x$ (\autoref{prop:coordenadas_normales}). The expansion of the metric in these coordinates~\eqref{eq:expansion_metrica_normales} gives:
	\[
	g_{ij}(z) = \delta_{ij} - \frac{1}{3} R_{ikjl}(x) z^k z^l + O(\|z\|^3).
	\]
	Let $\gamma$ be the geodesic from $x$ to $y$ with $\|z\| \leq \epsilon$. The length of $\gamma$ in normal coordinates satisfies:
	\[
	L(\gamma) = \int_0^1 \sqrt{g_{ij}(\gamma(t))\dot{\gamma}^i\dot{\gamma}^j} \, dt \leq \int_0^1 \|\dot{\gamma}(t)\|\bigl(1 + C'_\kappa \epsilon^2\bigr) \, dt,
	\]
	where we use that $|R_{ikjl}| \leq C\kappa$ by the sectional curvature bound. Since in normal coordinates $\gamma$ is identified with a segment in $\mathbb{R}^d$ (plus second-order corrections), we obtain $d_g(x,y) \leq \|x - y\|(1 + C'_\kappa \epsilon^2)$, which upon inverting gives the lower bound of~\eqref{eq:comparacion_metrica_local}.

	\textbf{Step 2: Probabilistic covering lemma.}
	We define the \emph{covering radius} of the sample $X_n$ as:
	\[
	r_n := \max_{q \in \mathcal{M}} \min_{x_i \in X_n} d_g(q, x_i).
	\]
	We claim that, under the condition $n\epsilon_n^d / \log n \to \infty$, with high probability:
	\begin{equation}\label{eq:radio_recubrimiento}
	r_n = O\!\left(\left(\frac{\log n}{n}\right)^{1/d}\right) = o(\epsilon_n).
	\end{equation}
	To prove this, we fix $\delta > 0$ and construct a $\delta$-covering of $\mathcal{M}$ (\autoref{def:numero_recubrimiento}). By the volume expansion of geodesic balls (a consequence of~\eqref{eq:expansion_metrica_normales}), there exists a constant $c_d > 0$ such that $\operatorname{Vol}(B_g(q, \delta)) \geq c_d \, \delta^d$ for $\delta$ sufficiently small. Since $p(x) \geq p_{\min} > 0$, the probability that \emph{no} sample point falls in $B_g(q, \delta)$ is:
	\[
	\bigl(1 - p_{\min} \, c_d \, \delta^d\bigr)^n \leq \exp(-n \, p_{\min} \, c_d \, \delta^d).
	\]
	The number of balls in a $\delta$-covering is $O(\delta^{-d})$ (by compactness of $\mathcal{M}$). Applying a union bound:
	\[
	\Pr[r_n > \delta] \leq O(\delta^{-d}) \cdot \exp(-n \, p_{\min} \, c_d \, \delta^d).
	\]
	Choosing $\delta = C\bigl(\frac{\log n}{n}\bigr)^{1/d}$ with $C$ sufficiently large, the exponent dominates the polynomial factor and $\Pr[r_n > \delta] \to 0$. The condition $n\epsilon_n^d/\log n \to \infty$ guarantees $r_n = o(\epsilon_n)$.

	\textbf{Step 3: Lower bound for $d_{G_n}$.}
	Let $x_i = z_0, z_1, \ldots, z_L = x_j$ be an optimal path in $G_n$, with $(z_l, z_{l+1}) \in E_n$ and therefore $\|z_l - z_{l+1}\| \leq \epsilon_n$. Applying Step~1 to each edge:
	\[
	\|z_l - z_{l+1}\| \geq d_g(z_l, z_{l+1})\bigl(1 - C_\kappa \epsilon_n^2\bigr).
	\]
	Summing over all edges and applying the triangle inequality of the geodesic distance (\autoref{prop:dg_metrica}):
	\begin{equation}\label{eq:cota_inferior_dGn}
	d_{G_n}(x_i, x_j) = \sum_{l=0}^{L-1} \|z_l - z_{l+1}\| \geq \bigl(1 - C_\kappa \epsilon_n^2\bigr) \sum_{l=0}^{L-1} d_g(z_l, z_{l+1}) \geq \bigl(1 - C_\kappa \epsilon_n^2\bigr) \, d_g(x_i, x_j).
	\end{equation}

	\textbf{Step 4: Upper bound for $d_{G_n}$.}
	Let $\gamma: [0, L_\gamma] \to \mathcal{M}$ be the minimizing geodesic from $x_i$ to $x_j$ parametrized by arc length, with $L_\gamma = d_g(x_i, x_j)$. We discretize $\gamma$ by choosing equispaced points $p_0 = x_i, p_1 = \gamma(h), p_2 = \gamma(2h), \ldots$ with step $h = \epsilon_n/3$.

	By Step~2, with high probability $r_n = o(\epsilon_n)$; in particular, for $n$ sufficiently large, $r_n < \epsilon_n/6$. For each $p_l$ ($l \neq 0, L_\gamma/h$), there exists a sample point $z_l \in X_n$ with $d_g(z_l, p_l) \leq r_n$. For the endpoints, we take $z_0 = x_i$ and $z_{\lceil L_\gamma/h \rceil} = x_j$.

	We verify connectivity: for consecutive edges, by the triangle inequality in $(\mathcal{M}, d_g)$ and the upper bound of Step~1:
	\[
	\|z_l - z_{l+1}\| \leq d_g(z_l, z_{l+1}) \leq d_g(z_l, p_l) + d_g(p_l, p_{l+1}) + d_g(p_{l+1}, z_{l+1}) \leq 2r_n + h < \epsilon_n,
	\]
	hence $(z_l, z_{l+1}) \in E_n$.

	Summing the weights of the edges of the constructed path:
	\[
	d_{G_n}(x_i, x_j) \leq \sum_l \|z_l - z_{l+1}\| \leq \sum_l d_g(z_l, z_{l+1}).
	\]
	Each term is bounded using the triangle inequality:
	\[
	d_g(z_l, z_{l+1}) \leq d_g(p_l, p_{l+1}) + 2r_n = h + 2r_n.
	\]
	The number of edges is $\lceil L_\gamma / h \rceil \leq L_\gamma/h + 1$, so that:
	\begin{equation}\label{eq:cota_superior_dGn}
	d_{G_n}(x_i, x_j) \leq (L_\gamma/h + 1)(h + 2r_n) = L_\gamma + 2L_\gamma \frac{r_n}{h} + h + 2r_n = d_g(x_i,x_j) + O\!\left(d_g(x_i,x_j) \cdot \frac{r_n}{\epsilon_n}\right) + O(\epsilon_n).
	\end{equation}

	\textbf{Step 5: Uniform convergence.}
	Combining the bounds~\eqref{eq:cota_inferior_dGn} and~\eqref{eq:cota_superior_dGn}:
	\[
	\bigl|d_{G_n}(x_i, x_j) - d_g(x_i, x_j)\bigr| \leq d_g(x_i, x_j) \cdot O\!\left(\max\!\left(\epsilon_n^2, \frac{r_n}{\epsilon_n}\right)\right) + O(\epsilon_n).
	\]
	Since $\mathcal{M}$ is compact, $d_g(x_i, x_j) \leq \operatorname{diam}(\mathcal{M}) < \infty$. Moreover, $\epsilon_n \to 0$ and $r_n/\epsilon_n \to 0$ by~\eqref{eq:radio_recubrimiento}. Therefore:
	\[
	\sup_{x_i, x_j \in X_n} \bigl|d_{G_n}(x_i, x_j) - d_g(x_i, x_j)\bigr| \to 0 \quad \text{as } n \to \infty. \qedhere
	\]
\end{proof}

\begin{remark}[Complementarity of the convergence theorems]\label{rem:complementariedad_convergencias}
\autoref{thm:convergencia_grafos_variedades_detallado} and \autoref{thm:convergencia_distancias_geodesicas} are \emph{complementary} results, not equivalent, that guarantee the fidelity of the discrete approximation from distinct perspectives:
\begin{itemize}
	\item \autoref{thm:convergencia_grafos_variedades_detallado} establishes \textbf{spectral} convergence: the graph Laplacian $L_n$ converges to the Laplace-Beltrami operator $\Delta_\mathcal{M}$, capturing the local differential structure (diffusion, curvature, eigenfunctions).
	\item \autoref{thm:convergencia_distancias_geodesicas} establishes \textbf{metric} convergence: the shortest-path distances $d_{G_n}$ converge to the geodesic distances $d_g$, capturing the global geometry (intrinsic separation between points).
\end{itemize}
The scaling conditions differ revealingly: $n\epsilon^{d+2} \to \infty$ for spectral convergence (more neighbors are needed to estimate second derivatives of the kernel) versus $n\epsilon_n^d / \log n \to \infty$ for metric convergence (it suffices to guarantee connectivity and uniform coverage). Together, these results justify that a sufficiently dense graph approximates the manifold both locally and globally.
\end{remark}

\begin{example}[Numerical approximation: sphere $S^2$]\label{ex:geodesicas_discretas_esfera}
	Consider the unit sphere $S^2 \subset \mathbb{R}^3$ with the induced metric. The geodesic distance between two points $p, q \in S^2$ is:
	\[
		d_g(p, q) = \arccos(\langle p, q \rangle)
	\]
	corresponding to the great-circle arc joining them.

	To illustrate the convergence, consider $p = (1, 0, 0)$ (pole) and $q = (0, 1, 0)$ (equator), with geodesic distance $d_g(p, q) = \pi/2 \approx 1.571$.

	\begin{center}
	\begin{tabular}{ccc}
		\toprule
		$n$ (points) & $\epsilon$ & $d_{G_n}(p, q)$ \\
		\midrule
		100 & 0.5 & 1.62 \\
		500 & 0.3 & 1.59 \\
		2000 & 0.2 & 1.575 \\
		10000 & 0.1 & 1.572 \\
		\bottomrule
	\end{tabular}
	\end{center}

	The table shows how the discrete distance $d_{G_n}$ converges to the geodesic $d_g = \pi/2$ as the sampling increases and the connectivity radius decreases.
\end{example}

	The convergence of discrete geodesic distances has direct implications for the design of neural networks on graphs:

\begin{enumerate}
	\item \textbf{Positional encodings:} Architectures such as GraphGPS \citep{rampasek2022graphgps} use shortest-path distances to construct positional encodings. The theorem guarantees that these encodings approximate intrinsic geodesic distances when the graph is a dense sampling of an underlying manifold.
	\item \textbf{Diffusion vs. attention:} Diffusion on graphs (local propagation) captures information at scale $O(\epsilon)$, while the geodesic distance enables long-range attention that respects the geometry. This justifies hybrid architectures that combine local message passing with global attention mechanisms weighted by geodesic distance.
	\item \textbf{Radius scaling:} The condition $\epsilon_n \sim (n / \log n)^{-1/d}$ of the theorem indicates that the connectivity radius must decrease more slowly than $n^{-1/d}$ to maintain connectivity, but fast enough to capture local geometry.
\end{enumerate}


\section{Gauges: Local Invariances and Coordinate Freedoms}\label{sec:gauges_invarianzas}

Gauge theory provides the mathematical framework for understanding and handling the coordinate freedoms and local invariances that arise in geometric architectures. In theoretical physics, gauge transformations represent internal symmetries that leave the equations of the system invariant; in \gls{GDL}, these ideas translate into neural architectures whose predictions are independent of arbitrary choices of local reference frames.

The mathematical treatment of these invariances rests on the theory of \textbf{principal bundles} (\autoref{subsec:fibrados_principales}), where each point of the base manifold carries an associated space of possible reference frames. The notions of \textbf{connection} (\autoref{subsec:conexiones_fibrados}) and \textbf{parallel transport} (\autoref{subsec:transporte_paralelo_gdl}) allow comparing quantities at distinct points intrinsically, and the \textbf{curvature} (\autoref{subsec:curvatura_estructura}) measures the obstruction to such transport being path-independent. These concepts underpin the gauge-equivariant architectures presented in \autoref{ch:cnn}.

\subsection{Theory of Principal Bundles}\label{subsec:fibrados_principales}

So far we have seen that many data domains possess the structure of a homogeneous space: the Orbit-Stabilizer Theorem (\ref{thm:orbita_estabilizador}) showed that domains such as the sphere $S^{n-1} \cong SO(n)/SO(n-1)$ and the hyperbolic space $\mathbb{H}^n \cong SO^+(n,1)/SO(n)$ are expressed as quotients $G/H$, and homogeneous spaces (\autoref{def:espacio_homogeneo}) formalize this structure on differentiable manifolds (\autoref{def:variedad_diferenciable}) with Lie groups (\autoref{def:grupo_lie}).

However, in a quotient $G/H$, concretely identifying the stabilizer of a point $p$ with the subgroup $H$ requires choosing a representative $g$ of the coset $gH$. This choice is inherently arbitrary: any other representative $g' = gh$ ($h \in H$) describes the same point $p = gH$, but defines a different identification. This arbitrariness is the \textbf{gauge freedom}, and in general there is no canonical way to fix a globally consistent choice over the entire manifold.

The mathematical solution is not to choose: instead of selecting an arbitrary representative at each point, a \textit{principal bundle} records \textbf{all} possible choices simultaneously. To each point $p \in \mathcal{M}$ the full group $G$ of possible reference frames is attached, forming a total space $P$ that encodes the structure of local symmetries intrinsically.

\begin{definition}[Principal Bundle]\label{def:fibrado_principal_formal}
A \textit{principal bundle} with structure group $G$ is a quadruple $(P, \mathcal{M}, \pi, G)$ where:
\begin{enumerate}
    \item $P$ is a differentiable manifold called the \textbf{total space}
    \item $\mathcal{M}$ is a differentiable manifold called the \textbf{base space} (or base manifold)
    \item $\pi: P \to \mathcal{M}$ is a surjective differentiable submersion (\autoref{def:submersion_suave}) called the \textbf{projection}
    \item $G$ is a Lie group that acts \textbf{freely} on the right on $P$, that is, if $u \cdot g = u$ for some $u \in P$, then $g = e$ (equivalently, $\mathrm{Stab}(u) = \{e\}$ for all $u \in P$, cf.\ \autoref{def:estabilizador})
\end{enumerate}
satisfying the following conditions:
\begin{enumerate}
    \item[(a)] \textbf{Fiber diffeomorphic to $G$}: for each $p \in \mathcal{M}$, the fiber $\pi^{-1}(p)$ is diffeomorphic to $G$. The group $G$ is the \textit{typical fiber} of the bundle, and consequently $\dim P = \dim \mathcal{M} + \dim G$. The diffeomorphism $\pi^{-1}(p) \cong G$ is not canonical: fixing it is equivalent to choosing a ``gauge'' (a reference frame) at the point $p$

    \item[(b)] \textbf{Action preserves fibers}: $\pi(u \cdot g) = \pi(u)$ for all $u \in P$, $g \in G$. Combined with the freedom of the action (condition 4 above), $G$ acts freely and transitively on each fiber, making each $\pi^{-1}(p)$ a \textit{torsor} (principal homogeneous space) of $G$: a space that ``looks like'' $G$ but without a distinguished identity element

    \item[(c)] \textbf{Local trivializations}: for each $p \in \mathcal{M}$, there exists an open neighborhood $U \ni p$ and a diffeomorphism $\phi_U: \pi^{-1}(U) \to U \times G$ such that $\phi_U(u \cdot g) = (\pi(u), \phi_2(u) \cdot g)$, where $\phi_2$ denotes the second component. Locally, the bundle is a product $U \times G$, but globally it may have a ``twisted'' topology. If $U$ and $V$ are two open sets with trivializations $\phi_U$, $\phi_V$, the \textit{transition functions} $g_{UV}: U \cap V \to G$ defined by $\phi_V \circ \phi_U^{-1}(p, g) = (p, g_{UV}(p) \cdot g)$ encode how the two trivializations are related and completely determine the global topology of the bundle
\end{enumerate}
For a systematic development of the theory of principal bundles, see~\citep{Kobayashi1996}.
\end{definition}

\begin{remark}[Submanifold structure of the fibers]\label{rem:fibras_subvariedad}
The condition that $\pi: P \to \mathcal{M}$ be a smooth submersion guarantees, by the Regular Level Set Theorem (\autoref{cor:fibras_submersion} of \autoref{ch:teoremas_complementarios}), that each fiber $\pi^{-1}(p)$ is a smooth submanifold of $P$ of dimension $\dim P - \dim \mathcal{M} = \dim G$. This justifies condition (a) of the definition: being submanifolds of the correct dimension on which $G$ acts freely and transitively, the diffeomorphism $\pi^{-1}(p) \cong G$ is determined by the choice of a point $u_0 \in \pi^{-1}(p)$, via the map $g \mapsto u_0 \cdot g$.
\end{remark}

\begin{proposition}[Homogeneous Spaces as Principal Bundles]\label{prop:homogeneo_fibrado}
Let $G$ be a Lie group and $H \leq G$ a closed subgroup. Then the quadruple $(G, G/H, \pi, H)$ is a principal bundle, where $\pi: G \to G/H$ is the canonical projection to the quotient and $H$ acts on $G$ by right multiplication.
\end{proposition}

\begin{proof}
We verify the conditions of \autoref{def:fibrado_principal_formal}:
\begin{enumerate}
	\item[(a)] For each coset $gH \in G/H$, the fiber $\pi^{-1}(gH) = \{gh : h \in H\} = gH$. The left translation $L_g: H \to gH$ defined by $h \mapsto gh$ is a diffeomorphism, hence each fiber is diffeomorphic to $H$.

	\item[(b)] The right action of $H$ on $G$ preserves fibers: for all $g \in G$ and $h \in H$, $\pi(g \cdot h) = (gh)H = gH = \pi(g)$, since $hH = H$.

	\item[(c)] The existence of local trivializations is equivalent to the existence of smooth local sections $s: U \to G$ of the quotient $\pi: G \to G/H$. These exist for every closed subgroup $H$ of a Lie group $G$; see~\citep[Ch.~I, Prop.~5.5]{Kobayashi1996}.
\end{enumerate}
\end{proof}

This proposition connects the two pillars, algebraic and geometric, of the theoretical framework. The homogeneous spaces obtained via the Orbit-Stabilizer Theorem~\ref{thm:orbita_estabilizador} are automatically bases of principal bundles: $S^{n-1} \cong SO(n)/SO(n-1)$ (\autoref{prop:esfera_cociente}), $\mathbb{H}^n \cong SO^+(n,1)/SO(n)$ (\autoref{prop:hiperbolico_cociente}), etc. In each case, the stabilizer $H$ is reinterpreted as the \textbf{fiber} of the bundle (the local symmetries that fix a point of the domain), while the orbit $G/H$ constitutes the \textbf{base} (the data space over which the signals are defined). Thus, the gauge freedoms that we formalize in this section are precisely the freedoms of choice of an element within each fiber $gH$.

\begin{figure}[htbp]
\centering
\includegraphics[width=0.7\textwidth]{figures/fiber_bundle.pdf}
\caption[Structure of a principal bundle]{Structure of a principal bundle. The total space $P$ projects onto the base manifold $\mathcal{M}$ via $\pi$. Each fiber $\pi^{-1}(x)$ is diffeomorphic to the structure group $G$, and local trivializations $\phi_\alpha$ allow representing the bundle as products $U_\alpha \times G$ in open regions.}
\label{fig:fiber_bundle}
\end{figure}

\begin{example}[Bundle of Orthonormal Frames]\label{ex:fibrado_marcos}
Let $(\mathcal{M}, g)$ be an $n$-dimensional Riemannian manifold. The \textit{bundle of orthonormal frames} $\mathrm{FM}(\mathcal{M})$ is a principal bundle $(P, \mathcal{M}, \pi, O(n))$ where:
\begin{itemize}
    \item The total space $P$ consists of all orthonormal frames: $P = \{(p, e_1, \ldots, e_n) : p \in \mathcal{M}, \{e_i\} \text{ orthonormal basis of } T_p\mathcal{M}\}$
    \item The projection $\pi: P \to \mathcal{M}$ sends $(p, e_1, \ldots, e_n) \mapsto p$
    \item The group $O(n)$ acts on the right: $(p, e_1, \ldots, e_n) \cdot A = (p, \sum_j A_{j1}e_j, \ldots, \sum_j A_{jn}e_j)$
\end{itemize}
Each fiber $\pi^{-1}(p)$ is diffeomorphic to $O(n)$ and represents all possible choices of orthonormal frame at the point $p$.
\end{example}

\begin{example}[Trivial Bundle]\label{ex:fibrado_trivial}
The simplest case is the \textit{trivial bundle} $P = \mathcal{M} \times G$ with projection $\pi(p, g) = p$ and action $(p, h) \cdot g = (p, hg)$. This bundle admits a global section $s: \mathcal{M} \to P$ given by $s(p) = (p, e)$, where $e$ is the identity of $G$. Not all principal bundles admit global sections; the existence of such sections is related to the topology of the base space.
\end{example}

In the context of \gls{GDL}, the base space $\mathcal{M}$ represents the domain where the data live (for example, the sphere $S^2$ for spherical data), while the group $G$ encodes the relevant local symmetries (for example, $SO(2)$ for rotations in the tangent plane). The fiber over each point contains all the possible ``local orientations'' or ``reference frames'' at that point.

\subsection{Connections on Principal Bundles}\label{subsec:conexiones_fibrados}

A connection on a principal bundle provides a way to ``transport'' information between different fibers consistently with the structure of the group. Mathematically, a connection specifies a decomposition of the tangent space to the bundle into ``horizontal'' and ``vertical'' components.

\begin{definition}[Vertical Space]\label{def:espacio_vertical}
Let $(P, \mathcal{M}, \pi, G)$ be a principal bundle. For each $u \in P$, the \textit{vertical space} $V_u P$ is the subspace of $T_u P$ tangent to the fiber:
\[
V_u P = \ker(d\pi_u) = \{X \in T_u P : d\pi_u(X) = 0\}
\]
\end{definition}

\begin{definition}[Fundamental Vector Field]\label{def:campo_vectorial_fundamental}
Let $(P, \mathcal{M}, \pi, G)$ be a principal bundle. For each $A \in \mathfrak{g}$ (\autoref{def:algebra_lie}) and each $u \in P$, the \textit{fundamental vector field} generated by $A$ at $u$ is:
\[
A^*_u = \frac{d}{dt}\bigg|_{t=0} u \cdot \exp(tA)
\]
where $\exp: \mathfrak{g} \to G$ is the exponential map (\autoref{def:mapa_exponencial}). The curve $t \mapsto u \cdot \exp(tA)$ remains entirely within the fiber $\pi^{-1}(\pi(u))$, since the action of $G$ preserves fibers (condition (b) of \autoref{def:fibrado_principal_formal}), whence $A^*_u \in V_u P$.
\end{definition}

\begin{proposition}[Canonical Isomorphism $V_u P \cong \mathfrak{g}$]\label{prop:vertical_lie_isomorfismo}
For each $u \in P$, the map $\sigma_u: \mathfrak{g} \to V_u P$ defined by $\sigma_u(A) = A^*_u$ is an isomorphism of vector spaces.
\end{proposition}

\begin{proof}
Consider the \textit{orbit map} $\Phi_u: G \to P$ defined by $\Phi_u(g) = u \cdot g$. Its differential at the identity is $\sigma_u = (d\Phi_u)_e: T_eG = \mathfrak{g} \to T_uP$, whose image is contained in $V_uP$ by \autoref{def:campo_vectorial_fundamental}.

\textbf{Linearity.} As the differential of a smooth map, $\sigma_u = (d\Phi_u)_e$ is a linear map.

\textbf{Injectivity.} Since the action of $G$ is free, $\Phi_u$ is injective. Moreover, $\Phi_u$ is an immersion (\autoref{def:inmersion_suave}): if $(d\Phi_u)_e(A) = 0$, the curve $t \mapsto u \cdot \exp(tA)$ would have zero velocity at $t = 0$, which by freedom of the action implies $A = 0$. Hence $\sigma_u$ is injective.

\textbf{Surjectivity.} Since $\pi$ is a submersion (\autoref{def:submersion_suave}), $d\pi_u$ is surjective and $\dim V_u P = \dim \ker(d\pi_u) = \dim P - \dim \mathcal{M} = \dim G = \dim \mathfrak{g}$. An injective linear map between spaces of equal finite dimension is an isomorphism.
\end{proof}

\begin{definition}[Connection on a Principal Bundle]\label{def:conexion_fibrado}
A \textit{connection} on the principal bundle $(P, \mathcal{M}, \pi, G)$ is a 1-form (\autoref{def:1_forma_diferencial}) $\omega \in \Omega^1(P, \mathfrak{g})$ with values in the Lie algebra $\mathfrak{g}$ of $G$ that satisfies:
\begin{enumerate}
    \item \textbf{Normalization}: $\omega(A^*_u) = A$ for all $A \in \mathfrak{g}$ and $u \in P$, where $A^*$ is the fundamental vector field generated by $A$
    \item \textbf{Equivariance}: $R_g^* \omega = \mathrm{Ad}_{g^{-1}} \omega$ for all $g \in G$, where $R_g: P \to P$ denotes the right action and $\mathrm{Ad}$ is the adjoint representation
\end{enumerate}
Equivalently, a connection determines a \textit{horizontal distribution} $H_u P = \ker(\omega_u) \subset T_u P$ complementary to the vertical space, such that $T_u P = H_u P \oplus V_u P$ and $H_{u \cdot g} P = (R_g)_* H_u P$.

For a detailed exposition, see~\citep{Taubes2011}.
\end{definition}

\textbf{Geometric Interpretation.} The horizontal distribution $H_u P$ specifies which directions in the total space are ``purely spatial'' (parallel to the base) versus ``purely gauge'' (along the fiber). A curve $\tilde{\gamma}$ in $P$ is \textit{horizontal} if $\dot{\tilde{\gamma}}(t) \in H_{\tilde{\gamma}(t)} P$ for all $t$; such curves represent transport ``without rotation'' of the reference frame.

\subsection{Parallel Transport}\label{subsec:transporte_paralelo_gdl}

Parallel transport is the mechanism that allows comparing objects (vectors, tensors, or more generally elements of the bundle) located at different points of the base manifold. This comparison is essential to define operations such as convolution on non-Euclidean manifolds.

\begin{definition}[Parallel Transport on Principal Bundles]\label{def:transporte_paralelo_formal}
Let $(P, \mathcal{M}, \pi, G)$ be a principal bundle with connection $\omega$. For a smooth path $\gamma: [0,1] \to \mathcal{M}$, the \textit{parallel transport} along $\gamma$ is the map
\[
\Gamma_\gamma: \pi^{-1}(\gamma(0)) \to \pi^{-1}(\gamma(1))
\]
defined as follows: given $u_0 \in \pi^{-1}(\gamma(0))$, there exists a unique \textit{horizontal lift} $\tilde{\gamma}: [0,1] \to P$ such that:
\begin{enumerate}
    \item $\pi \circ \tilde{\gamma} = \gamma$ (the lift projects onto $\gamma$)
    \item $\tilde{\gamma}(0) = u_0$ (initial condition)
    \item $\dot{\tilde{\gamma}}(t) \in H_{\tilde{\gamma}(t)} P$ for all $t$ (horizontality)
\end{enumerate}
Then $\Gamma_\gamma(u_0) = \tilde{\gamma}(1)$.
\end{definition}

\begin{theorem}[Properties of Parallel Transport]\label{thm:propiedades_transporte}
Parallel transport satisfies the following fundamental properties:
\begin{enumerate}
    \item \textbf{Equivariance}: $\Gamma_\gamma(u \cdot g) = \Gamma_\gamma(u) \cdot g$ for all $g \in G$
    \item \textbf{Composition}: $\Gamma_{\gamma_2 * \gamma_1} = \Gamma_{\gamma_2} \circ \Gamma_{\gamma_1}$, where $\gamma_2 * \gamma_1$ denotes the concatenation of paths
    \item \textbf{Inversion}: $\Gamma_{\gamma^{-1}} = (\Gamma_\gamma)^{-1}$, where $\gamma^{-1}(t) = \gamma(1-t)$
\end{enumerate}
\end{theorem}

\begin{proof}
\textbf{Equivariance.} Let $\tilde{\gamma}$ be the horizontal lift of $\gamma$ with $\tilde{\gamma}(0) = u$, and define $\tilde{\gamma}_g(t) = \tilde{\gamma}(t) \cdot g$. We verify the three conditions of \autoref{def:transporte_paralelo_formal}:
\begin{enumerate}
    \item \textbf{Projection}: since the action of $G$ preserves fibers, $\pi(\tilde{\gamma}_g(t)) = \pi(\tilde{\gamma}(t)) = \gamma(t)$.
    \item \textbf{Initial condition}: $\tilde{\gamma}_g(0) = \tilde{\gamma}(0) \cdot g = u \cdot g$.
    \item \textbf{Horizontality}: $\dot{\tilde{\gamma}}_g(t) = (R_g)_* \dot{\tilde{\gamma}}(t) \in (R_g)_* H_{\tilde{\gamma}(t)} P = H_{\tilde{\gamma}(t) \cdot g} P$, where the last equality uses the $G$-invariance of the horizontal distribution (\autoref{def:conexion_fibrado}).
\end{enumerate}
By uniqueness of the horizontal lift, $\tilde{\gamma}_g$ is the lift of $\gamma$ from $u \cdot g$, hence $\Gamma_\gamma(u \cdot g) = \tilde{\gamma}_g(1) = \tilde{\gamma}(1) \cdot g = \Gamma_\gamma(u) \cdot g$.

\textbf{Composition.} Let $\tilde{\gamma}_1$ be the horizontal lift of $\gamma_1$ from $u_0$, with endpoint $u_1 = \tilde{\gamma}_1(1)$, and let $\tilde{\gamma}_2$ be the horizontal lift of $\gamma_2$ from $u_1$. The concatenated curve
\[
\tilde{\gamma}(t) = \begin{cases} \tilde{\gamma}_1(2t), & t \in [0, \tfrac{1}{2}], \\ \tilde{\gamma}_2(2t - 1), & t \in [\tfrac{1}{2}, 1], \end{cases}
\]
satisfies: it projects onto $\gamma_2 * \gamma_1$, starts from $u_0$ and is horizontal on each segment. By uniqueness, it is the horizontal lift of $\gamma_2 * \gamma_1$ from $u_0$, and its endpoint is $\tilde{\gamma}_2(1) = \Gamma_{\gamma_2}(\Gamma_{\gamma_1}(u_0))$.

\textbf{Inversion.} Let $\tilde{\gamma}$ be the horizontal lift of $\gamma$ from $u_0$ with endpoint $u_1 = \tilde{\gamma}(1)$. Define $\tilde{\gamma}^{-}(t) = \tilde{\gamma}(1-t)$. Then: $\pi(\tilde{\gamma}^{-}(t)) = \gamma(1-t) = \gamma^{-1}(t)$; $\tilde{\gamma}^{-}(0) = u_1$; and $\dot{\tilde{\gamma}}^{-}(t) = -\dot{\tilde{\gamma}}(1-t) \in H_{\tilde{\gamma}(1-t)}P$, since the horizontal subspaces are vector spaces. By uniqueness, $\tilde{\gamma}^{-}$ is the lift of $\gamma^{-1}$ from $u_1$, hence $\Gamma_{\gamma^{-1}}(u_1) = \tilde{\gamma}^{-}(1) = u_0 = \Gamma_\gamma^{-1}(u_1)$.

For an alternative exposition, see~\citep[Ch.~II, Prop.~3.2]{Kobayashi1996}.
\end{proof}

\begin{figure}[htbp]
\centering
\includegraphics[width=0.85\textwidth]{figures/parallel_transport.pdf}
\caption[Parallel transport and holonomy]{Parallel transport along a curve $\gamma$ on a manifold. A tangent vector at $\gamma(0)$ is transported along the path keeping itself ``parallel'' according to the connection. Holonomy arises when the vector returns to its initial point after traversing a closed loop.}
\label{fig:parallel_transport}
\end{figure}

\textbf{Differential Equation of Transport.} In terms of the connection form, the horizontal lift $\tilde{\gamma}$ satisfies the differential equation:
\[
\omega(\dot{\tilde{\gamma}}(t)) = 0
\]
If we express $\tilde{\gamma}(t) = \tilde{\gamma}_0(t) \cdot g(t)$ where $\tilde{\gamma}_0$ is a local section and $g: [0,1] \to G$, then $g(t)$ satisfies:
\[
\frac{dg}{dt} + \omega(\dot{\gamma}(t)) g(t) = 0
\]
with formal solution $g(1) = \mathcal{P} \exp\left(-\int_0^1 \omega(\dot{\gamma}(t)) dt\right) g(0)$, where $\mathcal{P}$ denotes path ordering.

\subsection{Curvature and the Structure Equation}\label{subsec:curvatura_estructura}

The curvature of a connection measures the obstruction to the integrability of the horizontal distribution, or equivalently, how much the parallel transport depends on the path chosen between two points.

\begin{definition}[Curvature of a Connection]\label{def:curvatura_formal}
Let $\omega$ be a connection on the principal bundle $(P, \mathcal{M}, \pi, G)$. The \textit{curvature} $\Omega_\omega$ is the 2-form with values in $\mathfrak{g}$ defined via the exterior derivative (\autoref{def:derivada_exterior}) and the exterior product (\autoref{def:producto_exterior_k_formas}) by:
\begin{equation}\label{eq:curvatura_definicion}
\Omega_\omega = d\omega + \frac{1}{2}[\omega \wedge \omega]
\end{equation}
where $[\omega \wedge \omega](X, Y) = [\omega(X), \omega(Y)]$ uses the Lie bracket in $\mathfrak{g}$.

More explicitly, for vector fields $X, Y$ on $P$:
\[
\Omega_\omega(X, Y) = d\omega(X, Y) + [\omega(X), \omega(Y)]
\]
The curvature is a horizontal and equivariant 2-form~\citep{Chern2000}.
\end{definition}

\begin{theorem}[Cartan's Structure Equation]\label{thm:cartan_estructura}
The curvature $\Omega_\omega$ and the connection $\omega$ satisfy \textit{Cartan's structure equation}:
\begin{equation}\label{eq:estructura_cartan}
d\omega = \Omega_\omega - \frac{1}{2}[\omega \wedge \omega]
\end{equation}
Moreover, the curvature satisfies the \textit{Bianchi identity}:
\begin{equation}\label{eq:bianchi}
d\Omega_\omega = [\omega \wedge \Omega_\omega]
\end{equation}
\end{theorem}

\begin{proof}
The structure equation is simply a rearrangement of the definition of curvature. For the Bianchi identity, we apply $d$ to both sides of \eqref{eq:curvatura_definicion}:
\[
d\Omega_\omega = d(d\omega) + \frac{1}{2}d[\omega \wedge \omega] = \frac{1}{2}([d\omega \wedge \omega] - [\omega \wedge d\omega])
\]
Substituting $d\omega = \Omega_\omega - \frac{1}{2}[\omega \wedge \omega]$ and using the properties of the Lie bracket (Jacobi identity), we obtain $d\Omega_\omega = [\omega \wedge \Omega_\omega]$.
\end{proof}

\begin{definition}[Holonomy]\label{def:holonomia}
Let $\gamma: [0,1] \to \mathcal{M}$ be a closed loop based at $p$ (that is, $\gamma(0) = \gamma(1) = p$). The \textit{holonomy} of the connection $\omega$ along $\gamma$ is the element $g_\gamma \in G$ such that:
\[
\Gamma_\gamma(u) = u \cdot g_\gamma
\]
for any $u \in \pi^{-1}(p)$.

The \textit{holonomy group} $\mathrm{Hol}_p(\omega)$ is the subgroup of $G$ generated by the holonomies of all loops based at $p$.
\end{definition}

\textbf{Curvature as Infinitesimal Holonomy.} The curvature can be interpreted as the ``infinitesimal holonomy'': for an infinitesimal loop enclosing an area $dA$ in the base manifold, the holonomy is approximately $\exp(\Omega_\omega \cdot dA)$. If the curvature vanishes ($\Omega_\omega = 0$), the connection is \textit{flat} and parallel transport is independent of the path (for homotopic paths).

\subsection{Gauge Invariance in Neural Networks}\label{subsec:invarianza_gauge_redes}

The preceding concepts have direct implications for the design of equivariant neural architectures. Modern geometric architectures must respect gauge invariances to ensure that predictions do not depend on arbitrary choices of coordinates or reference frames.

\begin{definition}[Gauge Transformation]\label{def:transformacion_gauge}
A \textit{gauge transformation} is a smooth map $g: \mathcal{M} \to G$ that assigns to each point $p \in \mathcal{M}$ an element $g(p)$ of the structure group. The set of all gauge transformations forms the \textit{gauge group}:
\[
\mathcal{G} = C^\infty(\mathcal{M}, G)
\]
with the operation of pointwise multiplication: $(g_1 \cdot g_2)(p) = g_1(p) \cdot g_2(p)$.
\end{definition}

\begin{definition}[Gauge Equivariant Function]\label{def:funcion_gauge_equivariante_gdl}
Let $\rho: G \to GL(V)$ be a representation of $G$ on a vector space $V$. A function $f: P \to V$ on the principal bundle is \textit{gauge equivariant} (or of type $\rho$) if:
\begin{equation}
f(u \cdot g) = \rho(g^{-1}) f(u)
\end{equation}
for all $u \in P$ and $g \in G$.
\end{definition}

\begin{figure}[htbp]
\centering
\includegraphics[width=0.75\textwidth]{figures/gauge_transformation.pdf}
\caption[Gauge transformation on a graph]{Gauge transformation on a graph. A gauge transformation $g: V \to G$ independently rotates the reference frame at each vertex, transforming the features $f_i$ into $\rho(g_i)f_i$. The dashed arcs indicate the local rotation applied at each node.}
\label{fig:gauge_transformation}
\end{figure}

Gauge transformations act on the functions of the bundle, and a neural architecture $\Phi$ is gauge equivariant if it commutes with this action. Formally, if the output is a scalar (global prediction), the network is \textit{gauge invariant}:
\begin{equation}\label{eq:gauge_invariance_network}
\Phi(g \cdot f) = \Phi(f) \quad \text{for all } g: V \to G
\end{equation}
This guarantees that:

\begin{itemize}
    \item \textbf{Networks on graphs}: The output is invariant under reordering of nodes (permutation gauge)
    \item \textbf{Networks on manifolds}: Predictions are independent of the chosen parametrization or local chart
    \item \textbf{Equivariant networks}: Features transform coherently under changes of local reference frame
\end{itemize}

\textbf{Connection with Theoretical Physics.} Gauge theory in \gls{GDL} is analogous to gauge theories in physics, where the fundamental forces (electromagnetism, weak force, strong force) are described via connections on principal bundles with groups $U(1)$, $SU(2)$ and $SU(3)$ respectively. In \gls{GDL}, the gauge group encodes the local symmetries relevant to the learning problem.

The practical implementation of these ideas in concrete neural architectures, including Gauge Equivariant CNNs, is developed in detail in \autoref{sec:gauge_cnns} of \autoref{ch:cnn}.

\subsection{Examples of Gauge Equivariance in GDL}\label{subsec:ejemplos_gauge_gdl}

To illustrate the practical relevance of the concepts of gauge equivariance, we present concrete examples of how these ideas manifest in \gls{GDL} applications.

\begin{example}[Gauge Equivariance on 3D Triangulated Meshes]\label{ex:gauge_mallas_3d}
Consider a 3D surface represented as a triangular mesh $\mathcal{M} = (V, F)$, where $V$ is the set of vertices and $F$ the set of triangular faces. At each vertex $v \in V$, the tangent plane $T_v\mathcal{M}$ is well defined geometrically, but \textbf{there is no canonical choice of orthonormal basis} $(e_1^v, e_2^v)$ for this plane.

The gauge problem arises when we try to define vector fields or directional features on the mesh:
\begin{itemize}
    \item \textbf{Gauge choice}: For each vertex $v$, we choose an orthonormal basis $(e_1^v, e_2^v)$ of the tangent plane. This choice is arbitrary: any rotation $R_\theta \in SO(2)$ produces another equally valid basis.

    \item \textbf{Gauge transformation}: If we rotate the basis at vertex $v$ by an angle $\theta_v$, the components of a vector $\mathbf{u} = u_1 e_1^v + u_2 e_2^v$ transform as:
    \begin{equation}
    \begin{pmatrix} u_1' \\ u_2' \end{pmatrix} = \begin{pmatrix} \cos\theta_v & \sin\theta_v \\ -\sin\theta_v & \cos\theta_v \end{pmatrix} \begin{pmatrix} u_1 \\ u_2 \end{pmatrix}
    \end{equation}

    \item \textbf{Gauge equivariance requirement}: A neural network that processes vector features on the mesh must produce predictions that are \textbf{independent} of the choice of local bases. Mathematically, if $f: \mathcal{F}(\mathcal{M}) \to \mathcal{Y}$ is the learned function, it must satisfy:
    \begin{equation}
    f(\rho(g) \cdot \mathbf{x}) = \rho'(g) \cdot f(\mathbf{x})
    \end{equation}
    where $g = \{R_{\theta_v}\}_{v \in V}$ represents a gauge transformation (independent rotation at each vertex).
\end{itemize}
\end{example}

\begin{example}[Impact of the Gauge Choice on Predictions]\label{ex:impacto_gauge_predicciones}
To illustrate why gauge equivariance is crucial, consider a naive neural network that processes vector features on meshes without respecting gauge invariances.

\textbf{Problem:} Given a surface with a principal curvature field, we want to predict the direction of maximum curvature at each point.

\textbf{Non-gauge-equivariant architecture:} An \gls{MLP} that takes the components $(u_1, u_2)$ of vectors in local coordinates and predicts the components of the curvature direction.

\textbf{Observed failure:} If we change the orientation of the local frame at a vertex (a gauge transformation), the predictions change inconsistently:
\begin{itemize}
    \item Original frame: The network predicts direction $(\hat{d}_1, \hat{d}_2) = (0.8, 0.6)$
    \item After rotating the frame by $45^\circ$: The network predicts $(\hat{d}_1', \hat{d}_2') = (0.3, 0.9)$: a completely different direction!
\end{itemize}

\textbf{Gauge-equivariant architecture:} Networks such as Gauge Equivariant Mesh CNNs~\citep{cohen2019gaugeequivariantconvolutional} use convolutions that transform coherently under gauge changes. Their predictions satisfy:
\begin{equation}
\hat{\mathbf{d}}' = R_\theta \hat{\mathbf{d}}
\end{equation}
where $R_\theta$ is the same rotation applied to the local frame. Thus, although the \textit{numerical components} change, the predicted \textit{geometric direction} remains invariant.
\end{example}

\begin{example}[Gauge Equivariance for 3D Molecules]\label{ex:gauge_moleculas}
In the modeling of three-dimensional molecules, gauge equivariance manifests in a particularly important way for architectures that predict tensorial properties.

Consider a molecule with $n$ atoms at positions $\{\mathbf{r}_i\}_{i=1}^n \subset \mathbb{R}^3$. We want to predict the \textbf{polarizability tensor} $\boldsymbol{\alpha} \in \mathbb{R}^{3 \times 3}$, which describes how the electronic distribution responds to external electric fields.

\textbf{Fundamental physical property:} If we rotate the molecule by $R \in SO(3)$, the polarizability tensor transforms as:
\begin{equation}
\boldsymbol{\alpha}' = R \boldsymbol{\alpha} R^T
\end{equation}
This is a rank-2 tensor transformation under the group $SO(3)$.

\textbf{Gauge equivariant architectures:} Networks such as NequIP~\citep{batzner2022nequip} and Allegro represent intermediate features via \textbf{spherical tensors}, which transform according to irreducible representations of $SO(3)$. The polarizability tensor is constructed via equivariant tensor products:
\begin{equation}
\boldsymbol{\alpha} = \sum_{\ell=0,2} W_\ell \left( \mathbf{Y}^{(\ell)} \otimes \mathbf{Y}^{(\ell)} \right)
\end{equation}
where $\mathbf{Y}^{(\ell)}$ are spherical harmonics of order $\ell$ and $W_\ell$ are learned weights. This construction automatically guarantees the correct transformation under rotations.

\textbf{Practical advantage:} Predictions are physically consistent independently of the arbitrary orientation of the molecule in space, eliminating the need for data augmentation by rotations and significantly improving the efficiency of training.
\end{example}

\textbf{Synthesis of the examples.} The three preceding examples illustrate a common principle: when the data have local symmetries (gauge freedoms), neural architectures must respect these symmetries to produce physically meaningful predictions. The theory of principal bundles provides the rigorous mathematical framework for understanding and designing such architectures.


\section{Synthesis: The Unifying Framework of GDL}\label{sec:sintesis_marco_unificador}

The five fundamental G's of \gls{GDL} (Geometry, Groups, Graphs, Geodesics and Gauges) form a deeply interconnected conceptual framework: \textbf{Geometry} determines the structure of the domain and defines which \textbf{Groups} of symmetries are natural; \textbf{Graphs} discretize continuous geometries while preserving topological and spectral properties; \textbf{Geodesics} define optimal paths that respect the geometry; and \textbf{Gauges} ensure that the representations are intrinsic. From this perspective, the differences between neural architectures reflect choices about three geometric aspects: the \textbf{structure of the domain} (Euclidean spaces, graphs, manifolds), the \textbf{symmetry group} (trivial, translations, permutations) and the \textbf{connectivity} (global, fixed neighborhoods, adaptive). Traditional \glspl{ANN}, operating in $\mathbb{R}^d$ with trivial symmetries (\autoref{ch:neural_networks}), constitute the base case: they do not exploit any geometric structure, so that all invariance must be learned from the data. \gls{GDL} solves this by incorporating the symmetries directly into the architecture via the principle of equivariance (\autoref{subsec:equivarianza_invarianza}).

\subsection{Message Passing: The Unifying Framework}\label{subsec:message_passing_unificador}

\textbf{Message Passing Neural Networks} (\glspl{MPNN}) provide the unifying mathematical framework that underlies all \gls{GDL} architectures. Their conceptual foundation lies in the graph structure introduced in \autoref{def:grafo_neuronal} of \autoref{ch:neural_networks} and \autoref{def:vecindad_grafo} of \autoref{sec:grafos_discretizacion}, where the data are organized as nodes connected by edges that encode structural relationships.

This perspective reveals that the apparent differences between \glspl{ANN}, \glspl{CNN}, \glspl{Transformer} and \glspl{GNN} are manifestations of different graph topologies and connectivity strategies over the same fundamental paradigm: the propagation of information through structured neighborhoods.

\begin{definition}[Message Passing Neural Network]\label{def:mpnn_unificador}
An \textit{\gls{MPNN}} operates on a graph $\mathcal{G} = (V, E)$ with node features $\{h_v\}_{v \in V}$ and edge features $\{e_{uv}\}_{(u,v) \in E}$. The update at layer $\ell$ consists of two phases:

\begin{enumerate}
    \item \textbf{Message Passing}:
    \begin{equation}
    m_v^{(\ell+1)} = \sum_{u \in \mathcal{N}(v)} M_\ell(h_v^{(\ell)}, h_u^{(\ell)}, e_{vu})
    \end{equation}
      \item \textbf{Node Update}:
    \begin{equation}
    h_v^{(\ell+1)} = U_\ell(h_v^{(\ell)}, m_v^{(\ell+1)})
    \end{equation}
\end{enumerate}

where $M_\ell$ is the message function, $U_\ell$ is the update function, and $\mathcal{N}(v)$ is the neighborhood of node $v$ (\autoref{def:vecindad_grafo}). This formulation naturally generalizes discrete data structures. Note that this scheme operates exclusively over nodes and edges; its generalization to higher-order $k$-simplices, \textit{simplicial message passing}, is developed in \autoref{subsubsec:tdl_mpnn_general}.
\end{definition}

\begin{figure}[htbp]
\centering
\includegraphics[width=\textwidth]{figures/mpnn_unification.pdf}
\caption[Unification of architectures via Message Passing]{The four main Deep Learning architectures as instances of the \gls{MPNN} framework. \textbf{ANN}: complete graph with specific weights for each pair. \textbf{CNN}: regular grid with weights shared by translation. \textbf{Transformer}: complete graph with adaptive weights (attention). \textbf{GNN}: arbitrary graph with neighborhood aggregation. The differences reside in the graph topology and the symmetry group $G$.}
\label{fig:mpnn_unification}
\end{figure}

\subsubsection{ANNs as Trivial MPNNs}

Traditional \glspl{ANN} can be interpreted as the simplest case of \glspl{MPNN}, where each node connects to all the others with uniform weights:

\begin{proposition}[\gls{ANN} as an \gls{MPNN} on a uniform complete graph]\label{prop:ann_mpnn}
A fully connected layer $h^{(\ell+1)} = \sigma(W h^{(\ell)} + b)$ is equivalent to an \gls{MPNN} where:
\begin{itemize}
    \item \textbf{Graph}: Complete graph $K_d$ where each feature is a node
    \item \textbf{Message}: $M(h_v, h_u, e_{vu}) = W_{vu} h_u$ (specific weight for each pair)
    \item \textbf{Aggregation}: Sum over all input features
    \item \textbf{Update}: $U(h_v, m_v) = \sigma(m_v + b_v)$
\end{itemize}
\end{proposition}

\begin{proof}
Let $h^{(\ell)} \in \mathbb{R}^d$ be the feature vector at layer $\ell$. The fully connected layer computes:
\[
h_v^{(\ell+1)} = \sigma\!\left(\sum_{u=1}^{d} W_{vu} h_u^{(\ell)} + b_v\right)
\]
Interpreting each component as a node of the complete graph $K_d$, this expression coincides exactly with the message passing (\autoref{def:mpnn_unificador}) with message $M(h_v, h_u, e_{vu}) = W_{vu} h_u^{(\ell)}$, aggregation $m_v = \sum_{u \in \mathcal{N}(v)} M(\cdot) = \sum_{u=1}^d W_{vu} h_u^{(\ell)}$, and update $U(h_v, m_v) = \sigma(m_v + b_v)$.
\end{proof}

This perspective reveals why \glspl{ANN} lack \textbf{geometric inductive bias}: they treat all connections between features as equally important, ignoring any spatial, topological, or similarity structure that may exist in the data.

\subsubsection{Convolutions as Regular Message Passing}

\glspl{CNN} constitute the first nontrivial case of message passing, where the grid structure imposes a fundamental geometric inductive bias: translational equivariance.

\begin{proposition}[\gls{CNN} as an \gls{MPNN} on a regular grid]\label{prop:cnn_mpnn}
A convolutional layer over a regular grid $\mathbb{Z}^d$ with kernel of support $\mathcal{K}$ is equivalent to an \gls{MPNN} where:
\begin{itemize}
    \item \textbf{Graph}: Regular grid $\mathbb{Z}^d$ with edges $\mathcal{N}(v) = \{u \in \mathbb{Z}^d : v - u \in \mathcal{K}\}$
    \item \textbf{Message}: $M(h_v, h_u, e_{vu}) = W_{v-u} h_u$ (weights shared by relative position)
    \item \textbf{Aggregation}: Sum over the local neighborhood $m_v = \sum_{u \in \mathcal{N}(v)} W_{v-u} h_u$
    \item \textbf{Update}: $U(h_v, m_v) = \sigma(m_v + b)$
\end{itemize}
\end{proposition}

\begin{proof}
Let $h^{(\ell)}: \mathbb{Z}^d \to \mathbb{R}^{c_{\mathrm{in}}}$ be the signal at layer $\ell$ and $\{W_\delta\}_{\delta \in \mathcal{K}}$ the convolutional kernel with $W_\delta \in \mathbb{R}^{c_{\mathrm{out}} \times c_{\mathrm{in}}}$. The discrete convolution (\autoref{def:convolucion_discreta}) computes:
\[
h_v^{(\ell+1)} = \sigma\!\left(\sum_{\delta \in \mathcal{K}} W_\delta\, h_{v-\delta}^{(\ell)} + b\right) = \sigma\!\left(\sum_{u \in \mathcal{N}(v)} W_{v-u}\, h_u^{(\ell)} + b\right)
\]
where the change of variable $u = v - \delta$ transforms the sum over kernel displacements into a sum over neighbors of the grid graph. This expression coincides with the message passing (\autoref{def:mpnn_unificador}) with message $M(h_v, h_u, e_{vu}) = W_{v-u} h_u^{(\ell)}$, aggregation $m_v = \sum_{u \in \mathcal{N}(v)} M(\cdot)$ and update $U(h_v, m_v) = \sigma(m_v + b)$. The complete constructive equivalence, including the layer-by-layer correspondence by induction, is established in \autoref{thm:equivalencia_constructiva_cnn_gnn}.
\end{proof}

The key difference with respect to \glspl{ANN} is that the weights $W_{v-u}$ depend exclusively on the relative position, not on the absolute pair $(v, u)$. This weight sharing is precisely the implementation of translational equivariance as a geometric inductive bias (\autoref{theorem:convolution_translation_equivariance}), drastically reducing the number of parameters from $O(n^2 c^2)$ to $O(|\mathcal{K}| c^2)$.

\subsubsection{Attention as Adaptive Message Passing}

\glspl{Transformer} represent the complementary case: the communication weights are computed dynamically from the content of the nodes, not from their position.

\begin{proposition}[\gls{Transformer} as an \gls{MPNN} on an adaptive complete graph]\label{prop:attention_mpnn}
A self-attention layer (\autoref{def:self_attention_basica}) over a sequence of $n$ tokens is equivalent to an \gls{MPNN} where:
\begin{itemize}
    \item \textbf{Graph}: Complete graph $K_n$ where each token is a node
    \item \textbf{Message}: $M(h_i, h_j) = \alpha_{ij}\, V h_j$, with attention coefficients
    \[
    \alpha_{ij} = \frac{\exp\!\left(\frac{(Qh_i)^T(Kh_j)}{\sqrt{d_k}}\right)}{\sum_{m=1}^n \exp\!\left(\frac{(Qh_i)^T(Kh_m)}{\sqrt{d_k}}\right)}
    \]
    \item \textbf{Aggregation}: Weighted sum $m_i = \sum_{j=1}^n M(h_i, h_j)$
    \item \textbf{Update}: $U(h_i, m_i) = m_i$
\end{itemize}
\end{proposition}

\begin{proof}
Let $\{h_1, \ldots, h_n\} \subset \mathbb{R}^d$ be the input sequence and $Q, K, V \in \mathbb{R}^{d \times d_k}$ the projection matrices. The self-attention layer (\autoref{def:self_attention_basica}) computes:
\[
h_i^{(\ell+1)} = \sum_{j=1}^n \alpha_{ij}\, V h_j^{(\ell)}
= \sum_{j=1}^n \underbrace{\alpha_{ij}\, V h_j^{(\ell)}}_{M(h_i, h_j)}
\]
Identifying the complete graph $K_n$ with $\mathcal{N}(i) = \{1, \ldots, n\}$, this expression is exactly the message passing (\autoref{def:mpnn_unificador}) with message $M(h_i, h_j) = \alpha_{ij} V h_j$ and identity update $U(h_i, m_i) = m_i$. The softmax normalization guarantees $\alpha_{ij} \in (0,1)$ with $\sum_{j=1}^n \alpha_{ij} = 1$, so that the aggregation is a convex combination, unlike the uniform sum of \glspl{ANN} (\autoref{prop:ann_mpnn}) and \glspl{CNN} (\autoref{prop:cnn_mpnn}).
\end{proof}

The fundamental difference with respect to the previous architectures is that the weights $\alpha_{ij}$ are functions of the content $(h_i, h_j)$, neither fixed (\glspl{ANN}) nor dependent on the relative position (\glspl{CNN}). This adaptability enables global connectivity at the cost of quadratic complexity $O(n^2)$ in the sequence length. \autoref{ch:attention} develops the geometric extensions of this mechanism.

\subsubsection{TDL as Higher-Order Message Passing}\label{subsubsec:tdl_mpnn_general}

The three previous architectures (\glspl{ANN}, \glspl{CNN} and \glspl{Transformer}) operate exclusively over \textbf{nodes} ($0$-simplices) of a graph, propagating information along edges. \gls{TDL} breaks this restriction: it assigns signals to $k$-simplices of arbitrary order and propagates information simultaneously through the lower and upper adjacencies defined by the boundary operators. This extension incorporates a \textbf{topological inductive bias} that captures higher-order relationships inaccessible to standard message passing.

\begin{definition}[Simplicial signal]\label{def:senal_simplicial_gdl}
	A \textit{simplicial signal of order $k$} on a simplicial complex $\mathcal{K}$ is a function $\mathbf{x}: \mathcal{K}_k \to \mathbb{R}^d$ that assigns a feature vector $\mathbf{x}_\sigma \in \mathbb{R}^d$ to each $k$-simplex $\sigma \in \mathcal{K}_k$.
\end{definition}

\begin{definition}[Simplicial message passing]\label{def:mensaje_simplicial_gdl}
	\textit{Simplicial message passing} generalizes message passing on graphs (\autoref{def:mpnn_unificador}) via three types of communication for each $k$-simplex $\sigma$:
	\begin{enumerate}
		\item \textbf{Lower message} (via $\partial_k^T$): Information from $(k{-}1)$-simplices in the boundary of $\sigma$:
		\[
			\mathbf{m}_\sigma^{\mathrm{down}} = \bigoplus_{\tau \in \mathcal{B}^{\mathrm{down}}(\sigma)} M^{\mathrm{down}}\!\left(\mathbf{h}_\sigma, \mathbf{h}_\tau\right)
		\]
		where $\mathcal{B}^{\mathrm{down}}(\sigma) = \{\tau : [\partial_k^T \partial_k]_{\sigma,\tau} \neq 0\}$ are the $k$-simplices adjacent through shared faces.

		\item \textbf{Upper message} (via $\partial_{k+1}$): Information from $(k{+}1)$-simplices containing $\sigma$:
		\[
			\mathbf{m}_\sigma^{\mathrm{up}} = \bigoplus_{\rho \in \mathcal{B}^{\mathrm{up}}(\sigma)} M^{\mathrm{up}}\!\left(\mathbf{h}_\sigma, \mathbf{h}_\rho\right)
		\]
		where $\mathcal{B}^{\mathrm{up}}(\sigma) = \{\rho : [\partial_{k+1} \partial_{k+1}^T]_{\sigma,\rho} \neq 0\}$ are the $k$-simplices that share a $(k{+}1)$-simplex.

		\item \textbf{Update}:
		\[
			\mathbf{h}_\sigma' = U\!\left(\mathbf{h}_\sigma, \mathbf{m}_\sigma^{\mathrm{down}}, \mathbf{m}_\sigma^{\mathrm{up}}\right)
		\]
	\end{enumerate}
	where $\bigoplus$ denotes an aggregation function (sum, mean, maximum), $M^{\mathrm{down}}, M^{\mathrm{up}}$ are learnable message functions, and $U$ is the update function \citep{bodnar2021weisfeiler}.
\end{definition}

\begin{proposition}[\gls{TDL} as a generalized instance of the \gls{MPNN} framework]\label{prop:tdl_mpnn}
Simplicial message passing (\autoref{def:mensaje_simplicial_gdl}) is a generalized instance of the \gls{MPNN} framework (\autoref{def:mpnn_unificador}) where:
\begin{itemize}
    \item \textbf{Domain}: Simplicial complex $\mathcal{K}$, with the $k$-simplices $\sigma \in \mathcal{K}_k$ as ``nodes''
    \item \textbf{Neighborhood}: Double neighborhood $\mathcal{N}(\sigma) = \mathcal{B}^{\mathrm{down}}(\sigma) \cup \mathcal{B}^{\mathrm{up}}(\sigma)$, determined by the lower Laplacian $L_k^{\mathrm{down}} = \partial_k^T \partial_k$ and upper Laplacian $L_k^{\mathrm{up}} = \partial_{k+1} \partial_{k+1}^T$
    \item \textbf{Message}: Combination of lower and upper messages:
    \[
    M(\mathbf{h}_\sigma) = \left(\bigoplus_{\tau \in \mathcal{B}^{\mathrm{down}}(\sigma)} M^{\mathrm{down}}(\mathbf{h}_\sigma, \mathbf{h}_\tau),\; \bigoplus_{\rho \in \mathcal{B}^{\mathrm{up}}(\sigma)} M^{\mathrm{up}}(\mathbf{h}_\sigma, \mathbf{h}_\rho)\right)
    \]
    \item \textbf{Update}: $U(\mathbf{h}_\sigma, \mathbf{m}_\sigma^{\mathrm{down}}, \mathbf{m}_\sigma^{\mathrm{up}})$ learnable
\end{itemize}
In particular, for $k=0$ the standard \gls{MPNN} on graphs is recovered.
\end{proposition}

\begin{proof}
We verify the correspondence with the three components of the \gls{MPNN} framework. Let $\mathcal{K}$ be a simplicial complex and fix $k \geq 0$. We consider the graph $\mathcal{G}_k = (\mathcal{K}_k, E_k)$ whose nodes are the $k$-simplices and whose edges encode the adjacency given by $L_k$:

\textbf{(1) Message.} For each $k$-simplex $\sigma$, simplicial message passing (\autoref{def:mensaje_simplicial_gdl}) computes two messages:
\[
\mathbf{m}_\sigma^{\mathrm{down}} = \bigoplus_{\tau \in \mathcal{B}^{\mathrm{down}}(\sigma)} M^{\mathrm{down}}(\mathbf{h}_\sigma, \mathbf{h}_\tau), \qquad
\mathbf{m}_\sigma^{\mathrm{up}} = \bigoplus_{\rho \in \mathcal{B}^{\mathrm{up}}(\sigma)} M^{\mathrm{up}}(\mathbf{h}_\sigma, \mathbf{h}_\rho)
\]
Defining $\mathcal{N}(\sigma) = \mathcal{B}^{\mathrm{down}}(\sigma) \cup \mathcal{B}^{\mathrm{up}}(\sigma)$ and the joint message $\widetilde{M}(\mathbf{h}_\sigma, \mathbf{h}_\nu) = M^{\mathrm{down}}(\mathbf{h}_\sigma, \mathbf{h}_\nu)$ if $\nu \in \mathcal{B}^{\mathrm{down}}(\sigma)$ and $\widetilde{M}(\mathbf{h}_\sigma, \mathbf{h}_\nu) = M^{\mathrm{up}}(\mathbf{h}_\sigma, \mathbf{h}_\nu)$ if $\nu \in \mathcal{B}^{\mathrm{up}}(\sigma)$, this takes the form $m_\sigma = \bigoplus_{\nu \in \mathcal{N}(\sigma)} \widetilde{M}(\mathbf{h}_\sigma, \mathbf{h}_\nu)$, which corresponds to the message phase of \autoref{def:mpnn_unificador}.

\textbf{(2) Aggregation.} The operator $\bigoplus$ (sum, mean or maximum) plays the role of the permutation-invariant aggregation function of the \gls{MPNN} scheme.

\textbf{(3) Update.} The function $U(\mathbf{h}_\sigma, \mathbf{m}_\sigma^{\mathrm{down}}, \mathbf{m}_\sigma^{\mathrm{up}})$ corresponds directly to the update function $U_\ell$ of the \gls{MPNN} framework.

\textbf{Reduction to $k=0$.} For $0$-simplices (nodes), $\partial_0 = 0$ implies $L_0^{\mathrm{down}} = 0$, so $\mathcal{B}^{\mathrm{down}}(v) = \emptyset$ and the lower message vanishes. The upper message reduces to $\mathbf{m}_v^{\mathrm{up}} = \bigoplus_{u \in \mathcal{B}^{\mathrm{up}}(v)} M^{\mathrm{up}}(\mathbf{h}_v, \mathbf{h}_u)$, where $\mathcal{B}^{\mathrm{up}}(v) = \{u : [\partial_1 \partial_1^T]_{vu} \neq 0\} = \mathcal{N}(v)$ is the standard neighborhood of node $v$ in the graph $1$-skeleton. This recovers exactly the classical message passing on graphs.
\end{proof}

\begin{proposition}[Expressiveness of simplicial message passing]\label{prop:expresividad_simplicial}
	Simplicial message passing is strictly more expressive than message passing on graphs. There exist pairs of non-isomorphic simplicial complexes whose 1-skeletons are identical (and therefore indistinguishable by \glspl{MPNN}), but which simplicial networks distinguish correctly thanks to the higher-order topological information.
\end{proposition}

This result, proved by \citet{bodnar2021weisfeiler}, shows that the topological information encoded in the boundary operators $\partial_k$ provides a discriminative power that surpasses the limitations of the \gls{WL} test established in \autoref{subsubsec:test_weisfeiler_leman}.

This proposition closes the arc of unification of the \gls{MPNN} framework: while \glspl{ANN}, \glspl{CNN} and \glspl{Transformer} encode spatial or relational inductive \textit{biases} over nodes, \gls{TDL} introduces a \textbf{higher-order topological inductive bias} that allows capturing interactions between triangles, tetrahedra and arbitrary $k$-simplices. As shown in \autoref{prop:expresividad_simplicial}, this topological bias provides expressive power strictly superior to that of message passing on graphs, surpassing the \gls{WL} barrier \citep{bodnar2021weisfeiler, hajij2023topological}.

\subsection{Expressive Power and Synthesis of the Unifying Framework}\label{subsec:poder_expresivo_sintesis}

The expressive power of message-passing-based architectures is fundamentally limited by deep theoretical considerations:

\subsubsection{The Weisfeiler-Leman Test}\label{subsubsec:test_weisfeiler_leman}

The \textbf{Weisfeiler-Leman test} (\gls{WL}) is a classical algorithm originally proposed in 1968~\citep{weisfeiler1968reduction} to determine whether two graphs are isomorphic via an iterative process of \textbf{refined coloring}. Its importance in \gls{GDL} lies in the fact that it establishes a fundamental theoretical upper bound for the expressive power of \glspl{MPNN}.

\begin{definition}[Weisfeiler-Leman algorithm (1-WL)]\label{def:weisfeiler_leman}
Let $G = (V, E)$ be a graph with $n$ vertices. The \textit{order-1 Weisfeiler-Leman algorithm} (1-WL) consists of the following iterative coloring process:
\begin{enumerate}
    \item \textbf{Initialization} ($t=0$): Assign to each node $v \in V$ an initial color $c^{(0)}(v)$, typically identical for all nodes (or based on node labels if they exist).

    \item \textbf{Iterative refinement}: At each iteration $t \geq 0$, update the color of each node according to:
    \begin{equation}\label{eq:wl_update}
    c^{(t+1)}(v) = \mathrm{HASH}\left(c^{(t)}(v), \{\{c^{(t)}(u) : u \in \mathcal{N}(v)\}\}\right)
    \end{equation}
    where $\mathcal{N}(v)$ is the neighborhood of $v$, $\{\{\cdot\}\}$ denotes a multiset, and $\mathrm{HASH}$ is an injective function that assigns a new unique color to each distinct pair (old color, multiset of neighboring colors).

    \item \textbf{Convergence}: The algorithm terminates when the color partition stabilizes, that is, when $c^{(t+1)}(v) = c^{(t)}(v)$ for all $v \in V$. This occurs in at most $n$ iterations.

    \item \textbf{Final histogram}: The result of the algorithm is the \textit{color histogram} $\mathbf{h}(G) = (|\{v : c^{(T)}(v) = k\}|)_{k}$, which counts how many nodes have each final color.
\end{enumerate}
\end{definition}

\begin{definition}[WL-equivalence]\label{def:wl_equivalencia}
Two graphs $G_1$ and $G_2$ are \textit{WL-equivalent}, denoted $G_1 \equiv_{\mathrm{WL}} G_2$, if the 1-WL algorithm produces the same color histogram for both graphs.
\end{definition}

\begin{figure}[htbp]
\centering
\includegraphics[width=\textwidth]{figures/weisfeiler_leman.pdf}
\caption[Weisfeiler-Leman color refinement]{Example of the 1-WL color refinement test on a graph of 6 nodes. (a) Uniform initial coloring. (b) After the first iteration, the nodes are distinguished by degree: blue (degree~3), red (degree~4), green (degree~2). (c) The second iteration refines according to the multiset of neighboring colors: $v_3$ and $v_4$ (yellow) are separated from $v_6$ (red) because their neighborhoods have distinct color compositions.}
\label{fig:weisfeiler_leman}
\end{figure}

The deep connection between the WL test and graph networks was rigorously established by~\citet{xu2019how} and~\citet{morris2019weisfeiler}:

\begin{theorem}[WL-MPNN equivalence]\label{thm:wl_mpnn_equivalencia}
Let $\mathcal{G}$ be a class of graphs and let $\mathcal{A}$ be an \gls{MPNN} with injective aggregation and update functions. Then:
\begin{enumerate}
    \item \textbf{Upper bound}: $\mathcal{A}$ cannot distinguish graphs that are WL-equivalent. That is, if $G_1 \equiv_{\mathrm{WL}} G_2$, then $\mathcal{A}(G_1) = \mathcal{A}(G_2)$.

    \item \textbf{Optimality}: There exists an \gls{MPNN} (specifically, a Graph Isomorphism Network or GIN) that is \textit{as expressive as} the WL test. That is, if $G_1 \not\equiv_{\mathrm{WL}} G_2$, there exists a configuration of weights such that $\mathcal{A}(G_1) \neq \mathcal{A}(G_2)$.
\end{enumerate}
\end{theorem}

\begin{proof}[Proof sketch]
The key idea is the structural correspondence between one iteration of the WL algorithm and one layer of an \gls{MPNN}:

\textbf{WL $\rightarrow$ MPNN}: The WL refinement $c^{(t+1)}(v) = \mathrm{HASH}(c^{(t)}(v), \{\{c^{(t)}(u)\}\})$ maps directly to a message-aggregation-update layer:
\[
h_v^{(t+1)} = \phi\left(h_v^{(t)}, \bigoplus_{u \in \mathcal{N}(v)} \psi(h_u^{(t)})\right)
\]
If $\phi$ and $\oplus$ are sufficiently expressive (for example, sum over unique embeddings followed by an MLP), they exactly reproduce the behavior of $\mathrm{HASH}$.

\textbf{MPNN $\rightarrow$ WL}: Conversely, if two graphs are WL-equivalent, they have the same sequence of color partitions. Since \glspl{MPNN} can only aggregate neighborhood information, they cannot ``break'' the symmetry that the WL test preserves.

The complete proof requires showing that the sum aggregator over sufficiently unique embeddings is injective, which is established via the Universal Approximation Theorem for sets~\citep{zaheer2017deepsets}.
\end{proof}

\begin{example}[Non-isomorphic graphs indistinguishable by WL]\label{ex:grafos_wl_indistinguibles}
The simplest example of non-isomorphic graphs that the 1-WL test cannot distinguish are regular graphs of the same degree and size:

Consider two hexagonal graphs:
\begin{itemize}
    \item $G_1$: A cycle of 6 nodes (hexagon): $C_6$
    \item $G_2$: Two disjoint triangles: $K_3 \sqcup K_3$
\end{itemize}

Both graphs have 6 nodes, 6 edges, and all nodes have degree 2. The WL algorithm assigns:
\begin{enumerate}
    \item Iteration 0: All nodes receive the same initial color $c^{(0)} = 1$.
    \item Iteration 1: Each node sees neighbors with color 1, multiset $\{\{1, 1\}\}$. New color: $c^{(1)} = \mathrm{HASH}(1, \{\{1,1\}\}) = 2$ for all.
    \item Convergence: The partition does not change. Final histogram: $\mathbf{h}(G_1) = \mathbf{h}(G_2) = (6)$.
\end{enumerate}

However, $G_1 \not\cong G_2$: the cycle $C_6$ is connected while $K_3 \sqcup K_3$ has two connected components. This limitation shows that standard \glspl{MPNN} cannot detect the global connectivity of regular graphs.

\textbf{Implication for GDL}: If a task requires distinguishing molecules with similar global structures but different connectivities (for example, structural isomers), standard \glspl{MPNN} may fail systematically.
\end{example}

\subsubsection{Over-smoothing: Spectral Analysis}\label{subsubsec:oversmoothing}

\autoref{thm:wl_mpnn_equivalencia} delimits which graphs \glspl{MPNN} can distinguish, but says nothing about what happens when we stack many message-passing layers. The answer, known as \textit{\gls{oversmoothing}}~\citep{li2018deeper}, is that the node representations converge to a uniform value, losing all discriminative capacity. This phenomenon admits an elegant spectral explanation that connects directly with the tools developed in \autoref{subsec:analisis_espectral_grafos_gdl}.

\begin{definition}[Dirichlet energy on graphs]\label{def:energia_dirichlet_grafos}
	Let $\mathcal{G} = (V, E)$ be a graph with Laplacian $L_G$ and let $H \in \mathbb{R}^{n \times d}$ be a matrix of node features ($n = |V|$, $d$ channels). The \textit{Dirichlet energy} of $H$ is
	\[
		E(H) = \operatorname{tr}(H^T L_G H) = \sum_{c=1}^{d} \mathbf{h}_c^T L_G \mathbf{h}_c = \sum_{c=1}^{d} \sum_{(i,j) \in E} W_{ij}(h_{ic} - h_{jc})^2,
	\]
	where $\mathbf{h}_c$ denotes the $c$-th column of $H$.
\end{definition}

The last equality follows from the quadratic form of \autoref{prop:propiedades_algebraicas_laplaciano}. The Dirichlet energy measures the total variation of the signals $H$ over the graph: $E(H) = 0$ if and only if $H$ is constant on each connected component. This is the discrete version of the Dirichlet energy mentioned in \autoref{subsec:analisis_espectral_grafos_gdl}.

\begin{proposition}[Linear MPNN as discrete diffusion]\label{prop:mpnn_difusion}
	Consider an \gls{MPNN} (\autoref{def:mpnn_unificador}) with mean aggregation and no nonlinearity:
	\[
		h_v^{(\ell+1)} = (1 - \alpha) h_v^{(\ell)} + \alpha \sum_{u \in \mathcal{N}(v)} \frac{h_u^{(\ell)}}{\deg(v)},
	\]
	with $\alpha \in (0, 1]$. Let $\tilde{L} = D^{-1}L_G = I - D^{-1}A$ be the random-walk Laplacian, which shares its spectrum with the symmetric normalized Laplacian of \autoref{def:laplaciano_normalizado}. Then the update is written in matrix form as
	\begin{equation}\label{eq:difusion_discreta}
		H^{(\ell+1)} = (I - \alpha \tilde{L})\, H^{(\ell)}.
	\end{equation}
\end{proposition}

\begin{proof}
	For each node $v$, the update rule is:
	\begin{align*}
		h_v^{(\ell+1)} &= (1-\alpha) h_v^{(\ell)} + \alpha \sum_{u \sim v} \frac{h_u^{(\ell)}}{\deg(v)} = h_v^{(\ell)} - \alpha\left(h_v^{(\ell)} - \sum_{u \sim v} \frac{h_u^{(\ell)}}{\deg(v)}\right)\\
		&= h_v^{(\ell)} - \alpha (D^{-1}L_G\, H^{(\ell)})_v = h_v^{(\ell)} - \alpha (\tilde{L}\, H^{(\ell)})_v,
	\end{align*}
	where we use that $\tilde{L} = D^{-1}L_G = I - D^{-1}A$ and that $(D^{-1}A H)_v = \sum_{u \sim v} h_u / \deg(v)$.
\end{proof}

Equation~\eqref{eq:difusion_discreta} is the temporal discretization of the heat equation $\partial_t H = -\alpha \tilde{L} H$ on the graph. Iterating $\ell$ layers, $H^{(\ell)} = (I - \alpha \tilde{L})^\ell H^{(0)}$, which progressively diffuses the signals.

\begin{theorem}[Exponential convergence of over-smoothing]\label{thm:oversmoothing_convergencia}
	Let $\mathcal{G}$ be a connected graph with eigenvalues of the normalized Laplacian $0 = \lambda_1 < \lambda_2 \leq \cdots \leq \lambda_n$ (\autoref{prop:propiedades_espectrales_laplaciano_grafos}) and $\alpha \in (0, 1/\lambda_n)$. Then the Dirichlet energy of the representations after $\ell$ layers of the linear \gls{MPNN} satisfies
	\begin{equation}\label{eq:oversmoothing_bound}
		E\bigl(H^{(\ell)}\bigr) \leq (1 - \alpha\lambda_2)^{2\ell}\, E\bigl(H^{(0)}\bigr).
	\end{equation}
	In particular, $E(H^{(\ell)}) \to 0$ exponentially as $\ell \to \infty$: all representations converge to a single vector.
\end{theorem}

\begin{proof}
	Let $\{u_k\}_{k=1}^n$ be the orthonormal basis of eigenvectors of $L_{\mathrm{sym}} = D^{-1/2} L_G D^{-1/2}$ with eigenvalues $\lambda_k$ (\autoref{def:transformada_fourier_grafos}). We define $\mathbf{f}_c = D^{1/2}\mathbf{h}_c$ for each column of $H$, so that $\mathbf{h}_c^T L_G \mathbf{h}_c = \mathbf{f}_c^T L_{\mathrm{sym}} \mathbf{f}_c$. The Fourier transform decomposes $\mathbf{f}_c = \sum_{k=1}^n \hat{f}_{c,k}\, u_k$.

	The diffusion operator satisfies $D^{1/2}(I - \alpha \tilde{L})D^{-1/2} = I - \alpha L_{\mathrm{sym}}$, which has eigenvalues $1 - \alpha\lambda_k$. After $\ell$ iterations, the spectral coefficients of $\mathbf{f}_c^{(\ell)} = D^{1/2}\mathbf{h}_c^{(\ell)}$ are scaled by $(1-\alpha\lambda_k)^\ell \hat{f}_{c,k}$.

	The Dirichlet energy in the spectral domain is (\autoref{thm:interpretacion_frecuencial_autovectores}):
	\[
		E(H^{(\ell)}) = \sum_{c=1}^d \mathbf{f}_c^{(\ell)T} L_{\mathrm{sym}} \mathbf{f}_c^{(\ell)} = \sum_{c=1}^d \sum_{k=2}^n \lambda_k (1-\alpha\lambda_k)^{2\ell} |\hat{f}_{c,k}|^2,
	\]
	where the term $k=1$ disappears because $\lambda_1 = 0$. Since $\alpha < 1/\lambda_n$, we have $|1-\alpha\lambda_k| < 1$ for $k \geq 2$, and in particular $|1-\alpha\lambda_k| \leq 1-\alpha\lambda_2$ for all $k \geq 2$ (the maximum is attained at $\lambda_2$, the lowest nontrivial frequency). Therefore:
	\[
		E(H^{(\ell)}) \leq (1-\alpha\lambda_2)^{2\ell} \sum_{c=1}^d \sum_{k=2}^n \lambda_k |\hat{f}_{c,k}|^2
		= (1-\alpha\lambda_2)^{2\ell}\, E(H^{(0)}).
	\]
	Since $0 < \alpha\lambda_2 < 1$, the factor $(1-\alpha\lambda_2)^{2\ell} \to 0$ exponentially.
\end{proof}

\begin{remark}[Geometric interpretation of over-smoothing]\label{rem:oversmoothing_geometrico}
	The convergence rate $(1-\alpha\lambda_2)^{2\ell}$ is controlled by the \textbf{spectral gap} $\lambda_2$, the second eigenvalue of the Laplacian. Graphs with large $\lambda_2$ (such as expanders) converge faster to the uniform state, while poorly connected graphs ($\lambda_2 \approx 0$) better preserve the diversity of representations. This is the exact discrete version of heat diffusion on Riemannian manifolds (\autoref{rem:laplace_beltrami_gdl}), where the spectral gap of the Laplace-Beltrami operator controls the mixing rate.

	Among the strategies to mitigate \gls{oversmoothing} stand out residual connections (which preserve the original signal), layer normalization, and topological approaches such as Neural Sheaf Diffusion~\citep{bodnar2022neural}, where scalar diffusion is replaced by a diffusion on sheaves that allows each edge to transport information in different vector spaces.
\end{remark}

\subsubsection{Over-squashing: Topological Bottlenecks}\label{subsubsec:oversquashing}

If \gls{oversmoothing} arises from the loss of high-frequency signal, the complementary phenomenon of \textit{\gls{oversquashing}} is the inability to propagate information over long distances. Intuitively, the information from a distant node must be compressed exponentially in order to ``pass'' through the bottlenecks of the graph.

\begin{definition}[Jacobian sensitivity of an MPNN]\label{def:sensibilidad_mpnn}
	Consider an \gls{MPNN} with representations $h_v^{(\ell)} \in \mathbb{R}^d$. The \textit{sensitivity} of node $v$ with respect to node $u$ after $\ell$ layers is
	\[
		S_{vu}^{(\ell)} = \left\|\frac{\partial h_v^{(\ell)}}{\partial h_u^{(0)}}\right\|,
	\]
	where the norm is the operator norm of the Jacobian matrix $\partial h_v^{(\ell)} / \partial h_u^{(0)} \in \mathbb{R}^{d \times d}$.
\end{definition}

If $S_{vu}^{(\ell)} \approx 0$, node $v$ is essentially insensitive to the input of node $u$: any information that $u$ carries has been lost during propagation.

\begin{proposition}[Exponential decay of sensitivity]\label{prop:oversquashing_exponencial}
	Consider an \gls{MPNN} with Lipschitz message and update functions with constants $c_M$ and $c_U$ respectively, and let $\hat{A} = D^{-1}A$ be the random-walk transition matrix on $\mathcal{G}$. Then for any nodes $v, u \in V$:
	\begin{equation}\label{eq:oversquashing_bound}
		S_{vu}^{(\ell)} \leq (c_U \cdot c_M)^\ell\, [\hat{A}^\ell]_{vu}.
	\end{equation}
	In particular, if $d_\mathcal{G}(v, u) > \ell$ (\autoref{def:distancia_geodesica_grafo}), then $[\hat{A}^\ell]_{vu} = 0$ and the sensitivity is exactly zero.
\end{proposition}

\begin{proof}[Proof sketch]
	We proceed by induction on $\ell$. For $\ell = 1$ and nodes $v \neq u$,
	since $h_v^{(1)} = U(h_v^{(0)}, m_v^{(1)})$, the full chain rule is
	\[
		\frac{\partial h_v^{(1)}}{\partial h_u^{(0)}}
		= \frac{\partial U}{\partial h_v^{(0)}}
		  \cdot \underbrace{\frac{\partial h_v^{(0)}}{\partial h_u^{(0)}}}_{=\,0 \text{ if } v \neq u}
		+ \frac{\partial U}{\partial m_v}
		  \cdot \frac{\partial m_v}{\partial h_u^{(0)}}
		= \frac{\partial U}{\partial m_v}
		  \cdot \frac{\partial m_v}{\partial h_u^{(0)}}.
	\]
	The second factor is nonzero only if $u \in \mathcal{N}(v)$, in which case
	the Lipschitz constant of $M$ contributes $c_M$ and the degree
	normalization contributes the factor $1/\deg(v) = [\hat{A}]_{vu}$. Thus
	$S_{vu}^{(1)} \leq c_U c_M [\hat{A}]_{vu}$.

	For the inductive step, we apply the multivariable chain rule:
	\[
		\frac{\partial h_v^{(\ell+1)}}{\partial h_u^{(0)}} = \sum_{w \in \mathcal{N}(v) \cup \{v\}} \frac{\partial h_v^{(\ell+1)}}{\partial h_w^{(\ell)}} \cdot \frac{\partial h_w^{(\ell)}}{\partial h_u^{(0)}}.
	\]
	Bounding each factor by the Lipschitz constants and using the inductive hypothesis for $\partial h_w^{(\ell)} / \partial h_u^{(0)}$, we obtain the matrix product $\hat{A}^{\ell+1}$ that records all paths of length $\ell+1$ from $u$ to $v$.
\end{proof}

\begin{remark}[Curvature and over-squashing]\label{rem:curvatura_oversquashing}
	The bound~\eqref{eq:oversquashing_bound} reveals that the topology of the graph, encoded in $\hat{A}^\ell$, is the determining factor of \gls{oversquashing}. \citet{topping2022understanding} formalize this connection via the \textbf{discrete Ricci curvature} (Ollivier): edges $(u,v)$ with negative Ricci curvature correspond to bottlenecks where $[\hat{A}^\ell]_{vu}$ decays rapidly. This is the discrete version of the Jacobi equation in Riemannian geometry (\autoref{rem:significado_curvatura}): negative sectional curvature causes geodesics to diverge, reducing the ``volume of information'' that can propagate. Topping et al. propose \textit{rewiring} the graph via Ricci flow, adding edges in regions of negative curvature to alleviate the bottlenecks, an example of how differential geometry prescribes solutions to problems of \glspl{GNN}.
\end{remark}

\subsubsection{The k-WL Hierarchy and Extensions}\label{subsubsec:jerarquia_kwl}

\autoref{thm:wl_mpnn_equivalencia} establishes that \glspl{MPNN} are as expressive as the 1-WL test, which operates over individual nodes. A natural extension is to consider higher-order tests that operate over tuples of nodes.

\begin{definition}[Order-$k$ Weisfeiler-Leman test]\label{def:k_wl}
	The \textit{$k$-\gls{WL} test} operates over $k$-tuples of nodes $\bar{v} = (v_1, \ldots, v_k) \in V^k$ \citep{morris2019weisfeiler}. An initial coloring $c^{(0)}(\bar{v})$ is defined based on the isomorphism type of the subtuple, along with the iterative refinement:
	\[
		c^{(t+1)}(\bar{v}) = \mathrm{HASH}\!\left(c^{(t)}(\bar{v}),\, \bigl\{\!\!\bigl\{ \bigl(c^{(t)}(\bar{v}[w/1]), \ldots, c^{(t)}(\bar{v}[w/k])\bigr) : w \in V \bigr\}\!\!\bigr\}\right),
	\]
	where $\bar{v}[w/i]$ denotes the tuple obtained by replacing the $i$-th component of $\bar{v}$ by $w$.
\end{definition}

\begin{proposition}[Strict expressiveness hierarchy]\label{prop:jerarquia_kwl}
	The Weisfeiler-Leman tests form a strictly increasing hierarchy in expressive power:
	\[
		1\text{-WL} \subsetneq 2\text{-WL} \subsetneq 3\text{-WL} \subsetneq \cdots
	\]
	That is, for each $k \geq 1$ there exist graphs that $k$-WL does not distinguish but $(k+1)$-WL does. The computational complexity of the $k$-WL test is $O(n^k \log n)$ per iteration, where $n = |V|$.
\end{proposition}

\begin{proof}
	The proof of the strict separation requires techniques from finite model logic (in particular, the connection with the Immerman-Lander logic $C^k$) and exceeds the scope of this thesis; see~\citet{morris2019weisfeiler} for the details. The complexity $O(n^k \log n)$ follows directly from the iterative refinement: there are $n^k$ tuples, each iteration requires inspecting $n$ substitutions per tuple, and stabilization occurs in at most $O(n^k)$ iterations with cost $O(\log n)$ for the hash.
\end{proof}

There are various strategies to overcome the 1-WL barrier without incurring the exponential cost of $k$-WL:
\begin{itemize}
	\item \textbf{Simplicial message passing}: As shown by \autoref{prop:expresividad_simplicial}, operating over simplicial complexes ($k$-chains) allows surpassing 1-WL by incorporating higher-order topological information.
	\item \textbf{Structural encodings}: Adding positional encodings based on eigenvectors of the Laplacian (\autoref{def:transformada_fourier_grafos}) or random walks breaks the symmetry of the nodes.
	\item \textbf{Spectral Graph Transformers}: Architectures such as~\citet{kreuzer2021spectral} combine global attention with spectral encodings, connecting the ideas of \autoref{subsec:analisis_espectral_grafos_gdl} with the attention mechanisms of \autoref{ch:attention}.
\end{itemize}

\medskip

The three limitations analyzed (\gls{oversmoothing}, \gls{oversquashing} and the WL barrier) are manifestations of the fundamentally \textbf{local} character of message passing: each layer only propagates information one hop. \gls{oversmoothing} shows that too many layers destroy the discriminative signal; \gls{oversquashing} shows that too few layers prevent long-range communication; and the WL barrier shows that certain global patterns are undetectable locally. The geometric tools developed in this chapter (spectral analysis, discrete curvature, higher-order topology) not only \textit{diagnose} these limitations but also \textit{prescribe} concrete solutions.

From this perspective, we can synthesize the unifying framework that emerges from the preceding analysis. The unifying perspective of \gls{GDL} reveals a natural progression in the geometric sophistication of neural architectures: \glspl{ANN} operate in Euclidean spaces with trivial symmetries; \glspl{CNN} introduce \textbf{spatial inductive bias} via locality and translational equivariance; \glspl{Transformer} generalize this via \textbf{global adaptive connectivity}, where attention implements message passing with dynamic weights; and this paradigm extends naturally to arbitrary graphs with complex topologies.

This theoretical framework is predictive, not merely descriptive: it suggests design principles for new hybrid architectures based on the geometric foundations established. Chapters~\ref{ch:cnn} and~\ref{ch:attention} delve into these architectures from their specific geometric perspectives, providing the complete mathematical foundations for equivariant convolutions and geometric attention mechanisms. Geometry is not optional in Deep Learning: it is the fundamental organizing principle that determines which architectures are effective for which types of data and problems.
